% Generated by build_chapter.py from the GY1971 files. Do not edit by hand; edit the files and rebuild.
\chapter{The academic years 1971--72 and 1972--73}\label{ch:1971}
\noindent\textbf{Scope and source control.}
This chapter covers every one of the 45 requested MCCME records: 229--271,
286 and 288.  The packet contains 99 scanned pages, including overlapping
talks and material outside the requested records.  Their actual handwritten
boundaries govern the entries.  In particular, P000358--P000359 continue
Perelomov's preceding-year talk, while P000450--P000451 and the neighbouring
Buchstaber pages in the 1973 packet are not assigned to year headings.
Record 288 is retained because it was explicitly requested, but its clear
handwritten date, 19 February 1973 on P000480, corrects the catalogue's 1972
date.  Its identifiers therefore begin \texttt{1972-288}; this inclusion does
not constitute a complete companion for 1972--73.  The catalogue is cited as
\citet{mccme2026gelfand}; all substantive status statements are dated November
2026.

\paragraph{The year's mathematical lines.}
The autumn sequence moves from Kirillov's prequantization and polarizations to
Springer fibres, finite incidence geometry and microlocal analysis.  Bernstein's
Godbillon--Vey exposition and Atiyah's spectral-flow discussion connect
characteristic classes, foliations and index theory.  The Vinogradov lectures
organize differential operators through modules, jets and their universal
properties; the formal-geometry discussion links infinite jets with
cohomology of Lie algebras of vector fields.

The spring has a sustained algebraic \(K\)-theory thread: quadratic and
hermitian forms, elementary matrices and Steinberg relations, polynomial and
Laurent constructions, Vaserstein's comparison programme, Segal's categories,
and Quillen's finite-field calculation.  Parallel sequences on polarized
Hodge structures and their variation, sites and derived categories supply the geometric
and categorical language.  Scattering theory returns through trace-class
perturbations and stationary constructions.  Morava's formal-group picture
anticipates the height-based organization of complex bordism.  Dyson's
identity and the continued formal-group discussion close the main packet;
the requested Dobrushin outlier is only a brief field-theory report.

\paragraph{What became of the material.}
The occurrence index has 262 items.  Counts are occurrences, not numbers
of independent discoveries; splitting definitions, methods and claims changes
their granularity.  The 4 catalogue headings and 1 reading
plan have no invented mathematical items.  The 257 core-year occurrences and 5
Dobrushin occurrences are distinguished in the report.

\begin{center}
\begin{tabular}{lrrrrrrr}
\toprule
Class & A & B & C & D & E & F & Total\\
\midrule
Occurrences & 34 & 18 & 3 & 0 & 182 & 25 & 262\\
\bottomrule
\end{tabular}
\end{center}

The predominance of class E reflects serious exposition rather than
unrecorded novelty.  Publication and first circulation must be separated:
Quillen's finite-field calculation, the Adams-operation fibre and the
odd-prime cohomology formula were already announced at the 1970 Congress
\citep[\S\S2--3]{quillen1971cohomology}, although the detailed Annals paper
appeared in 1972.  Conversely, the spectral-flow programme is followed into
\citet[\S7]{atiyahpatodisinger1976spectral3}.  Morava's early circulated notes
are identified as unpublished; later published constructions are not
backdated to them.

\paragraph{Questions and critical qualifications.}
The 9 problem-index entries preserve definite questions even when
they were answered quickly.  The elementary-relations question is linked to
Vaserstein's answer later in the same seminar sequence.  Dyson's exponent-26
question has later root-system explanations.  Pt\'ak's rational-\(p\)
uncertainty was already overtaken by a 1971 paper before his talk; the exact
general fate of the minimum-over-all-norms and extremal questions was not
established by the searches.  Such a verification gap is stated expressly,
and is not evidence that a problem remains open in 2026.  The trace-class
equivalence question has a dedicated later classification paper, but its
full criterion is not quoted without a verified theorem text.

No class C or D item supplies both a securely assessed present frontier and
a justified bounded research contribution in this run.  Accordingly, the
development section contains editorial verifications: the parity obstruction
to nonzero spectral flow for the full de Rham family, three exact algorithms
for Dyson's coefficients through 100, and an exact recovery of the classical
rank-32 lattice bound.  Independent re-derivations and checks found no error;
finite computation is never promoted to an all-degree theorem.

\paragraph{Continuity and limits.}
Earlier links include prequantization (\texttt{1969-178.3--4}),
the distribution-to-foliation problem (\texttt{1970-206.1--2},
\texttt{GP-1970-26}), and the
site and sheaf expositions (\texttt{1968-151}, \texttt{1970-184},
\texttt{1970-226}).  A thematic connection does not by itself justify merging
problem identifiers.  The 127 correction rows distinguish errors
of transcription or translation, uncertain readings, catalogue provenance,
and mathematical errors faithfully preserved in the notes.  Neither a
class label nor a verified bibliographic URL substitutes for the explicit
publication-match and verification fields attached to each occurrence.
\clearpage

\GStalk{229}{1971-09-06}{A. A. Kirillov}{P000360-P000364}{Prequantization and polarizations}
The actual heading on P000360 begins a five-page exposition of Kostant’s geometric quantization: symplectic and Hamiltonian vector fields, integral curvature, line bundles with connection, the prequantum representation, and polarized sections. The notebook’s factors of two and signs do not fit one convention throughout. The supplied catalogue span P000358--P000360 includes the end of the preceding year’s Perelomov discussion and omits most of this talk; those boundary pages are not new occurrences of record 229.

\begin{GSitem}{1971-229.1}{E}{Symplectic form, Darboux coordinates, and Hamiltonian fields}
\GSfield{Kind; source}{definition; \GSpage{P000360}}
\GSfield{The notes}{The notes take a closed nondegenerate two-form \(\omega\) on a \(2n\)-manifold, write \(\omega=\sum dp_j\wedge dq_j\), and identify vector fields and one-forms. The coordinate identification and the names of locally Hamiltonian and Hamiltonian fields are not consistent with the displayed contraction.}
\GSfield{Statement}{On a smooth symplectic manifold the map \(X\mapsto\iota_X\omega\) is an isomorphism. A vector field preserves \(\omega\) exactly when \(d\iota_X\omega=0\); it is Hamiltonian, in a stated sign convention, when that one-form is globally exact. Darboux coordinates exist locally, and \(\omega^n/n!\) is a nonvanishing volume form.}
\GSfield{Reconstruction}{Use the usual unnormalised exterior product and state a contraction sign explicitly; this changes the notebook’s coordinate signs and absorbs its factor of two.}
\GSfield{Fate}{These are established symplectic preliminaries in \citet[geometric preliminaries]{kostant1970quantization}, published in 1970; the seminar is expository here.}
\GSfield{Present status}{As of November 2026, the distinction between local closedness and global exactness remains essential. The obstruction to a symplectic field being Hamiltonian is its class in \(H^1_{\mathrm{dR}}(M)\).}
\GSfield{Assessment}{literal statement: correct with missing hypotheses; publication match: partial; verification: read the result; development: annotated; confidence: high}
\end{GSitem}
\begin{GSitem}{1971-229.2}{E}{Poisson bracket and the normalization problem}
\GSfield{Kind; source}{construction; \GSpage{P000361}}
\GSfield{The notes}{The notebook sets \(dF=2\alpha_{\xi_F}\), defines \(\{F,G\}=\xi_FG\), and also identifies this with \(2\omega(\xi_F,\xi_G)\). It explicitly allows a wedge convention with \((dp\wedge dq)(\partial_p,\partial_q)=1/2\).}
\GSfield{Statement}{Fix \(\iota_{X_f}\omega=-df\) and \(\{f,g\}=X_fg\). Then \(\{f,g\}=\omega(X_f,X_g)\), the bracket is a Lie bracket and a derivation in each argument, and \([X_f,X_g]=X_{\{f,g\}}\). These identities must all use the same normalization.}
\GSfield{Reconstruction}{The displayed convention is an editorial normalization, not a claim about Kirillov’s exact intended sign.}
\GSfield{Fate}{The Poisson Lie algebra is part of the prequantization framework already developed in \citet{kostant1970quantization}. The discrepancy is internal to the recorded formulas, not a later mathematical obstruction.}
\GSfield{Present status}{As of November 2026, the construction is standard; the notebook’s simultaneous sign and normalization choices cannot be adopted literally without alteration.}
\GSfield{Assessment}{literal statement: false as written; publication match: partial; verification: read the result; development: annotated; confidence: high}
\end{GSitem}
\begin{GSitem}{1971-229.3}{E}{Integral symplectic forms and prequantum line bundles}
\GSfield{Kind; source}{theorem; \GSpage{P000361-P000362}}
\GSfield{The notes}{Local primitives \(\alpha_j\), differences \(df_{jk}\), integral triple overlaps, and transition functions \(e^{2\pi i f_{jk}}\) construct a line bundle.}
\GSfield{Statement}{A closed real two-form \(\omega\) is the curvature class of a Hermitian line bundle with unitary connection of curvature \(-2\pi i\omega\) if and only if \([\omega]\) lies in the image of \(H^2(M;\mathbb Z)\to H^2_{\mathrm{dR}}(M;\mathbb R)\). Local potentials and a good cover give the Čech construction.}
\GSfield{Fate}{This is the prequantization integrality criterion of \citet{kostant1970quantization}, published before the talk. The notes’ line-bundle construction is an exact conceptual match.}
\GSfield{Present status}{As of November 2026, the theorem is established. Torsion in integral cohomology is invisible to the de Rham class and contributes to the choice of prequantization.}
\GSfield{Assessment}{literal statement: correct with missing hypotheses; publication match: exact; verification: read the result; development: annotated; confidence: high}
\end{GSitem}
\begin{GSitem}{1971-229.4}{E}{Ambiguity of prequantization and flat twists}
\GSfield{Kind; source}{claim; \GSpage{P000362}}
\GSfield{The notes}{The ambiguity in the line-bundle construction is described using \(H^1(X,\mathbb C^*)\).}
\GSfield{Statement}{Isomorphism classes of line bundles with connection of a fixed curvature, when nonempty, form a torsor for flat line bundles with connection. In the unitary setting this group is \(H^1(M;\mathbb R/\mathbb Z)\), equivalently homomorphisms \(\pi_1(M)\to U(1)\); for complex flat connections use the locally constant sheaf \(\mathbb C^*\).}
\GSfield{Reconstruction}{The constant-sheaf interpretation and the distinction between unitary and general complex flat bundles are supplied explicitly.}
\GSfield{Fate}{The ambiguity is part of the classification in \citet{kostant1970quantization}. The sheaf and connection data are necessary to interpret the abbreviated notation.}
\GSfield{Present status}{As of November 2026, one must not confuse this torsor with the classification of bare smooth complex line bundles by \(H^2(M;\mathbb Z)\).}
\GSfield{Assessment}{literal statement: correct with missing hypotheses; publication match: partial; verification: read the result; development: annotated; confidence: high}
\end{GSitem}
\begin{GSitem}{1971-229.5}{E}{The prequantum representation by first-order operators}
\GSfield{Kind; source}{construction; \GSpage{P000362-P000363}}
\GSfield{The notes}{Compatible scalar first-order operators are expressed through a connection and multiplication. The notes give \(D_\varphi=\nabla_{\xi_\varphi}+\pi i\varphi\) and a Poisson commutator formula.}
\GSfield{Statement}{With \(\iota_{X_f}\omega=-df\) and \(F_\nabla=-2\pi i\omega\), the operators \(D_f=\nabla_{X_f}+2\pi i f\) satisfy \([D_f,D_g]=D_{\{f,g\}}\). They are first-order operators on smooth sections of the prequantum line bundle; suitable domain and completeness assumptions are needed to exponentiate to group representations.}
\GSfield{Reconstruction}{The coefficients in the modern formula follow the normalization fixed in item 1971-229.2. No literal equality with the notebook’s \(\pi i\) convention is asserted.}
\GSfield{Fate}{\citet{kostant1970quantization} develops this representation before the seminar. The notebook explicitly cites the 1970 lecture volume; it is not an unpublished new quantization theorem.}
\GSfield{Present status}{As of November 2026, the Lie-algebra representation is standard, while its analytic integration and the choice of a smaller quantum space are separate questions.}
\GSfield{Assessment}{literal statement: correct with missing hypotheses; publication match: partial; verification: read the result; development: annotated; confidence: high}
\end{GSitem}
\begin{GSitem}{1971-229.6}{E}{Real, complex, and mixed polarizations}
\GSfield{Kind; source}{definition; \GSpage{P000363}}
\GSfield{The notes}{A polarization is an involutive complex Lagrangian distribution \(F\). The cases \(F=\bar F\) and \(F\cap\bar F=0\) are called real and Kähler; mixed polarizations require a constant-rank condition. A sphere is said not to admit a global polarization.}
\GSfield{Statement}{A regular polarization is a complex rank-\(n\) involutive Lagrangian subbundle of \(T_\mathbb C M\). A Kähler polarization additionally has the appropriate definiteness of \(i\omega(v,\bar v)\). On \(S^2\), no nonsingular real line-field polarization exists, but the standard complex structure gives a Kähler polarization.}
\GSfield{Reconstruction}{The restriction of the sphere claim to the real case is inferred from its local context and is flagged rather than silently supplied.}
\GSfield{Fate}{Polarizations and their use in quantization are developed in \citet{kostant1970quantization}. The sphere remark is correct only for nonsingular real polarizations, not for all complex polarizations.}
\GSfield{Present status}{As of November 2026, positivity, regularity, and the possible singularities of a real polarization must be specified. Transversality to the conjugate distribution alone gives a compatible complex structure, not positivity.}
\GSfield{Assessment}{literal statement: correct with missing hypotheses; publication match: partial; verification: read the result; development: annotated; confidence: high}
\end{GSitem}
\begin{GSitem}{1971-229.7}{E}{Polarized sections and the Poisson normalizer}
\GSfield{Kind; source}{theorem; \GSpage{P000364}}
\GSfield{The notes}{The quantum space is written \(\Gamma(L,F)=\{s:\nabla_Xs=0\text{ for }X\in F\}\). Leafwise constant functions form a locally maximal Poisson-commutative algebra \(A(F)\); its normalizer acts by first-order operators.}
\GSfield{Statement}{For a regular real polarization, locally in a foliated Darboux chart, \(A(F)\) is maximal Poisson-commutative. If \(X_f\) preserves \(F\), the prequantum operator \(D_f\) preserves polarized sections. This preservation condition is locally equivalent to \(\{f,A(F)\}\subset A(F)\). Global maximality requires enough global first integrals and is not automatic.}
\GSfield{Fate}{This is the standard polarization-preserving observable construction of \citet{kostant1970quantization}. The notebook itself qualifies maximal commutativity by \emph{locally}.}
\GSfield{Present status}{As of November 2026, the geometric statement is established. Global polarized sections may be too small or vanish, and half-form and Hilbert-space choices are additional structures, not supplied by this page.}
\GSfield{Assessment}{literal statement: correct with missing hypotheses; publication match: partial; verification: read the result; development: annotated; confidence: high}
\end{GSitem}
\begin{GSitem}{1971-229.8}{E}{Vertical polarization on a cotangent bundle}
\GSfield{Kind; source}{construction; \GSpage{P000364}}
\GSfield{The notes}{For the distribution \(dq=0\), functions of \(q\) act by multiplication and functions linear in \(p\) give first-order operators.}
\GSfield{Statement}{On \(T^*Q\) with its vertical real polarization, local polarized sections depend only on \(q\). The functions preserving this polarization are locally \(f(q,p)=a^j(q)p_j+b(q)\); their prequantum operators descend to first-order differential operators on \(Q\), with constants and signs fixed by the chosen connection.}
\GSfield{Fate}{This Schrödinger-model example is part of the earlier construction in \citet{kostant1970quantization}. The affine, rather than strictly homogeneous linear, dependence on \(p\) includes multiplication operators.}
\GSfield{Present status}{As of November 2026, the local assertion is standard; a global measure or half-density convention determines the additional divergence term in a symmetric operator realization.}
\GSfield{Assessment}{literal statement: correct with missing hypotheses; publication match: partial; verification: read the result; development: annotated; confidence: high}
\end{GSitem}

\GStalk{230}{}{1971–72}{P000365}{Academic-year heading}
This is the catalogue’s academic-year heading, retained as a record in its own right. The supplied range P000358--P000366 overlaps substantive records; no independent mathematical assertion is attached to the heading. P000365 is an ordering anchor before the 13 September talks, not a claim that its scan contains a separate handwritten year heading.


\GStalk{231}{1971-09-13}{I. M. Gelfand}{P000365-P000366}{Springer resolution, flags, and Hall polynomials}
Gelfand describes Springer’s resolution of the unipotent variety and the geometry of invariant complete flags. The finite-field flag counts connect Hall polynomials with the even cohomology of flag varieties; a centralizer formula bounds the dimensions of Springer fibres. The literal claim about arbitrary reductive groups over every field needs geometric and characteristic qualifications.

\begin{GSitem}{1971-231.1}{E}{Regular unipotents and the singular locus in type A}
\GSfield{Kind; source}{theorem; \GSpage{P000365}}
\GSfield{The notes}{For \(SL_n(\mathbb C)\), one Jordan block gives a nonsingular point of the variety \(R\) of unipotent elements; several blocks give singular points. The \(SL_2\) example is a quadratic cone with vertex at the identity.}
\GSfield{Statement}{For \(G=SL_n(\mathbb C)\), the unipotent variety is isomorphic by \(u\mapsto u-I\) to the nilpotent cone in \(\mathfrak{sl}_n\). Its smooth locus is the regular nilpotent orbit, equivalently the matrices with a single Jordan block. In \(SL_2\), writing \(a=1+x,d=1-x\) gives \(x^2+bc=0\).}
\GSfield{Fate}{The regularity framework precedes the seminar: \citet[Theorems 1.1--1.5]{steinberg1965regular}; the notebook applies it in type A. The complex quadratic cone is a complex surface, not literally a real two-dimensional cone.}
\GSfield{Present status}{As of November 2026, this is a standard established description. The restriction to \(\mathbb C\) avoids inseparability and bad-characteristic issues hidden by the notebook’s broader comment.}
\GSfield{Assessment}{literal statement: correct with missing hypotheses; publication match: partial; verification: read the result; development: annotated; confidence: high}
\end{GSitem}
\begin{GSitem}{1971-231.2}{E}{Springer’s resolution by invariant complete flags}
\GSfield{Kind; source}{construction; \GSpage{P000365}}
\GSfield{The notes}{Pairs \((p,g)\), where \(p\in B\backslash G\) and \(pgp^{-1}\in U\), define \(\widetilde R\to R\); the source is smooth and the generic fibre is one point. The margin calls \(B\backslash G\) projective space.}
\GSfield{Statement}{For \(G=SL_n(\mathbb C)\), let \(\mathcal B\) be the complete flag variety and \(\widetilde R=\{(u,F_\bullet):uF_i\subset F_i\}\). The projection to the unipotent variety is proper and birational, and \(\widetilde R\simeq G\times^B U\) is smooth. A regular unipotent has a unique invariant complete flag. The base is a projective variety, generally not a projective space.}
\GSfield{Fate}{This is the resolution in \citet{springer1969unipotent}, an account from the 1968 Bombay colloquium published in 1969. The original chapter was not obtained; its identity is corroborated by later primary papers, while the construction here is independently checked as an associated bundle.}
\GSfield{Present status}{As of November 2026, the type-A construction is established. An assertion over every field should be phrased as a geometric morphism after base change, not as a resolution by sets of rational points.}
\GSfield{Assessment}{literal statement: correct with missing hypotheses; publication match: exact; verification: indirect; development: annotated; confidence: high}
\end{GSitem}
\begin{GSitem}{1971-231.3}{E}{Invariant flag counts and Hall polynomials}
\GSfield{Kind; source}{theorem; \GSpage{P000365-P000366}}
\GSfield{The notes}{The number of flags over \(\mathbb F_q\) is polynomial in \(q\); for a unipotent \(g\), the fixed-flag count depends only on its Jordan partition. Green’s 1955 character paper is recommended.}
\GSfield{Statement}{The number of complete flags in \(\mathbb F_q^n\) is \(\prod_{j=1}^n(1+q+\cdots+q^{j-1})\). For a nilpotent operator of fixed Jordan type \(\lambda\), its invariant complete flags are counted by a polynomial in \(q\). It can be computed by successive Hall-number counts of submodules of the corresponding \(\mathbb F_q[t]\)-module; the total count need not be identified with one unspecified Hall polynomial.}
\GSfield{Fate}{The Hall-polynomial and finite-general-linear-group connection was already available in \citet{green1955characters}. A later geometric verification through affine decompositions is read in \citet[Introduction and \S1]{shimomura1980fixed}. Green’s full text was not successfully opened in this run.}
\GSfield{Present status}{As of November 2026, the counting assertion is established. The field-independent Jordan type and the choice of complete rather than partial flags must be fixed.}
\GSfield{Assessment}{literal statement: correct with missing hypotheses; publication match: partial; verification: read the result; development: annotated; confidence: high}
\end{GSitem}
\begin{GSitem}{1971-231.4}{E}{Flag counts and the even Poincaré polynomial}
\GSfield{Kind; source}{claim; \GSpage{P000366}}
\GSfield{The notes}{The notes compare the finite-field polynomial with the Poincaré polynomial after replacing \(q\) by \(t^2\).}
\GSfield{Statement}{The complex full flag variety has an affine Bruhat decomposition with cells \(\mathbb C^{\ell(w)}\), \(w\in S_n\). Thus its cohomology is concentrated in even degrees and \(P(t)=\sum_wt^{2\ell(w)}\), while its split finite-field model has \(\sum_wq^{\ell(w)}\) points. A corresponding comparison for invariant-flag varieties requires a compatible affine paving, not an arbitrary variety whose point count happens to be polynomial.}
\GSfield{Fate}{The full-flag comparison is classical. The later explicit invariant-flag decomposition of \citet[\S1]{shimomura1980fixed} substantiates the stronger type-A comparison; the paper credits Spaltenstein’s 1976 complete-flag result.}
\GSfield{Present status}{As of November 2026, the full-flag statement and the type-A Springer-fibre analogue are established. The notebook does not record a general theorem equating arbitrary point counts and Betti numbers.}
\GSfield{Assessment}{literal statement: correct with missing hypotheses; publication match: partial; verification: read the result; development: annotated; confidence: high}
\end{GSitem}
\begin{GSitem}{1971-231.5}{E}{Dimension of a Springer fibre and the centralizer}
\GSfield{Kind; source}{claim; \GSpage{P000366}}
\GSfield{The notes}{The fibre dimension is bounded above by \(\tfrac12(\dim Z_g-(n-1))\).}
\GSfield{Statement}{For unipotent \(u\in SL_n(\mathbb C)\), \(\dim\mathcal B_u=\tfrac12(\dim Z_{SL_n}(u)-(n-1))\). If the Jordan partition is \(\lambda\), this is \(\sum_i(i-1)\lambda_i\). In \(GL_n\), replace \(n-1\) by \(n\), since the centralizer dimension also increases by one.}
\GSfield{Fate}{The notebook’s upper bound belongs to Springer’s earlier geometry \citep{springer1969unipotent}; later fixed-flag work supplies sharper structural descriptions \citep[Introduction]{shimomura1980fixed}. The exact original locator for the upper bound was not retrieved, so the historical match is indirect.}
\GSfield{Present status}{As of November 2026, the equality in type A is established. The rank term is group-dependent; mixing a \(GL_n\) centralizer with the \(SL_n\) rank gives an incorrect half-integer formula.}
\GSfield{Assessment}{literal statement: correct with missing hypotheses; publication match: partial; verification: indirect; development: annotated; confidence: high}
\end{GSitem}

\GStalk{232}{1971-09-13}{B. I. Rosenfeld}{P000366}{Derivations of differential forms and an unidentified note by Rego}
The lower part of P000366 reports a note attributed to Rego in volume 272 of the Comptes rendus and recalls the Frölicher--Nijenhuis description of graded derivations. The named note and its proposed cohomology theory were not identified securely. The displayed calculus can nevertheless be reconstructed without inventing the missing result.

\begin{GSitem}{1971-232.1}{F}{The unlocated Rego note on cohomology and integrable distributions}
\GSfield{Kind; source}{remark; \GSpage{P000366}}
\GSfield{The notes}{The heading says Rosenfeld reports on Rego, \emph{C. R.}, volume 272 (1971), and associates cohomology with an integrable distribution. The author name is legible on the scan.}
\GSfield{Statement}{The surviving text does not specify the cochain groups, differential, equivalence relation on distributions, or an exact theorem. It therefore does not determine a unique cohomology theory or a definite open problem.}
\GSfield{Fate}{Searches combining Rego, volume 272, 1971, foliation/cohomology, and Frölicher--Nijenhuis terminology did not identify the cited note securely. The surrounding derivation calculus can be identified independently, but it does not justify assigning the missing theorem to a guessed author.}
\GSfield{Present status}{As of November 2026, the bibliographic identification remains unresolved in this bounded search. No assertion that the idea was undeveloped or remains open is warranted.}
\GSfield{Assessment}{literal statement: not determinable; publication match: none found; verification: unresolved lead; development: none; confidence: low}
\end{GSitem}
\begin{GSitem}{1971-232.2}{E}{Decomposition of graded derivations and the Frölicher–Nijenhuis bracket}
\GSfield{Kind; source}{theorem; \GSpage{P000366}}
\GSfield{The notes}{The notes describe derivations of the graded algebra of forms by vector-valued forms and insertion. For a vector-valued \(p\)-form \(L\), they write \(d_L=d\,i_L+(-1)^p i_Ld\).}
\GSfield{Statement}{Every degree-\(k\) graded derivation of \(\Omega^*(M)\) has a unique decomposition \(\mathcal L_K+i_L\), with \(K\in\Omega^k(M;TM)\), \(L\in\Omega^{k+1}(M;TM)\), and \(\mathcal L_K=[i_K,d]\). The commutator of Lie derivations defines \([K,H]_{FN}\) by \([\mathcal L_K,\mathcal L_H]=\mathcal L_{[K,H]_{FN}}\).}
\GSfield{Reconstruction}{The form degrees and graded commutator are supplied. For \(L=\operatorname{Id}_{TM}\), the notebook convention gives \(-d\), consistently with this sign conversion.}
\GSfield{Fate}{This calculus is due to Frölicher and Nijenhuis in 1956. It is explicitly attributed and reviewed in the author’s account \citet[Introduction]{michor1989derivations}; the original 1956 publisher text was not accessible here.}
\GSfield{Present status}{As of November 2026, the decomposition is established. The notebook’s \(d_L\) equals \((-1)^p\mathcal L_L\) in this convention; that is a convention difference, not by itself a mathematical error.}
\GSfield{Assessment}{literal statement: correct with missing hypotheses; publication match: partial; verification: read the result; development: annotated; confidence: high}
\end{GSitem}
\begin{GSitem}{1971-232.3}{E}{When an induced Lie derivation is a differential}
\GSfield{Kind; source}{claim; \GSpage{P000366}}
\GSfield{The notes}{The short entry links the construction with integrability and a cohomology, without specifying a degree or a torsion identity.}
\GSfield{Statement}{For a vector-valued one-form \(N\), \(\mathcal L_N^2=\tfrac12\mathcal L_{[N,N]_{FN}}\). Hence vanishing Nijenhuis torsion gives a degree-one differential. An arbitrary graded Lie derivation does not automatically square to zero; a foliation gives a leafwise de Rham differential by a different, explicitly specified construction.}
\GSfield{Reconstruction}{The vector-valued one-form formulation is a mathematically valid interpretation, not a recovered missing theorem.}
\GSfield{Fate}{The precise graded identity belongs to the Frölicher--Nijenhuis calculus reviewed in \citet[\S\S1--2]{michor1989derivations}; \citet{michor1987remarks} gives a broader account. The relation of this identity to the unidentified Rego announcement is only a reconstruction, not an exact publication match.}
\GSfield{Present status}{As of November 2026, the torsion-zero implication is standard. Which of these complexes Rosenfeld intended cannot be decided from this page alone.}
\GSfield{Assessment}{literal statement: correct with missing hypotheses; publication match: partial; verification: read the result; development: annotated; confidence: medium}
\end{GSitem}

\GStalk{233}{1971-09-13}{Vlastimil Pták}{P000367}{Critical exponents of finite-dimensional normed spaces}
The talk asks how many consecutive norm-preserving powers of a contraction force spectral radius one. It records exact Hilbert and polyhedral examples, a two-dimensional counterexample, questions about \(\ell_p^n\), the minimum over all norms, and a constrained spectral-radius extremum. A 1971 paper already answers the rational-\(p\) finiteness question more generally than the notebook acknowledges; the exact later fate of the other stated questions was not established.

\begin{GSitem}{1971-233.1}{E}{The critical exponent of a finite-dimensional norm}
\GSfield{Kind; source}{definition; \GSpage{P000367}}
\GSfield{The notes}{The notes define an exponent \(q\) by \(\|A\|=\|A^q\|=1\Rightarrow r(A)=1\), together with a counterexample for \(q-1\). The dimension sign is visually ambiguous between finite dimension and \(\le\infty\).}
\GSfield{Statement}{For a finite-dimensional real or complex normed space \(E\), let \(q(E)\) be the least positive integer \(q\), if it exists, such that every contraction \(A\) with \(\|A^q\|=1\) has spectral radius one; otherwise set \(q(E)=\infty\). For \(q>1\), minimality gives a contraction \(B\) with \(\|B^{q-1}\|=1\) and \(r(B)<1\). In dimension one, \(q(E)=1\), and the notebook’s additional requirement \(\|B\|=1\) cannot be imposed on a predecessor witness.}
\GSfield{Reconstruction}{Finite dimension and the dimension-one exception are explicit editorial qualifications.}
\GSfield{Fate}{The definition and equivalent eventual-norm-drop formulation are given in \citet[p.~672]{kirzhnertabachnikov1971critical}, referring to earlier work of Pták and Mařík. It predates this September talk.}
\GSfield{Present status}{As of November 2026, the finite-dimensional invariant is well defined with the stated convention. An infinite-dimensional extension requires care: spectral radius, spectral values, and norm-attainment arguments cannot simply be copied.}
\GSfield{Assessment}{literal statement: correct with missing hypotheses; publication match: partial; verification: read the result; development: annotated; confidence: high}
\end{GSitem}
\begin{GSitem}{1971-233.2}{E}{Critical exponent of a Hilbert space}
\GSfield{Kind; source}{theorem; \GSpage{P000367}}
\GSfield{The notes}{The table gives \(q(\ell_2^n)=n\).}
\GSfield{Statement}{For a real or complex Hilbert space of finite dimension \(n\), every contraction \(A\) with \(\|A^n\|=1\) has spectral radius one, and the exponent is sharp. For \(n\ge2\), a nilpotent shift has \(\|A^{n-1}\|=1\) and \(A^n=0\); dimension one is immediate.}
\GSfield{Fate}{This was Pták’s 1962 result, cited on p.~672 of \citet{kirzhnertabachnikov1971critical}. \citet[\S\S1--2]{ptak1976isometric} gives a later proof by isometric parts. The 1976 date is not the date the theorem first became known.}
\GSfield{Present status}{As of November 2026, the exact exponent is established. The Hilbert geometry is essential; an arbitrary norm on the same vector space may have a different exponent or no finite one.}
\GSfield{Assessment}{literal statement: correct with missing hypotheses; publication match: exact; verification: read the result; development: annotated; confidence: high}
\end{GSitem}
\begin{GSitem}{1971-233.3}{E}{Critical exponents of the one and infinity norms}
\GSfield{Kind; source}{theorem; \GSpage{P000367}}
\GSfield{The notes}{The table gives \(q(\ell_1^n)=q(\ell_\infty^n)=n^2-n+1\).}
\GSfield{Statement}{For the usual \(\ell_1\) and \(\ell_\infty\) norms in dimension \(n\), the critical exponent is \(n^2-n+1\), with the predecessor-witness convention separately understood for \(n=1\). Duality identifies the two exponents because \(\|(A^*)^j\|=\|A^j\|\) and \(r(A^*)=r(A)\).}
\GSfield{Fate}{\citet[p.~672]{kirzhnertabachnikov1971critical} explicitly records these established values and their earlier references. Thus they were not new in September 1971; the historical source for the exact maximum-norm result is earlier work of Mařík and Pták.}
\GSfield{Present status}{As of November 2026, the exact values are established. They demonstrate why dimension alone does not determine the exponent.}
\GSfield{Assessment}{literal statement: correct with missing hypotheses; publication match: exact; verification: read the result; development: annotated; confidence: high}
\end{GSitem}
\begin{GSitem}{1971-233.4}{E}{A two-dimensional norm with no finite critical exponent}
\GSfield{Kind; source}{claim; \GSpage{P000367}}
\GSfield{The notes}{The notebook attributes an example to Danzer and sketches a unit ball with rotations by small angles; it asserts \(q(E)=\infty\) in dimension two.}
\GSfield{Statement}{There exists a real two-dimensional normed space for which, for every integer \(m\), a contraction \(A_m\) has \(\|A_m^m\|=1\) but \(r(A_m)<1\). A bare rotation has spectral radius one and is not by itself such a witness; the norm geometry and the contracting linear map must be specified.}
\GSfield{Fate}{\citet[example, pp.~674--675]{kirzhnertabachnikov1971critical} gives an explicit two-dimensional example using a piecewise algebraic unit sphere and diagonal contracting maps. This verifies the existence claim before the talk, but it does not identify the notebook’s Danzer diagram or its attribution.}
\GSfield{Present status}{As of November 2026, the existence of infinite critical exponent in dimension two is established. The exact Danzer construction remains an unresolved bibliographic detail, not an open existence problem.}
\GSfield{Assessment}{literal statement: correct with missing hypotheses; publication match: partial; verification: read the result; development: annotated; confidence: medium}
\end{GSitem}
\begin{GSitem}{1971-233.5}{C}{Finiteness for finite-dimensional lp spaces}
\GSfield{Kind; source}{question; \GSpage{P000367}}
\GSfield{The notes}{The notes say existence of a critical exponent for \(\ell_p^n\) is unknown except for \(p=2^m\), and call it useful for numerical analysis.}
\GSfield{Statement}{For each fixed dimension \(n\) and \(1\le p<\infty\), ask whether \(q(\ell_p^n)<\infty\), specifying real or complex scalars. A verified partial answer is affirmative for every rational \(p\): if the unit sphere lies in a real algebraic zero set not containing the origin, a finite exponent exists, and rational-\(p\) spheres satisfy that criterion.}
\GSfield{Fate}{\citet[Theorem 1 and Corollaries 2--3]{kirzhnertabachnikov1971critical}, received 14 January 1970 and printed May--June 1971, proves rational-\(p\) finiteness over both scalar fields. The notebook’s state-of-knowledge claim is already obsolete on its catalogue date. Searches through November 2026 did not verify an exact all-real-\(p\) successor or optimal general exponents.}
\GSfield{Present status}{As of November 2026, rational-\(p\) finiteness is established. The status of the remaining irrational parameters and sharp quantitative values was not established by this bounded search; this is not a certification that they are open today.}
\GSfield{Assessment}{literal statement: false as written; publication match: exact; verification: read the result; development: annotated; confidence: high; problem \hyperref[prob:GP-1971-01]{GP-1971-01}}
\GSfield{Worth developing}{A focused follow-up is worthwhile to identify the exact irrational-\(p\) literature and sharpen an explicitly chosen finite-dimensional case; the algebraic proof itself is already published.}
\end{GSitem}
\begin{GSitem}{1971-233.6}{C}{Does the minimum critical exponent tend to infinity with dimension?}
\GSfield{Kind; source}{question; \GSpage{P000367}}
\GSfield{The notes}{The notes define \(M(n)\) as the minimum over norms, report \(M(n)\le n\), examples below \(n\), and a subsequence on the scale \(\sqrt n\), and ask whether \(\liminf M(n)=\infty\).}
\GSfield{Statement}{Fix the scalar field and define \(M(n)=\min\{q(E):\dim E=n\}\), using \(q(E)=\infty\) when necessary. The Hilbert norm gives \(M(n)\le n\). The definite question is whether, for every integer \(K\), only finitely many dimensions admit a norm with \(q(E)\le K\). The notebook’s subsequence claim is a separate reported bound whose precise construction was not located.}
\GSfield{Fate}{Pták’s Hilbert result verifies the basic upper bound \citep{ptak1976isometric}. Bilingual searches of critical exponents, minimum over norms, and \(M(n)\) did not establish the exact later fate. The 1995 literature on \emph{exponential numbers} uses a different orbit-growth invariant and cannot be substituted for this contraction exponent.}
\GSfield{Present status}{As of November 2026, the historical question is precise, and some special-space bounds are established, but its exact current resolution remains unverified here. Class C records that partial development and unresolved identification, not a proof of present-day openness.}
\GSfield{Assessment}{literal statement: not determinable; publication match: partial; verification: indirect; development: annotated; confidence: medium; problem \hyperref[prob:GP-1971-02]{GP-1971-02}}
\GSfield{Worth developing}{It is worth first recovering the precise norm construction behind the reported \(\sqrt n\) subsequence; without that bibliographic step a new asymptotic claim would risk duplicating or misidentifying existing work.}
\end{GSitem}
\begin{GSitem}{1971-233.7}{F}{The twelve-dimensional direct-sum example is not securely reconstructible}
\GSfield{Kind; source}{claim; \GSpage{P000367}}
\GSfield{The notes}{The scan writes a sum rendered \(\ell_4^2\oplus_{\ell_1}\ell_8^2\), followed by dimension 12 and critical exponent 11.}
\GSfield{Statement}{With the conventional notation \(\ell_p^n\), the displayed sum has dimension four, not twelve. Swapping the roles of superscript and subscript would give Hilbert spaces of dimensions four and eight, but this reading is not established by the scan and its exponent must not be inferred. A unique normed space and claimed exact exponent cannot presently be assigned.}
\GSfield{Fate}{The direct-sum construction was not located securely in the searched Pták literature. The preceding discussion of \(M(n)<n\) makes it plausible that a twelve-dimensional example was intended, but plausibility does not resolve the notation.}
\GSfield{Present status}{As of November 2026, the example remains unidentified. Neither the exponent 11 for the literal four-dimensional space nor a proposed repaired twelve-dimensional space is claimed verified.}
\GSfield{Assessment}{literal statement: not determinable; publication match: none found; verification: unresolved lead; development: annotated; confidence: low}
\end{GSitem}
\begin{GSitem}{1971-233.8}{C}{The constrained power-norm extremum}
\GSfield{Kind; source}{problem; \GSpage{P000367}}
\GSfield{The notes}{The page defines \(C(E,r,\rho)=\max\{\|T^r\|:\|T\|\le1,\ r(T)\le\rho\}\), with \(0<\rho<1\).}
\GSfield{Statement}{\ours\ For finite-dimensional \(E\), integer \(r\ge1\), and \(0\le\rho<1\), the admissible set is compact, so this maximum exists. If \(q(E)<\infty\) and \(r\ge q(E)\), then \(C(E,r,\rho)<1\) for every fixed \(\rho<1\). Conversely, if \(r<q(E)\), some \(\rho<1\) gives \(C(E,r,\rho)=1\). Determining exact values and maximizers is a quantitative problem depending on the norm, dimension, \(r\), and \(\rho\).}
\GSfield{Reconstruction}{\ours\ Compactness, continuity of spectral radius, and attainment justify the maximum and strict inequality; these hypotheses and deductions are supplied explicitly.}
\GSfield{Fate}{The extremal viewpoint is part of Pták’s programme; the critical-exponent results in \citet{ptak1976isometric} concern its qualitative threshold. The compactness implication stated here is an elementary editorial deduction, not a claim that the full optimization has been solved. Searches did not identify a uniform exact solution for arbitrary norms.}
\GSfield{Present status}{As of November 2026, the qualitative finite-dimensional relation to \(q(E)\) is clear. Exact general values and the complete later classification of extremizers were not verified in this run; no blanket open-problem claim is made.}
\GSfield{Assessment}{literal statement: correct with missing hypotheses; publication match: partial; verification: read the result; development: annotated; confidence: medium; problem \hyperref[prob:GP-1971-03]{GP-1971-03}}
\GSfield{Worth developing}{For a specified low-dimensional norm and fixed spectral-radius bound, a certified optimization or sharp estimate could be a bounded contribution; the unrestricted problem is too broad for a defensible new result here.}
\end{GSitem}

\GStalk{235}{1971-09-13}{I. M. Gelfand}{P000368}{The affine Radon transform over a finite field}
Gelfand derives inversion by counting affine hyperplanes through one or two points and identifies a parallel-class range condition. A final remark connects finite-field subspace counts to Gaussian binomial coefficients and their \(q\to1\) limit. This heading precedes Kirillov’s P000369 heading despite the opposite order of catalogue record numbers.

\begin{GSitem}{1971-235.1}{B}{Affine finite-field Radon inversion by incidence counting}
\GSfield{Kind; source}{theorem; \GSpage{P000368}}
\GSfield{The notes}{For \(f:\mathbb F_q^n\to\mathbb C\), \(\varphi(p)\) sums \(f\) over a \emph{plane}. Counting planes through one or two points yields an inversion identity.}
\GSfield{Statement}{Let \(n\ge1\), let \(\mathcal H\) be all affine hyperplanes of \(V=\mathbb F_q^n\), and set \(Rf(H)=\sum_{x\in H}f(x)\). Put \(b=(q^{n-1}-1)/(q-1)\) and \(C=\sum_Vf\). Then \(\sum_{H\ni x}Rf(H)=q^{n-1}f(x)+bC\), so \(f(x)=q^{1-n}(\sum_{H\ni x}Rf(H)-bC)\). \ours\ The count uses \((q^n-1)/(q-1)\) hyperplanes through one point and \(b\) through two distinct points.}
\GSfield{Reconstruction}{\ours\ \emph{Plane} means hyperplane, as forced by the numerical coefficients; function values are taken in characteristic zero so division by \(q\) is valid.}
\GSfield{Fate}{The identity is verified directly from the incidence counts on the page. Related finite integral geometry was developed in \citet[Theorems 1--2]{graevpavlenko1972radon}, submitted February 1972, for projective spaces over finite commutative rings; that is a related generalization, not an exact publication of this affine formula.}
\GSfield{Present status}{As of November 2026, the inversion formula is established and elementary. The exact first publication of this notebook formulation was not identified. The fate classification records the later developed subject, without assigning priority.}
\GSfield{Assessment}{literal statement: correct with missing hypotheses; publication match: partial; verification: read the result; development: annotated; confidence: medium}
\end{GSitem}
\begin{GSitem}{1971-235.2}{B}{The range condition from parallel classes}
\GSfield{Kind; source}{claim; \GSpage{P000368}}
\GSfield{The notes}{The sum of \(\varphi(p)\) over each parallel class is the same constant \(C\), called a finite analogue of a Paley--Wiener condition.}
\GSfield{Statement}{Over \(\mathbb C\), the image of \(R\) consists exactly of functions \(\varphi\) on affine hyperplanes for which \(\sum_{H\parallel H_0}\varphi(H)\) is independent of the direction \(H_0\). \ours\ Necessity follows by partitioning \(V\). For sufficiency, there are \(a=(q^n-1)/(q-1)\) parallel classes, each of size \(q\); the \(a-1\) independent equal-sum constraints leave dimension \(aq-(a-1)=q^n\), equal to the domain dimension, and the preceding inversion proves injectivity.}
\GSfield{Reconstruction}{\ours\ The sufficiency proof and its dimension count are supplied here; the page explicitly records only the equal-sum condition.}
\GSfield{Fate}{This is an editorial completion of the notebook’s elementary range observation. The projective transform in \citet{graevpavlenko1972radon} supplies related later development but has different range geometry; no exact source for this affine characterization was verified.}
\GSfield{Present status}{As of November 2026, the range characterization is proved by the displayed finite-dimensional argument. It is not an unresolved Paley--Wiener problem.}
\GSfield{Assessment}{literal statement: correct with missing hypotheses; publication match: thematic only; verification: read the result; development: annotated; confidence: medium}
\end{GSitem}
\begin{GSitem}{1971-235.3}{E}{Gaussian binomial coefficients and the limit q equals one}
\GSfield{Kind; source}{remark; \GSpage{P000368}}
\GSfield{The notes}{A remark attributed to Rota interprets \(C_k^n(q)\) as the number of \(k\)-dimensional subspaces and states its limit is \(\binom nk\).}
\GSfield{Statement}{The number of \(k\)-subspaces of \(\mathbb F_q^n\) is \(\genfrac{[}{]}{0pt}{}{n}{k}_q=\prod_{j=0}^{k-1}(q^{n-j}-1)/(q^{k-j}-1)\), a polynomial in \(q\) with integer coefficients. Its polynomial value at \(q=1\), equivalently the limit of this rational expression, is \(\binom nk\). It is a polynomial specialization, not a finite field with one element.}
\GSfield{Fate}{This finite-vector-space combinatorics is already developed in \citet{goldmanrota1970finite}, published September 1970, whose bibliography also records their 1969 subspace-count paper. Publisher metadata and references were checked; the collected full text exceeded the retrieval limit.}
\GSfield{Present status}{As of November 2026, the identity is established. It explains the numerical coefficients in the preceding Radon argument.}
\GSfield{Assessment}{literal statement: correct with missing hypotheses; publication match: partial; verification: publisher metadata checked; development: annotated; confidence: high}
\end{GSitem}

\GStalk{234}{1971-09-13}{A. A. Kirillov}{P000369}{Wavefront sets and oscillatory integrals}
The one-page talk introduces the wavefront set, conormal Lagrangians, pseudodifferential kernels, and oscillatory representations. Hörmander’s 1971 paper provides a close contemporary source. The notebook omits localization in its wavefront criterion and overstates the converse from wavefront containment to membership in a Lagrangian distribution class.

\begin{GSitem}{1971-234.1}{E}{Microlocal regularity and the wavefront set}
\GSfield{Kind; source}{definition; \GSpage{P000369}}
\GSfield{The notes}{For a distribution \(A\), the notes test \(\langle A,\varphi e^{it\psi}\rangle=o(t^{-N})\) for every \(N\), with \(d\psi\) prescribed at one point and arbitrary compactly supported \(\varphi\). Smooth functions have empty wavefront.}
\GSfield{Statement}{For \(u\in\mathcal D^{\prime}(M)\), a nonzero covector \((x_0,\xi_0)\) is outside \(\operatorname{WF}(u)\) if, in a coordinate chart, there is a cutoff equal to one near \(x_0\) whose Fourier transform after multiplication by \(u\) decreases faster than every power in an open conic neighbourhood of \(\xi_0\). The wavefront is closed, conic, and projects to the singular support.}
\GSfield{Fate}{\citet[\S2.5]{hormander1971fourier} gives the contemporary definition and calculus. The article was printed on 1 June 1971 and developed material circulated in the 1969 Nordic summer-school lectures, before this September talk.}
\GSfield{Present status}{As of November 2026, the definition is standard. The notebook’s condition requires sufficiently small support and uniform control in neighbouring phase directions; requiring arbitrary support while prescribing \(d\psi\) only at one point is not a valid literal criterion.}
\GSfield{Assessment}{literal statement: false as written; publication match: exact; verification: read the result; development: annotated; confidence: high}
\end{GSitem}
\begin{GSitem}{1971-234.2}{E}{Conormal bundles as Lagrangians}
\GSfield{Kind; source}{construction; \GSpage{P000369}}
\GSfield{The notes}{A submanifold \(N\subset M\) gives the normal bundle of annihilating covectors, described as a Lagrangian of half the cotangent dimension.}
\GSfield{Statement}{For an embedded smooth submanifold \(N\), its conormal bundle \(N^*M=\{(x,\xi):x\in N,\xi|_{T_xN}=0\}\) is Lagrangian in \(T^*M\). The wavefront of a conormal distribution lies in \(N^*M\setminus0\); the Dirac density along a smooth submanifold has precisely this wavefront when its density is nonvanishing.}
\GSfield{Fate}{The conormal model is a basic case of the oscillatory constructions in \citet[\S1 and \S3]{hormander1971fourier}. It was established before the seminar.}
\GSfield{Present status}{As of November 2026, the assertion is standard; \emph{normal} here means conormal covectors, without choosing a Riemannian metric.}
\GSfield{Assessment}{literal statement: correct with missing hypotheses; publication match: partial; verification: read the result; development: annotated; confidence: high}
\end{GSitem}
\begin{GSitem}{1971-234.3}{E}{Pseudodifferential kernels are conormal to the diagonal}
\GSfield{Kind; source}{claim; \GSpage{P000369}}
\GSfield{The notes}{A pseudodifferential operator is described as a distribution kernel on \(M\times M\) whose wavefront is constructed from the diagonal.}
\GSfield{Statement}{The Schwartz kernel of a properly supported pseudodifferential operator of finite order belongs locally to a conormal distribution class for the diagonal; in particular \(\operatorname{WF}(K)\subset N^*\Delta\setminus0\). Wavefront containment alone is not a characterization of the symbol class or of finite operator order.}
\GSfield{Fate}{The symbol hypotheses and oscillatory kernel definition are present in \citet[Definition 1.1.1 and the pseudodifferential sections]{hormander1971fourier}. The notebook records the geometric support of singularities but omits their quantitative order.}
\GSfield{Present status}{As of November 2026, the inclusion is established. A converse needs symbol estimates or equivalent iterated regularity conditions.}
\GSfield{Assessment}{literal statement: correct with missing hypotheses; publication match: partial; verification: read the result; development: annotated; confidence: high}
\end{GSitem}
\begin{GSitem}{1971-234.4}{E}{Oscillatory parametrization of a Lagrangian distribution}
\GSfield{Kind; source}{claim; \GSpage{P000369}}
\GSfield{The notes}{The page states that every distribution with wavefront in a Lagrangian is given by an oscillatory integral, and writes a phase-gradient description of \(\Lambda\).}
\GSfield{Statement}{A nondegenerate phase \(\phi(x,\theta)\), homogeneous of degree one in \(\theta\), parametrizes \(\Lambda=\{(x,d_x\phi):d_\theta\phi=0\}\). Integrating \(e^{i\phi}a\) against a symbol \(a\) yields a Lagrangian distribution. Membership in \(I^m(M,\Lambda)\) entails amplitude estimates and is strictly stronger than the bare inclusion \(\operatorname{WF}(u)\subset\Lambda\). Locally, members of this class have such representations modulo smooth functions.}
\GSfield{Reconstruction}{The stationary equations \(d_\theta\phi=0\) and the symbol class are supplied; they are essential, not cosmetic.}
\GSfield{Fate}{\citet[\S3]{hormander1971fourier} establishes the phase and amplitude calculus. The notebook collapses the quantitative class condition into a wavefront condition, so its literal universal converse is not the theorem of that paper.}
\GSfield{Present status}{As of November 2026, the appropriately qualified theorem is established. No new classification problem is inferred from the missing hypotheses.}
\GSfield{Assessment}{literal statement: false as written; publication match: partial; verification: read the result; development: annotated; confidence: high}
\end{GSitem}

\GStalk{236}{1971-09-20}{I. N. Bernstein}{P000370}{The Godbillon–Vey class}
Bernstein explains the three-form attached to a cooriented codimension-one foliation, a homogeneous nonzero example, and the relation to formal vector-field cohomology. The claim of stability under deformation requires correction: arbitrary smooth families can change the Godbillon--Vey class. Atiyah’s talk starts later on the same page.

\begin{GSitem}{1971-236.1}{E}{The Godbillon–Vey three-class}
\GSfield{Kind; source}{construction; \GSpage{P000370}}
\GSfield{The notes}{For an integrable defining one-form \(\omega\), choose \(\omega_1,\omega_2\) with \(d\omega=\omega\wedge\omega_1\), \(d\omega_1=\omega\wedge\omega_2\); the closed form \(\Omega=\omega\wedge\omega_1\wedge\omega_2\) defines a choice-independent class.}
\GSfield{Statement}{For a smooth cooriented codimension-one foliation, take a nowhere-vanishing defining form \(\omega\) and a one-form \(\eta\) with \(d\omega=\omega\wedge\eta\). The form \(\eta\wedge d\eta\) is closed and its de Rham class is independent of the choices. With the notebook convention, \(\Omega=-\eta\wedge d\eta\); this overall sign must be fixed when comparing numerical invariants.}
\GSfield{Fate}{\citet[\S1]{godbillonvey1971invariant} introduced the class in the issue dated 12 July 1971. Bernstein’s 20 September report is therefore exposition of a published result, not a later 1972 discovery.}
\GSfield{Present status}{As of November 2026, the smooth construction is established. Compactness is needed to integrate over a fundamental cycle, not to define the cohomology class.}
\GSfield{Assessment}{literal statement: correct with missing hypotheses; publication match: exact; verification: read the result; development: annotated; confidence: high}
\end{GSitem}
\begin{GSitem}{1971-236.2}{E}{Vanishing for a foliation by fibres}
\GSfield{Kind; source}{claim; \GSpage{P000370}}
\GSfield{The notes}{The class is zero when the foliation is a fibration.}
\GSfield{Statement}{For a cooriented codimension-one foliation given by the fibres of a smooth submersion onto an oriented one-dimensional manifold, a nowhere-zero closed defining one-form can be chosen by pulling back a one-form on the base. Taking \(\eta=0\) makes the Godbillon--Vey form zero.}
\GSfield{Fate}{This is a special case of \citet[\S2(i)]{godbillonvey1971invariant}: a nonsingular closed Pfaffian defining form gives zero class. The result predates the talk.}
\GSfield{Present status}{As of November 2026, the fibre-foliation assertion is established with the indicated global defining-form hypothesis.}
\GSfield{Assessment}{literal statement: correct with missing hypotheses; publication match: partial; verification: read the result; development: annotated; confidence: high}
\end{GSitem}
\begin{GSitem}{1971-236.3}{B}{Deformation invariance versus foliated concordance}
\GSfield{Kind; source}{claim; \GSpage{P000370}}
\GSfield{The notes}{The notebook says \ru{устойчив при деформ.}, stable under deformations, with no qualification on the deformation.}
\GSfield{Statement}{The Godbillon--Vey class is natural and invariant under a foliated concordance on \(M\times[0,1]\) transverse to the boundary. It need not be constant along an arbitrary smooth parameter family of foliations on \(M\). Thurston’s polygon construction yields foliations of \(S^3\) whose integrated classes vary as \(4\pi\operatorname{Area}(P)\), up to orientation convention.}
\GSfield{Fate}{The original \citet[\S1(iii)]{godbillonvey1971invariant} uses differentiable homotopy of foliations, which must be interpreted in its foliation-theoretic sense. \citet[p.~513]{thurston1972noncobordant}, communicated 15 December 1971 and published in 1972, demonstrates variable values. This refutes an unrestricted reading of the notebook.}
\GSfield{Present status}{As of November 2026, unrestricted deformation invariance is false. A path of integrable plane fields is weaker than an integrable foliation on the product giving a concordance.}
\GSfield{Assessment}{literal statement: false as written; publication match: exact; verification: read the result; development: annotated; confidence: high}
\end{GSitem}
\begin{GSitem}{1971-236.4}{E}{The nonzero homogeneous example on a compact SL2 quotient}
\GSfield{Kind; source}{construction; \GSpage{P000370}}
\GSfield{The notes}{The page sketches a compact quotient of \(SL_2(\mathbb R)\) by a discrete subgroup and the foliation by triangular-subgroup orbits, asserting dense trajectories and a nonzero class.}
\GSfield{Statement}{For a cocompact lattice \(\Gamma\subset SL_2(\mathbb R)\), the foliation on \(\Gamma\backslash SL_2(\mathbb R)\) associated to a two-dimensional Borel subgroup has a Godbillon--Vey form which is a nonzero invariant volume form. Hence its integral, with an orientation, is nonzero. Left/right quotient and invariant-form conventions must be made compatible.}
\GSfield{Fate}{This is Roussarie’s example recorded explicitly in \citet[\S2(iii)]{godbillonvey1971invariant}. Thus the nonzero example was published before Bernstein’s report. The page’s separate density assertion is not required for computing the class.}
\GSfield{Present status}{As of November 2026, the homogeneous nonvanishing is established. The word \emph{always} about density is meaningful only with the compact quotient and the intended connected orbit foliation fixed; it is not a statement for every discrete subgroup.}
\GSfield{Assessment}{literal statement: correct with missing hypotheses; publication match: partial; verification: read the result; development: annotated; confidence: high}
\end{GSitem}
\begin{GSitem}{1971-236.5}{E}{Transverse pseudogroups and formal vector-field characteristic classes}
\GSfield{Kind; source}{method; \GSpage{P000370}}
\GSfield{The notes}{Local submersions \(\varphi_i\) give transition maps on intervals; because they form a pseudogroup rather than a single finite-dimensional group, the talk turns to \(W_1\), formal vector fields on the line.}
\GSfield{Statement}{A codimension-one foliation admits a Haefliger cocycle \(\varphi_j=\psi_{ji}\circ\varphi_i\) on sufficiently small overlaps, where the \(\psi_{ji}\) are local diffeomorphisms of intervals. Formal transverse jets yield a characteristic homomorphism from the appropriate continuous formal-vector-field cohomology to de Rham cohomology of the foliated manifold.}
\GSfield{Reconstruction}{The term Haefliger cocycle and the explicit factorization specify the modern meaning of the sketch.}
\GSfield{Fate}{The original \citet[\S3]{godbillonvey1971invariant} constructs a universal Lie coalgebra for this purpose. The broader formal-coordinate construction is developed in the contemporary \citet{gelfandkazhdan1971geometry} and is discussed in record 243.}
\GSfield{Present status}{As of November 2026, this is an established source of secondary foliation classes. The raw composition \(\varphi_j\varphi_i\) on the page is not defined; the overlap factorization is the correct relation.}
\GSfield{Assessment}{literal statement: correct with missing hypotheses; publication match: partial; verification: read the result; development: annotated; confidence: high}
\end{GSitem}
\begin{GSitem}{1971-236.6}{E}{The continuous cohomology of W1}
\GSfield{Kind; source}{theorem; \GSpage{P000370}}
\GSfield{The notes}{The page says the cohomology vanishes except for \(H^3(W_1;\mathbb R)=\mathbb R\), identifying that generator with the preceding class.}
\GSfield{Statement}{For the Lie algebra \(W_1\) of formal vector fields in one variable with its formal topology and trivial coefficients, continuous cohomology is one-dimensional in degrees 0 and 3 and zero in other degrees. The degree-three generator induces the Godbillon--Vey class under the characteristic map.}
\GSfield{Fate}{The formal-vector-field computation belongs to \citet{gelfandfuchs1970formal}; \citet[\S3]{godbillonvey1971invariant} identifies the foliation class. The missing degree-zero copy is an omitted standard qualification.}
\GSfield{Present status}{As of November 2026, the cohomology statement is established for continuous cochains. It does not assert the same result for arbitrary discontinuous algebraic cochains.}
\GSfield{Assessment}{literal statement: correct with missing hypotheses; publication match: partial; verification: read the result; development: annotated; confidence: high}
\end{GSitem}

\GStalk{237}{1971-09-20}{Michael Atiyah}{P000370-P000373}{Spectral flow and an index on a mapping torus}
Atiyah describes the integer obtained by following eigenvalues of a self-adjoint elliptic family and identifies it with the index of a suspended operator. A circle example, an eigenvalue-orbit argument, and remarks about diffeomorphisms and higher signatures follow. The main idea subsequently appeared in the third Atiyah--Patodi--Singer paper; the notebook’s derivative estimate and unitary conjugation sign need repair.

\begin{GSitem}{1971-237.1}{E}{Discrete spectrum of a self-adjoint elliptic operator}
\GSfield{Kind; source}{theorem; \GSpage{P000370}}
\GSfield{The notes}{On a compact manifold, a self-adjoint elliptic differential operator has a discrete real spectrum, with positive and negative eigenvalues.}
\GSfield{Statement}{On a closed smooth manifold, a formally self-adjoint elliptic differential operator of positive order on a Hermitian vector bundle has a self-adjoint realization with compact resolvent. Its eigenvalues are real, have finite multiplicities, and accumulate only at infinity. There need not be infinitely many eigenvalues of both signs; boundary problems require boundary conditions.}
\GSfield{Fate}{These spectral facts are older elliptic theory, used as preliminaries in \citet[\S2]{atiyahpatodisinger1976spectral3}; the 1976 paper is a verification source, not their first publication.}
\GSfield{Present status}{As of November 2026, the assertion is established with closedness or an appropriate elliptic boundary realization. Two-sided unbounded spectrum is an additional condition relevant to nontrivial loop examples.}
\GSfield{Assessment}{literal statement: correct with missing hypotheses; publication match: partial; verification: read the result; development: annotated; confidence: high}
\end{GSitem}
\begin{GSitem}{1971-237.2}{A}{Spectral flow for a closed self-adjoint family}
\GSfield{Kind; source}{definition; \GSpage{P000370-P000371}}
\GSfield{The notes}{The eigenvalues move with \(t\), and the net number passing zero defines an integer. The family is periodic or closes by \(D_{t+1}=UD_tU^{-1}\).}
\GSfield{Statement}{A continuous loop of self-adjoint Fredholm operators, in a topology supporting continuous spectral projections near zero, has an integer spectral flow counting negative-to-positive crossings with sign. For smooth elliptic families, it is invariant under homotopies of loops. For open paths, endpoint conditions or a specified endpoint convention are required.}
\GSfield{Fate}{\citet[\S7, pp.~93--94]{atiyahpatodisinger1976spectral3} defines and studies this invariant; the paper notes earlier joint work with Lusztig. The present September 1971 notebook is evidence of earlier seminar circulation; the published paper was received in January 1975 and appeared in 1976.}
\GSfield{Present status}{As of November 2026, spectral flow is established and applies beyond compact-resolvent differential operators. It must not be identified with an unsigned count of crossings.}
\GSfield{Assessment}{literal statement: correct with missing hypotheses; publication match: exact; verification: read the result; development: annotated; confidence: high}
\end{GSitem}
\begin{GSitem}{1971-237.3}{A}{The unit spectral-flow example on the circle}
\GSfield{Kind; source}{construction; \GSpage{P000371}}
\GSfield{The notes}{The example is \(D_t=-i\,d/dx+t\), with \(Uf=e^{ix}f\), and the recorded flow is \(\pm1\).}
\GSfield{Statement}{On the circle of length \(2\pi\), with periodic domain \(H^1\), the eigenvalues are \(n+t\), \(n\in\mathbb Z\). Over any unit interval with invertible endpoints, exactly one eigenvalue crosses upward, giving spectral flow \(+1\). Multiplication by \(e^{ix}\) satisfies \(UD_tU^{-1}=D_t-1\), so closure as written requires \(U=e^{-ix}\) or reversing the conjugation.}
\GSfield{Fate}{The same spectral model \(n+t\) is given in \citet[p.~94]{atiyahpatodisinger1976spectral3}. Its presence here is an early explicit example; the conjugation discrepancy is a faithful sign error in the notebook, independently checked by differentiating \(e^{-ix}f\).}
\GSfield{Present status}{As of November 2026, the example is exact after a consistent orientation and gluing convention. The alternative \(\pm1\) records convention, whereas the displayed conjugation formula must actually be changed.}
\GSfield{Assessment}{literal statement: false as written; publication match: exact; verification: read the result; development: annotated; confidence: high}
\end{GSitem}
\begin{GSitem}{1971-237.4}{A}{Computing spectral flow from geometric data}
\GSfield{Kind; source}{problem; \GSpage{P000371}}
\GSfield{The notes}{The stated problem is to compute spectral flow from the manifold, operator, and gluing \(U\). The proposed answer is the index of \(P=\partial_t+D_t\) on the mapping torus.}
\GSfield{Statement}{Let \(D_t\) be a smooth family of self-adjoint first-order elliptic operators on the fibres of a smooth bundle over \(S^1\), acting on a compatibly glued Hermitian vector bundle. Then the suspended operator \(\partial_t+D_t\), with a connection in the parameter direction, is elliptic and \(\operatorname{sf}(D_t)=\operatorname{ind}(\partial_t+D_t)\). Equivalently spectral flow is the topological index of the family symbol. An arbitrary Hilbert-space unitary need not define such a geometric mapping torus.}
\GSfield{Fate}{\citet[Theorem 7.4 and pp.~94--95]{atiyahpatodisinger1976spectral3} proves precisely the topological formula and the first-order suspension equality, including a remark on nontrivial fibre bundles over the circle. This is a strong A identification, with 1971 circulation distinct from 1976 publication.}
\GSfield{Present status}{As of November 2026, the geometric computation problem is solved under the stated elliptic and gluing hypotheses. The general Fredholm-loop invariant remains meaningful when a local differential mapping-torus model is unavailable.}
\GSfield{Assessment}{literal statement: correct with missing hypotheses; publication match: exact; verification: read the result; development: annotated; confidence: high; problem \hyperref[prob:GP-1971-04]{GP-1971-04}}
\end{GSitem}
\begin{GSitem}{1971-237.5}{A}{Analytic eigenvalue branches and the orbit heuristic}
\GSfield{Kind; source}{method; \GSpage{P000371-P000372}}
\GSfield{The notes}{Real-analytic dependence is used to label eigenvalues globally, and periodicity permutes their branches. Finite orbits and branches travelling towards the two infinities are used to read the integer. The bound \(|\lambda^{\prime}(t)|<C|\lambda(t)|\) is displayed.}
\GSfield{Statement}{For an analytic self-adjoint elliptic family with common domain and uniform elliptic estimates, local real-analytic eigenvalue branches can be continued along the real parameter. Their derivatives obey an estimate of the form \(|\lambda^{\prime}(t)|\le C(1+|\lambda(t)|)\) on compact parameter intervals. The notebook’s homogeneous bound forbids a nonzero derivative at a zero crossing and is contradicted by its own branch \(\lambda(t)=t\).}
\GSfield{Fate}{The later rigorous treatment in \citet[\S7]{atiyahpatodisinger1976spectral3} defines spectral flow through spectral graphs and perturbation, without requiring the informal orbit classification to carry the proof. The branch sketch is an antecedent, not a complete published argument.}
\GSfield{Present status}{As of November 2026, the corrected analytic estimates support the spectral picture; the exact four-type orbit description needs hypotheses and is retained only as a heuristic. No general monotonicity of eigenvalues is implied.}
\GSfield{Assessment}{literal statement: false as written; publication match: partial; verification: read the result; development: annotated; confidence: high}
\end{GSitem}
\begin{GSitem}{1971-237.6}{A}{The diagonal scalar-ODE model for the suspended index}
\GSfield{Kind; source}{method; \GSpage{P000372}}
\GSfield{The notes}{Assuming eigenvectors independent of \(t\), the solutions are \(u_n(t)=c_n\exp(-\int\lambda_n)\). The notes use orbit counts for \(\ker P\) and \(\ker P^*\), and write \(\int_0^N\lambda_n(t)dt=0\) for a finite orbit.}
\GSfield{Statement}{In a diagonal model, \((\partial_t+\lambda(t))u=0\) gives \(u(t)=c\exp(-\int_0^t\lambda(s)ds)\). On a periodic orbit of length \(N\), a nonzero periodic solution exists exactly when the scalar monodromy is one; for real \(\lambda\), this is \(\int_0^N\lambda=0\). It is a condition, not an automatic identity. Such finite orbits contribute equally to kernel and cokernel. Infinite unfolded orbits require explicit square-integrability and gluing conditions.}
\GSfield{Fate}{The notebook labels the calculation heuristic and acknowledges extra terms from moving eigenvectors. The rigorous equality is \citet[pp.~94--95]{atiyahpatodisinger1976spectral3}; the displayed scalar model is only a special-case explanation.}
\GSfield{Present status}{As of November 2026, the elementary ODE calculation is valid with its boundary conditions, but does not by itself prove the general spectral-flow theorem.}
\GSfield{Assessment}{literal statement: correct with missing hypotheses; publication match: partial; verification: read the result; development: annotated; confidence: high}
\end{GSitem}
\begin{GSitem}{1971-237.7}{A}{A diffeomorphism and the obstruction from a natural operator family}
\GSfield{Kind; source}{claim; \GSpage{P000373}}
\GSfield{The notes}{A diffeomorphism \(T:X\to X\), a metric, and a natural self-adjoint operator determine a family on the mapping torus. If \(T\) preserves a metric, the family is constant and its flow vanishes.}
\GSfield{Statement}{For a diffeomorphism-equivariant elliptic construction, compatible endpoint identification and a path of metrics give a closed operator family. If the spectral-flow class is independent of that path and an invariant metric exists, one may use the constant metric, so the invariant is zero. Consequently a verified nonzero value obstructs an invariant metric; the conclusion is conditional on the operator family actually having nonzero possible flow.}
\GSfield{Fate}{The mapping-torus extension appears in \citet[remark after Theorem 7.4]{atiyahpatodisinger1976spectral3}. The notebook’s geometric observation is a correct application of homotopy invariance, but its next choice \(d+\delta\) requires a separate audit.}
\GSfield{Present status}{As of November 2026, the implication is valid for a suitably specified natural elliptic family. It does not establish that every suggested natural operator yields an effective obstruction.}
\GSfield{Assessment}{literal statement: correct with missing hypotheses; publication match: partial; verification: read the result; development: annotated; confidence: high}
\end{GSitem}
\begin{GSitem}{1971-237.8}{F}{The ungraded de Rham operator has zero spectral flow}
\GSfield{Kind; source}{claim; \GSpage{P000373}}
\GSfield{The notes}{The example offered for the preceding construction is \(d+\delta\) on forms. No chirality restriction or odd-signature modification is written.}
\GSfield{Statement}{On all differential forms, \(D_t=d_t+d_t^*\) anticommutes with degree parity \(\Gamma\omega=(-1)^{\deg\omega}\omega\). If the endpoint gluing preserves parity, a closed family is unitarily conjugate to its negative; spectral-flow invariance and sign reversal give \(\operatorname{sf}(D_t)=0\). A potentially nontrivial signature construction uses a different restricted or chirality-modified operator.}
\GSfield{Reconstruction}{\ours\ The parity-symmetry deduction is supplied here and proved in the chapter’s verification subsection. It diagnoses the literal recorded example; it does not reconstruct an unrecorded operator restriction.}
\GSfield{Fate}{The general spectral-flow theorem is \citet[\S7]{atiyahpatodisinger1976spectral3}; its signature discussion in \S6 uses the appropriate signature operator. The notebook does not record the missing modification. The chapter’s grading-symmetry verification explains why its ungraded example cannot supply a nonzero obstruction.}
\GSfield{Present status}{As of November 2026, the vanishing is an elementary consequence of parity symmetry. The intended nontrivial example cannot be recovered merely by citing the general index theorem.}
\GSfield{Assessment}{literal statement: false as written; publication match: partial; verification: read the result; development: annotated; confidence: medium}
\end{GSitem}
\begin{GSitem}{1971-237.9}{F}{The incomplete higher-signature example}
\GSfield{Kind; source}{remark; \GSpage{P000373}}
\GSfield{The notes}{The page mentions characters \(\pi_1(X)\to\mathbb Z\to S^1\), \(X^{4k+1}\times S^1\), Novikov’s invariant, and Hirzebruch’s \(L\)-polynomial.}
\GSfield{Statement}{A family of flat line-bundle twists over a character circle can be paired with an odd elliptic symbol through an index theorem. The notes do not specify a signature-type operator, its grading, the twisting bundle on the product, or the precise cohomology pairing. Therefore no unique higher-signature formula follows from the recorded data.}
\GSfield{Fate}{\citet[\S\S5--7]{atiyahpatodisinger1976spectral3} provides related index formulas for flat bundles and signature operators. The match here is thematic, unlike the exact identification of item 1971-237.4. A dimensional or grading repair must not be silently attributed to Atiyah’s unrecorded speech.}
\GSfield{Present status}{As of November 2026, the general subject is developed, but this particular compressed example remains incompletely identifiable. The literal \(d+\delta\) interpretation has the zero-flow limitation just explained.}
\GSfield{Assessment}{literal statement: not determinable; publication match: thematic only; verification: read the result; development: annotated; confidence: low}
\end{GSitem}

\GStalk{238}{1971-09-20}{Pierre Deligne; translated by Yu. I. Manin}{P000373}{Motives: a two-line notice}
The entry says only that Deligne, with Manin translating, spoke about motives and the various cohomology theories of algebraic varieties. It records no theorem, construction, or conjecture sufficiently specific for a publication-level identification.

\begin{GSitem}{1971-238.1}{F}{Motives and different cohomology theories: no recorded result}
\GSfield{Kind; source}{remark; \GSpage{P000373}}
\GSfield{The notes}{The entry reads Deligne, in Manin’s translation, on the theory of motives and different cohomology theories of algebraic varieties.}
\GSfield{Statement}{The text supports only this subject identification. It gives no category of motives, equivalence relation on algebraic cycles, realization functor, or conjectural comparison that could be made into a uniquely attributable theorem.}
\GSfield{Fate}{No specific publication is assigned. Attaching the Weil conjectures or a later motivic theory would fill a gap with information not contained in the notes.}
\GSfield{Present status}{As of November 2026, the entry remains a subject notice; this is insufficient evidence for any claim about the fate of a particular mathematical problem.}
\GSfield{Assessment}{literal statement: not determinable; publication match: none found; verification: unresolved lead; development: none; confidence: high}
\end{GSitem}

\GStalk{239}{1971-10-25}{S. I. Gelfand}{P000373-P000375}{Schubert cells and invariant characteristic forms}
The talk compares the Schubert decomposition of a complex Grassmannian with invariant differential forms, curvature polynomials, and Schur functions. A separate intervention by I. M. Gelfand asks broadly about a defining-characteristic group algebra for \(SL_2(\mathbb F_q)\); the intended classification target is not specified. Vinogradov’s talk starts on the lower part of P000375.

\begin{GSitem}{1971-239.1}{E}{Schubert cells in a complex Grassmannian}
\GSfield{Kind; source}{construction; \GSpage{P000373-P000374}}
\GSfield{The notes}{The notes choose \(a_1\le\cdots\le a_n\le N\), a complete flag \(E_j\), and jumps \(a_i+i\) for \(\dim(F\cap E_j)\).}
\GSfield{Statement}{For \(\operatorname{Gr}(n,\mathbb C^N)\), choose \(0\le a_1\le\cdots\le a_n\le N-n\). The corresponding Schubert cell consists of planes whose intersection-dimension function jumps precisely at \(b_i=a_i+i\); it is an affine space of complex dimension \(\sum_i a_i\). The closure is defined by the associated rank inequalities and is a projective Schubert variety.}
\GSfield{Fate}{The notebook explicitly cites \citet{horrocks1957schubert}, published in 1957. The geometry and its indexing are also checked in \citet[\S4]{tamvakis2004connection}; his quotient-bundle convention is transposed from the usual subspace convention and should not be copied without translation.}
\GSfield{Present status}{As of November 2026, the affine decomposition is established. Requiring equality only at a few selected flag dimensions, without the complete jump data, does not specify the open cell.}
\GSfield{Assessment}{literal statement: correct with missing hypotheses; publication match: partial; verification: read the result; development: annotated; confidence: high}
\end{GSitem}
\begin{GSitem}{1971-239.2}{E}{Schubert cycles give an integral homology basis}
\GSfield{Kind; source}{theorem; \GSpage{P000374}}
\GSfield{The notes}{The cell decomposition is used to compare bases in homology and cohomology.}
\GSfield{Statement}{The closures of the complex Schubert cells give a \(\mathbb Z\)-basis of the homology of the Grassmannian. The affine cells have even real dimensions, so the cellular boundary maps vanish and odd homology is zero. A cohomology basis is obtained by Poincaré duality, with complementary rather than equal complex dimensions.}
\GSfield{Fate}{This is classical Schubert calculus underlying \citet{horrocks1957schubert}; \citet[\S4]{tamvakis2004connection} makes the cycle/cohomology conventions explicit. The seminar is expository.}
\GSfield{Present status}{As of November 2026, the integral assertion is established, stronger than a real de Rham basis. Open noncompact cells themselves are not ordinary compact integration cycles.}
\GSfield{Assessment}{literal statement: correct with missing hypotheses; publication match: partial; verification: read the result; development: annotated; confidence: high}
\end{GSitem}
\begin{GSitem}{1971-239.3}{E}{Basic invariant forms on the homogeneous quotient}
\GSfield{Kind; source}{method; \GSpage{P000374}}
\GSfield{The notes}{The Grassmannian is \(U(N)/(U(n)\times U(N-n))\). Forms are lifted to \(U(N)\); the first descent condition is written as vanishing when the listed arguments lie in the subgroup tangent space.}
\GSfield{Statement}{For a principal \(H\)-bundle \(G\to G/H\), a differential form descends precisely when it is \(H\)-invariant and horizontal: insertion of any one vertical tangent vector gives zero. Vanishing only when all arguments are vertical is insufficient in degree greater than one. Invariant forms on the Grassmannian are the \(H\)-invariant alternating tensors on \(\mathfrak g/\mathfrak h\).}
\GSfield{Fate}{This is the basic-form construction behind the Chern--Weil treatment in \citet[\S3]{tamvakis2004connection}, tracing the classical work of Chern, Weil, and Cartan. The descent criterion is older than the talk.}
\GSfield{Present status}{As of November 2026, the criterion is established. The notebook’s compressed first condition must be read in the stronger insertion sense for its later construction to work.}
\GSfield{Assessment}{literal statement: false as written; publication match: partial; verification: read the result; development: annotated; confidence: high}
\end{GSitem}
\begin{GSitem}{1971-239.4}{E}{The curvature matrix and invariant characteristic forms}
\GSfield{Kind; source}{construction; \GSpage{P000374}}
\GSfield{The notes}{The Maurer--Cartan entries \(\theta_{jk}\) are used in \(\Theta_{jk}=-(2\pi i)^{-1}\sum_s\theta_{js}\wedge\overline{\Theta}_{ks}\); the second factor is marked graphically uncertain. Conjugation-invariant polynomials in this matrix descend to the Grassmannian.}
\GSfield{Statement}{For the tautological Hermitian bundle or its dual, the induced connection has a curvature matrix whose entries are sums of wedges of two off-diagonal Maurer--Cartan one-forms. In a compatible convention the displayed second factor must be \(\overline\theta_{ks}\), so the result has degree two. A degree-\(r\) invariant polynomial in normalized curvature is a closed \(2r\)-form with connection-independent cohomology class.}
\GSfield{Reconstruction}{The lower-case second factor is a mathematical reconstruction. The case uncertainty is preserved in the corrections rather than called a proven transcription error.}
\GSfield{Fate}{This is classical Chern--Weil theory, explicitly reviewed in \citet[\S3 and \S4]{tamvakis2004connection}. The scan does not conclusively distinguish upper- and lower-case theta in the uncertain factor; degree counting gives the mathematical emendation, not paleographic certainty.}
\GSfield{Present status}{As of November 2026, the curvature construction is established. The sign depends on whether the tautological bundle, quotient, or dual is used; that choice must be made before identifying its Chern or Schur class.}
\GSfield{Assessment}{literal statement: correct with missing hypotheses; publication match: partial; verification: read the result; development: annotated; confidence: high}
\end{GSitem}
\begin{GSitem}{1971-239.5}{E}{The alternant quotient as a Schur polynomial and character}
\GSfield{Kind; source}{theorem; \GSpage{P000374-P000375}}
\GSfield{The notes}{A quotient of determinants with exponents \(\lambda_n,\lambda_{n-1}+1,\ldots,\lambda_1+n-1\) is said to be a polynomial and the character of the representation with highest weight \(\lambda_1\ge\cdots\ge\lambda_n\ge0\).}
\GSfield{Statement}{For a partition \(\lambda\) of length at most \(n\), \(s_\lambda(x)=\det(x_i^{\lambda_j+n-j})/\det(x_i^{n-j})\) extends to a symmetric polynomial, including repeated eigenvalues. It is the character of the irreducible polynomial representation of \(GL_n\) with highest weight \(\lambda\). Substituting a matrix of commuting even-degree forms is interpreted through the polynomial, not by dividing differential forms.}
\GSfield{Fate}{The representation-theoretic identity is classical and is the algebraic input of \citet{horrocks1957schubert}. \citet[\S2]{tamvakis2004connection} verifies the Schur functor and invariant-polynomial interpretation directly.}
\GSfield{Present status}{As of November 2026, the formula is established. The notebook’s reversed order of determinant columns gives the same quotient because numerator and denominator acquire the same sign.}
\GSfield{Assessment}{literal statement: correct; publication match: partial; verification: read the result; development: annotated; confidence: high}
\end{GSitem}
\begin{GSitem}{1971-239.6}{E}{Schur curvature forms and the dual Schubert basis}
\GSfield{Kind; source}{claim; \GSpage{P000375}}
\GSfield{The notes}{The notes substitute curvature into the Schur polynomial, put \(a_i=N-n-\lambda_i\), and assert integral one over the corresponding Schubert cell and zero over the others; form degree is equated with cell dimension.}
\GSfield{Statement}{In the standard subspace convention, let \(S\) be the rank-\(n\) tautological bundle and let \(\lambda\) fit the \(n\times(N-n)\) rectangle. The Schur class \(s_\lambda(S^*)\) is the Schubert cohomology class \(\sigma_\lambda\) of degree \(2|\lambda|\). If \(\lambda^\vee_i=N-n-\lambda_{n+1-i}\), then \(\int_{[\Omega_\mu]}\sigma_\lambda=\delta_{\mu,\lambda^\vee}\) for \(|\mu|=n(N-n)-|\lambda|\). Integrate over closed Schubert cycles, with consistent indexing and real degree.}
\GSfield{Reconstruction}{The dual bundle and standard partition convention are fixed here explicitly; the notebook’s uncertain \(N^{\prime}\) is not interpreted as a new ambient dimension.}
\GSfield{Fate}{This is the Schur-form/Schubert-class relation of \citet{horrocks1957schubert}; a complete modern formulation is \citet[Theorem 1, \S4]{tamvakis2004connection}. The notebook’s displayed indexing conflates the cycle Poincaré-dual to a class with the cycle on which that class evaluates to one.}
\GSfield{Present status}{As of November 2026, the corrected duality statement is established. The literal formula is generally dimensionally inconsistent unless the partition and cycle are complemented appropriately.}
\GSfield{Assessment}{literal statement: false as written; publication match: partial; verification: read the result; development: annotated; confidence: high}
\end{GSitem}
\begin{GSitem}{1971-239.7}{F}{Gelfand’s unspecified request to describe a modular group algebra}
\GSfield{Kind; source}{problem; \GSpage{P000375}; I. M. Gelfand}
\GSfield{The notes}{I. M. Gelfand asks to describe the group ring of \(SL_2(\mathbb F_q)\) with coefficients in the algebraic closure of \(\mathbb F_q\).}
\GSfield{Statement}{The object is \(\overline{\mathbb F}_q[SL_2(\mathbb F_q)]\), in the defining characteristic. The word \emph{describe} does not specify whether one seeks simple modules, blocks, the Jacobson radical, a presentation, indecomposables, or a Morita-equivalent algebra. These are different questions. Since the characteristic divides the group order, a semisimple character decomposition alone does not describe this algebra.}
\GSfield{Fate}{The object is identifiable but the requested invariant is not. Searches of modular \(SL_2\) group-algebra descriptions did not turn this wording into one uniquely matchable problem, so no theorem about simple representations is substituted for a full algebra classification.}
\GSfield{Present status}{As of November 2026, this occurrence remains an underdetermined request, not a certified unsolved problem. It receives no new problem-index identifier because no definite target is stated.}
\GSfield{Assessment}{literal statement: not determinable; publication match: none found; verification: unresolved lead; development: annotated; confidence: medium}
\end{GSitem}

\GStalk{240}{1971-10-25}{A. M. Vinogradov}{P000375-P000377}{Differential operators as functors over a ring}
Vinogradov defines order by iterated commutators with multiplication and introduces the left and right module structures underlying jets, cojets, and higher derivations. His 1972 announcement gives the coherent version over a commutative unital ring. The notebook’s apparent permission to take an arbitrary noncommutative ring cannot be carried through unchanged.

\begin{GSitem}{1971-240.1}{E}{Algebraic order of a differential operator}
\GSfield{Kind; source}{definition; \GSpage{P000375-P000376}}
\GSfield{The notes}{Vinogradov proposes an algebraic theory over any ring, explicitly allowing a noncommutative ring. For an additive map \(A:M\to N\), the notes define \((kA)m=kA(m)\), \((k^+A)m=A(km)\), and \(\delta_k=k^+-k\); order at most \(n\) means every \((n+1)\)-fold iterated \(\delta\) vanishes.}
\GSfield{Statement}{Let \(K\) be a commutative unital algebra over a commutative base \(R\), and let \(M,N\) be unital \(K\)-modules. An \(R\)-linear map \(A:M\to N\) has order at most \(n\) when \(\delta_{k_0}\cdots\delta_{k_n}A=0\) for all \(k_j\in K\). Order zero means \(K\)-linearity. Additive operators correspond to the base \(R=\mathbb Z\). For a noncommutative ring these same formulas do not furnish the ensuing commutative two-module calculus without additional choices.}
\GSfield{Fate}{The commutator definition predates this lecture: \citet[\S16.8]{grothendieck1967egaiv4} was published in 1967. The closely matching presentation in \citet[\S1]{vinogradov1972logic}, received 19 January 1972, explicitly assumes a commutative ring with identity; that hypothesis is essential to its subsequent functors.}
\GSfield{Present status}{As of November 2026, the commutative definition is standard. Noncommutative differential calculi require a specified replacement; the notebook’s unrestricted wording is not a theorem extending the entire construction.}
\GSfield{Assessment}{literal statement: false as written; publication match: exact; verification: read the result; development: annotated; confidence: high}
\end{GSitem}
\begin{GSitem}{1971-240.2}{A}{Agreement with differential operators on vector bundles}
\GSfield{Kind; source}{theorem; \GSpage{P000376}}
\GSfield{The notes}{The notes state that the algebraic definition agrees with the classical one.}
\GSfield{Statement}{For a smooth finite-dimensional manifold \(X\), take \(K=C^\infty(X,\mathbb R)\), \(R=\mathbb R\), and \(M=\Gamma(E),N=\Gamma(F)\) for smooth finite-rank vector bundles. The \(\mathbb R\)-linear operators satisfying the iterated-commutator criterion of order at most \(n\) are exactly the usual linear differential operators of that order. This assertion concerns section modules, not arbitrary algebraic modules over \(C^\infty(X)\).}
\GSfield{Fate}{\citet[\S1]{vinogradov1972logic} states this comparison as part of the announced programme; \citet[Theorem 2.2]{vinogradov2015logic} restates it. The classical differential operators themselves were already familiar; Class A concerns this lecture’s published algebraic presentation.}
\GSfield{Present status}{As of November 2026, the comparison is established under the stated smooth hypotheses. Identification of every universal algebraic form with an ordinary smooth form requires a further category restriction, discussed in occurrence 1971-241.3.}
\GSfield{Assessment}{literal statement: correct with missing hypotheses; publication match: exact; verification: read the result; development: annotated; confidence: high}
\end{GSitem}
\begin{GSitem}{1971-240.3}{A}{Cojet corepresentation and its evaluation operator}
\GSfield{Kind; source}{construction; \GSpage{P000376}}
\GSfield{The notes}{The right module \(\operatorname{Diff}_n^+(M,N)\) is written as \(\operatorname{Hom}_K(M,\operatorname{Diff}_n^+N)\), with \(\operatorname{Diff}_n^+N=\operatorname{Diff}_n^+(K,N)\). The latter object is called cojets.}
\GSfield{Statement}{For commutative unital \(K/R\), put \(C_n(N)=\operatorname{Diff}_{K/R,n}(K,N)\) with scalar action \((a\cdot B)(b)=B(ab)\). There is a natural \(K\)-module isomorphism \(\operatorname{Diff}_n^+(M,N)\cong\operatorname{Hom}_K(M,C_n(N))\), given by \(A\mapsto[m\mapsto(b\mapsto A(bm))]\). Its inverse is composition with \(\epsilon(B)=B(1)\), an operator of order at most \(n\).}
\GSfield{Fate}{This exact cojet formulation is printed in \citet[\S1]{vinogradov1972logic}. The November 8 continuation supplies its proof in six elementary lemmas, recorded separately as occurrences 1971-242.1--6.}
\GSfield{Present status}{As of November 2026, the natural isomorphism is established. The precomposition action is indispensable: replacing it by pointwise multiplication would make the displayed currying map generally nonlinear over \(K\).}
\GSfield{Assessment}{literal statement: correct with missing hypotheses; publication match: exact; verification: read the result; development: annotated; confidence: high}
\end{GSitem}
\begin{GSitem}{1971-240.4}{E}{Jets represent differential operators in the target variable}
\GSfield{Kind; source}{construction; \GSpage{P000376}}
\GSfield{The notes}{The notes pair the cojet formula with \(\operatorname{Diff}_n(M,N)=\operatorname{Hom}(J^n(M),N)\) and announce existence of jets.}
\GSfield{Statement}{For commutative \(K/R\), a universal map \(j_n:M\to J^n_{K/R}(M)\) of order at most \(n\) induces \(\operatorname{Hom}_K(J^n_{K/R}(M),N)\cong\operatorname{Diff}_{K/R,n}(M,N)\) naturally in \(N\). Here the operator module uses multiplication on its values. Existence follows from the quotient construction written explicitly in the next lecture; no finite generation or projectivity is required for algebraic module representability.}
\GSfield{Fate}{Principal parts and their universal property are older, in \citet[\S16.8]{grothendieck1967egaiv4}. \citet[\S1]{vinogradov1972logic} places the two universal properties side by side. The construction is expository as a representability fact.}
\GSfield{Present status}{As of November 2026, the algebraic universal property is established. Being a finite-rank geometric jet bundle is a stronger smooth or regularity statement.}
\GSfield{Assessment}{literal statement: correct with missing hypotheses; publication match: exact; verification: read the result; development: annotated; confidence: high}
\end{GSitem}
\begin{GSitem}{1971-240.5}{E}{The algebra generated by the two scalar actions}
\GSfield{Kind; source}{claim; \GSpage{P000376}}
\GSfield{The notes}{The notes suggest constructing \(J^n(K)\) as the algebra generated by \(k\) and \(k^+\) on \(\operatorname{Diff}_nK\); a lower index in the formula remains uncertain.}
\GSfield{Statement}{The canonical algebraic construction is \(J^n_{K/R}(K)=(K\otimes_R K)/I^{n+1}\), where \(I=\ker(K\otimes_RK\to K)\). The two factors give the two scalar actions and their differences generate \(I\). An action of this quotient on an operator module need not be faithful, so identifying the universal algebra with its image on one chosen module needs a separate argument.}
\GSfield{Reconstruction}{The diagonal-ideal quotient is a mathematical reconstruction, not a claimed reading of the crossed notation.}
\GSfield{Fate}{This is the principal-parts construction of \citet[\S16.8]{grothendieck1967egaiv4}; the jet quotient is also given in \citet[Proposition 7.1]{vinogradov2015logic}. Neither source justifies an unrestricted faithfulness claim for the notebook’s chosen representation.}
\GSfield{Present status}{As of November 2026, the universal quotient is established; the exact intended lower index and whether equality with an operator image was meant remain uncertain.}
\GSfield{Assessment}{literal statement: correct with missing hypotheses; publication match: partial; verification: read the result; development: annotated; confidence: medium}
\end{GSitem}
\begin{GSitem}{1971-240.6}{A}{Recursive multiderivation functors}
\GSfield{Kind; source}{construction; \GSpage{P000377}}
\GSfield{The notes}{Before the 1 November heading, the notes define \(D_1(N)=\operatorname{Der}(K,N)\), explain that it is a left submodule but only a subset of \(\operatorname{Diff}_1^+N\), and put \(D_i(N)=D\{D_{i-1}(N)\subset(\operatorname{Diff}_1^+)^{i-1}N\}\).}
\GSfield{Statement}{For a subset \(S\) of a \(K\)-module \(Q\), write \(D(S\subset Q)\) for the \(R\)-derivations \(K\to Q\) whose values lie in \(S\). Starting with \(D_0(N)=N\), the notebook’s recursive definition uses the indicated precomposition module at each stage. It cannot be replaced by repeated application of the ordinary left-module functor \(D\). In the smooth geometric category it produces multivector fields with coefficients in \(N\).}
\GSfield{Fate}{The recursive construction and the warning that the intermediate subset is not a right submodule appear in \citet[\S2]{vinogradov1972logic}; the later explanation is \citet[\S5.2]{vinogradov2015logic}.}
\GSfield{Present status}{As of November 2026, this construction is established in the commutative setting. Its representing objects are forms; the functor values are multiderivations, a distinction blurred in the next lecture.}
\GSfield{Assessment}{literal statement: correct with missing hypotheses; publication match: exact; verification: read the result; development: annotated; confidence: high}
\end{GSitem}

\GStalk{241}{1971-11-01}{A. M. Vinogradov}{P000377-P000383}{Jets, higher derivations, and algebraic differential complexes}
The second lecture develops universal jet modules, the recursive derivation functors, symbols, and de Rham and Spencer constructions. The characteristic-two discussion must distinguish skew symmetry from alternation. A literal example using only the prime field \(\mathbb F_2\) cannot support nonzero derivations; nearby dual-number and group-algebra examples are mathematically different.

\begin{GSitem}{1971-241.1}{A}{Repeated commutators recover the Hamilton–Jacobi symbol}
\GSfield{Kind; source}{method; \GSpage{P000377}}
\GSfield{The notes}{The 1 November entry calls \(\delta_f\cdots\delta_f\mathcal D\), with the same function everywhere, the left side of a Hamilton--Jacobi equation, and asks the audience to calculate it.}
\GSfield{Statement}{For a scalar differential operator \(A=\sum_{|\alpha|\le m}a_\alpha(x)\partial^\alpha\) over a smooth real manifold, use the symbol convention \(\sigma_m(A)(x,\xi)=\sum_{|\alpha|=m}a_\alpha(x)\xi^\alpha\), without factors of \(i\). Let \(m_f\) denote multiplication by \(f\), and define \(\delta_fA=[A,m_f]=A\circ m_f-m_f\circ A\). Then \(\delta_f^mA\) is multiplication by \(m!\,\sigma_m(A)(x,df_x)\). Thus the leading characteristic phase equation is \(\sigma_m(A)(x,df_x)=0\). Repeated commutators over a general ring need not permit division by \(m!\).}
\GSfield{Reconstruction}{\ours{} The precise normalization follows by applying the product rule \(m\) times: all lower-order terms disappear and each top-order derivative contributes the corresponding symmetric product of \(df\), with \(m!\) permutations.}
\GSfield{Fate}{\citet[\S7]{vinogradov1972logic} publishes the symbol/Hamiltonian part of this programme in 1972. The principal-symbol and Hamilton--Jacobi concepts are older; Class A records the algebraic reformulation anticipated by this exercise, not a claim of priority for those concepts.}
\GSfield{Present status}{As of November 2026, the scalar smooth identity is established and elementary. A Hamilton--Jacobi equation for a system or an inhomogeneous problem requires an additional specified principal-symbol equation.}
\GSfield{Assessment}{literal statement: correct with missing hypotheses; publication match: partial; verification: read the result; development: annotated; confidence: high}
\end{GSitem}
\begin{GSitem}{1971-241.2}{E}{The tensor-quotient construction of universal jets}
\GSfield{Kind; source}{construction; \GSpage{P000377-P000378}}
\GSfield{The notes}{The notes write \(J^s(M)=(K\otimes_{\mathbb Z}M)/\Pi_s\), where \(\Pi_s\) is generated by \((\bar k_0-k_0)\cdots(\bar k_s-k_s)\) applied to tensors.}
\GSfield{Statement}{For commutative \(K/R\), give \(K\otimes_R M\) the two actions \(a(b\otimes m)=ab\otimes m\) and \(\bar a(b\otimes m)=b\otimes am\). Let \(\Pi_s\) be the submodule generated by every \((\bar a_0-a_0)\cdots(\bar a_s-a_s)(b\otimes m)\). The quotient with \(j_s(m)=[1\otimes m]\) represents order-\(s\) operators: \(A\) corresponds to the unique homomorphism \([b\otimes m]\mapsto bA(m)\).}
\GSfield{Fate}{This is the older principal-parts universal construction, in \citet[\S16.8]{grothendieck1967egaiv4}; the lecture’s exact module form is restated in \citet[Proposition 7.1]{vinogradov2015logic}. The \(\mathbb Z\)-tensor in the notes correctly corresponds to additive, absolute operators.}
\GSfield{Present status}{As of November 2026, the quotient construction works without smoothness, finite generation, or freeness. The base ring must remain fixed when one compares it with relative differential operators.}
\GSfield{Assessment}{literal statement: correct with missing hypotheses; publication match: exact; verification: read the result; development: annotated; confidence: high}
\end{GSitem}
\begin{GSitem}{1971-241.3}{A}{Which universal forms recover ordinary smooth forms?}
\GSfield{Kind; source}{claim; \GSpage{P000378-P000381}}
\GSfield{The notes}{The representing modules \(\Lambda_i\) are identified with ordinary differential forms, while all \(K\)-modules are apparently allowed.}
\GSfield{Statement}{An identification with the ordinary smooth cotangent and exterior-form modules requires the geometric category over \(C^\infty(X)\), or an equivalent smooth-module universal property. Universal algebraic Kähler differentials representing derivations into every algebraic module can be larger. The distinction does not invalidate the differential-operator comparison for finite-rank bundle sections.}
\GSfield{Fate}{\citet[Proposition 7.2 and Example 7.1]{vinogradov2015logic} explicitly separates the two categories; its example warns that algebraic differentiation of \(e^x\) need not satisfy the ordinary smooth chain rule in the unrestricted universal module. This clarifies a suppressed hypothesis in the 1971 account.}
\GSfield{Present status}{As of November 2026, the geometric-category reconstruction is established. Calling the unrestricted representing objects ordinary smooth forms is not literally correct.}
\GSfield{Assessment}{literal statement: false as written; publication match: partial; verification: read the result; development: annotated; confidence: high}
\end{GSitem}
\begin{GSitem}{1971-241.4}{A}{Skew multiderivations and the distinction from covariant forms}
\GSfield{Kind; source}{construction; \GSpage{P000378-P000379}}
\GSfield{The notes}{The repeated derivation calculation yields \(D(a)(b)+D(b)(a)=0\). The notes then call \(D_k(M)\) \emph{\(k\)-forms with values in \(M\)}. In the two-coordinate calculation, several coefficients multiplying \(b\) are followed by \(\partial_x\) rather than \(\partial_y\).}
\GSfield{Statement}{The recursive objects are multiderivations. If 2 is invertible in the base, their skew symmetry gives alternating maps, and \(D_k(M)\cong\operatorname{Hom}_K(\bigwedge^k\Omega^1_{K/R},M)\). Over a smooth real manifold, with geometric modules, these are multivector fields with coefficients in \(M\); \(\Lambda_k\), the representing covariant object, is the form module. In characteristic two, skew symmetry alone does not impose alternation.}
\GSfield{Fate}{\citet[\S2]{vinogradov1972logic} publishes the recursive functors; \citet[\S5.2]{vinogradov2015logic} explicitly calls their smooth realizations multivector fields. The variance distinction repairs the notebook’s terminology.}
\GSfield{Present status}{As of November 2026, the characteristic-zero interpretation is established. The characteristic-two examples must be examined with their exact functor and differential, rather than silently importing ordinary exterior powers.}
\GSfield{Assessment}{literal statement: false as written; publication match: exact; verification: read the result; development: annotated; confidence: high}
\end{GSitem}
\begin{GSitem}{1971-241.5}{F}{The claimed nonzero second functor over the prime field}
\GSfield{Kind; source}{claim; \GSpage{P000379}}
\GSfield{The notes}{The bottom line reads \(K=\mathbb Z_2\), \(D_2\ne0\), followed by a parenthesis referring to second de Rham cohomology.}
\GSfield{Statement}{If \(\mathbb Z_2\) denotes the field \(\mathbb F_2\), the literal claim is false for the stated additive-derivation functors: every derivation from \(\mathbb F_2\) to a unital module is zero, so \(D_i(M)=0\) for every \(i\ge1\). A nonreduced algebra such as \(\mathbb F_2[C_2]\), considered on the next page, is a different object. Nonzero \(D_2\) is also not the same assertion as nonzero second cohomology.}
\GSfield{Reconstruction}{\ours{} A derivation satisfies \(D(1)=D(1\cdot1)=2D(1)\), hence \(D(1)=0\); additivity then makes it zero on the prime field. Induction in the notebook’s definition gives the higher vanishing.}
\GSfield{Fate}{No matching published statement was identified in the examined 1972 announcement or later survey. Their definitions suffice to reject the literal prime-field example; the neighbouring group-algebra example is a possible source of compression, not a justified silent substitution.}
\GSfield{Present status}{As of November 2026, the literal prime-field assertion is refuted. The intended alternative example and the parenthetical identification remain unverified.}
\GSfield{Assessment}{literal statement: false as written; publication match: none found; verification: unresolved lead; development: annotated; confidence: medium}
\end{GSitem}
\begin{GSitem}{1971-241.6}{F}{The group-ring/de Rham comparison is not identified}
\GSfield{Kind; source}{claim; \GSpage{P000380}}
\GSfield{The notes}{For the group algebra of the two-element group over \(\mathbb Z_2\), the notes claim all de Rham cohomology groups are \(\mathbb Z_2\), then assert that de Rham cohomology of group rings equals group cohomology. They contrast this with the diagonal-ideal definition of classical forms.}
\GSfield{Statement}{The claim requires a specified complex and coefficient convention. For ordinary exterior Kähler forms of \(A=\mathbb F_2[C_2]\cong\mathbb F_2[y]/(y^2)\), one obtains \(H^0_{\rm dR}(A/\mathbb F_2)=\mathbb F_2\), \(H^1_{\rm dR}(A/\mathbb F_2)=\mathbb F_2\), and zero in higher degrees. Thus this ordinary complex cannot realize the stated nonzero group in every degree. A different recursively defined characteristic-two complex cannot be identified solely from that numerical claim.}
\GSfield{Reconstruction}{\ours{} For the stated ordinary complex, \(\Omega^1=A\,dy\), \(\Omega^i=0\) for \(i\ge2\), and \(d(a+by)=b\,dy\). Its kernel and cokernel give the displayed calculation. This is an editorial diagnostic, not a computation of the unidentified alternative complex.}
\GSfield{Fate}{English and Russian searches for the group-ring comparison, and examination of \citet{vinogradov1972logic,vinogradov2015logic}, did not locate a theorem matching the unrestricted assertion. Its relation to later algebraic de Rham constructions remains unidentified.}
\GSfield{Present status}{As of November 2026, no exact publication match for the notebook’s alternative theory is claimed. The ordinary de Rham interpretation is false; this does not establish that every conceivable corrected comparison fails.}
\GSfield{Assessment}{literal statement: not determinable; publication match: none found; verification: unresolved lead; development: annotated; confidence: low}
\end{GSitem}
\begin{GSitem}{1971-241.7}{A}{Representation of multiderivations by universal forms}
\GSfield{Kind; source}{theorem; \GSpage{P000379-P000381}}
\GSfield{The notes}{The lecture announces existence of \(\Lambda_i(K)\) with \(D_i(M)=\operatorname{Hom}_K(\Lambda_i(K),M)\).}
\GSfield{Statement}{In the commutative algebraic setting, the multiderivation functors are represented by universal modules. Over a base in which 2 is invertible, the standard alternating interpretation gives \(\Lambda_i(K)=\bigwedge^i\Omega^1_{K/R}\), with \(\Lambda_0=K\). Over general rings the exact skew-versus-alternating convention must be retained. Universal objects for a restricted smooth category are taken in that category.}
\GSfield{Fate}{The announcement \citet[\S5]{vinogradov1972logic} gives the representing-module formulation; \citet[\S7.2]{vinogradov2015logic} makes category dependence explicit. Kähler differentials and exterior powers themselves are older constructions.}
\GSfield{Present status}{As of November 2026, the universal-property formulation is established. It does not identify multiderivations themselves with covariant forms, nor erase the characteristic-two distinction.}
\GSfield{Assessment}{literal statement: correct with missing hypotheses; publication match: exact; verification: read the result; development: annotated; confidence: high}
\end{GSitem}
\begin{GSitem}{1971-241.8}{A}{The exterior differential from a natural transformation}
\GSfield{Kind; source}{construction; \GSpage{P000380-P000382}}
\GSfield{The notes}{The inclusion \(D_i\to D_{i-1}(\operatorname{Diff}_1^+)\) is translated through universal objects into a first-order operator \(d:\Lambda_{i-1}\to\Lambda_i\). The notes check that consecutive compositions vanish.}
\GSfield{Statement}{In the commutative setting with a consistent form convention, the natural transformation is represented by a map \(J^1(\Lambda_{i-1})\to\Lambda_i\); composing with the universal jet gives \(d\). For ordinary alternating forms, \(d(a_0\,da_1\wedge\cdots\wedge da_i)=da_0\wedge da_1\wedge\cdots\wedge da_i\), which gives \(d^2=0\). The differential is first order over \(K\), generally not \(K\)-linear.}
\GSfield{Fate}{This natural-transformation construction is published in \citet[\S5]{vinogradov1972logic}; the representative/jet mechanism is restated in \citet[\S7.2]{vinogradov2015logic}. Class A concerns the categorical construction rather than the classical existence of \(d\).}
\GSfield{Present status}{As of November 2026, the de Rham complex is established. Its cohomology depends on the ring, base, and module category; the disputed group-ring claim is not a consequence of representability alone.}
\GSfield{Assessment}{literal statement: correct with missing hypotheses; publication match: exact; verification: read the result; development: annotated; confidence: high}
\end{GSitem}
\begin{GSitem}{1971-241.9}{A}{Finite and stable algebraic Spencer complexes}
\GSfield{Kind; source}{construction; \GSpage{P000382-P000383}}
\GSfield{The notes}{Composition of operators and the functor inclusions give \(D_{i+1}(\operatorname{Diff}_k^+M)\to D_i(\operatorname{Diff}_{k+1}^+M)\). The finite and direct-limit chains are called a Spencer resolution, followed by the more cautious statement \emph{this is a complex}.}
\GSfield{Statement}{For each total order \(r\), the terms \(D_i(\operatorname{Diff}_{r-i}^+M)\) form a complex ending in \(\operatorname{Diff}_r^+M\to M\), with augmentation evaluation at 1. Taking the filtered union of finite-order operator modules gives the stable complex. For arbitrary commutative rings these complexes need not be exact; extra regularity and category assumptions are required to call them resolutions.}
\GSfield{Fate}{\citet[\S3]{vinogradov1972logic} explicitly states that its finite complex is exact in the last two terms, that the order-one complex is exact, and that higher complexes need not be exact. The smooth antecedent is \citet{spencer1969overdetermined}; the new feature is the functorial algebraic framework.}
\GSfield{Present status}{As of November 2026, the complexes and their possible nonzero homology are established. The notebook’s universal word \emph{resolution} must be qualified; its later word \emph{complex} is the valid unrestricted assertion.}
\GSfield{Assessment}{literal statement: false as written; publication match: exact; verification: read the result; development: annotated; confidence: high}
\end{GSitem}
\begin{GSitem}{1971-241.10}{F}{An index over an arbitrary ring is only a programme heading}
\GSfield{Kind; source}{problem; \GSpage{P000383}}
\GSfield{The notes}{The concluding headings include \emph{the index over any ring}. No index group, admissible class of operators, or finiteness condition is given.}
\GSfield{Statement}{A mathematical index theorem here would need a defined class of operators, a target group for the index, hypotheses making the relevant kernel/cokernel or homological classes meaningful, and a comparison with a symbol construction. The notebook does not specify these data, so it does not pose one unique general theorem.}
\GSfield{Fate}{\citet[\S4 and concluding remarks of \S7]{vinogradov1972logic} relates Spencer homology to kernels and cokernels under acyclicity assumptions and announces a later algebraic index generalization. That is a concrete partial programme, but not a verified all-rings index theorem matching the bare heading.}
\GSfield{Present status}{As of November 2026, the occurrence remains insufficiently specified for a separate problem-index entry or a claim of current openness. No blanket general index theorem is inferred from later uses of the same terminology.}
\GSfield{Assessment}{literal statement: not determinable; publication match: thematic only; verification: read the result; development: annotated; confidence: medium}
\end{GSitem}
\begin{GSitem}{1971-241.11}{A}{Hamiltonian formalism over a commutative algebra}
\GSfield{Kind; source}{method; \GSpage{P000383}}
\GSfield{The notes}{The final list also promises a Hamiltonian formalism \emph{for every ring}.}
\GSfield{Statement}{For a commutative algebra \(K/R\), the order filtration on scalar differential operators has commutative associated graded algebra. The commutator induces a degree-minus-one bracket \(\{\sigma_p(A),\sigma_q(B)\}=\sigma_{p+q-1}([A,B])\). In the smooth real case this identifies with the usual Poisson bracket on fibre-polynomial functions on \(T^*X\), up to the chosen symbol sign convention. The corresponding statement for an arbitrary noncommutative ring is not supplied by these formulas.}
\GSfield{Reconstruction}{The associated-graded bracket is the precise published reconstruction of an otherwise terse programme heading.}
\GSfield{Fate}{\citet[\S7]{vinogradov1972logic} publishes precisely the symbol and Hamiltonian programme for commutative rings. This gives a close mathematical match to the heading once its ambient category is corrected.}
\GSfield{Present status}{As of November 2026, the commutative symbol bracket is established. It does not provide one unrestricted Hamiltonian calculus on all noncommutative rings.}
\GSfield{Assessment}{literal statement: correct with missing hypotheses; publication match: exact; verification: read the result; development: annotated; confidence: high}
\end{GSitem}

\GStalk{242}{1971-11-08}{A. M. Vinogradov}{P000384-P000386}{Six lemmas on differential operators and cojets}
The third lecture proves the elementary identities behind the differential-operator functors: pre- and post-composition, the first-order product identity, addition of orders under composition, and the universal evaluation/currying adjunction. These are the technical foundations of the 1972 announcement, with a commutative base ring understood.

\begin{GSitem}{1971-242.1}{A}{Lemma 1: functoriality in the source and target}
\GSfield{Kind; source}{theorem; \GSpage{P000384}}
\GSfield{The notes}{With fixed \(N\), \(\operatorname{Diff}_n^+(M,N)\) is contravariant in \(M\); with fixed \(M\), the pointwise module \(\operatorname{Diff}_n(M,N)\) is covariant in \(N\). One displayed composition is typed \(D\circ f:M_1\to M_2\).}
\GSfield{Statement}{For commutative unital \(K/R\), a \(K\)-linear map \(f:M_1\to M_2\) gives \(D\mapsto D\circ f\) from \(\operatorname{Diff}_n^+(M_2,N)\) to \(\operatorname{Diff}_n^+(M_1,N)\). A \(K\)-linear map \(g:N_1\to N_2\) gives \(D\mapsto g\circ D\) on the pointwise operator modules. Both maps preserve order and the stated scalar action.}
\GSfield{Fate}{These elementary functorial properties underpin \citet[\S1]{vinogradov1972logic}. Class A concerns the organized cojet presentation whose proof the November lecture supplies.}
\GSfield{Present status}{As of November 2026, the functoriality is established. The target in the displayed first composition must be \(N\), a type error rather than a change to the theorem.}
\GSfield{Assessment}{literal statement: false as written; publication match: exact; verification: read the result; development: annotated; confidence: high}
\end{GSitem}
\begin{GSitem}{1971-242.2}{E}{Lemma 2: first-order operators split into a derivation and multiplication}
\GSfield{Kind; source}{theorem; \GSpage{P000384}}
\GSfield{The notes}{The notes prove \(A(ab)=aA(b)+bA(a)-abA(1)\) for \(A\in\operatorname{Diff}_1(K,M)\), its converse, and the corollary that \(A-A(1)\) is a derivation.}
\GSfield{Statement}{An \(R\)-linear map \(A:K\to M\) has order at most one exactly when this product formula holds for every \(a,b\in K\). Equivalently, \(A(a)=aA(1)+D(a)\) for a unique \(R\)-derivation \(D\). The subtraction in the corollary means the multiplication map \(a\mapsto aA(1)\), not a constant map.}
\GSfield{Fate}{This is an elementary consequence of the pre-1971 commutator calculus \citep[\S16.8]{grothendieck1967egaiv4}; the derivation functor is part of \citet[\S2]{vinogradov1972logic}. It is expository rather than a new first-order decomposition.}
\GSfield{Present status}{As of November 2026, the equivalence holds over any commutative unital base, without dividing by 2 and without smoothness. The corrected reading of the subtraction is necessary for additivity.}
\GSfield{Assessment}{literal statement: correct with missing hypotheses; publication match: exact; verification: read the result; development: annotated; confidence: high}
\end{GSitem}
\begin{GSitem}{1971-242.3}{E}{Lemma 3: orders add under composition}
\GSfield{Kind; source}{theorem; \GSpage{P000384}}
\GSfield{The notes}{For \(A\in\operatorname{Diff}_m(M_1,M_2)\) and \(B\in\operatorname{Diff}_n(M_2,M_3)\), the notes prove \(BA\in\operatorname{Diff}_{m+n}\) by induction and a commutator product rule. Here \(\delta_k\) has the opposite sign from P376.}
\GSfield{Statement}{For commutative \(K/R\), the composition of operators of orders at most \(m,n\) has order at most \(m+n\). With \(\delta_aT=T\circ a-a\circ T\), one has \(\delta_a(BA)=(\delta_aB)A+B(\delta_aA)\). The opposite convention gives the same order condition and the same formal product rule.}
\GSfield{Fate}{The filtered composition law is part of the earlier algebraic operator calculus \citep[\S16.8]{grothendieck1967egaiv4}. The lecture supplies its elementary inductive verification; \citet[\S3]{vinogradov1972logic} uses it to define Spencer differentials.}
\GSfield{Present status}{As of November 2026, the order bound is established. Equality need not hold: principal symbols can vanish on composition. The local sign switch does not invalidate the argument.}
\GSfield{Assessment}{literal statement: correct; publication match: exact; verification: read the result; development: annotated; confidence: high}
\end{GSitem}
\begin{GSitem}{1971-242.4}{A}{Lemma 4: evaluation on cojets is a differential operator}
\GSfield{Kind; source}{theorem; \GSpage{P000385}}
\GSfield{The notes}{The canonical map \(C:\operatorname{Diff}_n^+(K,N)\to N\), \(C(A)=A(1)\), is called an operator of order \(n\). The proof transfers its iterated commutators to those of \(A\).}
\GSfield{Statement}{With the precomposition \(K\)-module structure on the source, \(C\) has order at most \(n\). Composition with a \(K\)-linear \(f:M\to\operatorname{Diff}_n^+(K,N)\) therefore defines an order-at-most-\(n\) operator \(C\circ f:M\to N\), and this assignment is \(K\)-linear when the resulting operator module uses precomposition.}
\GSfield{Fate}{This evaluation map is the universal cojet operator in \citet[\S1]{vinogradov1972logic}. The proof given on the scan is a correct direct verification, using its chosen commutator sign.}
\GSfield{Present status}{As of November 2026, the order bound is established. Exact order \(n\) is not asserted: it can drop for particular rings or modules.}
\GSfield{Assessment}{literal statement: correct with missing hypotheses; publication match: exact; verification: read the result; development: annotated; confidence: high}
\end{GSitem}
\begin{GSitem}{1971-242.5}{A}{Lemma 5: currying an operator into a cojet-valued homomorphism}
\GSfield{Kind; source}{theorem; \GSpage{P000385}}
\GSfield{The notes}{The map sends \(A\) to \(g_A(m)(k)=A(km)\); the notes check order, additivity, and scalar linearity. The first check is called restriction to a submodule.}
\GSfield{Statement}{For \(A\in\operatorname{Diff}_n^+(M,N)\), the map \(k\mapsto A(km)\) has order at most \(n\) because it is the composite of \(A\) with the \(K\)-linear map \(K\to M\), \(k\mapsto km\). The assignment \(m\mapsto g_A(m)\) is \(K\)-linear into the precomposition cojet module, and \(A\mapsto g_A\) is a homomorphism between the indicated operator and Hom modules. The map \(K\to M\) need not be injective.}
\GSfield{Fate}{This is the explicit forward direction of the universal cojet isomorphism published in \citet[\S1]{vinogradov1972logic}. The scan contains the three checks rather than merely announcing the universal property.}
\GSfield{Present status}{As of November 2026, the construction is established for arbitrary unital modules. The word restriction should be replaced by precomposition with a module map; no torsion-free hypothesis on \(M\) is necessary.}
\GSfield{Assessment}{literal statement: correct with missing hypotheses; publication match: exact; verification: read the result; development: annotated; confidence: high}
\end{GSitem}
\begin{GSitem}{1971-242.6}{A}{Lemma 6: the two cojet maps are mutually inverse}
\GSfield{Kind; source}{theorem; \GSpage{P000385-P000386}}
\GSfield{The notes}{The lecture finishes by computing both composites and concludes \(\operatorname{Diff}_n^+N=\operatorname{Diff}_n^+(K,N)\) as the concrete universal cojet object.}
\GSfield{Statement}{For \(A:M\to N\), evaluation of its curried map gives \(g_A(m)(1)=A(m)\). Conversely, if \(g:M\to\operatorname{Diff}_n^+(K,N)\) is \(K\)-linear, then the curried form of \(C\circ g\) satisfies \((C\circ g)(km)=g(km)(1)=(k\cdot g(m))(1)=g(m)(k)\). Thus the two natural module maps are inverse.}
\GSfield{Fate}{The resulting isomorphism is exactly that stated in \citet[\S1]{vinogradov1972logic}. The November 1971 notebook provides an earlier detailed seminar proof; the paper was received in January 1972.}
\GSfield{Present status}{As of November 2026, the proof is complete under the commutative-ring and fixed-base conventions already stated. The last equality is a realization of a universal object, not an assertion that the two scalar actions coincide.}
\GSfield{Assessment}{literal statement: correct; publication match: exact; verification: read the result; development: annotated; confidence: high}
\end{GSitem}

\GStalk{243}{1971-11-15}{I. M. Gelfand}{P000387-P000394}{Formal coordinate geometry and characteristic classes of foliations}
Gelfand develops the infinite frame bundle, its formal vector-field coframing, relative Lie algebra cohomology, and characteristic classes; Bernstein intervenes on inverse-limit manifolds. The final application is Bott’s obstruction to integrability, already discussed in the preceding year. The printed truncation threshold and the homogeneous degrees of Chern generators need correction. Arnold’s new heading begins later on P000394 and is excluded from this record.

\begin{GSitem}{1971-243.1}{E}{Infinitesimal homogeneous spaces and infinitesimal torsors}
\GSfield{Kind; source}{definition; \GSpage{P000387}}
\GSfield{The notes}{A Lie algebra acts by vector fields; the space is homogeneous when evaluation at every point is onto, and principal when evaluation is an isomorphism. The preceding group definition requires a unique group element carrying one point to another.}
\GSfield{Statement}{Let \(a:\mathfrak g\to\mathfrak X(X)\) be a Lie algebra homomorphism. Transitivity means that \(a_x:\mathfrak g\to T_xX\) is surjective for each \(x\). An infinitesimal torsor has each \(a_x\) invertible and the inverse coframing smooth in the chosen topological category. This definition is infinitesimal: a global integrated action can still have discrete stabilizers.}
\GSfield{Fate}{The principal-space language was already published by \citet[\S3]{gelfandkazhdan1971geometry}, received 5 June 1971. Its precise topological formulation is given in \citet[Definition 2.1]{bernsteinrosenfeld1973homogeneous}. The 15 November lecture explicitly refers to the earlier announcement.}
\GSfield{Present status}{As of November 2026, this is an established definition in formal geometry. The terms principal Lie-algebra space and principal group space should not be interchanged without a globalization and stabilizer condition.}
\GSfield{Assessment}{literal statement: correct with missing hypotheses; publication match: exact; verification: read the result; development: annotated; confidence: high}
\end{GSitem}
\begin{GSitem}{1971-243.2}{E}{Finite-dimensional Lie algebras integrate locally, not automatically globally}
\GSfield{Kind; source}{claim; \GSpage{P000387}}
\GSfield{The notes}{The underlined answer to when one can pass to a group is: when the Lie algebra is finite-dimensional. No completeness hypothesis is stated.}
\GSfield{Statement}{A smooth action of a finite-dimensional real Lie algebra on a finite-dimensional manifold integrates locally. A global action with these infinitesimal generators requires their completeness; finite-dimensionality by itself is insufficient. Even a complete infinitesimal torsor need not be a free torsor for the simply connected integrating group.}
\GSfield{Reconstruction}{\ours\ For the global claim, take \(X=(0,1)\) and the action of \(\mathbb R\) generated by \(\partial_x\). Evaluation is an isomorphism everywhere, but its integral curve \(x+t\) leaves \(X\) in finite time. For the freeness issue, translations by \(\mathbb R\) on \(\mathbb R/\mathbb Z\) have isomorphic evaluation maps and stabilizer \(\mathbb Z\).}
\GSfield{Fate}{This is a qualification of classical Lie integration, explicitly stated as local integration in \citet[Introduction, pp.~108--109]{bernsteinrosenfeld1973homogeneous}. It is not a new globalization theorem in the notebook.}
\GSfield{Present status}{As of November 2026, the local reading is correct; the unrestricted global reading is false. No classification of globalizations is inferred from the shorthand.}
\GSfield{Assessment}{literal statement: correct with missing hypotheses; publication match: partial; verification: read the result; development: annotated; confidence: high}
\end{GSitem}
\begin{GSitem}{1971-243.3}{E}{Formal vector fields and the origin-fixing formal coordinate group}
\GSfield{Kind; source}{definition; \GSpage{P000387-P000388}}
\GSfield{The notes}{\(W_n\) consists of formal vector fields. The notes say that its group does not exist because one cannot move; the subalgebra \(L_0\) of fields with zero constant term does integrate and has codimension \(n\). Hamiltonian and volume-preserving variants are mentioned.}
\GSfield{Statement}{For \(n\geq1\), let \(W_n=\operatorname{Der}^{\mathrm{cont}}_{\mathbb R}\mathbb R[[x_1,\ldots,x_n]]\) with its coefficient topology and continuous bracket. The subalgebra \(L_0=\{\sum a_i\partial_i:a_i(0)=0\}\) has vector-space codimension \(n\). The group \(G_n^0\) of invertible formal substitutions fixing the origin is the inverse limit of its finite-jet groups and has Lie algebra \(L_0\) in that category. A general formal series cannot be translated by a nonzero real constant by an unrestricted numerical substitution.}
\GSfield{Fate}{These objects are used in \citet[\S1]{gelfandfuchs1970formal} and \citet[\S3]{gelfandkazhdan1971geometry}; the formal-group and transitive-subalgebra conventions are made explicit in \citet[\S\S4--5]{bernsteinrosenfeld1973homogeneous}. They precede the lecture.}
\GSfield{Present status}{As of November 2026, the distinction between full formal vector fields and the origin-fixing formal group remains essential. The phrase that no group exists is category dependent; it is not a theorem excluding formal groupoids or every infinite-dimensional integration concept.}
\GSfield{Assessment}{literal statement: correct with missing hypotheses; publication match: exact; verification: read the result; development: annotated; confidence: high}
\end{GSitem}
\begin{GSitem}{1971-243.4}{E}{Automorphisms of the Lie algebra of formal vector fields}
\GSfield{Kind; source}{theorem; \GSpage{P000388}}
\GSfield{The notes}{The notebook states that the Lie algebra of the automorphism group of \(W_n\) is \(L_0\). It compares this with quotienting an ordinary Lie algebra by its centre.}
\GSfield{Statement}{Over a characteristic-zero field, every continuous Lie algebra automorphism of \(W_n\) is induced by an invertible origin-fixing formal coordinate substitution. Thus, over \(\mathbb R\), the automorphism group with the formal-coordinate, finite-jet group structure has Lie algebra \(L_0\). This is a statement about that topological group, not an unrestricted identification with all algebraic derivations of \(W_n\).}
\GSfield{Fate}{\citet[Theorem 1 and \S3, Theorem 2(1)]{rudakov1969automorphisms} proved preservation of the formal filtration and the coordinate-change description; the paper was received 18 December 1968 and published in Russian in 1969. Theorem 2 also treats the divergence-free, Hamiltonian and contact series. This is prior work, not an unpublished 1971 conjecture.}
\GSfield{Present status}{As of November 2026, the continuous formal-coordinate classification is established. The notebook comparison with quotient by the centre describes inner automorphisms in the classical case; it does not identify the entire automorphism algebra for arbitrary finite-dimensional Lie algebras.}
\GSfield{Assessment}{literal statement: correct with missing hypotheses; publication match: exact; verification: read the result; development: annotated; confidence: high}
\end{GSitem}
\begin{GSitem}{1971-243.5}{E}{The formal frame bundle as a principal Wn-space}
\GSfield{Kind; source}{construction; \GSpage{P000388}}
\GSfield{The notes}{For an arbitrary \(n\)-manifold \(M\), the space \(S(M)\) of jets of local coordinate systems is declared a principal \(W_n\)-space. Equivalence of charts is written as all derivatives of their transition map being zero.}
\GSfield{Statement}{For a smooth \(n\)-manifold \(M\), let \(S(M)\) consist of infinite jets at \(0\) of local diffeomorphisms \((\mathbb R^n,0)\to(M,x)\), for all \(x\). It has a canonical smooth \(W_n\)-valued coframing satisfying the Maurer--Cartan equation, equivalently a principal infinitesimal \(W_n\)-space structure. Two charts define the same jet when their coordinate difference is flat, or their transition has the identity infinite jet. The transition itself cannot have all derivatives zero, since its first derivative is invertible.}
\GSfield{Fate}{\citet[\S3 and \S4, Proposition 1]{gelfandkazhdan1971geometry} announced this construction before the lecture; \citet[Theorem 4.1]{bernsteinrosenfeld1973homogeneous} give a detailed construction and proof in the finite-coordinate category.}
\GSfield{Present status}{As of November 2026, the formal frame construction is established. The chart-equivalence sentence is a mathematical slip faithfully reproduced by the transcription; correcting it does not require guessing any cancelled words.}
\GSfield{Assessment}{literal statement: false as written; publication match: exact; verification: read the result; development: annotated; confidence: high}
\end{GSitem}
\begin{GSitem}{1971-243.6}{F}{Other substantive principal Wn-spaces}
\GSfield{Kind; source}{question; \GSpage{P000388}}
\GSfield{The notes}{The final sentence asks whether there are other substantive principal homogeneous spaces, after constructing \(S(M)\). It gives no meaning for substantive or for equivalence of examples.}
\GSfield{Statement}{The identifiable topic is the possible existence or classification of principal infinitesimal \(W_n\)-spaces beyond formal frame spaces. A definite classification problem requires specifying the category, local or global equivalence, permitted quotients and what counts as arising from a manifold. These data are absent, so no unique stronger proposition is reconstructed.}
\GSfield{Fate}{The later theory of \citet[\S\S2, 4--7]{bernsteinrosenfeld1973homogeneous} treats homogeneous spaces and geometric structures systematically. Searches for principal \(W_n\)-spaces and classification did not identify a theorem matching this qualitative question.}
\GSfield{Present status}{As of November 2026, the particular intended question remains unidentified. This is a limit of the notebook formulation and the checked literature, not evidence that a precise classification problem is open.}
\GSfield{Assessment}{literal statement: not determinable; publication match: thematic only; verification: unresolved lead; development: annotated; confidence: low}
\end{GSitem}
\begin{GSitem}{1971-243.7}{E}{The characteristic chain map from Lie algebra cochains}
\GSfield{Kind; source}{construction; \GSpage{P000389-P000390}}
\GSfield{The notes}{A skew multilinear form on \(W_n\) is evaluated on the vectors corresponding to tangent vectors of \(S(M)\). The resulting map is said to commute with the differentials. The printed exponent \(s+t=1\) is also visible in the scan.}
\GSfield{Statement}{For a principal infinitesimal \(\mathfrak g\)-space with Maurer--Cartan coframing \(\theta\), a continuous Chevalley--Eilenberg cochain \(c\) gives the differential form \(\Phi(c)(v_1,\ldots,v_q)=c(\theta v_1,\ldots,\theta v_q)\). With matching bracket and Maurer--Cartan signs, \(d\Phi(c)=\Phi(d_{\mathrm{CE}}c)\), giving \(H^*_{\mathrm{cont}}(\mathfrak g)\to H^*_{\mathrm{dR}}(S)\). For \(W_n\), continuous scalar cochains depend on finitely many jet coordinates.}
\GSfield{Reconstruction}{Using the convention \(d\theta+\tfrac12[\theta,\theta]=0\) and numbering arguments from 1, the bracket term is \(\sum_{s<t}(-1)^{s+t}c([\xi_s,\xi_t],\xi_1,\ldots,\widehat\xi_s,\ldots,\widehat\xi_t,\ldots,\xi_{q+1})\). Opposite action conventions change the displayed sign consistently. No literal equality sign belongs in this exponent.}
\GSfield{Fate}{This is the cohomological construction in \citet[\S4]{gelfandkazhdan1971geometry}, whose note is explicitly cited on the notebook page. A complete statement and proof are in \citet[\S3, Theorem 3.1]{bernsteinrosenfeld1973homogeneous}.}
\GSfield{Present status}{As of November 2026, the characteristic map is established. The malformed exponent must be emended together with a consistent left/right convention; its exact intended sign cannot be settled from the two horizontal strokes alone.}
\GSfield{Assessment}{literal statement: false as written; publication match: exact; verification: read the result; development: annotated; confidence: high}
\end{GSitem}
\begin{GSitem}{1971-243.8}{A}{Bernstein’s finite-coordinate differential calculus on jet towers}
\GSfield{Kind; source}{definition; \GSpage{P000389-P000390}; I. N. Bernstein (intervention)}
\GSfield{The notes}{After the heading I. Bernstein, an infinite-dimensional manifold is defined as an inverse limit of finite-dimensional manifolds with onto maps; functions are the union of the finite-stage function rings. The de Rham differential is then written by its vector-field formula.}
\GSfield{Statement}{For the jet towers in question one may use inverse systems whose sufficiently high stages are compatible affine bundles over a fixed finite stage. Their inverse limits admit charts modelled on \(\mathbb R^{\mathbb N}\). Scalar functions and forms are locally pulled back from a finite stage; vector fields and the exterior differential are defined compatibly. The local finite-stage condition is more general than requiring a single global stage for every smooth function. Arbitrary surjective inverse systems alone do not guarantee this manifold structure.}
\GSfield{Fate}{Bernstein’s intervention is developed with Rosenfeld in \citet[\S1, Definitions 1.1--1.10, Proposition 1.1 and Example 1.1]{bernsteinrosenfeld1973homogeneous}, received 29 March 1973 and published in 1973. That paper explicitly distinguishes its smooth finite-coordinate category from an arbitrary inverse limit.}
\GSfield{Present status}{As of November 2026, the construction is established for the formal jet spaces used here. The omitted local and bundle hypotheses are material; the statement should not be exported as a definition of all infinite-dimensional smooth manifolds.}
\GSfield{Assessment}{literal statement: correct with missing hypotheses; publication match: exact; verification: read the result; development: annotated; confidence: high}
\end{GSitem}
\begin{GSitem}{1971-243.9}{E}{The formal Lie algebra is an inverse limit as a vector space}
\GSfield{Kind; source}{claim; \GSpage{P000390}; I. N. Bernstein (intervention)}
\GSfield{The notes}{The Lie algebra is described parenthetically as a projective limit of finite-dimensional ones with the corresponding topology. Taken literally as finite-dimensional Lie algebra quotients, this excludes the main example.}
\GSfield{Statement}{In the formal-vector-field setting the underlying topological vector space is an inverse limit of finite-dimensional jet spaces, and the Lie bracket is continuous. The finite coefficient truncations of \(W_n\) are not Lie algebra quotients: order filtrations are not ideals because a constant vector field lowers vanishing order. No projective system of these truncated Lie algebras is required.}
\GSfield{Reconstruction}{\ours\ In one variable, \([\partial_x,x^{k+1}\partial_x]=(k+1)x^k\partial_x\). Thus the fields vanishing to order at least \(k+1\) are not stable under bracketing with all of \(W_1\), and the proposed quotient bracket is not well defined.}
\GSfield{Fate}{\citet[\S0.1]{gelfandfuchs1970formal} use the topological vector-space formulation, and \citet[\S2, preceding Definition 2.1]{bernsteinrosenfeld1973homogeneous} specify that \(\mathfrak g\) is modelled on \(\mathbb R^k\) or \(\mathbb R^{\mathbb N}\) as a vector space.}
\GSfield{Present status}{As of November 2026, the continuous Lie algebra formulation is standard. The literal finite-dimensional Lie-quotient reading is an error; finite jet vector spaces suffice for the subsequent continuous-cochain construction.}
\GSfield{Assessment}{literal statement: false as written; publication match: exact; verification: read the result; development: annotated; confidence: high}
\end{GSitem}
\begin{GSitem}{1971-243.10}{A}{Relative cochains and descent of basic forms}
\GSfield{Kind; source}{construction; \GSpage{P000390-P000392}; I. N. Bernstein (intervention)}
\GSfield{The notes}{The relative complex consists of cochains \(c\) for which both \(c\) and \(dc\) vanish when an argument lies in \(\mathfrak p\). A form on the base is obtained by lifting tangent vectors; independence of the point in the fibre is asserted.}
\GSfield{Statement}{The relative complex is \(C^*(\mathfrak g,\mathfrak p)=\{c:i_Xc=i_Xdc=0\text{ for every }X\in\mathfrak p\}\), equivalently horizontal and \(\mathfrak p\)-invariant cochains. For a compatible bundle \(P\to X\) of principal infinitesimal spaces whose fibres are connected principal \(\mathfrak p\)-spaces, the characteristic forms descend and give \(C^*(\mathfrak g,\mathfrak p)\to\Omega^*(X)\). With a disconnected structure group \(H\), one must require invariance under all of \(H\), as well as horizontality for its Lie algebra.}
\GSfield{Fate}{The relative-complex definition is classical. The homogeneous-space descent in the intervention is formulated precisely by \citet[\S3, Proposition 3.2, Theorem 3.2 and Remark 3.2]{bernsteinrosenfeld1973homogeneous}; Proposition 3.2 explicitly assumes connected fibres.}
\GSfield{Present status}{As of November 2026, the construction is established with those hypotheses. Horizontality gives independence of lifts, while infinitesimal invariance gives independence only within a connected fibre component; these are different steps.}
\GSfield{Assessment}{literal statement: correct with missing hypotheses; publication match: exact; verification: read the result; development: annotated; confidence: high}
\end{GSitem}
\begin{GSitem}{1971-243.11}{E}{Relative Lie algebra cohomology need not equal the cohomology of a homogeneous space}
\GSfield{Kind; source}{claim; \GSpage{P000391}; I. N. Bernstein (intervention)}
\GSfield{The notes}{The line after the relative-complex definition identifies its cohomology with that of the quotient in the classical case, without compactness or other comparison hypotheses.}
\GSfield{Statement}{For a Lie group \(G\) and closed subgroup \(H\), relative cochains describe suitable \(G\)-invariant differential forms on \(G/H\), including invariance under disconnected isotropy components when necessary. Their inclusion in all differential forms is not generally a quasi-isomorphism. The characteristic construction gives a homomorphism on cohomology; equality requires an additional comparison theorem.}
\GSfield{Reconstruction}{\ours\ Let \(G=(\mathbb R,+)\) and \(H=\{0\}\). The abelian one-dimensional Lie algebra has \(H^1(\mathfrak g;\mathbb R)=\mathbb R\), represented by its nonzero linear functional. Its corresponding invariant form on \(G\) is \(dx\), but \(dx=d(x)\), so \(H^1_{\mathrm{dR}}(\mathbb R)=0\).}
\GSfield{Fate}{\citet[Theorem 3.2 and Remark 3.2]{bernsteinrosenfeld1973homogeneous} assert the homomorphism, not a general isomorphism. The notebook’s unrestricted equality is therefore not part of the verified published result.}
\GSfield{Present status}{As of November 2026, the unrestricted equality is false, already in dimension one. This is an elementary correction, not an unresolved comparison problem.}
\GSfield{Assessment}{literal statement: false as written; publication match: partial; verification: read the result; development: annotated; confidence: high}
\end{GSitem}
\begin{GSitem}{1971-243.12}{E}{Removing linear frames leaves a space homotopy equivalent to the manifold}
\GSfield{Kind; source}{theorem; \GSpage{P000392}; I. N. Bernstein (intervention)}
\GSfield{The notes}{\(S(M)\) is said to be homotopy equivalent to the first-frame bundle, while \(S(M)/GL(n,\mathbb R)\) is homotopy equivalent to \(M\). A marginal line mentions a map \(M\to S(M)\); later, choices are associated with a connection.}
\GSfield{Statement}{The projection from formal frames to first frames has contractible higher-jet fibres. Similarly, the bundle \(S(M)/GL(n,\mathbb R)\to M\) has contractible fibre, and is a homotopy equivalence for a smooth paracompact manifold. A connection provides a formal exponential section of this quotient. A section of the full frame space would also choose a frame of \(TM\) and generally does not exist.}
\GSfield{Fate}{The formal-coordinate construction is in \citet[\S3]{gelfandkazhdan1971geometry}. The contractible higher-jet-fibre argument is made explicit in \citet[\S8, Theorems 8.3--8.4]{bernsteinrosenfeld1973homogeneous}.}
\GSfield{Present status}{As of November 2026, the two homotopy statements are established. The unclear marginal word is left unresolved; if the marginal map means a section, its target must be the quotient unless a global framing has been assumed.}
\GSfield{Assessment}{literal statement: correct with missing hypotheses; publication match: exact; verification: read the result; development: annotated; confidence: high}
\end{GSitem}
\begin{GSitem}{1971-243.13}{E}{The truncated characteristic polynomial ring in relative Wn-cohomology}
\GSfield{Kind; source}{theorem; \GSpage{P000392,P000394}; I. N. Bernstein (intervention)}
\GSfield{The notes}{The relative cohomology is described by generators \(\psi_1,\ldots,\psi_n\), written as having degree \(2n\), with products of degree greater than \(2n\) zero. The later page cites Gelfand and Fuks, Izvestiya, 1971.}
\GSfield{Statement}{With real coefficients and continuous cochains, the group-relative characteristic complex gives \[H^*(W_n,GL(n,\mathbb R))\simeq\mathbb R[c_1,\ldots,c_n]/I_{>2n},\qquad |c_i|=2i.\] Here \(I_{>2n}\) is the ideal generated by monomials of cohomological degree greater than \(2n\). Group-relative means horizontality for \(\mathfrak{gl}_n\) and invariance under the full linear group. The polynomial generators have different degrees; the quotient is essential.}
\GSfield{Fate}{The underlying formal-vector-field computation was published by \citet{gelfandfuchs1970formal}, received 4 November 1969 and published in Russian in 1970. Its English printing carries 1971 copyright, explaining the notebook date. The exact relative statement is \citet[Theorem 10.2]{bernsteinrosenfeld1973homogeneous}; see also \citet[\S2, pp.~125--126]{bott1976characteristic}.}
\GSfield{Present status}{As of November 2026, the truncated relative ring is established. It is not the untruncated cohomology of \(BGL_n\), and replacing every generator degree by \(2n\) loses the result.}
\GSfield{Assessment}{literal statement: false as written; publication match: exact; verification: read the result; development: annotated; confidence: high}
\end{GSitem}
\begin{GSitem}{1971-243.14}{E}{Chern–Weil classes from formal frames and the role of a connection}
\GSfield{Kind; source}{construction; \GSpage{P000392}; I. N. Bernstein (intervention)}
\GSfield{The notes}{The images of the generators are linked to Chern classes of the tangent bundle and to Pontryagin classes. Forms are canonical upstairs; nonuniqueness appears when a connection is chosen to obtain forms on \(M\).}
\GSfield{Statement}{Normalize the formal curvature generators so that their characteristic images are \(c_i(TM\otimes_{\mathbb R}\mathbb C)\) in real cohomology. Odd Chern classes of a complexified real bundle vanish over \(\mathbb R\), and \(p_i(TM)=(-1)^i c_{2i}(TM\otimes\mathbb C)\) in the standard convention. Pulling a fixed universal cocycle back along a formal exponential section gives a differential-form representative; changing the connection changes the representative by an exact form, not its cohomology class.}
\GSfield{Fate}{This is the classical Chern--Weil construction expressed through \citet{gelfandkazhdan1971geometry}. The normal-bundle version is identified in \citet[Theorem 10.5]{bernsteinrosenfeld1973homogeneous}, whose normalization sends its generator \(c_{2i}\) directly to \(p_i\). \citet[\S4]{bott1976characteristic} verifies the corresponding characteristic map.}
\GSfield{Present status}{As of November 2026, the characteristic classes are established and independent of the auxiliary connection. The notebook’s Chern/Pontryagin wording should be read with complexification and normalization stated, not as two unrelated sets of real tangent-bundle classes.}
\GSfield{Assessment}{literal statement: correct with missing hypotheses; publication match: exact; verification: read the result; development: annotated; confidence: high}
\end{GSitem}
\begin{GSitem}{1971-243.15}{A}{Formal transverse frames of a foliation}
\GSfield{Kind; source}{construction; \GSpage{P000393-P000394}; I. N. Bernstein (intervention)}
\GSfield{The notes}{For a codimension-\(n\) foliation, transverse formal coordinates define \(S_{\mathcal F}(M)\); the notes stress that this is not itself supplied with the claimed natural principal Lie-algebra action. Adapted full coordinates define \(S^{\prime}_{\mathcal F}(M)\), a principal \(W_{N,n}\)-space, with a leaf-tangent subalgebra \(L\) and quotient \(W_n\).}
\GSfield{Statement}{Let \(\mathcal F\) be a regular smooth foliation of codimension \(n\) on an \(N\)-manifold. In formal coordinates \((y,z)\), with \(y\in\mathbb R^n\) transverse, \(W_{N,n}\) consists of fields whose \(y\)-components depend only on \(y\). The ideal \(L\) of fields tangent to the standard leaves has quotient \(W_{N,n}/L\simeq W_n\). Adapted formal frames give a homogeneous pair with isotropy \(L\), hence \(C^*(W_n)=C^*(W_{N,n},L)\to\Omega^*(S_{\mathcal F}(M))\). Passing to full linear-group basic cochains yields the normal-bundle characteristic map on \(M\), since \(S_{\mathcal F}(M)/GL_n\to M\) has contractible higher-jet fibres.}
\GSfield{Fate}{The characteristic-class programme was announced by Bernstein and Rosenfeld in 1972, as recorded in reference [3] and the introduction of their detailed \citet[\S\S6, 8--10]{bernsteinrosenfeld1973homogeneous}. That full paper, received 29 March 1973, verifies the homogeneous-space and characteristic-map mechanism. The present notebook records Bernstein’s 1971 intervention.}
\GSfield{Present status}{As of November 2026, this construction is established for regular smooth foliations. The transverse quotient and the adapted full-frame principal space have different roles; merely postulating a \(W_n\)-action on the former would omit the actual descent argument.}
\GSfield{Assessment}{literal statement: correct with missing hypotheses; publication match: exact; verification: read the result; development: annotated; confidence: high}
\end{GSitem}
\begin{GSitem}{1971-243.16}{E}{Bott vanishing for the real normal bundle}
\GSfield{Kind; source}{theorem; \GSpage{P000392-P000394}; I. N. Bernstein (intervention)}
\GSfield{The notes}{The explicit formula kills products of the \(\psi_i\) of degree \(>2n\). The following sentence appears to change this to \(\geq2n\), although the additional stroke is uncertain; the result is named Bott’s theorem.}
\GSfield{Statement}{For a regular \(C^\infty\) foliation of real codimension \(n\) with normal bundle \(\nu=TM/T\mathcal F\), every monomial in the real Pontryagin classes of \(\nu\) having total cohomological degree strictly greater than \(2n\) vanishes. Equivalently, the corresponding real Chern classes of \(\nu\otimes\mathbb C\) satisfy the truncated polynomial relations. This gives no blanket vanishing in degree exactly \(2n\), and does not assert vanishing of integral torsion classes.}
\GSfield{Fate}{This is the prior obstruction of \citet{bott1970integrability}, from the Berkeley symposium of July 1968, printed in 1970. The formal-cohomology deduction is explicit in \citet[Theorem 10.2, Theorem 10.5 and Corollary 10.1]{bernsteinrosenfeld1973homogeneous} and is revisited by \citet[\S4]{bott1976characteristic}.}
\GSfield{Present status}{As of November 2026, the strict-degree real-cohomology theorem is established. The scan does not justify silently adopting the stronger boundary inequality; the preceding clear formula supplies the mathematically correct bound.}
\GSfield{Assessment}{literal statement: correct with missing hypotheses; publication match: exact; verification: read the result; development: annotated; confidence: high}
\end{GSitem}
\begin{GSitem}{1971-243.17}{E}{Characteristic obstruction to homotoping a distribution to a foliation}
\GSfield{Kind; source}{problem; \GSpage{P000394}; I. N. Bernstein (intervention)}
\GSfield{The notes}{The last sentence before Arnold’s new heading calls the vanishing relation an obstruction to being homotopically integrable.}
\GSfield{Statement}{Given a smooth codimension-\(n\) subbundle \(D\subset TM\), ask whether it is homotopic through subbundles of the same rank to the tangent distribution of a regular foliation. A necessary condition is that every real Pontryagin monomial of \(TM/D\) of degree \(>2n\) vanish. The statement is an obstruction, not a sufficient criterion.}
\GSfield{Reconstruction}{\ours\ A homotopy \(D_t\) defines a vector bundle over \(M\times[0,1]\); its quotient bundles at the endpoints are isomorphic. Their Pontryagin classes therefore agree. If \(D_1\) is integrable, Bott vanishing for \(TM/D_1\) gives the stated necessary condition for \(TM/D_0\).}
\GSfield{Fate}{This is the same integrability problem and Bott obstruction already indexed as GP-1970-26, occurrences 1970-206.1 and 1970-206.2. The notebook supplies a new formal-cohomology derivation of the known obstruction \citep{bott1970integrability}; its characteristic interpretation is verified in \citet[\S4]{bott1976characteristic}.}
\GSfield{Present status}{As of November 2026, the necessary condition is established and recurs here. No claim of a complete classification of integrable homotopy classes is made on the basis of the literature checked for this occurrence.}
\GSfield{Assessment}{literal statement: correct with missing hypotheses; publication match: exact; verification: read the result; development: annotated; confidence: high; problem \hyperref[prob:GP-1970-26]{GP-1970-26}}
\end{GSitem}

\GStalk{244}{1971-11-15}{V. I. Arnold}{P000394}{An unidentified characteristic invariant of a complex foliation}
The final paragraph of P000394 is explicitly headed Arnold. It concerns the foliation associated with a rational differential equation on the complex projective plane and asks for a geometric, possibly integral or rational, interpretation of an unnamed invariant. The request to recover rotation numbers suggests characteristic-class methods, but neither a class nor an evaluation cycle is specified. The preceding Bott discussion belongs to record 243. The catalogue span P000393--P000396 therefore substantially overstates the identifiable Arnold material.

\begin{GSitem}{1971-244.1}{F}{Geometric meaning and arithmetic of an unnamed foliation invariant}
\GSfield{Kind; source}{question; \GSpage{P000394}}
\GSfield{The notes}{Arnold writes \(dw/dz=P(z,w)/Q(z,w)\), describes a foliation on \(\mathbb{CP}^{2}\) away from several points, and mentions an invariant on a three-dimensional complex manifold described as ``point--vector''. He asks for its geometric characterization and integrality or rationality.}
\GSfield{Statement}{The equation, after homogenization and removal of its singular locus, determines a holomorphic line distribution on an open subset of \(\mathbb{CP}^{2}\). The notes do not specify the auxiliary threefold, the characteristic cocycle, its degree, a cycle on which to evaluate it, or its normalization. Consequently no definite integrality or rationality assertion can be extracted.}
\GSfield{Reconstruction}{The interpretation as a holomorphic line distribution is standard reconstruction. Identifying the threefold with a particular frame bundle, or the invariant with a Bott or Godbillon--Vey class, would add information absent from the notebook.}
\GSfield{Fate}{Targeted searches combining Arnold, the 1971 date, foliation invariants, rotation numbers, projective plane, and rationality did not identify an exact publication. The thematic proximity to characteristic classes does not establish a match.}
\GSfield{Present status}{As of November 2026: unidentified from the surviving statement. This is not evidence that a well-defined problem remains open; recovering the actual cocycle is prerequisite to further work.}
\GSfield{Assessment}{literal statement: not determinable; publication match: none found; verification: unresolved lead; development: annotated; confidence: low}
\end{GSitem}
\begin{GSitem}{1971-244.2}{F}{Recovering rotation numbers from the proposed invariant}
\GSfield{Kind; source}{problem; \GSpage{P000394}}
\GSfield{The notes}{The final exercise asks to define rotation numbers as invariants of the same kind, and the preceding sentence calls the construction a generalization of Poincar\'e rotation numbers.}
\GSfield{Statement}{For comparison, an orientation-preserving circle homeomorphism \(f\), with lift \(F:\mathbb R\to\mathbb R\) satisfying \(F(x+1)=F(x)+1\), has rotation number \(\rho(f)=\lim_{n\to\infty}(F^n(x)-x)/n\pmod{\mathbb Z}\). The requested realization of this number by the unnamed foliation invariant is not specified well enough to formulate an equivalence.}
\GSfield{Reconstruction}{The displayed modern definition is supplied solely to distinguish the classical invariant from the unidentified proposed realization. It is not claimed to be the notebook's missing construction.}
\GSfield{Fate}{The classical rotation number predates the seminar, but a publication matching the proposed characteristic-class realization was not identified by the searches recorded in the research log. No cross-year problem identifier was reused merely because earlier foliation items mention \(\mathbb{CP}^{2}\).}
\GSfield{Present status}{As of November 2026: the exercise's intended map from the auxiliary geometry to \(\mathbb R/\mathbb Z\) remains unidentifiable in these notes.}
\GSfield{Assessment}{literal statement: not determinable; publication match: none found; verification: unresolved lead; development: annotated; confidence: low}
\end{GSitem}

\GStalk{245}{}{-}{P000395}{Year heading: 1972}
This is the catalogue's year-divider record, not an additional mathematical talk. The next explicit notebook heading is ``14 February. L. Vaserstein'' at the top of P000395. The assigned page is the location of this transition; it does not transfer Arnold's preceding paragraph or Vaserstein's ensuing mathematics to the heading.


\GStalk{246}{1972-02-14}{L. N. Vaserstein}{P000395-P000399}{Quadratic form parameters and the beginnings of hermitian K-theory}
Vaserstein develops quadratic forms over rings with involution, including the form parameter needed in characteristic two, hyperbolic forms, Grothendieck--Witt and Witt groups, and stable classical groups. Examples over fields and the integers lead to remarks about classification and group rings. The final page introduces Steinberg-type central extensions and a simplicial higher-theory programme. The identifiable block starts at the dated heading on P000395 and ends before Kazhdan's heading on P000400. Several assertions are correct only after separating nonsingular forms, stable equivalence, and isometry classes.

\begin{GSitem}{1971-246.1}{E}{Quadratic polynomials, involutions, and twisted duals}
\GSfield{Kind; source}{definition; \GSpage{P000395}}
\GSfield{The notes}{The talk begins with \(q(x)=\sum a_{ij}x_ix_j\) over a commutative ring, and says that when \(2\) is invertible one may use symmetric matrices uniquely. Over a ring with involution it uses \(f(u\mu,v\lambda)=\mu^*f(u,v)\lambda\), and corrects the left/right module description of \(V^*=\operatorname{Hom}_A(V,A)\).}
\GSfield{Statement}{For an associative unital ring \(A\) with anti-involution \(*\), give \(V^*\) the right action \((f\cdot a)(v)=a^*f(v)\). Then right-linear maps \(L:V\to V^*\) correspond to sesquilinear pairings \(b(u,v)=L(u)(v)\). Over a commutative ring with \(2\in A^\times\), formal homogeneous quadratic polynomials on \(A^n\) have unique symmetric coefficient matrices; arbitrary polynomial functions over a finite ring need not determine their polynomial expressions.}
\GSfield{Reconstruction}{The twisted right action makes the notebook's successive corrections coherent. The qualification ``formal polynomial'' avoids an unnecessary assertion about polynomial functions.}
\GSfield{Fate}{These are the algebraic conventions underlying the form-ring theory, already used before the talk in Vaserstein's 1970 paper. A precise modern statement appears in \citet[\S1]{bakvavilov2000structure}.}
\GSfield{Present status}{As of November 2026: standard definitions. The crossed-out module adjectives are not evidence for a different duality convention.}
\GSfield{Assessment}{literal statement: correct with missing hypotheses; publication match: exact; verification: read the result; development: annotated; confidence: high}
\end{GSitem}
\begin{GSitem}{1971-246.2}{E}{The form parameter and quadratic forms as a quotient}
\GSfield{Kind; source}{definition; \GSpage{P000396}}
\GSfield{The notes}{A central \(\varepsilon\) satisfies \(\varepsilon\varepsilon^*=1\), and an additive subgroup \(\Lambda\) lies between \(\{x-\varepsilon x^*\}\) and \(\{x=-\varepsilon x^*\}\), with \(a^*\Lambda a\subseteq\Lambda\). Quadratic forms are sesquilinear forms modulo the subgroup defined by \(X=-\varepsilon X^*\) and \(X(v,v)\in\Lambda\).}
\GSfield{Statement}{Fix \((A,*,\varepsilon,\Lambda)\) with precisely those conditions. On a right \(A\)-module \(V\), identify sesquilinear forms \(Q,Q'\) when \(X=Q-Q'\) satisfies \(X(u,v)+\varepsilon X(v,u)^*=0\) and \(X(v,v)\in\Lambda\) for all \(u,v\). This quotient retains both the associated \(\varepsilon\)-hermitian pairing and the diagonal refinement modulo \(\Lambda\).}
\GSfield{Reconstruction}{The matrix star is interpreted as conjugate transpose; otherwise the displayed quotient has no invariant meaning.}
\GSfield{Fate}{The framework is that of Bak's form rings, whose pre-1972 origin is documented in \citet[introduction and \S1]{bakvavilov2000structure}. The latter gives the exact bounds on \(\Lambda\), the conjugation condition, and the corresponding quadratic module.}
\GSfield{Present status}{As of November 2026: established form-ring language, particularly essential when \(2\) is not invertible.}
\GSfield{Assessment}{literal statement: correct; publication match: exact; verification: read the result; development: annotated; confidence: high}
\end{GSitem}
\begin{GSitem}{1971-246.3}{E}{Hermitianization and the failure of the diagonal alone at maximal parameter}
\GSfield{Kind; source}{claim; \GSpage{P000396}}
\GSfield{The notes}{The notes define \((u,v)_q=Q(u,v)+\varepsilon Q(v,u)^*\) and \(|u|_q=Q(u,u)+\Lambda\). The final line says that for maximal \(\Lambda\) the quadratic form can be recovered from its \(|u|_q\), with ``linearization'' in parentheses.}
\GSfield{Statement}{Both \(b_q\) and \(q_\Lambda\) are independent of the representative. At \(\Lambda=\Lambda_{\max}\), the quotient quadratic form is determined by \(b_q\), not generally by \(q_\Lambda\). For a field with trivial involution, \(\varepsilon=-1\), and \(\Lambda=A\), every diagonal refinement takes values in \(A/\Lambda=0\), while nonzero alternating pairings exist.}
\GSfield{Reconstruction}{The replacement of the diagonal expression by the hermitianized pairing is a mathematical emendation. The counterexample diagnoses the literal assertion without attributing it to the speaker's intended formulation.}
\GSfield{Fate}{The two pieces of data and their relation are explicit in \citet[\S1, especially Lemma 1.1]{bakvavilov2000structure}. The scan confirms the diagonal symbol in the disputed sentence; it is not an English translation error.}
\GSfield{Present status}{As of November 2026: the corrected statement is standard; the displayed maximal-parameter assertion is false literally.}
\GSfield{Assessment}{literal statement: false as written; publication match: partial; verification: read the result; development: annotated; confidence: high}
\end{GSitem}
\begin{GSitem}{1971-246.4}{E}{Hyperbolic quadratic modules and nonsingularity}
\GSfield{Kind; source}{construction; \GSpage{P000397}}
\GSfield{The notes}{The example is called Hamiltonian or hyperbolic: \(H(U)=U\oplus U^*\), with a representative \(Q(u+f,u'+f')=f(u')\). A triangular matrix and the alternating matrix \(\left(\begin{smallmatrix}0&1\\-1&0\end{smallmatrix}\right)\) are displayed. Nonsingularity means that the induced map to the dual is an isomorphism.}
\GSfield{Statement}{For finite projective \(U\), the evaluation pairing defines a hyperbolic \(\Lambda\)-quadratic module on \(U\oplus U^*\). With the compatible ordering of summands, its hermitianized matrix is \(\left(\begin{smallmatrix}0&I\\\varepsilon I&0\end{smallmatrix}\right)\), up to exchanging the two blocks. The alternating matrix in the notes is the specialization \(\varepsilon=-1\).}
\GSfield{Reconstruction}{Finite projectivity supplies the canonical bidual isomorphism and hence nonsingularity. The sign is restored as a parameter rather than silently fixing the entire talk to the symplectic case.}
\GSfield{Fate}{See \citet[\S1, hyperbolic quadratic module]{bakvavilov2000structure}; Vaserstein's own earlier treatment of stable unitary and orthogonal groups is \citet{vaserstein1970stabilization}.}
\GSfield{Present status}{As of November 2026: standard hyperbolic construction; the displayed minus sign cannot represent all the form parameters introduced on P000396.}
\GSfield{Assessment}{literal statement: correct with missing hypotheses; publication match: exact; verification: read the result; development: annotated; confidence: high}
\end{GSitem}
\begin{GSitem}{1971-246.5}{E}{Grothendieck--Witt group versus Witt quotient}
\GSfield{Kind; source}{definition; \GSpage{P000397}}
\GSfield{The notes}{Nonsingular forms on finite projective modules form a semigroup under orthogonal sum. Its group completion is first called the Witt group, then explicitly corrected to Witt--Grothendieck; the Witt group is the quotient by Hamiltonian forms. The notation is \(K_0O(A)\), with the form parameters suppressed.}
\GSfield{Statement}{The group completion of isometry classes of nonsingular \(\Lambda\)-quadratic forms is \(GW_0(A,*,\varepsilon,\Lambda)\). Its quotient by the subgroup generated by hyperbolic forms is the corresponding Witt group. Witt cancellation, in its familiar form, holds for finite-dimensional nonsingular quadratic spaces over a field of characteristic different from two; analogous involutive versions require their own hypotheses.}
\GSfield{Reconstruction}{The notebook's correction is retained. The characteristic restriction is supplied for the classical cancellation statement, without asserting an unrestricted theorem for every parameter over every division ring.}
\GSfield{Fate}{These are classical distinctions in the theory of forms, rather than a new higher \(K\)-theory result. The stable group framework had already appeared in \citet{vaserstein1970stabilization}; \citet[\S1]{bakvavilov2000structure} fixes the parameter conventions used here.}
\GSfield{Present status}{As of November 2026: established. Conflating group completion with the hyperbolic quotient would change subsequent examples.}
\GSfield{Assessment}{literal statement: correct with missing hypotheses; publication match: partial; verification: indirect; development: annotated; confidence: high}
\end{GSitem}
\begin{GSitem}{1971-246.6}{E}{Hyperbolic complements and the finiteness restriction}
\GSfield{Kind; source}{theorem; \GSpage{P000397}}
\GSfield{The notes}{Theorem 1 says that every finite projective module with a quadratic form is an orthogonal summand of \(H(A^n)\). A following remark says that allowing cancellation by arbitrary modules makes the group zero, illustrated by \(a+b+a+b+\cdots\).}
\GSfield{Statement}{For a nonsingular quadratic module \((P,q)\), the diagonal in \((P,q)\perp(P,-q)\) is a lagrangian, giving a hyperbolic module; choosing \(P\oplus P'\cong A^n\) supplies a finite hyperbolic complement. If the category instead admits countable orthogonal sums, the isomorphism \(X\oplus\bigoplus_{m\ge1}X\cong\bigoplus_{m\ge1}X\) forces \([X]=0\) in its group completion.}
\GSfield{Reconstruction}{Nonsingularity is inherited from the immediately preceding definition and is indispensable: a singular form cannot be a nonsingular orthogonal summand of a nonsingular form. The infinite-sum remark is an Eilenberg swindle, not a statement that ordinary finite-projective Witt groups vanish.}
\GSfield{Fate}{Both arguments belong to the standard stable algebraic framework; compare the finite-projective definitions in \citet[\S12]{bass1964stable}. No new general vanishing theorem is being announced.}
\GSfield{Present status}{As of November 2026: valid with the stated category restrictions. The finite-projective and arbitrary-module group completions must remain distinct.}
\GSfield{Assessment}{literal statement: correct with missing hypotheses; publication match: partial; verification: indirect; development: annotated; confidence: high}
\end{GSitem}
\begin{GSitem}{1971-246.7}{E}{Field and integral examples: stable invariants are not complete isometry invariants}
\GSfield{Kind; source}{claim; \GSpage{P000398}}
\GSfield{The notes}{Examples over a field with \(\varepsilon=\pm1\) are followed by \(K_0O(A)=\mathbb Z\). Over \(\mathbb Z\), the notes discuss dimension and signature, and for \(\varepsilon=-1\) distinguish \(\Lambda=\mathbb Z\) from \(2\mathbb Z\), mentioning the Arf invariant.}
\GSfield{Statement}{Nonsingular alternating forms over a field are classified by their even dimension, yielding \(GW_0^{\mathrm{alt}}\cong\mathbb Z\). The corresponding blanket statement for symmetric forms is false: over \(\mathbb R\), positive and negative indices give \(GW_0(\mathbb R)\cong\mathbb Z^2\). Integral rank and signature do not classify definite even unimodular lattices up to isometry. For a unimodular alternating integral pairing, a quadratic refinement modulo two induces a nonsingular quadratic space on the reduction modulo two, whose Arf invariant distinguishes its two symplectic equivalence classes in positive even dimension.}
\GSfield{Reconstruction}{The examples are separated by parameter and equivalence relation. The Arf statement is about the induced mod-two refinement, not an invariant of the underlying alternating bilinear pairing alone.}
\GSfield{Fate}{The form-parameter distinction is explicit in \citet[\S1]{bakvavilov2000structure}. The general framework predates the talk. The Arf discussion is thematically related to earlier occurrence 1970-228.2, but the notebook poses no recurrent question requiring a problem-index identification.}
\GSfield{Present status}{As of November 2026: the corrected examples are standard; the universal \(\mathbb Z\) claim and an unrestricted isometry classification by rank and signature are false.}
\GSfield{Assessment}{literal statement: false as written; publication match: partial; verification: indirect; development: annotated; confidence: high}
\end{GSitem}
\begin{GSitem}{1971-246.8}{F}{An unspecified K-group of a finite group ring}
\GSfield{Kind; source}{question; \GSpage{P000398}}
\GSfield{The notes}{The notes say that little is known for \(A=\mathbb Z[g]\), and give \(g=\mathbb Z_p\times\mathbb Z_q\) as an example where \(K_0\) is unknown.}
\GSfield{Statement}{A precise computation would require specifying whether \(K_0\) means ordinary projective \(K_0\), its reduced class group, or the hermitian group \(K_0O\) of this talk; in the last case one must also fix the group-ring involution, \(\varepsilon\), and \(\Lambda\). The roles and assumptions on \(p,q\) are not supplied.}
\GSfield{Reconstruction}{No particular group-ring conjecture is substituted for these omitted choices. In particular, a theorem about projective ideal class groups would not automatically settle the hermitian question.}
\GSfield{Fate}{Searches of Vaserstein's publications and the surrounding algebraic \(K\)-theory literature did not identify a unique computation with these parameters. The supplied prior index contains no exact recurrence.}
\GSfield{Present status}{As of November 2026: the historical ``unknown'' cannot be converted into a present-day open-problem claim without identifying the functor.}
\GSfield{Assessment}{literal statement: not determinable; publication match: none found; verification: unresolved lead; development: annotated; confidence: low}
\end{GSitem}
\begin{GSitem}{1971-246.9}{E}{The mass-formula lower bound for rank-32 even unimodular lattices}
\GSfield{Kind; source}{remark; \GSpage{P000398}}
\GSfield{The notes}{A numerical aside mentions approximately \(80\cdot10^6\) forms in 32 variables with determinant one.}
\GSfield{Statement}{Let \(\Omega_{32}\) be the integral isometry classes of positive-definite even unimodular lattices of rank 32. The classical mass \(m_{32}=\sum_{L\in\Omega_{32}}|\operatorname{Aut}L|^{-1}\) satisfies \(|\Omega_{32}|\ge\lceil2m_{32}\rceil=80,618,466\), since every such lattice has the two automorphisms \(\pm1\). This is a lower bound, not an exact count.}
\GSfield{Reconstruction}{Positive definiteness, evenness, integral isometry, and the lower-bound interpretation are restored from the mass formula. The notebook gives only the rounded number and determinant-one condition, so the identification is strong but not textually complete. An independent exact-arithmetic check of the mass normalization and the counting inequality is given in the editorial verification \ref{verification:mass}; this checks the historical numerical identification and is not a new lattice-classification result.}
\GSfield{Fate}{The number is the classical Minkowski--Siegel mass bound. An exact retrospective check is \citet[\S6, comparison following Corollary 17]{king2003mass}, which prints 80,618,466 and obtains the stronger lower bound 1,162,109,024. Its preprint first appeared in 2000 and the paper in 2003; that later improvement is not attributed to the 1972 speaker.}
\GSfield{Present status}{As of November 2026: the recovered lower bound is established and has long been surpassed. No assertion of a complete modern enumeration is made.}
\GSfield{Assessment}{literal statement: correct with missing hypotheses; publication match: partial; verification: read the result; development: annotated; confidence: medium}
\end{GSitem}
\begin{GSitem}{1971-246.10}{E}{Matrix-ring invariance and the choice of duality}
\GSfield{Kind; source}{theorem; \GSpage{P000398}}
\GSfield{The notes}{The notes record equality of the \(K\)-groups of \(A\) and of its matrix rings.}
\GSfield{Statement}{For a unital associative ring and \(r\ge1\), Morita equivalence yields \(K_n(M_r(A))\cong K_n(A)\) for ordinary algebraic \(K\)-theory. In a hermitian version the equivalence must transport the involution, duality, and form parameter; an arbitrary unrelated involution on the matrix ring is not covered.}
\GSfield{Reconstruction}{Ordinary Morita invariance is the secure statement. The duality requirement prevents importing it unchanged into every hermitian theory mentioned in the talk.}
\GSfield{Fate}{The low-degree stable-algebra form is already in \citet[\S\S12--14]{bass1964stable}; the higher-theory comparison is subsequently incorporated into \citet{vaserstein1973foundations}.}
\GSfield{Present status}{As of November 2026: established, with the stated compatibility qualification for forms.}
\GSfield{Assessment}{literal statement: correct with missing hypotheses; publication match: partial; verification: read the result; development: annotated; confidence: high}
\end{GSitem}
\begin{GSitem}{1971-246.11}{E}{Stable orthogonal groups and hermitian K1}
\GSfield{Kind; source}{definition; \GSpage{P000398-P000399}}
\GSfield{The notes}{The stable group \(O(A)\) is formed from automorphism groups of \(H(A^n)\); \(K_1O(A)\) is its abelianization, with an elementary subgroup playing the role of commutators.}
\GSfield{Statement}{Fix the form ring and use stabilization by a hyperbolic plane. The stable unitary group \(U(A,\Lambda)\) has the stable elementary subgroup \(EU(A,\Lambda)\); under the standard stable form-ring conventions their quotient is hermitian \(K_1\). Assertions about a single finite-rank group require separate stability bounds and cannot be read off from this direct limit.}
\GSfield{Reconstruction}{The letter \(O\) is understood in the notebook's broad parameter-dependent sense, including symplectic and unitary specializations, rather than only the usual orthogonal group of a field.}
\GSfield{Fate}{This is the subject of \citet[\S1, Theorems 1.4 and 1.4O]{vaserstein1970stabilization}, received 29 January 1969 and published in the Russian original in 1970 (English translation copyright 1971). It therefore predates the talk. The general form-parameter convention is made explicit in \citet[\S\S1--2]{bakvavilov2000structure}.}
\GSfield{Present status}{As of November 2026: standard stable hermitian \(K_1\); no unstated finite-rank stability theorem is attributed to the lecture.}
\GSfield{Assessment}{literal statement: correct with missing hypotheses; publication match: exact; verification: read the result; development: annotated; confidence: high}
\end{GSitem}
\begin{GSitem}{1971-246.12}{E}{Perfect groups, universal central extensions, and second homology}
\GSfield{Kind; source}{theorem; \GSpage{P000399}}
\GSfield{The notes}{The middle of P000399 records \(H_2(G,\mathbb Z)\to\widetilde G\to G\) with the condition \([G,G]=G\), following a Steinberg-type extension for the elementary classical group.}
\GSfield{Statement}{A group \(G\) admits a universal central extension if and only if it is perfect. Its kernel is naturally \(H_2(G;\mathbb Z)\). For the stable elementary linear group of an associative unital ring \(A\), the Steinberg group \(\operatorname{St}(A)\to E(A)\) is its universal central extension, with central kernel \(K_2(A)\cong H_2(E(A);\mathbb Z)\).}
\GSfield{Reconstruction}{The general perfect-group theorem is read separately from the proposed hermitian model. The verified stable linear case is a comparator; it does not by itself prove the same presentation theorem for every form parameter or every finite-rank classical group.}
\GSfield{Fate}{The relation with perfect groups and the classical Steinberg--Milnor construction predates the lecture. \citet[historical introduction, p.95; \S6, Theorem 6.1]{vaserstein1976foundations} gives the central-extension theorem and dates the field case to Steinberg's 1962 work; the modern general formulation is also treated in \citet[III, \S5]{weibel2013kbook}.}
\GSfield{Present status}{As of November 2026: established for perfect groups and for the stated stable linear group. The additional hermitian and simplicial claims are separated in the next item.}
\GSfield{Assessment}{literal statement: correct with missing hypotheses; publication match: partial; verification: read the result; development: annotated; confidence: high}
\end{GSitem}
\begin{GSitem}{1971-246.13}{F}{The incompletely specified higher hermitian simplicial model}
\GSfield{Kind; source}{construction; \GSpage{P000399}}
\GSfield{The notes}{The page writes \(0\to K_2O(A)\to\operatorname{St}O(A)\to EO(A)\to0\), calls it a universal central extension, and proposes \(K_nO(A)=\pi_{n-1}O(A)\). A question asks which finite subsets should be simplices; triangular unipotent matrices and their left and right translates are sketched. The ordinary linear analogue appears in the margin.}
\GSfield{Statement}{A precise higher hermitian construction must specify \((A,*,\varepsilon,\Lambda)\), the stable elementary unitary subgroup, its generators and root relations, the admissible triangular subgroups, and the simplicial realization whose homotopy groups are intended. The surviving sketches do not determine all these data or establish comparison with a particular modern Grothendieck--Witt spectrum. The notation \(\pi_{n-1}O(A)\) cannot mean homotopy of the discrete group alone.}
\GSfield{Reconstruction}{The general perfect-group theorem is retained separately in 1971-246.12. The ordinary Volodin model explains the method, but importing its linear comparison theorem does not identify the complete hermitian theory suggested here.}
\GSfield{Fate}{The prior stable form-group framework is \citet{vaserstein1970stabilization}; the linear simplicial method was published in \citet{volodin1971extraordinary}. Searches of Vaserstein's subsequent foundations papers found the ordinary associative-ring comparison \citep{vaserstein1973foundations,vaserstein1976foundations}, not a uniquely matching all-parameter hermitian statement for this sketch.}
\GSfield{Present status}{As of November 2026: the proposed model remains underdetermined by the notes. This classification concerns identification, not a claim that modern hermitian K-theory remains undeveloped.}
\GSfield{Assessment}{literal statement: not determinable; publication match: thematic only; verification: unresolved lead; development: annotated; confidence: medium}
\end{GSitem}

\GStalk{247}{1972-02-21}{D. A. Kazhdan}{P000400-P000403}{Elementary relations and Laurent polynomial K-theory}
Kazhdan reviews low algebraic \(K\)-groups and Steinberg relations, then poses the universal-relation question answered in Vaserstein's later March talks. Most of the lecture constructs the Laurent polynomial boundary and its splitting: first through projective modules and polynomial cokernels, then through truncation of two-sided matrices. The notebook itself repairs the matrix class to finite bandwidth. The date and heading occur at P000400; Gabrielov begins below the divider on P000403.

\begin{GSitem}{1971-247.1}{E}{Projective K0 and elementary-matrix K1}
\GSfield{Kind; source}{definition; \GSpage{P000400}}
\GSfield{The notes}{The opening formulas are \(K_0(A)\) as the Grothendieck group of finitely generated projectives and \(K_1(A)=GL(\infty,A)/E(A)\). A triangular sketch is followed by the explicit statement that \(E(A)\) is generated by all elementary matrices. A Helmholtz colour-subtraction analogy illustrates group completion.}
\GSfield{Statement}{For a unital associative ring \(A\), \(K_0(A)\) is the group completion of isomorphism classes of finitely generated projective right modules under direct sum. Set \(GL(A)=\varinjlim GL_n(A)\) under \(g\mapsto\operatorname{diag}(g,1)\), and let \(E(A)\) be generated by \(1+ae_{ij}\), \(i\ne j\). Then \(E(A)=[GL(A),GL(A)]\), so \(K_1(A)\) is the abelianization. A single upper triangular subgroup is not \(E(A)\).}
\GSfield{Reconstruction}{The later generator sentence governs the earlier schematic drawing. The colour analogy is treated as pedagogy, not as a claim of a mathematical identification with physical colour spaces.}
\GSfield{Fate}{The exact stable definitions and commutator result occur in \citet[\S12]{bass1964stable} and \citet[Proposition 0]{basshellerswan1964whitehead}, well before the lecture.}
\GSfield{Present status}{As of November 2026: standard definitions, with finite generation and stabilization retained.}
\GSfield{Assessment}{literal statement: correct; publication match: exact; verification: read the result; development: annotated; confidence: high}
\end{GSitem}
\begin{GSitem}{1971-247.2}{E}{Steinberg relations and the incomplete displayed presentation}
\GSfield{Kind; source}{definition; \GSpage{P000400-P000401}}
\GSfield{The notes}{Two displayed relations are additivity and \([E_{ij}(a),E_{jk}(b)]=E_{ik}(ab)\). The four-factor line above them has overwritten signs and right side \(1\). The abstract group with the relations is called \(St(A)\).}
\GSfield{Statement}{In the stable Steinberg group use generators \(x_{ij}(a)\), \(i\ne j\), with \(x_{ij}(a)x_{ij}(b)=x_{ij}(a+b)\), \([x_{ij}(a),x_{kl}(b)]=1\) when \(j\ne k\) and \(i\ne l\), and \([x_{ij}(a),x_{jk}(b)]=x_{ik}(ab)\) for distinct \(i,j,k\), with a consistent commutator convention. The four factors with adjacent roots do not in general multiply to the identity: with inverse signs in the first pair they yield \(1+ab e_{ik}\).}
\GSfield{Reconstruction}{The missing disjoint-root relation and distinct-index conditions are supplied. The overwritten first signs are left diplomatically unresolved; the erroneous identity right side is a separate mathematical defect.}
\GSfield{Fate}{See \citet[\S6]{vaserstein1976foundations} for the explicit stable Steinberg presentation and its low-degree role. These relations and Milnor's \(K_2\) are antecedent theory, not Kazhdan's new theorem.}
\GSfield{Present status}{As of November 2026: the completed presentation is standard; the notes' two displayed relations alone do not specify it fully.}
\GSfield{Assessment}{literal statement: false as written; publication match: partial; verification: read the result; development: annotated; confidence: high}
\end{GSitem}
\begin{GSitem}{1971-247.3}{E}{K2 as the central kernel of the stable Steinberg group}
\GSfield{Kind; source}{definition; \GSpage{P000400-P000401}}
\GSfield{The notes}{The notes write \(K_2(A)=\ker(St(A)\to E(A))\) and describe it informally as relations other than the standard ones.}
\GSfield{Statement}{For a unital associative ring, send \(x_{ij}(a)\) to \(1+ae_{ij}\). The stable Steinberg group is the universal central extension of the perfect group \(E(A)\), and \(K_2(A)\cong H_2(E(A),\mathbb Z)\). The kernel is a group of relation classes modulo Steinberg relations, not a free list of extra relators.}
\GSfield{Reconstruction}{Centrality, universality, and the stable range are made explicit. No analogous universal-extension assertion is made for every finite-rank \(St_n(A)\).}
\GSfield{Fate}{The low-degree theory was available before February 1972 in Milnor's introduction; the publisher dates its edition to January 1972. The precise extension theorem was read in \citet[Theorem 6.1]{vaserstein1976foundations}; see also \citet{milnor1972introduction} for historical publication metadata.}
\GSfield{Present status}{As of November 2026: established stable algebraic \(K_2\), independent of the later proposed higher models.}
\GSfield{Assessment}{literal statement: correct; publication match: exact; verification: read the result; development: annotated; confidence: high}
\end{GSitem}
\begin{GSitem}{1971-247.4}{E}{The determinant identifies K1 of a field}
\GSfield{Kind; source}{theorem; \GSpage{P000400}}
\GSfield{The notes}{The determinant-one subgroup of \(GL(4,\mathbb C)\) is said to be generated by elementary matrices. The final \(\mathbb C\) in \(K_1(\mathbb C)=\mathbb C\) has an uncertain upper mark, followed by ``and this holds for any field''.}
\GSfield{Statement}{For a commutative field \(F\), elementary row operations give \(SL_n(F)=E_n(F)\) for \(n\ge2\), and determinant induces \(K_1(F)\cong F^\times\). For a division ring the replacement is the abelianization of its unit group, rather than an ordinary determinant into its underlying additive group.}
\GSfield{Reconstruction}{The multiplicative group symbol is supplied mathematically; the scan does not justify claiming that the uncertain mark is definitely a legible star.}
\GSfield{Fate}{This is the familiar field case of the stable Whitehead group; see \citet[Proposition 0 and the division-ring corollary after Theorem 1]{basshellerswan1964whitehead}.}
\GSfield{Present status}{As of November 2026: established. The additive group of \(F\) is not the target.}
\GSfield{Assessment}{literal statement: correct with missing hypotheses; publication match: exact; verification: read the result; development: annotated; confidence: high}
\end{GSitem}
\begin{GSitem}{1971-247.5}{B}{Are universally valid elementary relations consequences of the Steinberg relations?}
\GSfield{Kind; source}{problem; \GSpage{P000401}}
\GSfield{The notes}{The page explicitly labels as unsolved the question whether a relation valid in every ring follows from the Steinberg relations. It also mentions an unidentified problem in discrete group theory as a trivial case.}
\GSfield{Statement}{A precise augmented formulation is as follows. Let \(R=\mathbb Z\langle x_s:s\in S\rangle\) be the free associative unital ring, with augmentation \(\epsilon:R\to\mathbb Z\) sending every variable to zero. Then \(\ker(K_2(R)\xrightarrow{\epsilon_*}K_2(\mathbb Z))=0\). Consequently a stable word in elementary generators that is the identity as a universal matrix and whose Steinberg lift becomes the identity after augmentation is already a consequence of the Steinberg relations. In particular this applies when every entry has zero constant term. Allowing arbitrary constant relations without this normalization introduces the nonzero group \(K_2(\mathbb Z)\).}
\GSfield{Reconstruction}{The augmentation condition is reconstructed from the subsequent use of free rings without identity. It expresses reduced \(K_2\)-vanishing without identifying the naive Steinberg kernel of an arbitrary nonunital ring with that reduced group. The cryptic discrete-group subproblem is not identifiable.}
\GSfield{Fate}{Vaserstein announces the answer on P000413. \citet[second theorem, p.231]{vaserstein1973foundations} states free nonunital acyclicity; \citet[Theorem 28.1 and Corollary 28.1]{vaserstein1976foundations} gives precisely \(K_n(\mathbb Z)\cong K_n(\mathbb Z\langle S\rangle)\) and its reduced consequence, with Gersten's proof cited for higher degrees. The supplied earlier-year index contains no exact recurrence. The match is partial because the February wording does not spell out the augmentation convention.}
\GSfield{Present status}{As of November 2026: the augmented universal-relation statement is resolved. The literal unnormalized statement must be qualified; it is not a remaining open problem.}
\GSfield{Assessment}{literal statement: not determinable; publication match: partial; verification: read the result; development: annotated; confidence: medium; problem \hyperref[prob:GP-1971-20]{GP-1971-20}}
\end{GSitem}
\begin{GSitem}{1971-247.6}{E}{Embedding K0 as a Laurent K1 summand}
\GSfield{Kind; source}{construction; \GSpage{P000401}}
\GSfield{The notes}{For \(B=A[t,t^{-1}]\), an automorphism of a finite projective module is embedded into an automorphism of a free module by adjoining the identity on a complement. Multiplication by \(t\) on \(P\otimes_A B\) defines \(K_0(A)\to K_1(B)\).}
\GSfield{Statement}{For \(P\oplus Q\cong A^r\), take the class of \(t\,1_{P_B}\oplus1_{Q_B}\) in \(K_1(B)\). It is independent of the complement and free identification, additive in \(P\), and defines a split injection \(K_0(A)\hookrightarrow K_1(B)\). The indeterminate is central, also when \(A\) is noncommutative.}
\GSfield{Reconstruction}{The projective-automorphism definition explains why the construction does not require \(P\) itself to be free.}
\GSfield{Fate}{The construction is the map \(h\) of \citet[pp.62--63 and \S5]{basshellerswan1964whitehead}. Their regular-ring decomposition predates the talk; the summand construction and index remain meaningful more generally, with Nil terms retained.}
\GSfield{Present status}{As of November 2026: standard part of the fundamental theorem for Laurent polynomial \(K\)-theory.}
\GSfield{Assessment}{literal statement: correct; publication match: exact; verification: read the result; development: annotated; confidence: high}
\end{GSitem}
\begin{GSitem}{1971-247.7}{E}{Polynomial cokernels and the normalization of the Laurent index}
\GSfield{Kind; source}{construction; \GSpage{P000402}}
\GSfield{The notes}{Write a Laurent invertible matrix \(L=t^{-n}L'\) with polynomial entries in \(L'\); its cokernel is said to be a finite projective \(A\)-module. The page calls this the inverse construction but does not subtract the shift contribution.}
\GSfield{Statement}{If \(g\in GL_r(A[t,t^{-1}])\), choose \(n\ge0\) with \(t^ng\in M_r(A[t])\). The Laurent index is \(\partial[g]=[\operatorname{coker}(t^ng:A[t]^r\to A[t]^r)]-nr[A]\). The cokernel is finite projective as an \(A\)-module. Increasing \(n\) by one adds \(r[A]\) to the first term, so the difference is independent of \(n\), and \(\partial h[P]=[P]\).}
\GSfield{Reconstruction}{The subtractive term is required; record 249 supplies it explicitly. The formula is written for stable classes, not claimed to turn two nonisomorphic polynomial cokernels into isomorphic modules.}
\GSfield{Fate}{This is the index construction behind \citet[\S5]{basshellerswan1964whitehead}; the modern version is also given in \citet[III, \S3]{weibel2013kbook}. The following week's return to it is a clarification of an old construction.}
\GSfield{Present status}{As of November 2026: established after normalization. Using only the cokernel class would assign a nonzero answer to the identity when one chooses \(n>0\).}
\GSfield{Assessment}{literal statement: false as written; publication match: partial; verification: read the result; development: annotated; confidence: high}
\end{GSitem}
\begin{GSitem}{1971-247.8}{B}{The K1 summand in Laurent K2}
\GSfield{Kind; source}{theorem; \GSpage{P000401-P000402}}
\GSfield{The notes}{Beside \(K_0(A)\subset K_1(B)\), the notes assert that \(K_1(A)\) is a direct summand of \(K_2(B)\), and leave construction of the forward map as an exercise.}
\GSfield{Statement}{For a unital associative ring \(A\) and central \(t\), the Laurent fundamental theorem has a natural split summand \(K_1(A)\subset K_2(A[t,t^{-1}])\); in full generality \(K_2(B)\cong K_2(A)\oplus K_1(A)\oplus NK_2(A)\oplus NK_2(A)\), where \(NK_2(A)=\ker(K_2(A[t])\to K_2(A))\).}
\GSfield{Reconstruction}{The full modern decomposition is supplied to prevent suppressing the two Nil terms. The notebook asserts only the indicated summand.}
\GSfield{Fate}{\citet{gersten1971laurent} had already proved a related regular-ring surjectivity result (received June 1970, printed October 1971), not the exact general decomposition. The general statement is among the explicit Laurent formulas of \citet[third theorem, p.231]{vaserstein1973foundations}; record 253 returns to it and cites Wagoner in the notes. The class records the checked general publication rather than attributing priority solely to this lecture.}
\GSfield{Present status}{As of November 2026: established; omission of Nil terms is justified only under additional regularity or vanishing hypotheses.}
\GSfield{Assessment}{literal statement: correct; publication match: exact; verification: read the result; development: annotated; confidence: medium}
\end{GSitem}
\begin{GSitem}{1971-247.9}{B}{Truncation modulo finite matrices as a boundary construction}
\GSfield{Kind; source}{construction; \GSpage{P000402-P000403}}
\GSfield{The notes}{Two-sided row-and-column-finite matrices are truncated to one side. A marginal correction restricts them to finitely many diagonals parallel to the main diagonal. A lifted elementary relation has image \(1\) modulo finite matrices and is then described as a finite matrix representing \(K_1(A)\).}
\GSfield{Statement}{For two-sided finite-band matrices, compression to the nonnegative half-line is multiplicative modulo the ideal of finitely supported matrices: the defect factors through only finitely many indices at the cut. Laurent multiplication matrices lie in this algebra. Lifting a Steinberg relation through the compression produces an invertible \(1+F\) with \(F\) finitely supported, hence a stable \(K_1(A)\)-class. It is \(L-1\), not \(L\), that has finite support.}
\GSfield{Reconstruction}{This reconstructs the algebraic Toeplitz boundary and makes explicit the finite-support defect. Independence of lifts is part of the boundary theorem; the notebook's sketches alone are not a proof.}
\GSfield{Fate}{The general suspension and exact-sequence treatment is published by \citet[first--third theorems]{vaserstein1973foundations}. Prior Laurent constructions appear in \citet{basshellerswan1964whitehead,gersten1971laurent}, but the retrieved earlier sources do not establish priority for this exact general Toeplitz implementation. Class B records the verified later formulation by another author, with the match qualified because the notebook provides only a sketch.}
\GSfield{Present status}{As of November 2026: standard after imposing the notebook's own bandwidth correction and reading the last matrix as identity plus a finite matrix.}
\GSfield{Assessment}{literal statement: false as written; publication match: partial; verification: read the result; development: annotated; confidence: medium}
\end{GSitem}

\GStalk{248}{1972-02-21}{A. M. Gabrielov}{P000403-P000405}{Polarized Hodge structures and the Hodge--Riemann signs}
Gabrielov explicitly reports Griffiths's work on variation of Hodge structure, with pointers to Griffiths's 1970 survey and a Deligne Bourbaki report. This first installment develops the geometry of one smooth projective variety: Hodge decomposition, Lefschetz decomposition, the primitive polarization, and the abstract linear-algebraic definition. These are prior results presented as preparation for the family theory of record 251. A sign inconsistency in the positive-definiteness formula is visible in the scan.

\begin{GSitem}{1971-248.1}{E}{An integral polarization and Hodge decomposition}
\GSfield{Kind; source}{definition; \GSpage{P000403}}
\GSfield{The notes}{For a nonsingular complex projective \(X\) of dimension \(d\), the notes introduce \(H^n(X,\mathbb Z)\), its complexification, and a closed integral \((1,1)\)-form \(\eta\), called a polarization. They write \(H^n(X,\mathbb C)=\bigoplus_{p+q=n}H^{p,q}\).}
\GSfield{Statement}{For smooth projective \(X/\mathbb C\), take \(\eta=c_1(L)\) for an ample line bundle, represented by a K\"ahler form. Then \(H^n(X,\mathbb C)=\bigoplus H^{p,q}(X)\), with \(\overline{H^{p,q}}=H^{q,p}\). The integral lattice for a pure integral Hodge structure is \(H^n(X,\mathbb Z)/\mathrm{torsion}\); complexification kills torsion.}
\GSfield{Reconstruction}{An arbitrary integral closed \((1,1)\)-class need not be positive. Ampleness/K\"ahler positivity and the torsion-free lattice are the missing qualifications.}
\GSfield{Fate}{These are the prerequisites of the already published survey \citet[\S2]{griffiths1970periods} and the pure-Hodge definitions of \citet[\S2.1]{deligne1971hodge}. Griffiths's text derives from a 1968 address and appeared in March 1970.}
\GSfield{Present status}{As of November 2026: established classical Hodge theory under the stated compactness and positivity hypotheses.}
\GSfield{Assessment}{literal statement: correct with missing hypotheses; publication match: exact; verification: read the result; development: annotated; confidence: high}
\end{GSitem}
\begin{GSitem}{1971-248.2}{E}{Hard Lefschetz, primitive cohomology, and orthogonal decomposition}
\GSfield{Kind; source}{theorem; \GSpage{P000404}}
\GSfield{The notes}{The operator \(\eta\) raises bidegree by \((1,1)\). Primitive cohomology is the kernel of a high power of \(\eta\), and the notes decompose cohomology into powers of \(\eta\) times primitive pieces.}
\GSfield{Statement}{Let \(L=\eta\smile-\). For \(0\le n\le d\), \(L^{d-n}:H^n(X,\mathbb Q)\to H^{2d-n}(X,\mathbb Q)\) is an isomorphism. Put \(P^n=\ker(L^{d-n+1}:H^n\to H^{2d-n+2})\). Then \(H^n=\bigoplus_{j\ge0}L^jP^{n-2j}\), and this decomposition respects Hodge types. With a K\"ahler metric the associated primitive Lefschetz components are orthogonal for the Hodge inner product.}
\GSfield{Reconstruction}{The degree range, coefficient field, and which inner product supplies orthogonality are stated explicitly; the formula for primitive degree is not extended beyond its intended range without reinterpretation.}
\GSfield{Fate}{\citet[\S2, pp.234--235]{griffiths1970periods} explicitly passes to primitive cohomology and gives the Lefschetz decomposition; this is classical input to the variation theory.}
\GSfield{Present status}{As of November 2026: established, not an instance of the Hodge conjecture.}
\GSfield{Assessment}{literal statement: correct with missing hypotheses; publication match: exact; verification: read the result; development: annotated; confidence: high}
\end{GSitem}
\begin{GSitem}{1971-248.3}{E}{Hodge--Riemann bilinear relations and a sign inconsistency}
\GSfield{Kind; source}{theorem; \GSpage{P000404-P000405}}
\GSfield{The notes}{The notes use \(\Psi(x,y)=(-1)^{n(n+1)/2}\int_X\eta^{d-n}\wedge x\wedge y\). They state \(i^{p-q}\Psi(x,\bar x)>0\) and later require \(\Psi(x,C\bar x)>0\), where \(C|_{H^{p,q}}=i^{p-q}\).}
\GSfield{Statement}{On primitive \(H^{p,q}\), \(p+q=n\le d\), let \(Q(x,y)=(-1)^{n(n-1)/2}\int_Xx\wedge y\wedge\eta^{d-n}\). Then \(Q(y,x)=(-1)^nQ(x,y)\), \(Q(H^{p,q},H^{r,s})=0\) unless \((r,s)=(q,p)\), and \(i^{p-q}Q(x,\bar x)>0\) for \(x\ne0\). With the notebook's \(\Psi=(-1)^nQ\), positivity is instead \(i^{q-p}\Psi(x,\bar x)>0\), equivalently \(\Psi(x,C\bar x)>0\).}
\GSfield{Reconstruction}{Two consistent conventions are displayed rather than altering one sign without explanation. The source's first and last positivity assertions differ in odd weight.}
\GSfield{Fate}{Compare the polarization conventions in \citet[\S2 and equations (9.8)--(9.9)]{griffiths1970periods} and \citet[Definition 2.1.16]{deligne1971hodge}. The scan confirms both the \(n(n+1)/2\) exponent and \(i^{p-q}\) in the first inequality.}
\GSfield{Present status}{As of November 2026: the bilinear relations are established; the literal combination of signs in the notes fails already in weight one.}
\GSfield{Assessment}{literal statement: false as written; publication match: partial; verification: read the result; development: annotated; confidence: high}
\end{GSitem}
\begin{GSitem}{1971-248.4}{E}{Abstract polarized Hodge structures and the Weil operator}
\GSfield{Kind; source}{definition; \GSpage{P000405}}
\GSfield{The notes}{The end replaces geometric cohomology by a finite-dimensional vector space with a real form, a \((p,q)\)-decomposition, the operator \(C\), and a bilinear form satisfying symmetry and positivity conditions.}
\GSfield{Statement}{A real Hodge structure of weight \(n\) is a finite-dimensional real vector space \(V\) with \(V_\mathbb C=\bigoplus_{p+q=n}V^{p,q}\), \(\overline{V^{p,q}}=V^{q,p}\). A polarization additionally has the type orthogonality and signed positivity stated in the preceding item. The Weil operator is real on \(V\), acts by \(i^{p-q}\) on \(V^{p,q}\), and satisfies \(C^2=(-1)^n\). An integral structure adds a lattice; invariance under \(C\) alone does not replace full Hodge-type orthogonality.}
\GSfield{Reconstruction}{The distinction between a real structure and an integral lattice is made explicit. This is a definition and does not assert that every abstract polarized Hodge structure comes from a smooth projective variety.}
\GSfield{Fate}{See \citet[\S2.1, especially Definition 2.1.16]{deligne1971hodge} and \citet[\S9]{griffiths1970periods}. Both were published before Gabrielov's report.}
\GSfield{Present status}{As of November 2026: standard linear-algebraic structure used in variations and period domains.}
\GSfield{Assessment}{literal statement: correct with missing hypotheses; publication match: exact; verification: read the result; development: annotated; confidence: high}
\end{GSitem}

\GStalk{249}{1972-02-28}{A. A. Kirillov}{P000405-P000407}{The normalized Laurent index and polynomial Nil groups}
This lecture revisits the preceding Laurent construction, explicitly noting that one must subtract the excess introduced by clearing negative powers. Gelfand repeats the projective-module map, and the final part states the fundamental exact sequence, the Nil summand, and Higman linearization. The catalogue date is preserved in metadata; the handwritten heading is clearly \(27/\mathrm{II}\), not \(28/\mathrm{II}\). The talk begins near the bottom of P000405 and ends at the divider before Morava on P000407.

\begin{GSitem}{1971-249.1}{E}{The shift-independent Laurent index revisited}
\GSfield{Kind; source}{construction; \GSpage{P000405-P000406}}
\GSfield{The notes}{Kirillov writes \(M=t^{-k}M_1=t^{-\ell}M_2\), observes that polynomial multiplication has zero kernel, and explicitly adds ``one must subtract the excess''. The same symbol \(k\) is also used for matrix size.}
\GSfield{Statement}{For \(g\in GL_r(A[t,t^{-1}])\) and \(N\) large enough, define \(\partial[g]=[\operatorname{coker}(t^Ng)]-Nr[A]\). The first cokernel is finite projective over \(A\). Multiplying by another \(t\) changes its class by \(r[A]\), proving independence of \(N\). Elementary changes preserve the resulting stable index; arbitrary Laurent elementary changes need not preserve a chosen polynomial cokernel before normalization.}
\GSfield{Reconstruction}{Separate symbols \(r,N\) prevent confusing size with denominator-clearing order. The positive-cokernel convention is chosen so that multiplication by \(t\) has index \(+[A]\); the analytic Fredholm convention would have the opposite sign.}
\GSfield{Fate}{This clarifies occurrence 1971-247.7. The index construction was already in \citet[\S5]{basshellerswan1964whitehead}; see \citet[III, \S3]{weibel2013kbook} for the general form.}
\GSfield{Present status}{As of November 2026: established; the new lecture's significance is the recorded repair of the missing normalization.}
\GSfield{Assessment}{literal statement: correct with missing hypotheses; publication match: exact; verification: read the result; development: annotated; confidence: high}
\end{GSitem}
\begin{GSitem}{1971-249.2}{E}{The projective-module splitting and Gelfand’s repetition}
\GSfield{Kind; source}{construction; \GSpage{P000406}; A. A. Kirillov; intervention by I. M. Gelfand}
\GSfield{The notes}{For \(N\oplus M=A^r\), the automorphism \(t\) on \(N_B\) and \(1\) on \(M_B\) gives a free \(B\)-matrix. The notes compute \(\operatorname{coker}(t:N[t]\to N[t])=N\). Gelfand repeats the diagram, whose displayed scalar-extended module is erroneously labelled \(A^r\).}
\GSfield{Statement}{The matrix for \(t1_{N_B}\oplus1_{M_B}\) is already polynomial after choosing an \(A\)-linear free identification \(N\oplus M\cong A^r\). Its polynomial cokernel is \(N\), so \(\partial h=1\) on \(K_0(A)\). Scalar extension yields \(N_B\oplus M_B\cong B^r\), not \(A^r\).}
\GSfield{Reconstruction}{This verifies one side of the splitting and does not assert that the two maps are inverses on all of \(K_1(B)\), which contains the other summands.}
\GSfield{Fate}{The map and its left inverse are those of \citet[pp.62--63 and \S5]{basshellerswan1964whitehead}. The two recorded explanations are part of the same construction, not two separate discoveries.}
\GSfield{Present status}{As of November 2026: established; the module label requires correction.}
\GSfield{Assessment}{literal statement: false as written; publication match: partial; verification: read the result; development: annotated; confidence: high}
\end{GSitem}
\begin{GSitem}{1971-249.3}{E}{The Laurent exact sequence and polynomial Nil summands}
\GSfield{Kind; source}{theorem; \GSpage{P000407}}
\GSfield{The notes}{The notes display \(0\to K_1(A)\to K_1(A[t])\oplus K_1(A[t^{-1}])\to K_1(A[t,t^{-1}])\to K_0(A)\to0\), and \(K_1(A[t])=K_1(A)\oplus\operatorname{Nil}A\). They say that Nil vanishes for regular rings.}
\GSfield{Statement}{For a unital ring, set \(NK_1(A)=\ker(K_1(A[t])\xrightarrow{t\mapsto0}K_1(A))\). Evaluation splits the polynomial group. The displayed Laurent sequence, with opposite signs on the two constant inclusions, is exact and its right boundary is split by \(h\). Thus \(K_1(A[t,t^{-1}])\cong K_1(A)\oplus K_0(A)\oplus NK_1(A)^{\oplus2}\). If \(A\) is left regular noetherian (every finite left module has finite projective dimension), then \(NK_1(A)=0\).}
\GSfield{Reconstruction}{The definition of regularity and the two copies of the Nil group are supplied. ``Functions on the suspension'' on P000405 is an index-theoretic analogy: algebraically the Laurent extension corresponds to product with the multiplicative group, not literally to a topological suspension.}
\GSfield{Fate}{\citet[Theorems 1--2]{basshellerswan1964whitehead} prove the regular version; \citet[III, Theorem 3.6 and \S3]{weibel2013kbook} gives the general formulation including Nil groups. Earlier item 1968-171.8 concerns related polynomial matrices, but no identical indexed problem was located.}
\GSfield{Present status}{As of November 2026: established, with regularity essential for deleting Nil terms.}
\GSfield{Assessment}{literal statement: correct with missing hypotheses; publication match: exact; verification: read the result; development: annotated; confidence: high}
\end{GSitem}
\begin{GSitem}{1971-249.4}{E}{Higman linearization and nilpotent matrices}
\GSfield{Kind; source}{method; \GSpage{P000407}}
\GSfield{The notes}{After removing the constant \(K_1\)-part, the notes write \(M=1+tN(t)\). Enlarging the matrix reduces higher powers to first order, so \(N\) becomes constant, and polynomial invertibility forces \(N^k=0\).}
\GSfield{Statement}{A polynomial matrix invertible over \(A[t]\) with constant term \(I\) becomes, after stabilization and elementary equivalence, \(I-tN\) for a finite nilpotent matrix \(N\) over \(A\). Conversely \(I-tN\) is polynomially invertible exactly when \(N\) is nilpotent, with inverse the finite geometric series. Nilpotent endomorphisms supply generators and relations for the Nil group; their raw set of matrices is not itself a group.}
\GSfield{Reconstruction}{The sign of \(N\) is conventional. Stable equivalence, rather than equality of the original matrices, is the content of the linearization.}
\GSfield{Fate}{This is the classical Higman trick, used in \citet[Proposition 2.1 and its proof]{basshellerswan1964whitehead} and explained in \citet[III, \S3]{weibel2013kbook}.}
\GSfield{Present status}{As of November 2026: established and still a useful computational model of polynomial Nil phenomena.}
\GSfield{Assessment}{literal statement: correct with missing hypotheses; publication match: exact; verification: read the result; development: annotated; confidence: high}
\end{GSitem}

\GStalk{250}{1972-02-28}{Jack Morava}{P000407-P000408}{Quillen’s plus construction and algebraic K-theory}
Morava explains the construction that kills a perfect normal subgroup while preserving homology, applies it to the discrete stable general linear group, and notes that degree zero has to be added separately. The talk compares this theory with Milnor's \(K_2\), Volodin's construction, and topological \(K\)-theory. It closes with Novikov's broad proposal for a Chern character. This is a report on Quillen's already circulating work, not the invention of the plus construction by Morava.

\begin{GSitem}{1971-250.1}{E}{Killing a perfect normal subgroup by the plus construction}
\GSfield{Kind; source}{theorem; \GSpage{P000407}}
\GSfield{The notes}{For \(H\triangleleft\pi_1(X)\), \(H=[H,H]\), the notes attach two- and three-cells, obtain \(\pi_1(X_H)=\pi_1(X)/H\), and preserve integral homology. They claim uniqueness up to homotopy and functoriality.}
\GSfield{Statement}{For a connected CW complex \(X\) and a perfect normal subgroup \(H\), there is a plus construction \(X\to X_H^+\) obtained by attaching cells in dimensions 2 and 3, with fundamental group quotient \(\pi_1(X)/H\). The map is acyclic: it induces homology isomorphisms with every local coefficient module over that quotient. The resulting map under \(X\) is unique up to homotopy. Ordinary untwisted homology alone is not the full characterization.}
\GSfield{Reconstruction}{CW hypotheses, local coefficients, and uniqueness of the map under \(X\) are supplied. The phrase ``homology ring'' is replaced by homology groups unless an additional multiplication is given.}
\GSfield{Fate}{Quillen's construction circulated before the lecture; \citet[IV, Theorem 1.5 and Lemma 1.6]{weibel2013kbook} dates the announcement to 1969 and gives the acyclic-map formulation. Vaserstein's 1973 announcement cites Quillen's 1970 ICM account as an existing definition.}
\GSfield{Present status}{As of November 2026: established. The local-coefficient condition is essential to the stated uniqueness formulation.}
\GSfield{Assessment}{literal statement: correct with missing hypotheses; publication match: partial; verification: read the result; development: annotated; confidence: high}
\end{GSitem}
\begin{GSitem}{1971-250.2}{E}{The discrete GL plus model and the missing degree zero}
\GSfield{Kind; source}{definition; \GSpage{P000408}}
\GSfield{The notes}{Morava stresses that \(GL(R)\) is discrete, kills its perfect elementary subgroup, and defines \(K_i(R)=\pi_i(BGL(R)^+)\). The crossed-out line about \(K_0\) is replaced by a note that the definition must change because the space is connected.}
\GSfield{Statement}{For a unital associative ring \(R\), define \(K_i(R)=\pi_i(BGL(R)^+)\) for \(i\ge1\), where plus kills the stable elementary subgroup. A representing space for connective algebraic \(K\)-theory is \(K_0(R)\times BGL(R)^+\), with the group of components \(K_0(R)\). Its fundamental group is \(GL(R)/E(R)\).}
\GSfield{Reconstruction}{The subscript range and component factor are explicitly restored. The source itself notices the degree-zero issue, so this is not counted as an uncorrected false theorem.}
\GSfield{Fate}{This is Quillen's plus definition as summarized in \citet[IV, \S1]{weibel2013kbook}; the comparison with the earlier low groups is included in \citet[first theorem]{vaserstein1973foundations}.}
\GSfield{Present status}{As of November 2026: standard connective algebraic \(K\)-theory; connectedness of \(BGL^+\) remains compatible with nonzero \(K_0\).}
\GSfield{Assessment}{literal statement: correct; publication match: exact; verification: read the result; development: annotated; confidence: high}
\end{GSitem}
\begin{GSitem}{1971-250.3}{E}{Agreement of Quillen K2 with the classical field group}
\GSfield{Kind; source}{theorem; \GSpage{P000408}}
\GSfield{The notes}{The notes attribute to Quillen the equality \(K_2^{\mathrm{Quillen}}(R)=k_2^{\mathrm{Milnor}}(R)\) for a field.}
\GSfield{Statement}{For a field \(F\), Quillen's \(\pi_2BGL(F)^+\) agrees with the classical Steinberg-kernel \(K_2(F)\), and Matsumoto's theorem identifies it with \((F^\times\otimes_{\mathbb Z}F^\times)/\langle a\otimes(1-a):a\ne0,1\rangle\). Agreement with the stable Steinberg-kernel definition holds for unital rings as well; the symbolic presentation stated here is the field case.}
\GSfield{Reconstruction}{The modern distinction between Milnor's original ring \(K_2\) and the Milnor symbol group of a field is made explicit; no assertion is extended to equality of higher Milnor and Quillen groups.}
\GSfield{Fate}{The low-degree identification belongs to the original plus theory; see \citet[IV, \S1]{weibel2013kbook} and the explicit \(K_2\)-identification in \citet[first theorem]{vaserstein1973foundations}.}
\GSfield{Present status}{As of November 2026: established in degree two. The field qualification does not justify replacing all higher algebraic \(K\)-groups by Milnor symbols.}
\GSfield{Assessment}{literal statement: correct; publication match: exact; verification: read the result; development: annotated; confidence: high}
\end{GSitem}
\begin{GSitem}{1971-250.4}{E}{Discrete and topological general linear groups give different theories}
\GSfield{Kind; source}{remark; \GSpage{P000408}}
\GSfield{The notes}{After insisting on discreteness, Morava says that a version for topological groups gives the usual \(K\)-functor for \(\mathbb R\) and \(\mathbb C\), and mentions Volodin's alternative definition.}
\GSfield{Statement}{Using \(GL_n(\mathbb C)\) or \(GL_n(\mathbb R)\) with their analytic topologies and group completion leads to the representing spaces \(\mathbb Z\times BU\) and \(\mathbb Z\times BO\) for connective complex and real topological \(K\)-theory. Using the same fields as discrete rings yields algebraic \(K\)-theory. The continuous inclusion from the discrete group induces a comparison map, not an equality of the theories.}
\GSfield{Reconstruction}{The precise spaces clarify the dependence on topology. Periodic topological \(K\)-theory additionally inverts Bott periodicity; it is not obtained by changing only the notation for \(K_0\).}
\GSfield{Fate}{The distinction is built into the classical classifying-space approach and is reviewed alongside algebraic models in \citet[IV, \S1]{weibel2013kbook}. Volodin's original construction was published in 1971, received July 1970 \citep{volodin1971extraordinary}.}
\GSfield{Present status}{As of November 2026: established comparison of different functors; their values must not be substituted for one another.}
\GSfield{Assessment}{literal statement: correct with missing hypotheses; publication match: partial; verification: read the result; development: annotated; confidence: high}
\end{GSitem}
\begin{GSitem}{1971-250.5}{E}{Cohomology on spaces from the K-theory representing space}
\GSfield{Kind; source}{construction; \GSpage{P000408}}
\GSfield{The notes}{A displayed definition reads \(K_i(R)(Y)=[S^i\wedge Y;X_{E(R)}]\), with a bracket labelled suspension.}
\GSfield{Statement}{For a pointed CW complex \(Y\), stable algebraic \(K\)-cohomology is defined by the connective \(K(R)\)-spectrum. Its reduced groups satisfy \(\widetilde K(R)^{-i}(Y)=[\Sigma^iY,\Omega^\infty K(R)]_*\) for \(i\ge0\), with the appropriate base component when \(\Sigma^iY\) is connected. The full zeroth space includes \(K_0(R)\). An arbitrary pointed space without compatible deloopings does not by itself define a generalized cohomology theory in every degree.}
\GSfield{Reconstruction}{Based homotopy classes, degree-sign convention, component information, and the spectrum structure are supplied. The page is a motivation for a cohomology theory, not a complete construction of its deloopings.}
\GSfield{Fate}{Quillen already announces Segal's coherent categorical cohomology theory in \citet[\S3, p.50]{quillen1971cohomology}, reporting at Nice in 1970 and printed in 1971. The detailed construction appears in \citet[\S\S1--4]{segal1974categories}, received in August 1972. Thus the later publication date does not make the seminar's basic assertion a new post-lecture development.}
\GSfield{Present status}{As of November 2026: standard representability once the infinite-loop-space structure is included.}
\GSfield{Assessment}{literal statement: correct with missing hypotheses; publication match: partial; verification: read the result; development: annotated; confidence: medium}
\end{GSitem}
\begin{GSitem}{1971-250.6}{F}{Novikov’s Chern-character programme}
\GSfield{Kind; source}{problem; \GSpage{P000408}; S. P. Novikov, as reported by Jack Morava}
\GSfield{The notes}{The closing sentence attributes to Novikov the problem of defining a Chern character in algebraic \(K\)-theory and says that here it reduces to computing \(H_*(BGL(R)^+;\mathbb Q)\).}
\GSfield{Statement}{The source does not specify the rings, target theory, grading, naturality requirements, or normalization of a desired character. Computing the rational homology of the plus space is meaningful and equals rational group homology of \(GL(R)\), but it does not uniquely specify a Chern-character map to a prescribed cohomology or homology theory.}
\GSfield{Reconstruction}{No later regulator, Dennis trace, cyclic-homology character, or motivic character is selected as the unique intended solution. Such choices answer distinct precise questions.}
\GSfield{Fate}{Searches for Novikov's algebraic Chern character, the 1972 seminar, Morava, and rational homology gave broad related programmes rather than an exact uniquely specified problem. The homology preservation itself is part of the plus construction \citep[IV, \S1]{weibel2013kbook}.}
\GSfield{Present status}{As of November 2026: the programme is historically recognizable but too broad for a definite problem-index entry or a claim of present-day openness.}
\GSfield{Assessment}{literal statement: not determinable; publication match: thematic only; verification: unresolved lead; development: annotated; confidence: medium}
\end{GSitem}

\GStalk{251}{1972-02-28}{A. M. Gabrielov}{P000409-P000410}{The holomorphic Hodge filtration and Gauss–Manin connection}
Gabrielov continues the exposition of variation of Hodge structure. The notes pass from the Hodge decomposition on a fibre to the relative de Rham complex, its spectral sequence, holomorphic subbundles, and the flat Gauss--Manin connection. The distinction between inverse-image sheaves and module pullback is explicitly emphasized. A strict inequality in the initial filtration and an increasing exact sequence later conflict with the intended decreasing-filtration convention; they are recorded as such rather than silently propagated.

\begin{GSitem}{1971-251.1}{E}{Holomorphic variation of the Hodge filtration}
\GSfield{Kind; source}{theorem; \GSpage{P000409}}
\GSfield{The notes}{The initial formula is \(F^p=\bigoplus_{p'>p}H^{p',q}\), and the notes emphasize that the filtration ``moves holomorphically''. They contrast it with the nonholomorphic behaviour of individual \((p,q)\)-pieces.}
\GSfield{Statement}{For a smooth projective morphism \(f:Y\to S\) of smooth complex varieties, the vector bundle \(\mathcal H^n=(R^nf_*\mathbb C)\otimes\mathcal O_S\) has holomorphic subbundles \(F^p\mathcal H^n\) whose fibre is \(\bigoplus_{r\ge p}H^{r,n-r}(Y_s)\). The Hodge decomposition gives a smooth splitting, generally not a holomorphic splitting into all \(H^{p,q}\).}
\GSfield{Reconstruction}{The modern decreasing convention uses \(r\ge p\), not \(r>p\). A globally shifted convention could be chosen, but it would require changing all later indices consistently.}
\GSfield{Fate}{The exact holomorphic-filtration statement is in \citet[\S2(b)]{griffiths1970periods}, published before the lecture. That survey uses an increasing reindexing, which helps explain but does not remove the notebook's mixed conventions.}
\GSfield{Present status}{As of November 2026: established variation-of-Hodge-structure theory; the literal index requires repair.}
\GSfield{Assessment}{literal statement: false as written; publication match: partial; verification: read the result; development: annotated; confidence: high}
\end{GSitem}
\begin{GSitem}{1971-251.2}{E}{Smooth families and the cohomological local system}
\GSfield{Kind; source}{construction; \GSpage{P000409}}
\GSfield{The notes}{The family is locally projective over a nonsingular base. ``Smooth'' is glossed as local triviality in the analytic sense, followed by a locally trivial holomorphic bundle with fibre \(H^n(Y_s,\mathbb C)\).}
\GSfield{Statement}{A smooth proper map of complex manifolds is locally a differentiable fibre bundle, and \(R^nf_*\mathbb C\) is a finite-rank local system. Its tensor product with \(\mathcal O_S\) is a holomorphic flat vector bundle. A holomorphic submersion is locally a projection near each point of the total space; this does not imply that the entire family is holomorphically a product over a neighbourhood in the base.}
\GSfield{Reconstruction}{The two senses of local triviality are separated. The notebook may mean the local submersion normal form; an assertion of isotriviality of all fibres is not adopted.}
\GSfield{Fate}{\citet[\S1 and \S2(a)]{griffiths1970periods} states the smooth proper projective setting and identifies the flat bundle with the cohomological local system.}
\GSfield{Present status}{As of November 2026: established. Varying complex structures are fully compatible with differentiable local triviality.}
\GSfield{Assessment}{literal statement: correct with missing hypotheses; publication match: partial; verification: read the result; development: annotated; confidence: high}
\end{GSitem}
\begin{GSitem}{1971-251.3}{E}{Inverse-image coefficients and the relative de Rham resolution}
\GSfield{Kind; source}{construction; \GSpage{P000409}}
\GSfield{The notes}{The source distinguishes the set-theoretic inverse-image sheaf \(f^\circ\mathcal O_S\) from the tensor pullback \(f^*\mathcal O_S=\mathcal O_Y\), and writes \(0\to f^\circ\mathcal O_S\to\mathcal O_Y\to\Omega^1_{Y/S}\to\cdots\).}
\GSfield{Statement}{For a holomorphic submersion, let \(f^{-1}\mathcal O_S\) denote the inverse-image sheaf of rings and \(f^*\mathcal O_S=f^{-1}\mathcal O_S\otimes_{f^{-1}\mathcal O_S}\mathcal O_Y\cong\mathcal O_Y\). Define \(\Omega^1_{Y/S}=\Omega_Y^1/\operatorname{im}(f^*\Omega_S^1)\) and exterior powers. The relative holomorphic Poincar\'e lemma makes \(\Omega^\bullet_{Y/S}\) a resolution of \(f^{-1}\mathcal O_S\) on the analytic site.}
\GSfield{Reconstruction}{The two pullback symbols are modernized without erasing the notebook's explicit distinction. The complex is a resolution in the analytic sheaf setting, not an exact sequence of global sections.}
\GSfield{Fate}{This is the relative de Rham construction used in \citet[\S2(b), Grothendieck's formulation]{griffiths1970periods}; it is already part of the theory being reported.}
\GSfield{Present status}{As of November 2026: standard, with the analytic topology and local nature of exactness stated.}
\GSfield{Assessment}{literal statement: correct; publication match: exact; verification: read the result; development: annotated; confidence: high}
\end{GSitem}
\begin{GSitem}{1971-251.4}{E}{Degeneration of the relative Hodge-to-de Rham spectral sequence}
\GSfield{Kind; source}{theorem; \GSpage{P000410}}
\GSfield{The notes}{The page calls the sequence Leray and writes \(E_1^{p,q}=R^qf_*\Omega^p_{Y/S}\), says all differentials vanish, and identifies the resulting filtration by holomorphic subbundles. It also displays \(\mathcal F^n=R^nf_*\Omega^0\) and \(0\to\mathcal F^{p-1}\to\mathcal F^p\to E_\infty^{p,n-p}\to0\).}
\GSfield{Statement}{For a smooth proper projective complex family, \(E_1^{p,q}=R^qf_*\Omega^p_{Y/S}\Rightarrow\mathbb R^{p+q}f_*\Omega^\bullet_{Y/S}\) degenerates at \(E_1\). The terms are locally free and compatible with base change. In decreasing indexing, \(0\to F^{p+1}\mathcal H^n\to F^p\mathcal H^n\to R^{n-p}f_*\Omega^p_{Y/S}\to0\), and \(F^n\mathcal H^n=f_*\Omega^n_{Y/S}\).}
\GSfield{Reconstruction}{The sequence is the relative Hodge-to-de Rham hypercohomology spectral sequence. The corrected indices follow the convention fixed in item 1971-251.1; the source's formulas combine incompatible conventions.}
\GSfield{Fate}{The degeneration and construction of the Hodge filtration are the prior Hodge-theoretic input described in \citet[\S2(b)]{griffiths1970periods}. The scan confirms the reversed direct-image and differential-form indices.}
\GSfield{Present status}{As of November 2026: established for the stated family; the displayed indices in the notes are false under their own initial decreasing convention.}
\GSfield{Assessment}{literal statement: false as written; publication match: partial; verification: read the result; development: annotated; confidence: high}
\end{GSitem}
\begin{GSitem}{1971-251.5}{E}{The Gauss–Manin connection}
\GSfield{Kind; source}{construction; \GSpage{P000410}}
\GSfield{The notes}{The last lines identify \(R^nf_*f^\circ\mathcal O_S\) with \(R^nf_*\mathbb C\otimes\mathcal O_S\), mark locally constant cohomology classes as sections, and display a connection into \(\Omega_S^1\otimes\mathcal H^n\).}
\GSfield{Statement}{The cohomological local system defines the flat holomorphic connection \(\nabla:\mathcal H^n\to\Omega_S^1\otimes\mathcal H^n\), characterized by having \(R^nf_*\mathbb C\) as horizontal sections. In the full variation theorem it satisfies Griffiths transversality \(\nabla F^p\subseteq\Omega_S^1\otimes F^{p-1}\).}
\GSfield{Reconstruction}{Flatness follows from the local-system construction. The transversality inclusion is supplied as context from the cited theorem; it is not claimed to be written on these two notebook pages.}
\GSfield{Fate}{\citet[\S2(a)--(c), equation (2.1)]{griffiths1970periods} provides the flat connection and infinitesimal period relation. Both predate the lecture.}
\GSfield{Present status}{As of November 2026: standard Gauss--Manin and Hodge-filtration structure; the notebook's final formula is its beginning, not a complete proof.}
\GSfield{Assessment}{literal statement: correct with missing hypotheses; publication match: exact; verification: read the result; development: annotated; confidence: high}
\end{GSitem}

\GStalk{252}{1972-03-06}{M. Sh. Birman}{P000411-P000412}{Scattering channels, trace ideals, and the spectral shift}
Birman starts with the Lax--Phillips exterior-wave picture, then gives the abstract definition of wave operators and asks for completeness. The principal results reported are the trace-class theorem, its invariance and resolvent-comparison extensions, the Birman--Krein determinant formula, and localized trace-ideal criteria for one or two Hilbert spaces. All have antecedents before this lecture. The compressed opening must distinguish conserved total energy from decay of a compressed or spatially localized evolution; the closing two-space statement must replace equality of domains in different spaces by a condition on the identification operator.

\begin{GSitem}{1971-252.1}{E}{Exterior waves and the Lax–Phillips interaction space}
\GSfield{Kind; source}{construction; \GSpage{P000411}}
\GSfield{The notes}{For the exterior Dirichlet wave equation \(u_{tt}=\Delta u\), the notes assume Huygens' principle and write the conserved energy \(\int_\Omega(|\nabla u|^2+|u_t|^2)\,dx\). Incoming and outgoing spaces \(K_-,K_+\) leave a complement \(K_0\), described as exponentially decaying solutions concentrated near the boundary.}
\GSfield{Statement}{For a smooth compact obstacle in odd-dimensional Euclidean space of dimension at least three, the Dirichlet wave group acts unitarily on its finite-energy space. Suitable incoming and outgoing subspaces \(D_-,D_+\) are orthogonal, and the interaction space is \(\mathcal K=\mathcal H\ominus(D_+\oplus D_-)\). The compressed evolution on \(\mathcal K\) is a contraction semigroup. Resonant states and local-energy decay belong to this reduced description. A nonzero finite-energy solution cannot decay in its full energy norm, which is conserved; uniform exponential local decay requires further geometric and analytic hypotheses not stated here.}
\GSfield{Reconstruction}{The exterior region, dimension, boundary regularity, and distinction between total and local or compressed energy are supplied. The page does not justify an exponential-decay theorem for every obstacle or every vector of \(K_0\).}
\GSfield{Fate}{\citet[Preface, pp.vii--ix; Introduction, pp.1--2]{laxphillips1967scattering} explicitly describes the incoming/outgoing spaces, odd-dimensional orthogonality, compressed semigroup, and conserved total energy. The preview contains these primary pages. The later \citet{laxphillips1969decaying} treats the exterior-wave resonances; its abstract was checked, not its full proof.}
\GSfield{Present status}{As of November 2026: the framework is established; the notebook's unrestricted exponential-decay wording remains too strong to use as a theorem.}
\GSfield{Assessment}{literal statement: false as written; publication match: partial; verification: read the result; development: annotated; confidence: medium}
\end{GSitem}
\begin{GSitem}{1971-252.2}{E}{Complex poles and the continued resolvent}
\GSfield{Kind; source}{remark; \GSpage{P000411}}
\GSfield{The notes}{The list of keywords includes characteristic functions, the second sheet of the resolvent, the scattering matrix, and complex poles of a self-adjoint operator; Gohberg--Krein and Lax--Phillips are mentioned.}
\GSfield{Statement}{For self-adjoint \(H\), the operator-valued resolvent \((H-z)^{-1}\) is analytic for \(\operatorname{Im}z\ne0\). Nonreal scattering poles concern a meromorphic continuation of suitably localized matrix elements or resolvents, or the spectrum of a non-self-adjoint compressed generator. The existence, domain, and sheets of continuation depend on the scattering problem; the phrase cannot mean nonreal eigenvalues of \(H\) itself.}
\GSfield{Reconstruction}{This distinguishes the objects in the keyword list without selecting a unique analytic continuation or a particular Gohberg--Krein theorem that the notes do not identify.}
\GSfield{Fate}{The exterior-wave resonances treated by \citet{laxphillips1969decaying} supply the identifiable prior setting. The bare phrase about complex poles has only a partial match to that theory, not the status of a separate spectral theorem.}
\GSfield{Present status}{As of November 2026: standard resonance terminology after the continued or compressed object has been specified.}
\GSfield{Assessment}{literal statement: correct with missing hypotheses; publication match: partial; verification: publisher metadata checked; development: annotated; confidence: medium}
\end{GSitem}
\begin{GSitem}{1971-252.3}{E}{Wave operators and the question of asymptotic completeness}
\GSfield{Kind; source}{definition; \GSpage{P000411}}
\GSfield{The notes}{Asymptotic closeness means \(\|e^{-itH_0}f-e^{-itH}g_\pm\|\to0\). The map \(f\mapsto g_\pm=U_\pm f\) is a partial isometry intertwining \(H_0\) and \(H\). The explicit question asks when its range exhausts the absolutely continuous subspace of \(H\).}
\GSfield{Statement}{For self-adjoint operators on one Hilbert space, define \(W_\pm(H,H_0)=\operatorname{s-lim}_{t\to\pm\infty}e^{itH}e^{-itH_0}P_{\mathrm{ac}}(H_0)\) when the limits exist. Their initial space is \(\mathcal H_{\mathrm{ac}}(H_0)\), their ranges lie in \(\mathcal H_{\mathrm{ac}}(H)\), and they intertwine the corresponding spectral calculi. Completeness means \(\operatorname{Ran}W_\pm=\mathcal H_{\mathrm{ac}}(H)\). Existence alone does not imply completeness.}
\GSfield{Reconstruction}{The absolutely continuous projection and strong operator topology are restored. The broad request is followed in the same lecture by sufficient criteria; it is not treated as one unspecified modern open problem.}
\GSfield{Fate}{The formulation and trace-class completeness result appear already in \citet[Theorem 1]{kato1957trace}; the two-space framework is explicit in \citet[\S1]{belopolskiiBirman1968wave}.}
\GSfield{Present status}{As of November 2026: standard definitions and a family of precise completeness theorems. No universal criterion is claimed by the notebook.}
\GSfield{Assessment}{literal statement: correct with missing hypotheses; publication match: exact; verification: read the result; development: annotated; confidence: high}
\end{GSitem}
\begin{GSitem}{1971-252.4}{E}{Orthogonal channels and multichannel completeness}
\GSfield{Kind; source}{definition; \GSpage{P000411-P000412}}
\GSfield{The notes}{The notes introduce \(H_0^{(1)},\ldots,H_0^{(n)}\), requiring the ranges of the channel maps \(U_+^{(i)}\) to be mutually orthogonal and to generate the absolutely continuous part of \(H\).}
\GSfield{Statement}{Let \(H_j\) be self-adjoint channel Hamiltonians on \(\mathcal H_j\), let \(J_j:\mathcal H_j\to\mathcal H\) be bounded identifications, and suppose the channel wave operators \(W_\pm^{(j)}=\operatorname{s-lim}_{t\to\pm\infty}e^{itH}J_je^{-itH_j}P_j\) exist as isometries on their designated channel spaces \(P_j\mathcal H_j\). Orthogonality and completeness mean \(\mathcal H_{\mathrm{ac}}(H)=\bigoplus_j\operatorname{Ran}W_\pm^{(j)}\). The channel spaces and identifications are part of the hypotheses.}
\GSfield{Reconstruction}{The direct-sum condition formalizes the definition in the notes; it does not prove completeness for an arbitrary collection of Hamiltonians or interactions.}
\GSfield{Fate}{This is the established multichannel language of the scattering theory surveyed by the lecture. The definition of two-space wave maps, used for each channel, is in \citet[\S1]{belopolskiiBirman1968wave}; the notes themselves do not state a general many-body completeness theorem.}
\GSfield{Present status}{As of November 2026: standard definition; actual completeness depends on the particular system and its channels.}
\GSfield{Assessment}{literal statement: correct with missing hypotheses; publication match: partial; verification: read the result; development: annotated; confidence: high}
\end{GSitem}
\begin{GSitem}{1971-252.5}{F}{Faddeev’s three-particle problem and the Lax–Phillips scheme}
\GSfield{Kind; source}{remark; \GSpage{P000412}}
\GSfield{The notes}{A single sentence says that Faddeev placed the three-particle problem in the Lax--Phillips scheme.}
\GSfield{Statement}{No interaction potentials, decay assumptions, threshold exclusions, channel decomposition, or map into a Lax--Phillips representation are specified. Thus the line identifies a subject and an attribution but does not determine an exact equivalence or completeness theorem.}
\GSfield{Reconstruction}{Faddeev's three-particle scattering theory is a firm historical lead. The additional assertion about a particular embedding into the Lax--Phillips scheme is left unresolved rather than replaced by a modern general many-body theorem.}
\GSfield{Fate}{The primary publication \citet{faddeev1963three} is verified as a 1963 monograph on the stated three-particle problem. Searches combining Faddeev, three particles, and Lax--Phillips did not locate the exact asserted embedding in a source that could be read here. The primary preface of \citet[p.ix]{laxphillips1967scattering} instead discusses Faddeev's three-body work as a steady-state approach distinct from the book's hyperbolic method; this does not exclude a later connection, but it does not verify the notebook's precise attribution.}
\GSfield{Present status}{As of November 2026: the broad attribution is recognizable; the specific formulation and its fate cannot be determined from this sentence.}
\GSfield{Assessment}{literal statement: not determinable; publication match: thematic only; verification: unresolved lead; development: annotated; confidence: medium}
\end{GSitem}
\begin{GSitem}{1971-252.6}{E}{Kato’s trace-class perturbation theorem}
\GSfield{Kind; source}{theorem; \GSpage{P000412}}
\GSfield{The notes}{The line ``Kato (1957)'' is followed by \(H-H_0\) nuclear.}
\GSfield{Statement}{Let \(H_0\) be self-adjoint and let \(V\) be a bounded self-adjoint trace-class operator. Then \(H=H_0+V\), on \(D(H_0)\), is self-adjoint and both wave operators \(W_\pm(H,H_0)\) exist and are complete. They implement unitary equivalence of the absolutely continuous parts, without requiring the whole spectra to be absolutely continuous.}
\GSfield{Reconstruction}{Boundedness, the common operator domain, and the restriction to the absolutely continuous parts are supplied. Nuclear means the Schatten trace ideal \(\mathfrak S_1\) in this Hilbert-space setting.}
\GSfield{Fate}{This is exactly \citet[Theorem 1, pp.260--261]{kato1957trace}, published in 1957. The paper itself acknowledges Rosenblum's contemporaneous related result; the lecture reports the theorem under Kato's name.}
\GSfield{Present status}{As of November 2026: established as the Kato--Rosenblum trace-class theorem.}
\GSfield{Assessment}{literal statement: correct with missing hypotheses; publication match: exact; verification: read the result; development: annotated; confidence: high}
\end{GSitem}
\begin{GSitem}{1971-252.7}{E}{Birman’s invariance principle and resolvent comparability}
\GSfield{Kind; source}{theorem; \GSpage{P000412}}
\GSfield{The notes}{The next line, dated 1962 and attributed to Birman, replaces \(H-H_0\) by \(\varphi(H)-\varphi(H_0)\) in the nuclear condition.}
\GSfield{Statement}{A precise sufficient version is: if \(H,H_0\) are self-adjoint and \((H-z)^{-1}-(H_0-z)^{-1}\in\mathfrak S_1\) for one nonreal \(z\), then the ordinary wave operators exist and are complete. More general invariance principles permit suitable monotone changes of spectral variable with additional regularity and localization hypotheses. Trace-class difference for an arbitrary function \(\varphi\) cannot suffice: a constant \(\varphi\) makes it zero for every pair.}
\GSfield{Reconstruction}{The verified resolvent theorem replaces the underqualified function shorthand. The constant-function observation shows why the missing hypotheses are substantive; no maximal class of admissible functions is inferred.}
\GSfield{Fate}{\citet[Theorem 2 and the following condition \(B_n\), pp.476--477]{birmanKrein1962wave} prove the stated resolvent-comparison result and describe further semibounded inverse-power conditions.}
\GSfield{Present status}{As of November 2026: established under the specified comparison hypothesis. The unrestricted functional-calculus formulation is false.}
\GSfield{Assessment}{literal statement: false as written; publication match: partial; verification: read the result; development: annotated; confidence: high}
\end{GSitem}
\begin{GSitem}{1971-252.8}{E}{The scattering operator and its spectral fibres}
\GSfield{Kind; source}{construction; \GSpage{P000412}}
\GSfield{The notes}{The source defines \(S=U_+^*U_-\) and calls \(S(\lambda)\) the matrix in the spectral decomposition of \(H_0\).}
\GSfield{Statement}{When both wave operators exist and are complete, \(S=W_+^*W_-\) is a unitary operator on \(\mathcal H_{\mathrm{ac}}(H_0)\). It commutes with the spectral measure of \(H_0\), hence in a direct-integral representation it acts by a measurable family of unitaries \(S(\lambda)\), defined almost everywhere. These fibres may act on infinite-dimensional multiplicity spaces; the word matrix does not imply finite dimension.}
\GSfield{Reconstruction}{Completeness, the almost-everywhere convention, and possible infinite spectral multiplicity are added.}
\GSfield{Fate}{This is the construction used explicitly in \citet[\S3, pp.476--477]{birmanKrein1962wave} and \citet[\S6]{belopolskiiBirman1968wave}.}
\GSfield{Present status}{As of November 2026: standard. Its determinant requires the extra trace-ideal conclusion in the following item.}
\GSfield{Assessment}{literal statement: correct with missing hypotheses; publication match: exact; verification: read the result; development: annotated; confidence: high}
\end{GSitem}
\begin{GSitem}{1971-252.9}{E}{The spectral shift and Birman–Krein determinant formula}
\GSfield{Kind; source}{theorem; \GSpage{P000412}}
\GSfield{The notes}{The displayed formulas are \(S(\lambda)-I\) nuclear, \(\det S(\lambda)=e^{-2\pi i\xi(\lambda)}\), and \(\operatorname{Tr}(\varphi(H)-\varphi(H_0))=\int\varphi'(\lambda)\xi(\lambda)\,d\lambda\); the latter is said to define \(\xi\).}
\GSfield{Statement}{If \(V=H-H_0\) is bounded self-adjoint and trace class, there is a unique real \(\xi\in L^1(\mathbb R)\) with \(\operatorname{Tr}(f(H)-f(H_0))=\int f'(\lambda)\xi(\lambda)\,d\lambda\) for smooth compactly supported \(f\), and \(\int\xi=\operatorname{Tr}V\). Almost everywhere on the absolutely continuous spectrum, \(S(\lambda)-I\in\mathfrak S_1\) and its Fredholm determinant equals \(\exp(-2\pi i\xi(\lambda))\). For merely resolvent-comparable pairs one uses a weighted-integrable spectral shift with an additional normalization; the determinant alone determines it only modulo integers.}
\GSfield{Reconstruction}{A sufficient test-function class and the normalization of the trace-class case are specified. The sign in the scan agrees with the convention \(S=W_+^*W_-\) and with the cited paper.}
\GSfield{Fate}{\citet[\S4, formulas (8)--(9), and \S5, Theorem 3, formula (11), p.477]{birmanKrein1962wave} gives the spectral-shift normalization and determinant formula, distinguishing trace-class perturbations from resolvent comparability.}
\GSfield{Present status}{As of November 2026: established; equality holds almost everywhere and the determinant is the Fredholm determinant, not a finite-matrix determinant unless multiplicity is finite.}
\GSfield{Assessment}{literal statement: correct with missing hypotheses; publication match: exact; verification: read the result; development: annotated; confidence: high}
\end{GSitem}
\begin{GSitem}{1971-252.10}{E}{Localized trace-class conditions for existence}
\GSfield{Kind; source}{theorem; \GSpage{P000412}}
\GSfield{The notes}{With \(D(H)=D(H_0)\), the notes require \((H-H_0)E^0(\Delta)\in\mathcal J_1\) for every bounded interval \(\Delta\), and conclude existence of wave operators.}
\GSfield{Statement}{Suppose self-adjoint \(H,H_0\) act on one Hilbert space with equal domains, and \((H-H_0)E_{H_0}(\Delta)\) extends to a trace-class operator for every bounded interval \(\Delta\). Then \(W_\pm(H,H_0)\) exist as isometries on \(\mathcal H_{\mathrm{ac}}(H_0)\). This one-sided condition yields existence; completeness follows, for example, if the corresponding reverse conditions also hold, so both inverse wave maps exist.}
\GSfield{Reconstruction}{The bounded extensions and the distinction between existence and completeness are stated. The spectral projections exhaust the absolutely continuous subspace and provide the increasing family used in the cited theorem.}
\GSfield{Fate}{This is the identity-identification specialization of \citet[Theorem 5.1, p.1126]{belopolskiiBirman1968wave}, received March 1968 and published that year.}
\GSfield{Present status}{As of November 2026: established; the lecture correctly labels this conclusion as existence, not automatically completeness.}
\GSfield{Assessment}{literal statement: correct with missing hypotheses; publication match: exact; verification: read the result; development: annotated; confidence: high}
\end{GSitem}
\begin{GSitem}{1971-252.11}{E}{Two-space wave operators and the identification operator}
\GSfield{Kind; source}{theorem; \GSpage{P000412}}
\GSfield{The notes}{For different Hilbert spaces the notes write \(U_\pm=\lim e^{itH}Je^{-itH_0}\), retain \(D(H)=D(H_0)\), and require localized \(JH_0-HJ\) to be nuclear and \((J^*J-I)E^0(\Delta)\) compact. The final remark says that choosing \(J=U_+\) is trivial, so one must approximately guess the identification.}
\GSfield{Statement}{Let \(H_0,H\) be self-adjoint on \(\mathcal H_0,\mathcal H\), and let \(J:\mathcal H_0\to\mathcal H\) be bounded with \(J D(H_0)\subset D(H)\). Suppose for every bounded interval \(\Delta\) that \((HJ-JH_0)E_{H_0}(\Delta)\) has a trace-class extension and \((J^*J-I)E_{H_0}(\Delta)\) is compact. Then \(W_\pm(H,H_0;J)=\operatorname{s-lim}e^{itH}Je^{-itH_0}P_{\mathrm{ac}}(H_0)\) exist and are isometric on \(\mathcal H_{\mathrm{ac}}(H_0)\). Completeness requires further hypotheses, such as suitable reverse comparison conditions.}
\GSfield{Reconstruction}{Domain equality across different spaces is replaced by the required domain compatibility. These hypotheses are a convenient stronger specialization of the theorem's exhausting projections. The last sentence is a method for choosing \(J\), not an existence proof using an operator not yet constructed.}
\GSfield{Fate}{\citet[Theorem 5.1, pp.1126--1127]{belopolskiiBirman1968wave} proves the existence and isometry assertion; its Theorem 5.2 supplies stronger conditions for unitarity onto the absolutely continuous part. The manuscript was received on 20 March 1968, four years before the talk.}
\GSfield{Present status}{As of November 2026: established two-Hilbert-space scattering theory, with domain and completeness issues separated.}
\GSfield{Assessment}{literal statement: false as written; publication match: partial; verification: read the result; development: annotated; confidence: high}
\end{GSitem}

\GStalk{253}{1972-03-06}{L. N. Vaserstein}{P000413-P000416}{Comparison of K-theories, suspension rings, and Laurent boundaries}
This substantial lecture compares Volodin's simplicial construction with the derived-functor theories of Gersten and Swan, announces the answer to the elementary-relation question, and develops infinite-matrix suspension, negative groups, relative groups, and polynomial Nil terms. Quillen's comparison is still presented as Vaserstein's expectation here. The final Laurent construction uses Whitehead's elementary factorization and one-sided compression. Many assertions reappear, with fuller definitions and proofs, in Vaserstein's 1973 announcement and 1976 article. The explicit Kazhdan heading at the bottom of P000416 ends this record.

\begin{GSitem}{1971-253.1}{E}{Volodin’s simplicial general linear group}
\GSfield{Kind; source}{construction; \GSpage{P000413}}
\GSfield{The notes}{The first theory is written \(K_n=\pi_{n-1}(GL(A))\), with a marginal attribution to Volodin and a reference to an abstract simplicial scheme. The first two groups are said to agree with Whitehead and Milnor.}
\GSfield{Statement}{For a unital ring \(A\), give the stable group \(GL(A)\) the simplicial structure whose finite simplices lie in a left coset of a simultaneously triangular unipotent subgroup, allowing finite permutations of the basis. If \(S(A)\) denotes its realization, Volodin's model has \(K_n^V(A)=\pi_{n-1}S(A)\) for \(n\ge1\), with components for \(n=1\). This is not ordinary homotopy of the discrete group, whose positive homotopy groups would vanish.}
\GSfield{Reconstruction}{The simplicial topology and degree shift are supplied from the named model; the displayed symbol \(GL(A)\) is shorthand for this extra structure.}
\GSfield{Fate}{The construction was published in \citet[\S\S1--3]{volodin1971extraordinary}, received in July 1970. The explicit scheme convention and comparisons are given in \citet[\S\S11--15]{vaserstein1976foundations}.}
\GSfield{Present status}{As of November 2026: an established model equivalent to Quillen's theory in the unital setting; the later comparison is separately indexed below.}
\GSfield{Assessment}{literal statement: correct with missing hypotheses; publication match: exact; verification: read the result; development: annotated; confidence: high}
\end{GSitem}
\begin{GSitem}{1971-253.2}{E}{Gersten–Swan derived functors and free simplicial resolutions}
\GSfield{Kind; source}{construction; \GSpage{P000413}}
\GSfield{The notes}{The source names Swan and ``Gershtein'', applies \(GL\) to a free-ring resolution, and observes that exactness is lost. It explicitly chooses rings without identity, \(F(A)=\mathbb Z\langle A\rangle_+\), the noncommutative polynomials without a constant term.}
\GSfield{Statement}{The free nonunital ring on the underlying set of \(A\) maps onto \(A\) by evaluation. Iterating the free-ring adjunction produces a simplicial free resolution. Applying stable linear or elementary-group functors gives the nonabelian derived-functor constructions associated with Gersten and Swan; a three-term sequence alone does not define all their derived groups. The unitization of \(F(A)\) is the unital free ring \(\mathbb Z\langle A\rangle\).}
\GSfield{Reconstruction}{The surname is normalized to Gersten. The full simplicial resolution is distinguished from the abbreviated chain in the notes, and the omitted identity is a substantive convention rather than a harmless typographical omission.}
\GSfield{Fate}{\citet[\S29, pp.142--143]{vaserstein1976foundations} identifies the antecedent constructions as Swan (1970) and Gersten (1971), and explains their equivalence via Keune's formulations. This verifies the framework retrospectively; the notebook is not the first definition.}
\GSfield{Present status}{As of November 2026: established antecedent constructions, with the later comparison stated using a suitable fibre-product functor in the next items.}
\GSfield{Assessment}{literal statement: correct with missing hypotheses; publication match: partial; verification: read the result; development: annotated; confidence: high}
\end{GSitem}
\begin{GSitem}{1971-253.3}{A}{Acyclicity of free rings in reduced K-theory}
\GSfield{Kind; source}{theorem; \GSpage{P000413}}
\GSfield{The notes}{The notes say \(K_1\) of a free ring is zero, and propose induction from an exact free presentation to define the higher groups.}
\GSfield{Statement}{For any set \(S\), the inclusion \(\mathbb Z\to R=\mathbb Z\langle S\rangle\) induces \(K_n(\mathbb Z)\cong K_n(R)\) for all integers \(n\), using nonconnective \(K\)-theory in negative degrees. Thus the augmented or reduced theory of \(R\), equivalently the chosen theory of its free nonunital augmentation ideal, vanishes. The unreduced statement \(K_1(R)=0\) is false: it contains \(K_1(\mathbb Z)\cong\{\pm1\}\).}
\GSfield{Reconstruction}{The reduced convention already indicated by ``without identity'' is made explicit. The full all-degree theorem is the published form of the lecture's acyclicity mechanism, not an assertion that the three written arrows prove it.}
\GSfield{Fate}{\citet[second theorem, p.231]{vaserstein1973foundations} announces the vanishing; \citet[Theorem 28.1 and Corollary 28.1, p.142]{vaserstein1976foundations} states the precise unital and reduced forms. The latter cites Gersten's higher-degree theorem and explicitly does not reproduce its long proof.}
\GSfield{Present status}{As of November 2026: established. The 1973 note was received on 14 November 1972; the 1976 article supplies the exact formulation and references for its proof.}
\GSfield{Assessment}{literal statement: correct with missing hypotheses; publication match: exact; verification: read the result; development: annotated; confidence: high}
\end{GSitem}
\begin{GSitem}{1971-253.4}{A}{The announced solution of the universal Steinberg-relation problem}
\GSfield{Kind; source}{claim; \GSpage{P000413}}
\GSfield{The notes}{Vaserstein says that he has solved the problem about Steinberg relations and that the solution follows from comparison with Volodin's functors.}
\GSfield{Statement}{For the augmented free ring \(R=\mathbb Z\langle S\rangle\), the kernel of \(K_2(R)\to K_2(\mathbb Z)\) vanishes. Hence a universally valid matrix relation whose Steinberg word specializes to the identity when all variables are zero is generated by Steinberg relations. This is the normalized precise version of the question recorded on P000401.}
\GSfield{Reconstruction}{The statement uses the augmentation convention established in item 1971-247.5. Neither vanishing of \(K_2\) for every ring nor a presentation of every finite-rank elementary group follows.}
\GSfield{Fate}{This occurrence is the notebook's answer to 1971-247.5. The free-ring theorem in \citet[second theorem]{vaserstein1973foundations} and \citet[Theorem 28.1]{vaserstein1976foundations} supports that answer, with the reconstructed normalization stated explicitly.}
\GSfield{Present status}{As of November 2026: resolved for the normalized universal question; it is retained in the problem index to connect the February question with the March answer.}
\GSfield{Assessment}{literal statement: not determinable; publication match: partial; verification: read the result; development: annotated; confidence: medium; problem \hyperref[prob:GP-1971-20]{GP-1971-20}}
\end{GSitem}
\begin{GSitem}{1971-253.5}{A}{Cone matrices, Morita invariance, and suspension}
\GSfield{Kind; source}{construction; \GSpage{P000413-P000414}}
\GSfield{The notes}{The ring \(MA\) has finitely many nonzero entries in total; \(\ell A\) has finitely many in each row and each column. The notes assert \(K_n(\ell A)=0\), define \(\mu A=\ell A/MA\), and obtain \(K_n(A)=K_n(MA)=K_{n+1}(\mu A)\). The proof is compared with Kuiper's theorem.}
\GSfield{Statement}{Let \(C(A)\) be the ring of countable row-and-column-finite matrices over a unital associative ring \(A\), and let \(M_\infty(A)\) be its ideal of finitely supported matrices. Put \(\Sigma A=C(A)/M_\infty(A)\). In the algebraic suspension construction, \(K_n(C(A))=0\), \(K_n(M_\infty(A))\cong K_n(A)\), and \(K_{n+1}(\Sigma A)\cong K_n(A)\), in all integer degrees after negative groups have been included.}
\GSfield{Reconstruction}{These are precisely the matrix classes in the 1973 announcement. No finite-band restriction is imposed on this cone ring: the bandwidth condition in record 247 applies to multiplicativity of a different two-sided truncation. Contractibility here means vanishing in the specified \(K\)-theory, not contractibility of a discrete matrix group as a topological space. The locally unital finite-matrix ideal permits the suspension exact sequence.}
\GSfield{Fate}{\citet[first theorem, p.231]{vaserstein1973foundations} gives these identities with the stated cone definition; \citet[\S\S18--21]{vaserstein1976foundations} supplies the asphericity, Morita, and suspension arguments. These are part of a contemporary algebraic suspension programme, not a claim that infinite matrices first entered \(K\)-theory in this lecture.}
\GSfield{Present status}{As of November 2026: established. The analytic analogy with Kuiper does not change the algebraic definition or provide the missing proof by itself.}
\GSfield{Assessment}{literal statement: correct with missing hypotheses; publication match: exact; verification: read the result; development: annotated; confidence: high}
\end{GSitem}
\begin{GSitem}{1971-253.6}{E}{Recovering K0 and Bass’s negative groups}
\GSfield{Kind; source}{construction; \GSpage{P000414}}
\GSfield{The notes}{The suspension formula yields the ordinary full \(K_0\), with a free module nonzero, by \(K_0(A)=K_1(\mu A)\). Iteration extends \(K_n\) to all integers and is said to coincide with Bass's groups.}
\GSfield{Statement}{The suspension model gives \(K_0(A)\cong K_1(\Sigma A)\) and \(K_{-r}(A)\cong K_0(\Sigma^r A)\) for \(r\ge1\). Here \(K_0\) is the full Grothendieck group of finitely generated projectives, not the quotient by free modules. The resulting negative groups agree with Bass's negative \(K\)-groups; for regular noetherian rings they vanish.}
\GSfield{Reconstruction}{The phrase about a free module excludes reduced \(\widetilde K_0\). The regular-noetherian vanishing is supplied as established context, not claimed to be written on P000414.}
\GSfield{Fate}{Bass's negative theory predates the talk (1968); its comparison with the suspension definition is explicitly recorded in \citet[first theorem and references 10--12]{vaserstein1973foundations} and proved in the framework of \citet[\S21]{vaserstein1976foundations}.}
\GSfield{Present status}{As of November 2026: standard nonconnective extension. Positive-degree polynomial invariance does not follow merely from the existence of these negative groups.}
\GSfield{Assessment}{literal statement: correct with missing hypotheses; publication match: exact; verification: read the result; development: annotated; confidence: high}
\end{GSitem}
\begin{GSitem}{1971-253.7}{A}{Polynomial Nil terms and the distinction from Karoubi–Villamayor theory}
\GSfield{Kind; source}{construction; \GSpage{P000414-P000415}}
\GSfield{The notes}{The notes define \(\operatorname{Nil}_n A=K_n(tA[t])\), contrast this with a theory that contracts \(tA[t]\), and use \(t(t-1)A[t]\to tA[t]\xrightarrow{t=1}A\) as Karoubi's loop construction.}
\GSfield{Statement}{For a unital ring, define \(NK_n(A)=\ker(K_n(A[t])\xrightarrow{t=0}K_n(A))\). The constant inclusion splits evaluation, so \(K_n(A[t])\cong K_n(A)\oplus NK_n(A)\). Interpreting the notebook's \(K_n(tA[t])\) requires the corresponding augmented or relative convention. Polynomial homotopy theories force the polynomial path ring to be contractible; ordinary algebraic \(K\)-theory can retain its Nil contribution. The kernel of evaluation at one on \(tA[t]\) is \(t(t-1)A[t]\).}
\GSfield{Reconstruction}{The modern \(NK\) notation distinguishes these groups from differently indexed categories of nilpotent endomorphisms. The equality with a nonunital absolute functor is not used as an assertion of excision for arbitrary ideals.}
\GSfield{Fate}{\citet[third theorem, p.231]{vaserstein1973foundations} states the polynomial and Laurent decompositions with these Nil terms. The distinction from polynomial-path \(K\)-theory is discussed in \citet[Introduction, p.96]{vaserstein1976foundations}; the underlying Karoubi--Villamayor theory had circulated before the lecture.}
\GSfield{Present status}{As of November 2026: established distinction; Nil groups are exactly the obstruction in this polynomial splitting, not an omission to be discarded for every ring.}
\GSfield{Assessment}{literal statement: correct with missing hypotheses; publication match: partial; verification: read the result; development: annotated; confidence: high}
\end{GSitem}
\begin{GSitem}{1971-253.8}{A}{Relative groups and the correction from kernels to fibre products}
\GSfield{Kind; source}{construction; \GSpage{P000413-P000414}}
\GSfield{The notes}{The source first writes \(K_{n+1}(A)=K_n(A'')\) after calling \(A''\) the kernel of a free presentation. It later proposes a pullback ring \(F\) giving an exact sequence \(\cdots\to K_n(A)\to K_n(B)\to K_n(F)\to K_{n-1}(A)\to\cdots\), with an ambiguous \(\mu B\)-label on an inclusion into \(\ell B\).}
\GSfield{Statement}{For a homomorphism \(f:A\to B\), define relative \(K\)-theory by the homotopy fibre of \(K(A)\to K(B)\); its groups fit \(\cdots\to K_n(A,B)\to K_n(A)\to K_n(B)\to K_{n-1}(A,B)\to\cdots\). In Vaserstein's ring model let \(F(B)\to B\) be the free nonunital presentation and put \(D_f=A\times_B F(B)\). His relative group is represented by \(K_n(D_f)\). For dimension shifting he uses \(D(A)=F(A)\times_A F(A)\), with \(K_n(D(A))\cong K_{n+1}(A)\).}
\GSfield{Reconstruction}{The fibre sequence fixes the indexing: a ring sitting after \(K_n(B)\) in the notes represents a shifted relative object. The published doubled fibre product is not simply the presentation kernel. The drawn inclusion cannot have the earlier quotient \(\mu B\) as its evident source; its uncertain letter is left unresolved, and the verified published model is supplied separately.}
\GSfield{Fate}{\citet[second theorem, p.231]{vaserstein1973foundations} gives the relative pullback model. \citet[\S20, Theorem 29.1 and its proof, pp.142--143]{vaserstein1976foundations} identifies the doubled fibre product and proves dimension shifting. General excision is not assumed.}
\GSfield{Present status}{As of November 2026: established relative and dimension-shift constructions. The literal kernel recursion and the unqualified pullback diagram are not sufficient definitions.}
\GSfield{Assessment}{literal statement: false as written; publication match: partial; verification: read the result; development: annotated; confidence: high}
\end{GSitem}
\begin{GSitem}{1971-253.9}{E}{The plus construction does not automatically make every group an H-space}
\GSfield{Kind; source}{claim; \GSpage{P000415}}
\GSfield{The notes}{For any group \(G\) with perfect commutator subgroup \(E\), the notes attach cells to \(BG\), kill \(E\), preserve homology, and say that this turns the result into a group or H-space.}
\GSfield{Statement}{The plus construction \(BG\to BG_E^+\) kills the specified perfect normal subgroup and preserves homology with the appropriate local coefficients. This operation alone does not provide an H-space multiplication for every such \(G\). For the stable general linear group, compatible block sums supply the additional multiplication, and Quillen's positive groups are \(\pi_n(BGL(A)^+)\).}
\GSfield{Reconstruction}{The general existence of plus is separated from the stable-linear source of its multiplication. Homotopy groups of a general plus space are not automatically ring \(K\)-groups without a ring and the relevant model.}
\GSfield{Fate}{Quillen's already circulated account \citep[\S3, p.50]{quillen1971cohomology} treats the stable general linear group and its coherent categorical structure. The general plus theorem is given with its correct hypotheses in \citet[IV, Theorem 1.5]{weibel2013kbook}.}
\GSfield{Present status}{As of November 2026: standard after this distinction; the universal H-space claim is an overstatement in the notebook.}
\GSfield{Assessment}{literal statement: false as written; publication match: partial; verification: read the result; development: annotated; confidence: high}
\end{GSitem}
\begin{GSitem}{1971-253.10}{F}{The assertion that everything extends to the unitary case}
\GSfield{Kind; source}{remark; \GSpage{P000415}}
\GSfield{The notes}{One emphatic sentence says that everything carries over to the unitary case.}
\GSfield{Statement}{The preceding discussion contains several inequivalent assertions: plus constructions, comparison of models, suspension, relative groups, and Laurent decompositions. A unitary or hermitian version must specify a ring with involution, symmetry and form parameter, degree conventions, and often additional hypotheses. None is supplied by this sentence.}
\GSfield{Reconstruction}{It is therefore not interpreted as a single theorem asserting all earlier formulas unchanged for every hermitian theory. It continues the programme of record 246, rather than settling every possible hermitian analogue.}
\GSfield{Fate}{The verified \citet{vaserstein1973foundations,vaserstein1976foundations} concern ordinary algebraic \(K\)-theory. Their linear comparison theorem does not verify this unrestricted unitary extension. The earlier form-ring work \citep{vaserstein1970stabilization} supports the setting but not the claimed all-degree package.}
\GSfield{Present status}{As of November 2026: the exact package intended remains underdetermined; no claim that hermitian \(K\)-theory itself was left undeveloped is made.}
\GSfield{Assessment}{literal statement: not determinable; publication match: thematic only; verification: unresolved lead; development: annotated; confidence: medium}
\end{GSitem}
\begin{GSitem}{1971-253.11}{A}{The higher Laurent fundamental theorem and its missing summands}
\GSfield{Kind; source}{theorem; \GSpage{P000415}}
\GSfield{The notes}{A theorem attributed to an uncertain surname resembling Wagoner is written \(K_2(A[t,t^{-1}])=K_2(A)\oplus K_1(A)\oplus ?\). It is asserted for all \(n\), with the other degrees following from this case.}
\GSfield{Statement}{For a unital associative ring and central Laurent variable, the nonconnective fundamental theorem is \[K_n(A[t,t^{-1}])\cong K_n(A)\oplus K_{n-1}(A)\oplus NK_n(A)\oplus NK_n(A),\qquad n\in\mathbb Z.\] For regular noetherian \(A\), the two Nil groups vanish. The constant inclusion is split by evaluation at \(t=1\); it does not make the remaining relative group zero.}
\GSfield{Reconstruction}{The circled final mark remains a question mark diplomatically; the two Nil terms are supplied from the precise theorem. The arrow from \(A\) means inclusion as constant polynomials, not the coefficient-extraction map from Laurent polynomials, which is not a ring homomorphism.}
\GSfield{Fate}{The full formula occurs in \citet[third theorem, p.231]{vaserstein1973foundations}. Its prior low-degree antecedents are \citet{basshellerswan1964whitehead,gersten1971laurent}. Wagoner is a plausible reading of the name, supported by the contemporary suspension literature, but the scan does not identify a precise cited paper; no priority claim rests on that reading.}
\GSfield{Present status}{As of November 2026: established, with two Nil terms for general rings. The claim that a written relative \(K_2\) is simply trivial requires correction.}
\GSfield{Assessment}{literal statement: false as written; publication match: partial; verification: read the result; development: annotated; confidence: medium}
\end{GSitem}
\begin{GSitem}{1971-253.12}{E}{Whitehead’s four factors and the Laurent K1-to-K2 class}
\GSfield{Kind; source}{construction; \GSpage{P000415}}
\GSfield{The notes}{The diagonal block \(\operatorname{diag}(a,a^{-1})\) is written as four elementary matrices, with two entries left as dots. The notes then place the diagonal blocks for \(gt\), \(t\), and \(g\) side by side, and say that their elementary decompositions give a relation, hence an element of \(K_2\).}
\GSfield{Statement}{For an invertible matrix \(a\), put \(U(b)=\left(\begin{smallmatrix}1&b\\0&1\end{smallmatrix}\right)\) and \(L(c)=\left(\begin{smallmatrix}1&0\\c&1\end{smallmatrix}\right)\). Then \(\operatorname{diag}(a,a^{-1})=U(a-1)L(1)U(a^{-1}-1)L(-a)\). If \(h(a)\) is the corresponding Steinberg lift and \(t\) is central, the product \(h(gt)h(t)^{-1}h(g)^{-1}\) maps to the identity and so defines a \(K_2(A[t,t^{-1}])\)-class. With compatible conventions this is the Laurent splitting from \(K_1(A)\).}
\GSfield{Reconstruction}{The two entries are uniquely recovered by multiplication. The inverses and multiplication order between the three lifts are supplied; without them the three displayed matrices do not themselves constitute a relation. Well-definedness on \(K_1\) is part of the established splitting theorem, not proved by the shorthand alone.}
\GSfield{Fate}{The elementary identity is Whitehead's classical stabilization device. The low-degree Laurent map is part of the theory discussed in \citet{milnor1972introduction,gersten1971laurent}; the full splitting is explicitly recorded by \citet[third theorem]{vaserstein1973foundations}. Milnor's book was published in January 1972 according to the publisher catalogue.}
\GSfield{Present status}{As of November 2026: established. The restored factors are an elementary verification of the written construction, not a new research contribution.}
\GSfield{Assessment}{literal statement: correct with missing hypotheses; publication match: partial; verification: read the result; development: annotated; confidence: high}
\end{GSitem}
\begin{GSitem}{1971-253.13}{A}{Compression of Laurent relations and the inverse boundary}
\GSfield{Kind; source}{method; \GSpage{P000416}}
\GSfield{The notes}{For the converse, Laurent multiplication is made into an infinite matrix, cut to one side, and considered modulo finite matrices. The quotient is written \(\ell A/\mu A\), although \(\mu A\) was previously the quotient ring.}
\GSfield{Statement}{Laurent polynomial matrices act as two-sided finite-band matrices. Compression to nonnegative indices is a ring homomorphism modulo the ideal of finitely supported matrices, because the multiplication defect has finite support. A \(K_2\)-relation can be lifted through the resulting extension; its boundary is a \(K_1(A)\)-class after Morita identification of the finite-matrix ideal. The quotient intended with the earlier notation is \(\ell A/MA=\mu A\).}
\GSfield{Reconstruction}{The method is the preceding record's truncation construction, now used in the reverse Laurent map. The word Laurent series cannot automatically mean arbitrary formal series: their matrices need not belong to the row-and-column-finite algebra being used. The safe algebraic domain here is Laurent polynomials or an explicitly enlarged topological matrix algebra, not supplied by the notes.}
\GSfield{Fate}{The precise suspension and exact-sequence context is \citet[first--third theorems]{vaserstein1973foundations} and \citet[\S\S18--20]{vaserstein1976foundations}. The notebook supplies an implementation sketch, so the publication match is partial.}
\GSfield{Present status}{As of November 2026: established as the algebraic boundary after the matrix class and ideal are fixed; the notation drift is not a new quotient construction.}
\GSfield{Assessment}{literal statement: false as written; publication match: partial; verification: read the result; development: annotated; confidence: medium}
\end{GSitem}
\begin{GSitem}{1971-253.14}{A}{Vaserstein’s explicit expectation of agreement with Quillen}
\GSfield{Kind; source}{conjecture; \GSpage{P000414}}
\GSfield{The notes}{After the pullback construction, the notebook says \ru{По мнению Л. В. теория Квиллена совпадает с этими}: in Vaserstein's opinion Quillen's theory agrees with these. The next line explicitly excludes Karoubi's polynomial theory.}
\GSfield{Statement}{For an associative ring \(A\) with identity, the stable Volodin--Vaserstein simplicial groups and the free-resolution groups of Swan and Gersten should be naturally isomorphic to \(K_n^Q(A)=\pi_nBGL(A)^+\) for every \(n\ge1\). The comparison does not assert agreement with Karoubi--Villamayor groups for arbitrary singular rings.}
\GSfield{Reconstruction}{The positive index range and unital hypothesis identify the exact theorem later published. The notation does not make an unrestricted assertion about all nonunital or polynomial-homotopy constructions.}
\GSfield{Fate}{The expectation is stated as a theorem in the comparison block assigned conditionally to record 257. \citet[first theorem, p.231]{vaserstein1973foundations} announces the agreement, and \citet[\S27, Theorem 27.2; \S29, Theorem 29.1]{vaserstein1976foundations} gives the detailed comparisons. The latter explicitly dates first announcement to Gelfand's seminar in spring 1972; the short article was received on 14 November 1972.}
\GSfield{Present status}{As of November 2026: resolved in the stated scope. GP-1971-21 links this conjectural wording, the suspension question in 1971-256.3, and the comparison theorem in 1971-257.2.}
\GSfield{Assessment}{literal statement: correct with missing hypotheses; publication match: exact; verification: read the result; development: annotated; confidence: high; problem \hyperref[prob:GP-1971-21]{GP-1971-21}}
\end{GSitem}

\GStalk{254}{1972-03-06}{D. A. Kazhdan}{P000416}{A residue as the trace of a compressed commutator}
Kazhdan's short contribution begins at the explicit heading below Vaserstein's Laurent boundary. For the pair \(k[[t]]\subset k((t))\), it asks for a residue description independent of a parameter, using complementary linear projections and a commutator trace. This is a recognizable instance of Tate's residue construction, published in 1968 from a 1967 lecture. The subscripts on the final projections are uncertain; the modern statement below fixes a convention and does not pretend to recover overwritten indices.

\begin{GSitem}{1971-254.1}{E}{Projections onto a Laurent-series lattice and a finite-rank commutator}
\GSfield{Kind; source}{construction; \GSpage{P000416}}
\GSfield{The notes}{The source sets \(K=k((t))\), \(A=k[[t]]\), and chooses maps \(p_1+p_2=1\), \(p_1A=0\), \(p_2K\subset A\). For \(f,g\in K\) it writes the trace of a commutator with uncertain projection subscripts and calls the operator finite rank.}
\GSfield{Statement}{Let \(P:K\to A\) be a \(k\)-linear projection with \(P|_A=1\), and define \(T_f=P\circ M_f|_A\), where \(M_f\) is multiplication by \(f\). The commutator \([T_f,T_g]\) is finite rank on \(A\): its two multiplication defects factor through the finite-dimensional quotients \((A+fA)/A\) and \((A+gA)/A\). Hence its ordinary finite-rank trace is defined. If one instead works on all of \(K\), Tate's finite-potent trace is the appropriate general construction; a full-space compressed commutator need not itself have finite rank.}
\GSfield{Reconstruction}{The stated conditions force \(p_2\) to be a projection onto \(A\); thus \(P=p_2\) is a natural modern convention. Compression to \(A\), the multiplication notation, and the distinction between finite-rank and finite-potent operators are supplied rather than read into the uncertain final subscripts.}
\GSfield{Fate}{\citet[\S1 and \S2, Theorem 1, pp.149--152]{tate1968residues} defines the trace and the abstract residue. The present compressed formula is a specialization of that theorem.}
\GSfield{Present status}{As of November 2026: established for the formal Laurent field over any ground field. No analytic Hilbert-space trace or characteristic-zero restriction is needed.}
\GSfield{Assessment}{literal statement: correct with missing hypotheses; publication match: partial; verification: read the result; development: annotated; confidence: high}
\end{GSitem}
\begin{GSitem}{1971-254.2}{E}{Tate’s coordinate-free formula for the residue}
\GSfield{Kind; source}{theorem; \GSpage{P000416}}
\GSfield{The notes}{The question asks for the residue of \(f\,dg\); the parenthetical phrase most plausibly says ``without a parameter''. The trace construction is the proposed answer.}
\GSfield{Statement}{With the preceding convention, \[\operatorname{res}_{t=0}(f\,dg)=\operatorname{Tr}_A([T_f,T_g]).\] The value is independent of the projection and of the chosen uniformizer of the discrete valuation ring \(A\). With the standard projection onto nonnegative powers it equals the coefficient of \(t^{-1}\) in \(f(t)g'(t)\). The normalization is \(\operatorname{res}(t^{-1}dt)=1\); on \(A\), the backward-shift/forward-shift commutator is the projection onto constants.}
\GSfield{Reconstruction}{The projection convention and commutator order fix the sign. Parameter independence is inferred as the intended question from Tate's exact matching construction, while the notebook's parenthesis remains diplomatically qualified. No global residue theorem is attributed to this short remark.}
\GSfield{Fate}{The coordinate-free theorem is \citet[Theorem 1 and property (R1), pp.151--152; Theorem 2, p.155]{tate1968residues}. Tate states that it grew from a Bowdoin lecture in summer 1967; publication was in 1968, well before Kazhdan's report.}
\GSfield{Present status}{As of November 2026: classical and fully resolved. The notebook is an economical exposition of the trace interpretation, not evidence for an orphan residue problem.}
\GSfield{Assessment}{literal statement: correct with missing hypotheses; publication match: partial; verification: read the result; development: annotated; confidence: high}
\end{GSitem}

\GStalk{255}{1972-03-13}{Jack Morava}{P000417-P000420}{Segal’s categories, coherent multiplication, and the sphere spectrum}
Morava reports Segal's construction of cohomology theories from categories with coherent sum. Gelfand's intervention spells out simplicial objects, realization, and nerves. The lecture contrasts strict commutative multiplication with coherent commutativity, constructs iterated deloopings, and treats projective modules and finite sets. The latter yields stable cohomotopy and a natural map to algebraic \(K\)-theory. The closing comparison with the image of \(J\) lacks enough information to identify an exact finite-field assertion. Segal's 1974 paper gives the full matching construction, but the basic programme was already announced by Quillen in the 1970 ICM report.

\begin{GSitem}{1971-255.1}{E}{The simplex category, simplicial objects, and realization}
\GSfield{Kind; source}{definition; \GSpage{P000417}; I. M. Gelfand, during Morava’s exposition}
\GSfield{The notes}{The auxiliary category has objects \([n]=\{0,1,\ldots,n\}\); the displayed monotonicity condition starts at \(\varphi(1)\). A boxed insertion says contravariant. The realization glues the pieces \(A_n\times\Delta^n\). Gelfand's name is written by this explanation.}
\GSfield{Statement}{Let \(\Delta\) have finite nonempty ordered sets \([n]\) as objects and all order-preserving maps as morphisms, so \(\varphi(0)\le\cdots\le\varphi(k)\). A simplicial object in \(\mathcal C\) is a functor \(\Delta^{\mathrm{op}}\to\mathcal C\). For a simplicial set \(A\), its geometric realization is \(\coprod_{n\ge0}A_n\times\Delta^n\) modulo the face and degeneracy identifications. It is a CW complex, generally not a simplicial complex without subdivision.}
\GSfield{Reconstruction}{The omitted endpoint in the monotonicity condition and the distinction between simplicial sets and complexes are restored. The named Gabriel--Zisman reference is contextual, not an additional mathematical assertion.}
\GSfield{Fate}{These constructions are given in \citet[\S1, pp.105--106]{segal1968classifying}, received in April 1967 and published in 1968; they predate this explanatory intervention.}
\GSfield{Present status}{As of November 2026: standard simplicial language, with the contravariance explicitly supported by the notebook's own insertion.}
\GSfield{Assessment}{literal statement: correct with missing hypotheses; publication match: exact; verification: read the result; development: annotated; confidence: high}
\end{GSitem}
\begin{GSitem}{1971-255.2}{E}{The nerve of a category and the group classifying space}
\GSfield{Kind; source}{construction; \GSpage{P000417-P000418}}
\GSfield{The notes}{The notes identify zero-simplices with objects, one-simplices with morphisms, and higher simplices with composable chains. A one-object groupoid gives \(BC_n=G^n\), and realization is compared with Milnor's \(BG\).}
\GSfield{Statement}{For a small category \(\mathcal C\), define \(N_n\mathcal C\) as strings of \(n\) composable arrows. Faces compose adjacent arrows or omit an endpoint; degeneracies insert identities. Set \(B\mathcal C=|N_\bullet\mathcal C|\). If \(\mathcal C\) has one object and automorphism group a discrete group \(G\), then \(N_n\mathcal C=G^n\) and \(B\mathcal C\) is a \(K(G,1)\). For a topological group one uses a topological nerve and suitable well-behaved realization; higher homotopy need not vanish.}
\GSfield{Reconstruction}{The source uses \(BC\) for both the simplicial object and its realization. Distinct notation prevents confusing these, and the discrete/topological distinction is retained.}
\GSfield{Fate}{\citet[\S\S2--3, pp.106--107]{segal1968classifying} gives exactly the categorical nerve, the one-object group example, and the comparison with Milnor's classifying space, including topological qualifications.}
\GSfield{Present status}{As of November 2026: standard. The one-object description does not identify topological \(BG\) with an Eilenberg--MacLane space unless the group is discrete.}
\GSfield{Assessment}{literal statement: correct with missing hypotheses; publication match: exact; verification: read the result; development: annotated; confidence: high}
\end{GSitem}
\begin{GSitem}{1971-255.3}{E}{Products of nerves and natural isomorphisms as homotopies}
\GSfield{Kind; source}{theorem; \GSpage{P000418}}
\GSfield{The notes}{The lemma states \(B(\alpha\times\beta)=B\alpha\times B\beta\). Direct sum induces multiplication, but \(M\oplus N\) and \(N\oplus M\) are isomorphic rather than literally equal. A shuffle picture explains homotopy commutativity.}
\GSfield{Statement}{The nerve preserves categorical products, and geometric realization preserves products in compactly generated spaces. A natural transformation of functors induces a homotopy between the corresponding maps of classifying spaces. Consequently the direct-sum functor on the groupoid of finitely generated projective modules makes its classifying space a homotopy-commutative and homotopy-associative H-space, with unit supplied by the zero module; the standard associativity and symmetry isomorphisms carry the coherent data.}
\GSfield{Reconstruction}{The ambient category of spaces and the natural-transformation mechanism are supplied. A chosen product triangulation or shuffle model realizes this mechanism; no strict equality of the two functors is asserted.}
\GSfield{Fate}{\citet[\S1 and Proposition 2.1, p.106]{segal1968classifying} proves product preservation and the homotopy induced by a natural transformation. The coherently additive categorical construction is developed in \citet[\S2]{segal1974categories}.}
\GSfield{Present status}{As of November 2026: established; the notebook correctly distinguishes equality of objects from canonical isomorphism.}
\GSfield{Assessment}{literal statement: correct with missing hypotheses; publication match: exact; verification: read the result; development: annotated; confidence: high}
\end{GSitem}
\begin{GSitem}{1971-255.4}{E}{Strict commutative multiplication and generalized Eilenberg–MacLane spaces}
\GSfield{Kind; source}{claim; \GSpage{P000419}}
\GSfield{The notes}{A remark says that strictly commutative and associative multiplication, rather than multiplication only up to homotopy, makes \(X\) a product of \(K(\pi_n,n)\). A nearby box says ``Boardman--Vogt 1968 unpublished''. No unit, inverses, or topological regularity are specified in the displayed implication.}
\GSfield{Statement}{A safe precise version is that a topological abelian group, after passage to its singular simplicial abelian group, has the weak homotopy type of a product of Eilenberg--MacLane spaces, noncanonically. For CW-type models this gives the corresponding homotopy equivalence. Such a conclusion for a raw binary associative commutative operation requires at least the relevant unit, group-completion, and topological hypotheses; a bare operation does not specify them.}
\GSfield{Reconstruction}{The statement is restricted to the standard abelian-group model and is a homotopy-type assertion, not a homeomorphism or a canonical splitting. The Boardman--Vogt note cannot be taken as evidence that no announcement had been published by 1972.}
\GSfield{Fate}{\citet[Introduction, p.293]{segal1974categories} explains that the topological-abelian-group construction yields ordinary-cohomology summands. Separately, \citet[p.1117]{boardmanvogt1968homotopy} already published their recognition announcement in November 1968, while promising full details elsewhere.}
\GSfield{Present status}{As of November 2026: the precise classical statement is established. The notebook's omitted structural assumptions and its unpublished label need qualification.}
\GSfield{Assessment}{literal statement: correct with missing hypotheses; publication match: partial; verification: read the result; development: annotated; confidence: medium}
\end{GSitem}
\begin{GSitem}{1971-255.5}{E}{Coherent commutativity and repeated delooping}
\GSfield{Kind; source}{theorem; \GSpage{P000419}}
\GSfield{The notes}{The theorem assumes multiplication commutative and associative ``in a canonical way'', constructs \(BX\) with the same structure, and iterates. It states \(B^k\mathcal C\simeq\Omega B^{k+1}\mathcal C\) for \(k\ge2\).}
\GSfield{Statement}{A special \(\Gamma\)-space \(A\) has \(A(0)\) contractible and Segal equivalences \(A(n)\simeq A(1)^n\), compatible with all structure maps. Its iterated classifying spaces form a connective spectrum. The map \(A(1)\to\Omega BA(1)\) is an equivalence when \(\pi_0A(1)\) is a group under multiplication; otherwise it is the group-completion map. Later deloopings are connected and have the required inverse structure. Equivalently, a grouplike coherently commutative H-space of CW type is an infinite loop space.}
\GSfield{Reconstruction}{The phrase ``canonical way'' is interpreted as full coherent structure, not merely one associating homotopy and one commuting homotopy. The group-completion condition accounts for why the initial category space is excluded from the notebook's eventual loop equivalences.}
\GSfield{Fate}{\citet[Theorem A]{boardmanvogt1968homotopy} states the grouplike recognition and repeated-classifying-space theorem. Quillen reports Segal's categorical cohomology construction in \citet[\S3, p.50]{quillen1971cohomology}; the detailed \(\Gamma\)-space version is \citet[Definitions 1.1--1.3 and Proposition 1.4]{segal1974categories}.}
\GSfield{Present status}{As of November 2026: established. The principal claim predates the lecture, despite the later appearance of Segal's full paper.}
\GSfield{Assessment}{literal statement: correct with missing hypotheses; publication match: exact; verification: read the result; development: annotated; confidence: high}
\end{GSitem}
\begin{GSitem}{1971-255.6}{E}{The represented cohomology theory and its grading}
\GSfield{Kind; source}{definition; \GSpage{P000419}}
\GSfield{The notes}{The notes write \(K_{\mathcal C}^{-k}(Y)=\lim_i[Y,\Omega^iB^{k+i}\mathcal C]\), equivalently a suspended-source formula, and say the limit is unnecessary after the displayed loop equivalences.}
\GSfield{Statement}{If \(E\) is the associated spectrum with spaces \(E_r\) and coherent structure maps, its reduced cohomology on a finite based CW complex is \(\widetilde E^j(Y)=\varinjlim_r[\Sigma^rY,E_{r+j}]_*\), with nonnegative indices after increasing \(r\). For an \(\Omega\)-spectrum the transition maps are bijections and one may use a single sufficiently high level. For general CW complexes the definition is the stable mapping group \([\Sigma^\infty Y,\Sigma^jE]\). In particular \(E^{-n}(\mathrm{pt})=\pi_nE\) in the unreduced theory.}
\GSfield{Reconstruction}{The displayed minus sign, the indexing of \(B\mathcal C\) as the nerve versus the first spectrum level, and the point/basepoint convention are not consistent enough to retain literally. The standard spectrum grading is supplied and kept separate from the notebook's \(K^{-k}\) notation.}
\GSfield{Fate}{The construction is announced in \citet[\S3]{quillen1971cohomology} and made precise in \citet[\S1 and Propositions 3.3--3.4]{segal1974categories}. The equivalence after delooping is Proposition 1.4 of the latter paper.}
\GSfield{Present status}{As of November 2026: standard represented generalized cohomology. Stabilization is removed only after an appropriate loop-spectrum replacement, not for every initial category nerve.}
\GSfield{Assessment}{literal statement: correct with missing hypotheses; publication match: partial; verification: read the result; development: annotated; confidence: medium}
\end{GSitem}
\begin{GSitem}{1971-255.7}{E}{Projective modules, automorphism groupoids, and algebraic K-theory}
\GSfield{Kind; source}{construction; \GSpage{P000418-P000419}}
\GSfield{The notes}{The category is explicitly restricted to finitely generated projective \(R\)-modules and isomorphisms. Its classifying space is written as a disjoint union of \(K(GL(M),1)\); the resulting K-functor agrees with Quillen's.}
\GSfield{Statement}{Choose a small skeleton of the groupoid \(\mathcal P(R)^{\simeq}\) of finitely generated projective right \(R\)-modules and isomorphisms. Then \(B\mathcal P(R)^{\simeq}\simeq\coprod_{[P]}B\operatorname{Aut}_R(P)\). Direct-sum group completion has zeroth space \(K_0(R)\times BGL(R)^+\), and its positive homotopy groups are Quillen's \(K_n(R)\). The sum ranges over isomorphism classes, and \(GL(M)\) here means the automorphism group of a projective module, not necessarily a free-module matrix group.}
\GSfield{Reconstruction}{A skeleton, group completion, and the notation for module automorphisms are specified. Taking all module morphisms instead would introduce a zero object and a contractible nerve; the source's restriction to isomorphisms is therefore essential and is preserved.}
\GSfield{Fate}{Quillen already states the projective-module categorical interpretation in \citet[\S3, p.50]{quillen1971cohomology}. \citet[\S2 and \S4, p.305]{segal1974categories} gives the group-completion model explicitly.}
\GSfield{Present status}{As of November 2026: established for unital associative rings. The uncompleted disjoint union itself is not yet the full K-theory space.}
\GSfield{Assessment}{literal statement: correct with missing hypotheses; publication match: exact; verification: read the result; development: annotated; confidence: high}
\end{GSitem}
\begin{GSitem}{1971-255.8}{E}{Finite sets and the Barratt–Priddy–Quillen theorem}
\GSfield{Kind; source}{theorem; \GSpage{P000419-P000420}}
\GSfield{The notes}{For the category of finite sets, the notes write \(B\mathcal C=\coprod_nB\Sigma_n\), and say the associated groups are stable cohomotopy. The entire union is called an H-space.}
\GSfield{Statement}{Use the groupoid of finite sets and bijections, with disjoint union. Its nerve is \(\coprod_{n\ge0}B\Sigma_n\). Its group-completion spectrum is equivalent to the sphere spectrum \(\mathbb S\); equivalently its zeroth infinite loop space is \(QS^0\simeq\mathbb Z\times B\Sigma_\infty^+\). Thus the associated cohomology is stable cohomotopy. The category of finite sets and all functions has a terminal object and would not give this result by its plain nerve.}
\GSfield{Reconstruction}{The morphisms are recovered from the displayed symmetric-group decomposition. Group completion and the missing component group \(\mathbb Z\) are made explicit.}
\GSfield{Fate}{\citet[Proposition 3.5 and \S4]{segal1974categories} proves the stated sphere-spectrum result and credits Barratt--Priddy--Quillen. Its references identify prior Barratt and Priddy announcements in 1971; the notebook itself calls the result known. The full Segal proof was received in August 1972 and published in 1974.}
\GSfield{Present status}{As of November 2026: classical Barratt--Priddy--Quillen theorem, with the correct groupoid and completion specified.}
\GSfield{Assessment}{literal statement: correct with missing hypotheses; publication match: exact; verification: read the result; development: annotated; confidence: high}
\end{GSitem}
\begin{GSitem}{1971-255.9}{E}{Linearization from finite sets to free modules}
\GSfield{Kind; source}{construction; \GSpage{P000420}}
\GSfield{The notes}{A finite set of size \(k\) is sent to \(R^k\), giving a map from stable cohomotopy to \(K(R)\). At a point this gives homomorphisms from stable homotopy groups of spheres to K-groups of every ring.}
\GSfield{Statement}{The functor \(S\mapsto R[S]\), with bijections acting by permutation matrices, is compatible with disjoint union and direct sum. After group completion it induces a natural spectrum map \(\mathbb S\to K(R)\), sending \(1\in\pi_0\mathbb S\cong\mathbb Z\) to \([R]\in K_0(R)\). Consequently it induces \(\pi_n^S\to K_n(R)\) in all nonnegative degrees and the corresponding natural transformation on represented cohomology.}
\GSfield{Reconstruction}{The map is interpreted as the spectrum map selected by the rank-one free module. For a general noncommutative ring no commutative ring-spectrum structure on \(K(R)\) is assumed.}
\GSfield{Fate}{This follows within Segal's categorical functorial construction \citep[\S\S2--3, Proposition 3.5]{segal1974categories}; the prior announcement of the machine is \citet[\S3]{quillen1971cohomology}. The actual finite-set functor is explicitly present in the notebook.}
\GSfield{Present status}{As of November 2026: standard linearization. Its image need not be the whole K-group.}
\GSfield{Assessment}{literal statement: correct with missing hypotheses; publication match: exact; verification: read the result; development: annotated; confidence: high}
\end{GSitem}
\begin{GSitem}{1971-255.10}{F}{The unresolved image-of-J comparison over a prime field}
\GSfield{Kind; source}{remark; \GSpage{P000420}}
\GSfield{The notes}{For \(R=\mathbb F_p\) and \(Y\) a point, a damaged relation sign connects \(K(R)^kY\) with ``image of the j-homomorphism''. Degrees, localization, and the map being compared are not supplied.}
\GSfield{Statement}{The preceding linearization gives a definite map \(\pi_n^S\to K_n(\mathbb F_p)\), but the sentence does not specify whether its image, a quotient, a prime-local summand, or the entire target is meant. Equality of every target K-group with the image of this map is false: in degree one the permutation matrices have determinant only \(\pm1\), whereas \(K_1(\mathbb F_5)=\mathbb F_5^\times\) has order four.}
\GSfield{Reconstruction}{The degree-one check separates a literal surjectivity reading from more refined relations with stable \(J\), which require prime and degree conventions. No one of those deeper formulations is selected as the intended missing theorem.}
\GSfield{Fate}{The source's link to Quillen's finite-field computation is clear \citep[\S3, Theorems 1--2]{quillen1971cohomology}; Segal's verified paper establishes the sphere-spectrum and linearization framework but does not supply this particular image-of-\(J\) claim. Searches combining Morava, Segal, finite-field K-theory and image \(J\) gave related theorems without recovering the lost qualifiers.}
\GSfield{Present status}{As of November 2026: the exact comparison remains unidentified from the notebook. This is not an assertion that the classical image-of-\(J\) problem is open.}
\GSfield{Assessment}{literal statement: not determinable; publication match: thematic only; verification: unresolved lead; development: annotated; confidence: medium}
\end{GSitem}

\GStalk{256}{1972-03-13}{L. N. Vaserstein}{P000420-P000422}{Simplicial K-theory, algebraic suspension, and comparisons}
The explicit Vaserstein heading in the lower part of P000420 begins a description of triangular simplices in the stable general linear group. The next page constructs an algebraic cone and suspension, states their degree-shifting property, returns to free-ring resolutions, and compares negative groups with Bass's. The first two lines of P000422 concern Karoubi's polynomial theory. The subsequent unheaded Quillen comparison block is allocated conditionally to catalogue record 257; there is no additional dated Kazhdan heading at that boundary. Two significant qualifications are needed: relative K-theory cannot generally be replaced by the K-theory of a bare kernel, and absence of zero divisors does not imply agreement with polynomial-homotopy K-theory.

\begin{GSitem}{1971-256.1}{E}{Triangular simplices on the stable general linear group}
\GSfield{Kind; source}{construction; \GSpage{P000420-P000421}}
\GSfield{The notes}{An abstract simplicial scheme is a collection of finite subsets closed under passage to subsets. On stable matrices the notes declare a finite set \(a_1,\ldots,a_n\) to be a simplex if simultaneous left and right multiplication by invertible matrices makes them upper unitriangular. Conjugates, permutation matrices, and a marginal reference to Volodin accompany \(K_i(A)=\pi_{i-1}GL(A)\).}
\GSfield{Statement}{For an associative unital ring \(A\), one equips the set \(GL(A)=\varinjlim_nGL_n(A)\) with the triangular simplicial structure, rather than its discrete topology. A finite subset contained in a translate of an admissible conjugate of a unitriangular subgroup determines a simplex. In the usual Volodin version the triangular subgroups are indexed by orderings of a finite basis. The realization \(S(A)\) of the specified stable complex satisfies \(K_i(A)\cong\pi_{i-1}S(A)\) for \(i\ge1\), with \(\pi_0\) carrying its stable group structure.}
\GSfield{Reconstruction}{The source's two-sided formulation and its marginal restrictions are recorded without claiming that each possible choice of conjugating matrices gives literally the same complex. The standard Volodin model is the precise comparator. The homotopy groups belong to a realization, not to an ordinary discrete general linear group.}
\GSfield{Fate}{The construction had appeared in \citet{volodin1971extraordinary}, received on 27 July 1970 and published in 1971. Its relationship with Vaserstein's relative formulation is specified in \citet[\S15, Theorem 15.1 and Corollary 15.1]{vaserstein1976foundations}; the later comparison with Quillen is a separate theorem, indexed under record 257.}
\GSfield{Present status}{As of November 2026: the simplicial models are established. The unreadable adjective in the margin is not used to impose a new matrix restriction.}
\GSfield{Assessment}{literal statement: correct with missing hypotheses; publication match: partial; verification: read the result; development: annotated; confidence: medium}
\end{GSitem}
\begin{GSitem}{1971-256.2}{E}{The algebraic cone and suspension rings}
\GSfield{Kind; source}{construction; \GSpage{P000420-P000421}}
\GSfield{The notes}{The notation \(\ell A\) denotes infinite matrices with finite rows and columns; \(MA\) denotes the finite matrices, and \(\mu A=\ell A/MA\) is called the ring of Fredholm operators.}
\GSfield{Statement}{Let \(C(A)\) be the ring of countably indexed matrices over \(A\) having finitely many nonzero entries in each row and each column, and let \(M_\infty(A)\) be its ideal of matrices with finite total support. Set \(\Sigma A=C(A)/M_\infty(A)\). Matrix multiplication is well-defined under these finiteness conditions. The invertible classes in the quotient are algebraic Fredholm classes; arbitrary elements of the quotient need not be invertible.}
\GSfield{Reconstruction}{The symbols \(C\) and \(\Sigma\) replace the source's \(\ell\) and \(\mu\). This purely algebraic quotient is analogous to a Calkin algebra. Even over \(\mathbb C\), row and column finiteness alone does not ensure that the matrix defines a bounded operator on \(\ell^2\).}
\GSfield{Fate}{The matrix cone is treated in \citet[\S\S3--4 and Chapter IV]{vaserstein1976foundations}. Its historical introduction credits earlier Fredholm-operator definitions of negative K-groups to Karoubi--Villamayor and Wagoner; the 1972 talk reuses that algebraic construction.}
\GSfield{Present status}{As of November 2026: standard algebraic suspension construction, with the ring and its group of invertible classes distinguished.}
\GSfield{Assessment}{literal statement: correct with missing hypotheses; publication match: exact; verification: read the result; development: annotated; confidence: high}
\end{GSitem}
\begin{GSitem}{1971-256.3}{A}{Suspension, vanishing of cone K-theory, and the index}
\GSfield{Kind; source}{theorem; \GSpage{P000421}}
\GSfield{The notes}{The framed theorem is \(K_iA=K_{i+1}\mu A\), in particular \(K_1\mu A=K_0A\), called the index. Its explanation is contractibility of \(GL(\ell A)\), or at least vanishing of \(K_i(\ell A)\). A separate sentence asks whether this holds for Quillen's functor.}
\GSfield{Statement}{For a unital associative ring \(A\), the algebraic cone is flasque: the stable sum swindle gives \(K_n(C(A))=0\). The extension \(M_\infty(A)\to C(A)\to\Sigma A\), together with Morita invariance and its applicable localization sequence, gives \[K_{n+1}(\Sigma A)\xrightarrow{\ \partial\ }K_n(A)\] as an isomorphism for \(n\ge0\), and for every integer \(n\) in nonconnective K-theory. For \(n=0\), \(\partial\) is the algebraic index.}
\GSfield{Reconstruction}{Contractibility refers to the chosen K-theory or simplicial model; it is not a claim that the underlying discrete ring or group is contractible. The special matrix ideal supplies the exactness required here, without asserting excision for arbitrary ideals. The exercise about finding the arrows is interpreted as supplying these boundary and comparison maps, not as an additional specified problem.}
\GSfield{Fate}{The degree shift and comparison are announced in \citet[the first theorem block, p.231]{vaserstein1973foundations}, received on 14 November 1972, and developed in \citet[Chapter IV, especially \S18; \S27]{vaserstein1976foundations}. The latter explicitly dates the comparison programme to Gelfand's seminar in spring 1972. The same degree shift holds in Quillen K-theory under the displayed hypotheses.}
\GSfield{Present status}{As of November 2026: resolved. The question about Quillen's theory is an instance of the model-comparison problem GP-1971-21.}
\GSfield{Assessment}{literal statement: correct with missing hypotheses; publication match: exact; verification: read the result; development: annotated; confidence: high; problem \hyperref[prob:GP-1971-21]{GP-1971-21}}
\end{GSitem}
\begin{GSitem}{1971-256.4}{F}{The unspecified Banach-ring generalization}
\GSfield{Kind; source}{remark; \GSpage{P000421}}
\GSfield{The notes}{One line says that the theorem generalizes to a Banach ring and gives the usual theorem. No topology, completed ideal, or K-theory convention is named.}
\GSfield{Statement}{A Banach-ring analogue must specify the category of modules or operators, its completion, the ideal being divided out, and whether K-theory uses norm-continuous paths or algebraic plus constructions. The preceding row-and-column-finite ring is not, as defined, a Banach algebra. Thus this line does not determine a single comparison or periodicity theorem.}
\GSfield{Reconstruction}{An analytic Fredholm-index theorem is a plausible context, but identifying it with a particular completed Calkin extension would supply data not present in the source. Algebraic and topological K-groups of a Banach algebra are not silently equated.}
\GSfield{Fate}{The verified matrix construction in \citet[Chapter IV]{vaserstein1976foundations} and the searches for Karoubi--Villamayor Banach-ring versions provide context, but no exact Banach theorem with the notebook's omitted choices was identified.}
\GSfield{Present status}{As of November 2026: unidentified at the required level of hypotheses; the ambiguity does not establish an open analytic problem.}
\GSfield{Assessment}{literal statement: not determinable; publication match: thematic only; verification: unresolved lead; development: annotated; confidence: low}
\end{GSitem}
\begin{GSitem}{1971-256.5}{A}{Free rings are K-acyclic only in the reduced convention}
\GSfield{Kind; source}{theorem; \GSpage{P000421}}
\GSfield{The notes}{A framed statement says \(K_i(\text{free ring})=0\), followed immediately by the warning that there is no contraction of the free ring.}
\GSfield{Statement}{If \(F\) is a free associative nonunital ring on a set \(X\), its unitization is \(\mathbb Z\langle X\rangle\). The relevant assertion is \[\ker\!\left(K_n(\mathbb Z\langle X\rangle)\longrightarrow K_n(\mathbb Z)\right)=0.\] Equivalently the inclusion of constants induces \(K_n(\mathbb Z)\cong K_n(\mathbb Z\langle X\rangle)\). This holds for all integers in the nonconnective extension. It does not assert \(K_0(\mathbb Z\langle X\rangle)=0\), nor a contraction by ring homomorphisms.}
\GSfield{Reconstruction}{The augmentation and reduced interpretation are supplied from the surrounding free nonunital resolution. No excision identification with an arbitrary ambient relative group is required. This is the same acyclicity used in occurrences 1971-253.3--1971-253.4.}
\GSfield{Fate}{See \citet[the second theorem block, p.231]{vaserstein1973foundations}, followed by \citet[Theorem 28.1 and Corollary 28.1]{vaserstein1976foundations}. The detailed article explains that the higher-degree argument uses Gersten's results and does not reproduce that long proof. The 1972 announcement and the 1973/1976 publications are therefore distinguished.}
\GSfield{Present status}{As of November 2026: established in the stated reduced sense. The free-ring vanishing supports the normalized universal-relation question GP-1971-20, but is stronger than that degree-two instance.}
\GSfield{Assessment}{literal statement: correct with missing hypotheses; publication match: exact; verification: read the result; development: annotated; confidence: high; problem \hyperref[prob:GP-1971-20]{GP-1971-20}}
\end{GSitem}
\begin{GSitem}{1971-256.6}{A}{Free resolutions and the required relative degree shift}
\GSfield{Kind; source}{method; \GSpage{P000421}}
\GSfield{The notes}{For \(0\to L'\to L\to A\to0\), with \(L\) free, the notes write \(K_iA=K_{i-1}L'\), refer to a long exact sequence, and say that induction gives the theories of Swan and Gersten.}
\GSfield{Statement}{Let \(F(A)\twoheadrightarrow A\) be a functorial free nonunital presentation. The homotopy fibre of \(K(F(A))\to K(A)\) has homotopy groups \(K_{n+1}(A)\), since the reduced K-space of \(F(A)\) is contractible. It must be modeled as relative K-theory. In Vaserstein's explicit ring recursion, the doubled fibre product \(D(A)=F(A)\times_A F(A)\), together with the split projection to \(F(A)\), supplies the functor used to shift degrees. The assertion for the bare kernel \(L'\) does not follow from free acyclicity alone.}
\GSfield{Reconstruction}{The relative fibre replaces the source's unrestricted use of \(K(L')\). This is the same excision issue as occurrence 1971-253.8. A long exact sequence is legitimate for the homotopy fibre by definition; identifying its terms with absolute K-groups of an ideal needs an additional theorem.}
\GSfield{Fate}{\citet[\S29, definition preceding Theorem 29.1 and Theorem 29.1]{vaserstein1976foundations} gives the precise doubled-ring recursion and comparison with Gersten and Swan. The comparison was announced in \citet{vaserstein1973foundations}, following the spring 1972 seminar.}
\GSfield{Present status}{As of November 2026: the corrected relative construction is established. The bare-kernel shortcut remains an unjustified literal statement in the notes.}
\GSfield{Assessment}{literal statement: correct with missing hypotheses; publication match: partial; verification: read the result; development: annotated; confidence: high}
\end{GSitem}
\begin{GSitem}{1971-256.7}{E}{Agreement with Bass negative groups and regular vanishing}
\GSfield{Kind; source}{theorem; \GSpage{P000421}}
\GSfield{The notes}{The final lines say that the groups with negative index also coincide with Bass's, and that all of them vanish for regular rings.}
\GSfield{Statement}{For a unital ring, the negative groups obtained from repeated algebraic suspension agree with Bass negative K-theory. If \(A\) is regular noetherian, then \(K_n(A)=0\) for all \(n<0\). Here regular means the homological regularity relevant to the fundamental theorem, not merely absence of zero divisors.}
\GSfield{Reconstruction}{The regular noetherian hypothesis is made explicit to avoid extending the assertion to arbitrary singular domains. The equality concerns negative K-groups, separately from the positive polynomial-homotopy comparison on P000422.}
\GSfield{Fate}{Bass's negative construction predates the lecture; its relationship with the Fredholm model is recalled in \citet[introduction, pp.95--96 and Chapter IV]{vaserstein1976foundations}. The modern nonconnective formulation is \citet[IV, \S10]{weibel2013kbook}; the regularity convention is already explicit in \citet{basshellerswan1964whitehead}.}
\GSfield{Present status}{As of November 2026: classical. Singular rings can have nonzero negative K-groups, so the regularity qualification carries mathematical content.}
\GSfield{Assessment}{literal statement: correct with missing hypotheses; publication match: exact; verification: read the result; development: annotated; confidence: high}
\end{GSitem}
\begin{GSitem}{1971-256.8}{B}{Karoubi–Villamayor comparison: a domain hypothesis is insufficient}
\GSfield{Kind; source}{claim; \GSpage{P000422}}
\GSfield{The notes}{The first two lines say that Karoubi's positive-index theory differs from Bass's, but coincides when there are no zero divisors: \ru{Нет делит. нуля}.}
\GSfield{Statement}{For a unital ring \(A\), polynomial-homotopy K-theory is \(KV_n(A)=\pi_{n-1}GL(A[\Delta^\bullet])\), \(n\ge1\). The natural comparison \(K_n(A)\to KV_n(A)\) is an isomorphism for regular noetherian rings. It is not an isomorphism for all integral domains. For example, the graded domain \(A=\mathbb F_p[t^2,t^3]\) has \(KV_2(A)=KV_2(\mathbb F_p)=0\), whereas \(K_2(A,(t^2,t^3))\ne0\) and hence \(K_2(A)\ne0\).}
\GSfield{Reconstruction}{The intended correct hypothesis is reconstructed as regularity, not read into the words ``no zero divisors''. \ours\ The diagnostic deduction is as follows. Fix a prime \(p\) and the perfect field \(k=\mathbb F_p\). The ring \(A=k[t^2,t^3]\subset k[t]\) is a domain, and \(t^2\mapsto s^2t^2,\ t^3\mapsto s^3t^3\) defines a polynomial homotopy from the augmentation followed by constants to the identity. Hence \(KV_2(A)=KV_2(k)=K_2(k)=0\). In Hesselholt--Nikolaus's Theorem A, take \(a=2,b=3,r=1\). The truncation set contains 1, and the first Witt coordinate maps its quotient onto \(k\), since the images of both Verschiebung maps \(V_2,V_3\) have zero first coordinate. Thus \(K_2(A,(t^2,t^3))\ne0\). The augmentation is split, so this relative group is a direct summand of \(K_2(A)\); indeed it is the whole group because \(K_2(k)=0\).}
\GSfield{Fate}{The distinction between polynomial paths and stable elementary equivalence predates the talk and is discussed in \citet[introduction, p.96]{vaserstein1976foundations}. The precise regular-ring comparison is \citet[IV, Theorem 11.8]{weibel2013kbook}. A direct modern counterexample to the literal domain claim follows from \citet[Theorem A]{hesselholtnikolaus2020cuspidal}: its Witt-vector quotient for \(a=2,b=3,r=1\) is nonzero. That paper first circulated on 20 March 2019 and was published in 2020; it is not attributed to the 1972 seminar.}
\GSfield{Present status}{As of November 2026: the regular-ring comparison is established, and the unrestricted domain claim is false. This is a correction of hypotheses, not a newly posed unresolved problem. The editorial counterexample deduction was independently checked against Theorem A, the split long exact sequence, and the polynomial contraction; it also applies in characteristics two and three.}
\GSfield{Assessment}{literal statement: false as written; publication match: partial; verification: read the result; development: annotated; confidence: high}
\end{GSitem}

\GStalk{257}{1972-03-20}{D. A. Kazhdan (catalogue attribution; notebook heading absent)}{P000422-P000423}{Quillen comparison block: conditional attribution}
This entry preserves the catalogue's record and date, but its notebook attribution is conditional. The unheaded block beginning ``Quillen's theory'' on P000422 follows Vaserstein's Karoubi comparison, and continues through the upper half of P000423. It explains the plus construction, states equality of the Quillen, Swan, Gersten, Volodin and Vaserstein theories, and sketches a quotient model. Morava is named only for an intervention. The next explicit date and speaker are 20 March and Gelfand below the divider on P000423; Kazhdan's explicit proof heading occurs only on P000424, in record 259. The theorem here belongs to the same comparison programme announced by Vaserstein, so the allocation is not evidence for an independent Kazhdan priority claim.

\begin{GSitem}{1971-257.1}{E}{The stable plus construction and its universal property}
\GSfield{Kind; source}{construction; \GSpage{P000422}}
\GSfield{The notes}{With \(G=GL(A)\), the notes identify \(E(A)\) with the perfect commutator subgroup, give \(BG\to BG^+\) with fundamental group \(G/E\) and unchanged homology, and state universality for maps to H-spaces. Morava adds that an H-space hypothesis can be dropped when the other map induces a homology isomorphism.}
\GSfield{Statement}{For a unital associative ring \(A\), \(E(A)=[GL(A),GL(A)]\) is perfect. The plus map \(BGL(A)\to BGL(A)^+\) kills \(E(A)\) and induces isomorphisms in homology with every local coefficient system pulled back from \(K_1(A)\). Its target is a homotopy-commutative H-space, and \(K_i(A)=\pi_iBGL(A)^+\), \(i\ge1\). The corresponding factorization property holds for connected H-space targets. A uniqueness comparison without an H-space assumption must retain the appropriate fundamental-group and acyclic-map hypotheses; an arbitrary integral homology equivalence alone is insufficient.}
\GSfield{Reconstruction}{Local coefficients and the killed subgroup are supplied to make the plus characterization precise. Morava's abbreviated correction is interpreted only within that setting, not as a blanket homology-Whitehead theorem for all spaces. The crossed equality sign does not alter the verified stable perfectness statement.}
\GSfield{Fate}{Quillen's construction was already reported in \citet[\S3, p.50]{quillen1971cohomology}, from the 1970 ICM lecture, printed in 1971. The precise characterization and universal properties are \citet[IV, Definition 1.1 and Theorems 1.5 and 1.8]{weibel2013kbook}. The explicit discussion continued in occurrence 1971-259.1.}
\GSfield{Present status}{As of November 2026: established. This exposition does not create a new comparison theorem merely because its catalogue attribution is uncertain.}
\GSfield{Assessment}{literal statement: correct with missing hypotheses; publication match: partial; verification: read the result; development: annotated; confidence: high}
\end{GSitem}
\begin{GSitem}{1971-257.2}{B}{Equality of the Quillen, Volodin, Swan, Gersten and Vaserstein groups}
\GSfield{Kind; source}{theorem; \GSpage{P000422}}
\GSfield{The notes}{The framed theorem states \(K_iA=K_i^QA\). Above the right-hand K-symbol the labels list Swan, Gersten, two Volodin constructions, and Vaserstein.}
\GSfield{Statement}{For an associative ring \(A\) with identity and \(i\ge1\), the stable simplicial K-groups of the indicated Volodin and Vaserstein constructions, the free-resolution constructions of Swan and Gersten, and \(\pi_iBGL(A)^+\) are naturally isomorphic, when the nonunital intermediary rings and relative constructions use the respective published definitions. This does not identify Karoubi--Villamayor's polynomial-homotopy theory with Quillen's for arbitrary rings.}
\GSfield{Reconstruction}{The unital hypothesis and the restriction to positive groups are stated explicitly. The list is treated as a comparison of constructions, not a claim that their underlying spaces are identical before completion or delooping. The preceding Karoubi warning is essential to the scope.}
\GSfield{Fate}{The precise match is \citet[the first theorem block]{vaserstein1973foundations}; the detailed proof is \citet[\S15, \S27, Theorem 27.2, and \S29, Theorem 29.1]{vaserstein1976foundations}. The latter records first announcement in Gelfand's seminar in spring 1972 and in Tbilisi in autumn 1972; the short paper was received on 14 November 1972. Class B is relative to the conditional catalogue speaker Kazhdan; it does not imply an independent announcement distinct from Vaserstein's occurrences 1971-253.14 and 1971-256.3.}
\GSfield{Present status}{As of November 2026: resolved for the specified models and rings. GP-1971-21 links the earlier question and this report without duplicating its mathematical priority.}
\GSfield{Assessment}{literal statement: correct with missing hypotheses; publication match: exact; verification: read the result; development: annotated; confidence: medium; problem \hyperref[prob:GP-1971-21]{GP-1971-21}}
\end{GSitem}
\begin{GSitem}{1971-257.3}{B}{Acyclic coinvariants and the quotient comparison}
\GSfield{Kind; source}{method; \GSpage{P000422-P000423}}
\GSfield{The notes}{A diagram compares realization before and after quotienting the cone's simplicial matrix space by \(GL(MA)\). The final quotient is labeled \(BG^+\), with two exclamation marks. Bars and initial B-symbols are overwritten, so the individual nodes are not all diplomatically secure.}
\GSfield{Statement}{The proof uses an equivariant simplicial model and its chain complex, not an unrestricted rule that taking a quotient commutes with geometric realization. In Vaserstein's formulation, if every finite subset of a ring \(R\) lies in a unital subring, the coinvariant chain complex \((C_*S(R))_{GL(R)}\), with his stable-matrix notation understood, has vanishing positive-degree homology. Applying this acyclicity to the cone comparison, together with the kernel on \(\pi_1\), identifies the resulting connected target as a plus construction.}
\GSfield{Reconstruction}{The explicitly published acyclicity theorem is supplied as the rigorous content behind the sketch. The point-set quotient, simplicial quotient, and their realizations remain distinct until the hypotheses and comparison maps are verified. No overwritten B-symbol is silently resolved into a different proof.}
\GSfield{Fate}{See \citet[Theorem 26.1, Lemma 27.1 and proof of Theorem 27.2, pp.139--142]{vaserstein1976foundations}. The detailed chain and subdivision argument justifies the spring 1972 programme later announced in \citet{vaserstein1973foundations}; it is more than a formal diagram chase.}
\GSfield{Present status}{As of November 2026: established in the stated locally unital setting, which includes unital rings and the relevant finite matrix ideals. The notebook gives a proof outline, not the required equivariant argument in full.}
\GSfield{Assessment}{literal statement: correct with missing hypotheses; publication match: partial; verification: read the result; development: annotated; confidence: medium}
\end{GSitem}
\begin{GSitem}{1971-257.4}{B}{The identity component and the missing K-zero factor}
\GSfield{Kind; source}{theorem; \GSpage{P000423}}
\GSfield{The notes}{The page begins \(GL(\mu A)=K_0A\oplus\cdots\), then identifies a simplicial quotient with the identity component of the realized suspension space. The final boxes state \(K_iA=\pi_iGL(\mu A)\) and an equality of its homology with stable group homology.}
\GSfield{Statement}{For the suspension model \(S(\Sigma A)\), \(\pi_0S(\Sigma A)\cong K_0(A)\), and its identity component is a model of \(BGL(A)^+\). Thus, after choices of component identifications, the full K-space has homotopy type \(K_0(A)\times BGL(A)^+\), while \(K_i(A)=\pi_iS(\Sigma A)\) at the identity basepoint for \(i\ge1\). The homology equality is \(H_i(S(\Sigma A)_0;\mathbb Z)\cong H_i(BGL(A);\mathbb Z)\). It is not generally the homology of all components together.}
\GSfield{Reconstruction}{The scan's own identity-component annotation is retained and explains the restriction; it is not presented as wholly absent from the source. The direct-sum ellipsis is replaced by a homotopy-type statement, without asserting a canonical splitting of H-spaces. The right-hand group homology means homology of a classifying space.}
\GSfield{Fate}{The suspension and plus identifications are \citet[the first theorem and the paragraph preceding the fourth theorem, p.231]{vaserstein1973foundations} and \citet[\S18 and Theorem 27.2]{vaserstein1976foundations}. The modern distinction between \(K_0(A)\times BGL(A)^+\) and its connected factor is explicit in \citet[IV, Definition 1.1.1]{weibel2013kbook}.}
\GSfield{Present status}{As of November 2026: established with the component convention. Dropping it changes even \(H_0\), so the notebook's margin is mathematically significant.}
\GSfield{Assessment}{literal statement: correct with missing hypotheses; publication match: partial; verification: read the result; development: annotated; confidence: medium}
\end{GSitem}

\GStalk{258}{1972-03-20}{I. M. Gelfand}{P000423-P000424}{Quillen’s K-groups of finite fields and Frobenius}
Below an explicit divider and the heading ``20/III. Gelfand'', the notes ask for the K-groups of a finite field, state the low-degree groups and Quillen's general answer, and compare the algebraic closure with periodic K-theory. A brief continuation on P000424 mentions Frobenius and contains an unclear asymptotic remark. Kazhdan's named proof begins immediately afterward and belongs to record 259. The principal result is expository: it was already announced in Quillen's Nice ICM lecture in 1970 and printed in 1971, before the detailed 1972 Annals paper.

\begin{GSitem}{1971-258.1}{E}{The first three K-groups of a finite field}
\GSfield{Kind; source}{theorem; \GSpage{P000423}}
\GSfield{The notes}{For a finite field \(F_q\), the notes give \(K_0(F_q)=\mathbb Z\), \(K_1(F_q)=F_q^*\), and say that \(K_2(F_q)=0\) is easy to show.}
\GSfield{Statement}{Let \(q=p^a\) be a prime power. Dimension induces \(K_0(\mathbb F_q)\cong\mathbb Z\), determinant induces \(K_1(\mathbb F_q)\cong\mathbb F_q^\times\cong\mathbb Z/(q-1)\), and \(K_2(\mathbb F_q)=0\). The dimension and determinant identifications are canonical; an isomorphism of the multiplicative group with a chosen cyclic group requires choosing a generator.}
\GSfield{Reconstruction}{The prime-power hypothesis and the distinction between canonical descriptions and generator choices are explicit. No unrecorded elementary proof of the K-two assertion is attributed to Gelfand.}
\GSfield{Fate}{The low-degree groups were already part of the theory of Bass, Steinberg and Milnor; Quillen's announcement explicitly says that his definition agrees with Bass and Milnor \citep[\S3, p.50]{quillen1971cohomology}. They also follow from its finite-field computation, stated at the 1970 ICM and printed in 1971.}
\GSfield{Present status}{As of November 2026: classical and completely determined.}
\GSfield{Assessment}{literal statement: correct; publication match: exact; verification: read the result; development: annotated; confidence: high}
\end{GSitem}
\begin{GSitem}{1971-258.2}{E}{Quillen’s complete finite-field calculation}
\GSfield{Kind; source}{theorem; \GSpage{P000423}}
\GSfield{The notes}{The even groups vanish, and the odd groups are written as multiplicative groups of finite extension fields. The overwritten odd index admits both \(2i-1\) and \(2i+1\) readings, while the exponent placement on \(F_{q^i}\) is also uncertain.}
\GSfield{Statement}{For every prime power \(q\) and every integer \(i\ge1\), \[K_{2i}(\mathbb F_q)=0,\qquad K_{2i-1}(\mathbb F_q)\cong\mathbb Z/(q^i-1).\] In particular the odd group is abstractly isomorphic to \(\mathbb F_{q^i}^{\times}\). If indexed instead by \(2i+1\), its order is \(q^{i+1}-1\), not \(q^i-1\). These abstract cyclic-group descriptions by themselves do not specify the natural maps under field extensions.}
\GSfield{Reconstruction}{The index is fixed by the verified theorem rather than by guessing the overwritten sign. The equation with a multiplicative field group is taken as an abstract group isomorphism, without a canonical or Frobenius-equivariant identification.}
\GSfield{Fate}{\citet[\S3, Theorem 1 and its consequences, p.51]{quillen1971cohomology} announces the Adams-fibre equivalence from which these groups follow. The first circulation was the Nice ICM in 1970; the proceedings were printed in 1971. \citet{quillen1972finitefields} is the later full Annals paper, whose publisher metadata was checked here. The proof mechanisms are covered in occurrence 1971-259.3, with Brauer lifting in 1971-259.5 and the Lang comparison in 1971-259.6.}
\GSfield{Present status}{As of November 2026: completely resolved. The 1972 publication date must not be used to classify this March report as a later solution.}
\GSfield{Assessment}{literal statement: correct with missing hypotheses; publication match: partial; verification: read the result; development: annotated; confidence: high}
\end{GSitem}
\begin{GSitem}{1971-258.3}{E}{K-theory of the algebraic closure and its apparent periodicity}
\GSfield{Kind; source}{theorem; \GSpage{P000423}}
\GSfield{The notes}{The displayed formula over \(\overline F_q\) gives vanishing positive even groups and identifies odd groups with \(\overline F_q^*\), followed by ``like the periodicity theorem''. The first odd index is overwritten and provisionally \(2i+1\).}
\GSfield{Statement}{For \(k=\overline{\mathbb F}_p\), \(K_0(k)=\mathbb Z\), and for every \(i\ge1\), \[K_{2i}(k)=0,\qquad K_{2i-1}(k)\cong\bigoplus_{\ell\ne p}\mathbb Q_\ell/\mathbb Z_\ell.\] Each odd group is abstractly isomorphic to \(k^\times\). This similarity of abstract groups is not a canonical integral Bott-periodicity theorem, and does not identify the Frobenius actions in different odd degrees.}
\GSfield{Reconstruction}{The explicit prime-to-p divisible group and the warning about Frobenius distinguish the actual calculation from an unrestricted periodicity assertion. Including \(i=1\) incorporates the separately displayed K-one group.}
\GSfield{Fate}{This is \citet[\S3, Theorem 2, p.51]{quillen1971cohomology}, already circulated in 1970 and published in 1971. The same theorem specifies the degree-dependent Frobenius action, which is lost if all odd groups are simply labeled \(k^\times\).}
\GSfield{Present status}{As of November 2026: completely known. The notes' periodicity analogy is useful at the level of abstract abelian groups but insufficient for naturality.}
\GSfield{Assessment}{literal statement: correct with missing hypotheses; publication match: partial; verification: read the result; development: annotated; confidence: high}
\end{GSitem}
\begin{GSitem}{1971-258.4}{E}{Frobenius invariants recover the finite-field orders}
\GSfield{Kind; source}{method; \GSpage{P000424}}
\GSfield{The notes}{The continuation says only that the Frobenius automorphism is used.}
\GSfield{Statement}{For \(k=\overline{\mathbb F}_p\), arithmetic Frobenius \(x\mapsto x^p\) acts on \(K_{2i-1}(k)\) as multiplication by \(p^i\). Thus \(x\mapsto x^q\), with \(q=p^a\), acts as multiplication by \(q^i\). Extension of scalars identifies \(K_{2i-1}(\mathbb F_q)\) with the subgroup of Galois invariants in \(K_{2i-1}(k)\); this is the kernel of multiplication by \(q^i-1\), cyclic of that order. For finite intermediate extensions the restriction-of-scalars map corresponds to the norm on these invariant groups.}
\GSfield{Reconstruction}{The action, invariant subgroup and norm statement are supplied from the matching theorem to explain the short reference to Frobenius. They are not asserted to have been written out in the seminar. Arithmetic Frobenius is fixed explicitly so that its inverse convention cannot reverse the action.}
\GSfield{Fate}{All the functorial assertions appear in \citet[\S3, Theorem 2, p.51]{quillen1971cohomology}. They sharpen the group-only formulas in the preceding two items and were already available before the lecture.}
\GSfield{Present status}{As of November 2026: established. In degree \(2i-1\), identifying the group with \(k^\times\) while retaining the usual Frobenius power \(q\) would be wrong unless \(i=1\).}
\GSfield{Assessment}{literal statement: correct with missing hypotheses; publication match: partial; verification: read the result; development: annotated; confidence: high}
\end{GSitem}
\begin{GSitem}{1971-258.5}{F}{An unclear statement about increasing degree}
\GSfield{Kind; source}{remark; \GSpage{P000424}}
\GSfield{The notes}{Before the Frobenius sentence, the notes say that as \(i\) increases the K-functor approaches an elementary K-functor, followed by an unreadable or abbreviated final word.}
\GSfield{Statement}{No limiting operation, comparison map, stable range, or precise elementary functor is specified. The assertion cannot be reconstructed uniquely as a theorem about increasing K-degree.}
\GSfield{Reconstruction}{The nearby finite-field formulas and periodicity analogy do not supply the missing notion of approximation. No equivalence with Milnor K-theory or a stability theorem is substituted for the unclear wording.}
\GSfield{Fate}{The verified Quillen announcement explains the surrounding calculation but does not identify this sentence. Searches of Quillen's finite-field and Segal-related accounts did not recover a matching degree-asymptotic assertion with the omitted parameters.}
\GSfield{Present status}{As of November 2026: unidentified from this brief continuation. It is not evidence for a presently open stabilization problem.}
\GSfield{Assessment}{literal statement: not determinable; publication match: none found; verification: unresolved lead; development: annotated; confidence: low}
\end{GSitem}

\GStalk{259}{1972-03-20}{D. A. Kazhdan}{P000424-P000426}{Quillen’s finite-field calculation: plus construction and Adams operations}
Kazhdan explains the proof of the finite-field calculation announced by Gelfand in record 258. The argument passes from a classifying space through a homology-preserving modification to a homotopy fibre of an Adams operation, using Brauer lifting and an étale comparison. Fuchs interrupts on P000424; Kazhdan resumes on P000425. The catalogue's P000423--P000425 interval starts too early and omits the proof's final page. Several apparent finite-rank statements actually require stabilization; the asserted identification of a fixed subgroup with maps into a homotopy fibre is a substantive error. These constructions had already been announced by Quillen at the 1970 ICM, printed in 1971; the later Annals date must not be mistaken for their first circulation.

\begin{GSitem}{1971-259.1}{E}{Stable plus construction for finite-field K-theory}
\GSfield{Kind; source}{construction; \GSpage{P000424-P000425}}
\GSfield{The notes}{For \(G_n=GL_n(\mathbb F_q)\), the notes enlarge \(K(G_n,1)\) to \(B_n\), kill the commutator subgroup, fill the resulting two-dimensional cycles, and write \(K_i(G_n)=\pi_i(B_n)\). The enlarged space is called an \(H\)-space.}
\GSfield{Statement}{Put \(G=\operatorname*{colim}_n GL_n(\mathbb F_q)\), and let \(E=[G,G]\), the stable elementary subgroup. The plus construction \(BG\to BG^+\) kills the perfect normal subgroup \(E\) and preserves integral homology, indeed homology with local coefficients pulled back from \(G/E\). For \(i\geq1\), \(K_i(\mathbb F_q)=\pi_i(BG^+)\). Block sum supplies the stable \(H\)-space structure. The construction can be made by attaching cells in dimensions two and three.}
\GSfield{Reconstruction}{The stable group, the perfectness condition, and the distinction between attaching three-cells and ``three-cycles'' are supplied. No claim that a fixed finite \(n\) computes all higher \(K\)-groups is retained.}
\GSfield{Fate}{The construction and its role in finite-field \(K\)-theory were already announced at the 1970 ICM and printed in 1971 \citep[Section 3, p. 50]{quillen1971cohomology}. The full Annals treatment appeared in November 1972 \citep{quillen1972finitefields}; the later publication does not make this a post-seminar discovery.}
\GSfield{Present status}{As of November 2026, the stable construction is standard. Fixed-rank spaces require a specified stability range, absent from the notebook.}
\GSfield{Assessment}{literal statement: correct with missing hypotheses; publication match: exact; verification: read the result; development: none; confidence: medium}
\end{GSitem}
\begin{GSitem}{1971-259.2}{E}{Adams operations represented on reduced complex K-theory}
\GSfield{Kind; source}{construction; \GSpage{P000425}}
\GSfield{The notes}{The power sum \(\sum_i\lambda_i^k\) of eigenvalues of a unitary matrix is a virtual character. It defines \(\psi^k\) on bundles and hence a self-map of \(BU\); the notes write \(K^0(X)=[X,BU]\), with a marginal qualification about quotienting trivial bundles.}
\GSfield{Statement}{For a pointed connected finite CW complex, \(\widetilde K^0(X)\cong[X,BU]_*\); unreduced \(K^0\) is represented by \(\mathbb Z\times BU\). For \(k\geq1\), Newton polynomials in the exterior-power operations define a natural ring endomorphism \(\psi^k\), characterized on line bundles by \(\psi^k[L]=[L^{\otimes k}]\). The character on a rank-\(n\) bundle is the indicated power sum.}
\GSfield{Reconstruction}{Reduced theory and a finite-CW domain are made explicit; a virtual representation need not be an actual bundle of nonnegative rank.}
\GSfield{Fate}{Adams operations were established before 1972 and are treated in \citet[Chapter II]{atiyah1967ktheory}. Kazhdan is using this existing apparatus in Quillen's calculation.}
\GSfield{Present status}{As of November 2026, the construction and normalization remain standard; the rank component must be included for unreduced theory.}
\GSfield{Assessment}{literal statement: correct with missing hypotheses; publication match: exact; verification: publisher metadata checked; development: none; confidence: high}
\end{GSitem}
\begin{GSitem}{1971-259.3}{E}{Finite-field K-theory as the fibre of an Adams operation}
\GSfield{Kind; source}{theorem; \GSpage{P000425-P000426}}
\GSfield{The notes}{The claim is \(E\psi^q=B_n=B\) ``for large \(n\)'', where \(E\psi^q\) is the fibre of \(\psi^q-1:BU\to BU\). Frobenius is identified with the topological Adams operation.}
\GSfield{Statement}{For a prime power \(q\), the stable space \(BGL(\mathbb F_q)^+\) has the homotopy type of \(\operatorname{hofib}(\psi^q-1:BU\to BU)\). Since \(\psi^q\) acts on \(\pi_{2j}(BU)\cong\mathbb Z\) by multiplication by \(q^j\), its fibre has \(\pi_{2j-1}\cong\mathbb Z/(q^j-1)\) and \(\pi_{2j}=0\) for \(j\geq1\).}
\GSfield{Reconstruction}{``Fibre'' means homotopy fibre, and the equality is a homotopy equivalence after stabilization. The long exact homotopy sequence explains the numerical result rather than assuming a fixed-rank identification.}
\GSfield{Fate}{The construction and its role in finite-field \(K\)-theory were already announced at the 1970 ICM and printed in 1971 \citep[Section 3, Theorem 1, p. 51]{quillen1971cohomology}. The full Annals treatment appeared in November 1972 \citep{quillen1972finitefields}; the later publication does not make this a post-seminar discovery.}
\GSfield{Present status}{As of November 2026, the equivalence and resulting finite-field groups are established; the manuscript's finite-rank shorthand remains unsafe without a stability bound.}
\GSfield{Assessment}{literal statement: correct with missing hypotheses; publication match: exact; verification: read the result; development: none; confidence: medium}
\end{GSitem}
\begin{GSitem}{1971-259.4}{E}{A homotopy fibre does not represent the bare fixed subgroup}
\GSfield{Kind; source}{claim; \GSpage{P000425}}
\GSfield{The notes}{The notes assert that \([K(X)]^k=\{\xi:\psi^k\xi=\xi\}\) is represented by \([X,E\psi^k]\), and call this easy to check.}
\GSfield{Statement}{Let \(F_k=\operatorname{hofib}(\psi^k-1:BU\to BU)\). For finite pointed CW complexes there is an exact sequence of abelian groups \[\widetilde K^{-1}(X)\xrightarrow{\psi^k-1}\widetilde K^{-1}(X)\longrightarrow[X,F_k]_*\longrightarrow\widetilde K^0(X)\xrightarrow{\psi^k-1}\widetilde K^0(X).\] Thus a map into \(F_k\) includes a chosen homotopy from \(\psi^k\xi\) to \(\xi\). For \(X=S^1\) and \(k>2\), the fixed subgroup in reduced degree zero is zero, whereas \([S^1,F_k]_*\cong\mathbb Z/(k-1)\).}
\GSfield{Reconstruction}{\ours{} The exact sequence and the circle counterexample are an editorial check of the literal assertion. They use \(\psi^k\) acting by \(k\) on \(\widetilde K^{-1}(S^1)\cong\widetilde K^0(S^2)\).}
\GSfield{Fate}{The distinction follows from the ordinary homotopy-fibre sequence and the pre-existing Adams-operation formalism in \citet{atiyah1967ktheory}. It is not an unresolved part of Quillen's theorem.}
\GSfield{Present status}{As of November 2026, the literal representability statement is false; retaining the null-homotopy data repairs it.}
\GSfield{Assessment}{literal statement: false as written; publication match: partial; verification: indirect; development: annotated; confidence: high}
\end{GSitem}
\begin{GSitem}{1971-259.5}{E}{Brauer lifting and Frobenius invariance}
\GSfield{Kind; source}{method; \GSpage{P000426}}
\GSfield{The notes}{Choose \(\mu:\overline{\mathbb F}_q^*\to\mathbb C^*\), form \(\varphi_\mu(g)=\sum_i\mu(\lambda_i)\), and invoke Brauer's theorem to obtain a virtual character. The equality of the multisets \(\{\lambda_i\}\) and \(\{\lambda_i^q\}\) gives Adams invariance; a map is then extended over the added cells.}
\GSfield{Statement}{Choose compatible lifts of roots of unity of order prime to \(p=\operatorname{char}\mathbb F_q\). Brauer lifting of the defining representation produces a stable map to complex \(K\)-theory. Frobenius invariance is the relation used to lift this map to the homotopy fibre of \(\psi^q-1\). For a unipotent matrix of size \(n\), the displayed sum of eigenvalue lifts equals \(n\), not \(1\); the corresponding reduced class has rank zero.}
\GSfield{Reconstruction}{The Brauer character, its stable realization, and the coherent homotopy required for the fibre lift must be distinguished. The manuscript's pointwise character argument alone does not construct that coherence or justify extension across every cell.}
\GSfield{Fate}{The construction and its role in finite-field \(K\)-theory were already announced at the 1970 ICM and printed in 1971 \citep[Section 3, pp. 50--51]{quillen1971cohomology}. The full Annals treatment appeared in November 1972 \citep{quillen1972finitefields}; the later publication does not make this a post-seminar discovery.}
\GSfield{Present status}{As of November 2026, the method is established. The claim \(\varphi_\mu=1\) on unipotents is incompatible with the displayed unnormalized rank-\(n\) formula.}
\GSfield{Assessment}{literal statement: false as written; publication match: partial; verification: read the result; development: none; confidence: medium}
\end{GSitem}
\begin{GSitem}{1971-259.6}{E}{The Lang map and étale comparison}
\GSfield{Kind; source}{method; \GSpage{P000426}}
\GSfield{The notes}{To compute the homology of \(K(G_n,1)\), the notes propose \(g\mapsto g^{-1}g^\sigma\), with \(\sigma\) Frobenius, followed by the spectral sequence of a fibration.}
\GSfield{Statement}{For a connected algebraic group \(G\) over \(\mathbb F_q\), the Lang map \(L(g)=g^{-1}F(g)\) is surjective and, for smooth \(G\), is a finite étale torsor under the finite group \(G(\mathbb F_q)\) acting on the left. Applied to general linear groups, it gives an étale analogue of the fibration used in the comparison with the Adams-operation fibre. A cohomological comparison requires coefficients prime to \(p\), followed by the separate treatment implicit in the finite-field theorem.}
\GSfield{Reconstruction}{The connectedness, smoothness and prime-to-characteristic qualifications are supplied. The spectral sequence is a method, not an explicit calculation recorded on this page.}
\GSfield{Fate}{The construction and its role in finite-field \(K\)-theory were already announced at the 1970 ICM and printed in 1971 \citep[Section 2, pp. 48--50]{quillen1971cohomology}. The full Annals treatment appeared in November 1972 \citep{quillen1972finitefields}; the later publication does not make this a post-seminar discovery.}
\GSfield{Present status}{As of November 2026, the method is standard; the manuscript does not specify enough terms or differentials to reconstruct its entire spectral sequence independently.}
\GSfield{Assessment}{literal statement: correct with missing hypotheses; publication match: partial; verification: read the result; development: none; confidence: medium}
\end{GSitem}

\GStalk{260}{1972-03-20}{D. B. Fuchs}{P000424}{When abelianizing the fundamental group preserves integral homology}
Fuchs inserts an existence theorem into Kazhdan's discussion: the commutator subgroup must be relatively perfect. He distinguishes existence from uniqueness and then invokes simple spaces and Whitehead's theorem. All of the identifiable intervention is on P000424, not throughout the catalogue's P000423--P000427 range.

\begin{GSitem}{1971-260.1}{B}{Relative perfectness and homology-preserving abelianization}
\GSfield{Kind; source}{theorem; \GSpage{P000424}}
\GSfield{The notes}{The condition is \([\pi_1X,\pi_1X]=[\pi_1X,[\pi_1X,\pi_1X]]\). It is said to be necessary and sufficient for a map \(X\to Y\) inducing integral homology isomorphisms with \(\pi_1Y\) abelian.}
\GSfield{Statement}{For a connected CW complex \(X\), write \(G=\pi_1X\) and \(N=[G,G]\). There is a CW complex \(Y\) and a map \(f:X\to Y\) inducing the quotient \(G\to G/N\) and an isomorphism on integral homology if and only if \([G,N]=N\). The space can be constructed by attaching two- and three-cells. This is relative perfectness, which is weaker than the perfect-kernel hypothesis \([N,N]=N\) for an acyclic plus construction.}
\GSfield{Reconstruction}{Surjectivity and the specified quotient on fundamental groups are made explicit. Ordinary integral homology equivalence must not be silently strengthened to an acyclic map with all local coefficients.}
\GSfield{Fate}{The exact criterion appears in \citet[Corollary 4.3]{ye2012unified}; its preprint circulated in 2011 and the journal article appeared in 2012. That later verification does not establish priority for the 1972 oral statement. The attached erratum concerns a different application, not this corollary.}
\GSfield{Present status}{As of November 2026, the criterion is established. Confusing relative perfectness with perfectness would wrongly reject the notebook's formula.}
\GSfield{Assessment}{literal statement: correct; publication match: exact; verification: read the result; development: none; confidence: high}
\end{GSitem}
\begin{GSitem}{1971-260.2}{F}{Simple targets and the limits of the uniqueness assertion}
\GSfield{Kind; source}{claim; \GSpage{P000424}}
\GSfield{The notes}{After saying the target \(Y\) need not be unique, Fuchs adds that a simple target is unique, citing Whitehead's theorem; he recalls that an \(H\)-space is simple.}
\GSfield{Statement}{A path-connected \(H\)-space is simple: its fundamental group is abelian and acts trivially on higher homotopy groups. A homology equivalence between connected simple CW complexes is a weak homotopy equivalence. This criterion proves uniqueness only after a comparison map between the proposed targets has been constructed; simplicity by itself is not the missing universal property of the preceding existence construction.}
\GSfield{Reconstruction}{The first two assertions are standard precise statements. The notes do not specify the category of target maps, a universal property, or the comparison map needed to infer uniqueness; that inference is therefore left unresolved rather than identified with the usual plus-construction uniqueness theorem.}
\GSfield{Fate}{The homotopy-theoretic ingredients are classical; the cited Whitehead principle is discussed in \citet[Chapters 10 and 15]{may1999concise}. No exact publication matching this particular uniqueness formulation was identified in the searches for Fuchs, simple targets and homology localization.}
\GSfield{Present status}{As of November 2026, the elementary ingredients are established, but the exact uniqueness assertion recoverable from this intervention remains under-specified; this is not a claim of an open problem.}
\GSfield{Assessment}{literal statement: not determinable; publication match: partial; verification: indirect; development: none; confidence: medium}
\end{GSitem}

\GStalk{261}{1972-03-20}{M. Sh. Birman}{P000427}{Trace-class perturbations and Gelfand’s spectral classification question}
Birman begins a stationary proof of the Kato--Rosenblum theorem and identifies two boundary-value facts used in it. Gelfand asks which prescribed spectral types can be realized with a trace-class difference, and what the spectrum in the corresponding operator quotient means. The next dated heading, Kazhdan on 27 March, begins at the bottom of this same page.

\begin{GSitem}{1971-261.1}{E}{The Kato–Rosenblum theorem}
\GSfield{Kind; source}{theorem; \GSpage{P000427}}
\GSfield{The notes}{If \(H,H_0\) are self-adjoint and \(V=H-H_0\) is nuclear, their absolutely continuous parts are unitarily equivalent.}
\GSfield{Statement}{Let \(H_0\) be self-adjoint on a separable Hilbert space and let \(V=V^*\in\mathcal S_1\) be bounded trace class. Set \(H=H_0+V\) on \(D(H_0)\). The wave operators \(W_\pm(H,H_0)\) exist and are complete, and intertwine the restrictions to the absolutely continuous spectral subspaces. Thus the absolutely continuous spectral measures and multiplicity functions agree up to unitary equivalence.}
\GSfield{Reconstruction}{A common-domain bounded-perturbation formulation is supplied. The conclusion does not assert equivalence of the singular spectral parts.}
\GSfield{Fate}{The trace-class theorem predates the talk. A stationary proof with precisely this hypothesis is given by \citet[Theorem 1, Section 3]{birmanentina1967stationary}; Section 5 identifies the stationary and time-dependent wave operators.}
\GSfield{Present status}{As of November 2026, the theorem is standard; trace-class control is a sufficient condition for complete scattering, not a classification of singular spectra.}
\GSfield{Assessment}{literal statement: correct; publication match: exact; verification: read the result; development: none; confidence: high}
\end{GSitem}
\begin{GSitem}{1971-261.2}{B}{Which spectral types admit trace-class-close realizations?}
\GSfield{Kind; source}{question; \GSpage{P000427}; I. M. Gelfand}
\GSfield{The notes}{Gelfand asks how to decide from a collection of spectral types whether realizations can be chosen whose difference is nuclear.}
\GSfield{Statement}{Given two unitary-equivalence classes of bounded self-adjoint operators \(A\) and \(B\) on separable Hilbert spaces of the same dimension, characterize the existence of a unitary \(U\) for which \(A-UBU^*\in\mathcal S_1\). The data must include spectral multiplicity, not merely the closed spectral sets. The Kato--Rosenblum theorem gives equality of the absolutely continuous parts as a necessary condition; it alone does not settle the question.}
\GSfield{Reconstruction}{Boundedness and the unitary-equivalence formulation make the phrase ``collection of spectral types'' precise. The notes do not say whether unbounded operators, a prescribed trace norm, or exact singular spectral measures are intended.}
\GSfield{Fate}{\citet{careypincus1976traceclass} published a paper specifically on unitary equivalence modulo trace class in 1976. Its bibliographic record is verified, but its full two-operator criterion was not retrieved. A narrower consequence is stated explicitly by \citet[slide 8]{xiaundateddiagonalization}: every bounded purely singular self-adjoint operator on a separable Hilbert space is diagonalizable modulo a self-adjoint trace-class perturbation of arbitrarily small trace norm. This verifies substantive later progress, not every quantifier of the notebook's classification question.}
\GSfield{Present status}{As of November 2026, the subject has a dedicated later classification literature and the purely singular diagonalization case is verified indirectly. This audit has not independently checked the complete necessary-and-sufficient criterion for a prescribed pair of spectral types, and does not assert that equality of absolutely continuous parts alone suffices.}
\GSfield{Assessment}{literal statement: correct with missing hypotheses; publication match: partial; verification: indirect; development: none; confidence: medium; problem \hyperref[prob:GP-1971-40]{GP-1971-40}}
\end{GSitem}
\begin{GSitem}{1971-261.3}{E}{Spectrum in the algebraic quotient by trace-class operators}
\GSfield{Kind; source}{question; \GSpage{P000427}; I. M. Gelfand}
\GSfield{The notes}{Gelfand asks: ``Spectrum in the quotient ring modulo the nuclear operators --- ?}
\GSfield{Statement}{For a bounded operator \(T\) on an infinite-dimensional Hilbert space, define the spectrum of its coset in the unital algebra \(\mathcal B(\mathcal H)/\mathcal S_1\) algebraically, using two-sided invertibility. This spectrum equals the Fredholm essential spectrum. Indeed an inverse modulo \(\mathcal S_1\) is an inverse modulo compact operators; conversely a Fredholm operator has a parametrix with finite-rank errors, hence trace-class errors.}
\GSfield{Reconstruction}{\ours{} This is one precise natural reading of the question, not a reconstruction of a proposed new invariant. The quotient is not a Banach algebra for the operator-norm quotient topology because \(\mathcal S_1\) is not operator-norm closed.}
\GSfield{Fate}{The argument is the classical Fredholm-parametrix criterion, already available before the talk; it is included as an elementary editorial clarification of what the quotient-ring spectrum can and cannot distinguish. The surrounding scattering context is documented in \citet{birmanentina1967stationary}.}
\GSfield{Present status}{As of November 2026, under this algebraic interpretation the answer is the essential spectrum. A finer trace-class invariant would require additional structure beyond ordinary quotient-algebra invertibility.}
\GSfield{Assessment}{literal statement: correct with missing hypotheses; publication match: thematic only; verification: indirect; development: annotated; confidence: high}
\end{GSitem}
\begin{GSitem}{1971-261.4}{E}{Sandwiched spectral derivatives and resolvent boundary values}
\GSfield{Kind; source}{theorem; \GSpage{P000427}}
\GSfield{The notes}{For Hilbert--Schmidt operators \(D_1,D_2\), the derivative of \(D_1E_\lambda D_2\) is nuclear almost everywhere, and \(D_1R_{\lambda+i\varepsilon}D_2\) has an almost-everywhere Hilbert--Schmidt boundary value as \(\varepsilon\downarrow0\).}
\GSfield{Statement}{Let \(H\) be self-adjoint with spectral resolution \(E\), and let \(D_1,D_2\in\mathcal S_2\). The function \(\lambda\mapsto D_1E(( -\infty,\lambda])D_2\) has bounded variation in \(\mathcal S_1\) and an almost-everywhere derivative in trace norm. For almost every real \(\lambda\), the sandwiched resolvents \(D_1(H-\lambda\mp i\varepsilon)^{-1}D_2\) have limits in \(\mathcal S_2\) as \(\varepsilon\downarrow0\). These statements do not assert an unsandwiched norm limit for the resolvent.}
\GSfield{Reconstruction}{The topology of convergence, the almost-everywhere qualification and the use of a sandwiched spectral measure are made explicit.}
\GSfield{Fate}{These are the auxiliary stationary-scattering facts developed in \citet[Section 2]{birmanentina1967stationary}. The notes correctly identify the Hilbert--Schmidt rigging that makes them usable.}
\GSfield{Present status}{As of November 2026, the facts remain standard tools for the stationary trace-class method.}
\GSfield{Assessment}{literal statement: correct; publication match: exact; verification: read the result; development: none; confidence: high}
\end{GSitem}
\begin{GSitem}{1971-261.5}{E}{A stationary formula for the intertwining operator}
\GSfield{Kind; source}{method; \GSpage{P000427}}
\GSfield{The notes}{An integral of the derivative of \((E_\lambda(I+VR_z^0)f,g)\), evaluated at \(z=\lambda\pm i0\), is proposed for \((W_\pm f,g)\). The perturbation \(V\) is factored into two Hilbert--Schmidt operators.}
\GSfield{Statement}{For a self-adjoint trace-class perturbation, factor \(V\) through \(\mathcal S_2\), use boundary values of the sandwiched resolvents to define stationary wave operators on the absolutely continuous subspace, and extend by zero on its orthogonal complement. The resulting operators are partial isometries with initial space \(\mathcal H_{\rm ac}(H_0)\) and final space \(\mathcal H_{\rm ac}(H)\).}
\GSfield{Reconstruction}{The notebook integral is a schematic weak formula: spectral-density differentiation, boundary evaluation and integration cannot be interchanged without the factorization estimates. The precise construction, rather than a falsely pointwise formula for every vector, is the modern statement.}
\GSfield{Fate}{The construction is explicit in \citet[Section 3, formula (3.17) and Theorem 1]{birmanentina1967stationary}, five years before this talk.}
\GSfield{Present status}{As of November 2026, the stationary construction and its equivalence to time-dependent scattering are established.}
\GSfield{Assessment}{literal statement: correct with missing hypotheses; publication match: partial; verification: read the result; development: none; confidence: high}
\end{GSitem}

\GStalk{262}{1972-03-27}{D. A. Kazhdan}{P000427-P000428}{Artin–Mazur étale homotopy and arithmetic symmetry}
Kazhdan outlines an algebraic construction of the homotopy type of a variety, its comparison over the complex numbers, and the action of field automorphisms. The completion paragraph runs onto P000428. The ordinary-space terminology in the notes suppresses the pro-homotopy category essential to Artin--Mazur's formulation.

\begin{GSitem}{1971-262.1}{E}{Étale homotopy type and complex comparison}
\GSfield{Kind; source}{construction; \GSpage{P000427}}
\GSfield{The notes}{An algebraic \(X\) over any field produces a complex \(C_X\); over \(\mathbb C\) it is the completion of the ordinary complex space.}
\GSfield{Statement}{A locally noetherian scheme has an Artin--Mazur \(\acute{e}\)tale homotopy type in a pro-homotopy category, constructed from connected components of \(\acute{e}\)tale hypercovers. For a connected pointed scheme \(X\) of finite type over \(\mathbb C\), the canonical comparison from the classical analytic homotopy type to \(X_{\acute{e}t}\) induces an isomorphism after profinite completion. If \(X\) is additionally geometrically unibranch---in particular, if \(X\) is smooth or normal---then \(X_{\acute{e}t}\) itself is isomorphic to the profinite completion of the classical type in that pro-homotopy category.}
\GSfield{Reconstruction}{The hypotheses on the scheme and its base point are supplied from the primary theorem. Finite type and geometric unibranchness have distinct roles in the two comparison statements. The notebook's informal ``homotopy type'' is not identified with an ordinary integral CW homotopy type.}
\GSfield{Fate}{The construction circulated in the Harvard seminar of 1965--1966 and appeared as \citet{artinmazur1969etale}. The precise comparison is Theorem 12.9 and Corollary 12.10, pp. 142--143, read in the full primary text.}
\GSfield{Present status}{As of November 2026, the construction is a foundation of étale homotopy theory. The notebook's shorthand is recognizable but not a complete theorem statement.}
\GSfield{Assessment}{literal statement: correct with missing hypotheses; publication match: exact; verification: read the result; development: none; confidence: high}
\end{GSitem}
\begin{GSitem}{1971-262.2}{E}{Galois action on geometric étale homotopy}
\GSfield{Kind; source}{construction; \GSpage{P000427}}
\GSfield{The notes}{Field automorphisms preserving \(X\) act on \(C_X\); for a variety defined over \(\mathbb Q\), the group is written \(\operatorname{Gal}(\mathbb C\mid\mathbb Q)\).}
\GSfield{Statement}{If \(X\) is defined over a field \(k\) with separable closure \(\bar k\), functoriality gives a homotopy-coherent action of \(\operatorname{Gal}(\bar k/k)\) on the geometric étale homotopy type of \(X_{\bar k}\). For pointed homotopy groups, a geometric point and base-point transport are required; the fundamental-group action is naturally outer until those choices are made.}
\GSfield{Reconstruction}{The arithmetic Galois group and the distinctions between a functorial action, an action on a chosen model, and an outer action on a fundamental group are made explicit. The huge abstract automorphism group of \(\mathbb C\) is not treated as a group of continuous automorphisms of the analytic space.}
\GSfield{Fate}{The arithmetic action is already discussed in the introduction, especially p. 4, of \citet{artinmazur1969etale}; its functorial construction belongs to their 1965--1966 seminar work published in 1969. The pointed-versus-outer-action distinction is our explicit formulation of the base-point issue.}
\GSfield{Present status}{As of November 2026, Galois actions on geometric étale homotopy types are established; the model and base-point qualifications remain essential.}
\GSfield{Assessment}{literal statement: correct with missing hypotheses; publication match: partial; verification: read the result; development: none; confidence: medium}
\end{GSitem}
\begin{GSitem}{1971-262.3}{E}{Completion is represented in a pro-homotopy category}
\GSfield{Kind; source}{construction; \GSpage{P000428}}
\GSfield{The notes}{The notes collect maps from \(X\) to spaces with finite homotopy groups and say that a compatible mapping functor is represented by \(\widehat X\).}
\GSfield{Statement}{Profinite homotopy completion is characterized by maps into finite homotopy types, ordinarily meaning spaces with finitely many components, finitely many nonzero homotopy groups and each homotopy group finite. The universal object belongs to a pro-homotopy category. Maps from a test object into an inverse system and maps from an inverse system into a test object have different limit/colimit formulas; they must not be conflated.}
\GSfield{Reconstruction}{The finite-type test category and pro-representability are a reconstruction. The notebook's product of compatible \(\operatorname{Hom}(Y,X_i)\) describes maps to a given pro-system, not by itself the universal construction of the completion of \(X\).}
\GSfield{Fate}{The pro-category and completion are constructed in \citet[Sections 2--3, pp. 20--34]{artinmazur1969etale}; the complex comparison is Theorem 12.9. This is prior foundational work, not a conjectural completion procedure first proposed in the seminar.}
\GSfield{Present status}{As of November 2026, the pro-categorical formulation is standard. Ordinary representability by one cell complex is not justified by the text.}
\GSfield{Assessment}{literal statement: correct with missing hypotheses; publication match: partial; verification: read the result; development: none; confidence: medium}
\end{GSitem}

\GStalk{264}{1972-03-27}{I. M. Gelfand}{P000428}{A reading plan for the seminar’s common categorical language}
Gelfand asks participants to learn categories and functors from the first chapter of Grothendieck's Tohoku paper, then topoi and descent from the material accompanying Mumford. This is a reading plan, with no separately stated theorem. It precedes Birman's continuation on the page, although its catalogue record number is later. The plan connects the earlier site and sheaf discussions (for example 1967-143 and 1970-226) to Geronimus's April lectures.


\GStalk{263}{1972-03-27}{M. Sh. Birman}{P000428-P000430}{Stationary and time-dependent wave operators}
Birman continues the trace-class scattering proof, explains completeness through reverse wave operators and the chain rule, and derives a resolvent identity controlling spectral densities. He explicitly refers to the 1967 Izvestiya paper. A final sketch gestures toward locally trace-class perturbations and unbounded operators; its missing domain hypotheses are retained as a limitation.

\begin{GSitem}{1971-263.1}{E}{Time-dependent wave operators, intertwining and completeness}
\GSfield{Kind; source}{theorem; \GSpage{P000428-P000429}}
\GSfield{The notes}{The strong limits \(W_\pm=\lim_{t\to\pm\infty}e^{itH}e^{-itH_0}P_0\) are defined, with \(P_0\) the absolutely continuous projection. The notes state \(W_\pm H_0=HW_\pm\) and say completeness is equivalent to existence of the reverse wave operator.}
\GSfield{Statement}{For self-adjoint \(H,H_0\), if the indicated strong limit exists, it is an isometry on \(\mathcal H_{\rm ac}(H_0)\), vanishes on its orthogonal complement, and intertwines bounded Borel functions of the operators. It is complete precisely when the corresponding reverse wave operator exists and has adjoint \(W_\pm(H,H_0)^*\). Under the common-domain trace-class perturbation hypothesis, existence holds in both directions and hence completeness follows.}
\GSfield{Reconstruction}{The adjoint relation and the spectral-projection meaning of the unbounded intertwining identity are supplied. Existence alone is not completeness for arbitrary nonsymmetric hypotheses.}
\GSfield{Fate}{These are established scattering-theory facts; \citet[Sections 1 and 5]{birmanentina1967stationary} give the stationary/time-dependent comparison used in the talk.}
\GSfield{Present status}{As of November 2026, the statements are standard with the stated absolutely continuous projections and domains.}
\GSfield{Assessment}{literal statement: correct with missing hypotheses; publication match: exact; verification: read the result; development: none; confidence: high}
\end{GSitem}
\begin{GSitem}{1971-263.2}{E}{The chain rule for wave operators}
\GSfield{Kind; source}{theorem; \GSpage{P000429}}
\GSfield{The notes}{The displayed composition is \(W_\pm(H_3,H_2)W_\pm(H_2,H_1)=W_\pm(H_3,H_1)\).}
\GSfield{Statement}{Let \(H_1,H_2,H_3\) be self-adjoint on the same Hilbert space, and use the ordinary wave operators \(W_\pm(H_j,H_i)=\operatorname*{s-lim}_{t\to\pm\infty}e^{itH_j}e^{-itH_i}P_{\rm ac}(H_i)\). If \(W_\pm(H_3,H_2)\) and \(W_\pm(H_2,H_1)\) exist for a fixed sign, then \(W_\pm(H_3,H_1)\) exists and equals their product. The range of the second factor lies in \(\mathcal H_{\rm ac}(H_2)\), which permits insertion of the intermediate projection. Completeness of both factors implies completeness of the product.}
\GSfield{Reconstruction}{A fixed common sign is required. The insertion of the intermediate absolutely continuous projection is the point that makes the strong-limit multiplication legitimate.}
\GSfield{Fate}{The chain rule belongs to the general wave-operator framework summarized in \citet[Section 1]{birmanentina1967stationary}; it is not a new consequence first obtained in 1972.}
\GSfield{Present status}{As of November 2026, the chain rule remains valid in this setting and is frequently used to reduce perturbation comparisons to intermediate operators.}
\GSfield{Assessment}{literal statement: correct; publication match: exact; verification: read the result; development: none; confidence: high}
\end{GSitem}
\begin{GSitem}{1971-263.3}{E}{Resolvent factorization and spectral localization of the stationary formula}
\GSfield{Kind; source}{method; \GSpage{P000429-P000430}}
\GSfield{The notes}{With \(Q_z=I+VR_z^0\), Birman writes \(R_z^0-R_{\bar z}^0=Q_z^*(R_z-R_{\bar z})Q_z\). The upper integration bound \(\sigma\) gives \((E_\sigma W_\pm f,g)\), and the construction is independent of the Hilbert--Schmidt factorization of \(V\).}
\GSfield{Statement}{Let \(H_0\) be self-adjoint, \(V=V^*\in\mathcal S_1\), and \(H=H_0+V\) on \(D(H_0)\). For \(z\notin\mathbb R\), put \(R_z=(H-z)^{-1}\), \(R_z^0=(H_0-z)^{-1}\), and \(Q_z=I+VR_z^0\). The resolvent identity gives \(R_zQ_z=R_z^0\). Combining this with \(R_z-R_{\bar z}=(z-\bar z)R_z^*R_z\) gives \(R_z^0-R_{\bar z}^0=Q_z^*(R_z-R_{\bar z})Q_z\). After Hilbert--Schmidt sandwiching and almost-everywhere boundary passage, this identifies the absolutely continuous spectral densities and makes the stationary partial isometry respect spectral projections. Uniqueness of the resulting weak form yields independence of the factorization.}
\GSfield{Reconstruction}{The formula is checked with the notes' resolvent convention. A density statement is made only after rigging; it is not an operator-norm identity for unsandwiched boundary resolvents.}
\GSfield{Fate}{The argument is the stationary proof in \citet[Section 3, Theorem 1]{birmanentina1967stationary}, the paper explicitly named in the notebook.}
\GSfield{Present status}{As of November 2026, the derivation is established and the resolvent identity is correct with this convention.}
\GSfield{Assessment}{literal statement: correct; publication match: exact; verification: read the result; development: none; confidence: high}
\end{GSitem}
\begin{GSitem}{1971-263.4}{E}{Birman’s local trace-class criterion with domain control}
\GSfield{Kind; source}{claim; \GSpage{P000430}}
\GSfield{The notes}{The ``more general variant'' asks for \(E(\Delta)(H-H_0)E_0(\Delta)\in\mathcal S_1\), followed by a comparison of domains of growing functions of \(H\) and \(H_0\). The inclusion sign and the functions are not fully resolved.}
\GSfield{Statement}{Let \(H,H_0\) be self-adjoint operators on a separable Hilbert space, with spectral projections \(E,E_0\). Assume that, for every bounded interval \(\Delta\), the bounded operator \[H E(\Delta)E_0(\Delta)-E(\Delta)H_0E_0(\Delta)\] belongs to the trace class. Assume also that there exist real Borel functions \(f,f_0\), bounded on compact intervals and satisfying \(|f(\lambda)|,|f_0(\lambda)|\to\infty\) as \(|\lambda|\to\infty\), such that \(D(f(H))=D(f_0(H_0))\). Then both time-dependent wave operators exist and are complete. The domain equality is a sufficient form of the mutual-subordination hypothesis in Birman's theorem.}
\GSfield{Reconstruction}{The scan does contain the equality of the two functional-calculus domains and the statement that the functions grow at infinity. Borel measurability, local boundedness and absolute growth are made explicit. The localized difference is written as the bounded expression above so that no undefined subtraction on unequal operator domains is needed. Birman's Definition 4.2 allows one domain inclusion for subordination; mutual subordination follows from the displayed equality.}
\GSfield{Fate}{This is the local criterion of \citet[Definition 4.2 and Theorem 4.4, pp. 897--900]{birman1968local}, received 23 January 1968 and published that year. Its hypotheses are more general than a globally trace-class difference and explain precisely the final lines of the 27 March lecture.}
\GSfield{Present status}{As of November 2026, this is an established scattering theorem. The notebook's compressed domain condition can be reconstructed from the 1968 source; no separate open extension is justified.}
\GSfield{Assessment}{literal statement: not determinable; publication match: exact; verification: read the result; development: none; confidence: high}
\end{GSitem}

\GStalk{265}{1972-04-03}{A. Yu. Geronimus}{P000430-P000436}{Sites, sheaves, derived cohomology and the classifying topos}
Following a minimum reading programme on topology, descent and simplicial methods, Geronimus develops covering families, sheaves, adjunction, derived cohomology and a Leray spectral sequence. He then constructs a classifying site for group objects and applies the language to Quillen's finite-field cohomology calculation. An exercise refers broadly to Artin on formal and algebraic objects, without stating a theorem. The exposition revisits exact earlier occurrences such as 1967-143.2, 1967-143.5 and 1970-226.2--3; no new problem identifier is created for those definitions.

\begin{GSitem}{1971-265.1}{E}{Grothendieck covering families and fibre products}
\GSfield{Kind; source}{definition; \GSpage{P000430-P000431}}
\GSfield{The notes}{Coverings are families \(\{X_\alpha\to X\}\), stable under isomorphisms, composition and arbitrary base change. The base-change axiom explicitly requires existence of \(X_\alpha\times_X U\); a universal diagram defines the fibre product.}
\GSfield{Statement}{A pretopology on a category \(T\) assigns covering families satisfying the three stated axioms. A fibre product \(U\times_XV\) represents pairs of arrows to \(U,V\) with equal composites to \(X\). The associated Grothendieck topology consists of sieves containing a covering family. The distinction between a pretopology and its topology matters when different covering families generate the same sheaves.}
\GSfield{Reconstruction}{The modern term pretopology and the associated-sieve construction are supplied; the manuscript's axioms themselves are accurate.}
\GSfield{Fate}{The framework was circulated in SGA 4 in 1963--64 and published in 1972 \citep{sga41972topos}; a precise current formulation is \citet[Tag 00VH]{stacks2026sites}.}
\GSfield{Present status}{As of November 2026, the definitions are standard and recur from the seminar's earlier years.}
\GSfield{Assessment}{literal statement: correct; publication match: exact; verification: read the result; development: none; confidence: high}
\end{GSitem}
\begin{GSitem}{1971-265.2}{E}{Presheaves, natural transformations and the functor category}
\GSfield{Kind; source}{definition; \GSpage{P000431-P000432}}
\GSfield{The notes}{A presheaf is a contravariant functor \(T\to A\), with \(A\) abelian. A commuting square defines a morphism of functors; the notes observe a possible set-theoretic size issue.}
\GSfield{Statement}{For a small site \(T\), an abelian presheaf is a functor \(T^{\rm op}\to\mathbf{Ab}\). A morphism consists of maps \(t_U:F(U)\to G(U)\) commuting with every restriction map. These form an abelian functor category with kernels and cokernels computed objectwise. For a large site, one fixes a universe or a suitable small presentation before asserting that morphism collections are sets.}
\GSfield{Reconstruction}{The coefficient category is restricted to abelian groups in the statements used below. Mere existence of infinite direct sums in an arbitrary abelian category is not a complete substitute for all limits and exactness conditions needed in the subsequent sheaf theory.}
\GSfield{Fate}{This is the categorical language of \citet[Chapter I]{grothendieck1957homological} and \citet[Expos\'e I]{sga41972topos}. The seminar is building shared terminology.}
\GSfield{Present status}{As of November 2026, these constructions are standard; the manuscript's warning about size is appropriate.}
\GSfield{Assessment}{literal statement: correct with missing hypotheses; publication match: exact; verification: publisher metadata checked; development: none; confidence: high}
\end{GSitem}
\begin{GSitem}{1971-265.3}{E}{The sheaf condition uses products over an arbitrary cover}
\GSfield{Kind; source}{definition; \GSpage{P000431}}
\GSfield{The notes}{The sheaf condition is displayed with \(\bigoplus_\alpha F(X_\alpha)\) and \(\bigoplus_{\alpha,\beta}F(X_\alpha\times_X X_\beta)\), although the covering index is an arbitrary set.}
\GSfield{Statement}{For every covering \(\{U_i\to U\}\), an abelian presheaf is a sheaf when \[0\longrightarrow F(U)\longrightarrow\prod_iF(U_i)\xrightarrow{\operatorname{pr}_1^*-\operatorname{pr}_2^*}\prod_{i,j}F(U_i\times_UU_j)\] is exact. Products and direct sums coincide for finite covers, but direct sums impose an unwarranted finite-support condition on infinite covers.}
\GSfield{Reconstruction}{\ours{} The products are an emendation of the mathematical statement, not a silent change to the diplomatic transcription. For an infinite discrete space covered by its points, a global locally constant function need not have finite support.}
\GSfield{Fate}{The exact equalizer formulation is \citet[Tag 00VM]{stacks2026sites}; it belongs to the pre-existing sheaf formalism of SGA 4.}
\GSfield{Present status}{As of November 2026, the arbitrary-cover statement in the notebook is false with direct sums and correct with products.}
\GSfield{Assessment}{literal statement: false as written; publication match: exact; verification: read the result; development: annotated; confidence: high}
\end{GSitem}
\begin{GSitem}{1971-265.4}{E}{Sheafification, adjunction and the abelian category of sheaves}
\GSfield{Kind; source}{theorem; \GSpage{P000432}}
\GSfield{The notes}{The full subcategory \(S_T\) of sheaves is abelian. The proof uses an associated-sheaf functor \(\#\), adjoint to the inclusion \(i:S_T\to P_T\). Adjointness is defined by a natural Hom-set isomorphism.}
\GSfield{Statement}{For abelian presheaves on a small site, sheafification \(a:P_T\to S_T\) is an exact left adjoint to the inclusion: \(\operatorname{Hom}(aF,G)\cong\operatorname{Hom}(F,iG)\), naturally in both variables. The category \(S_T\) is abelian; kernels are computed as presheaf kernels, whereas cokernels must be sheafified. A Hom-set bijection without naturality would not define an adjunction.}
\GSfield{Reconstruction}{The direction of the adjunction and exactness of sheafification are explicit. The conclusion is stated for abelian groups, the coefficient category actually needed later.}
\GSfield{Fate}{These results are in the classical site formalism \citep{sga41972topos}; \citet[Tag 03CM, Lemmas 18.3.1--2]{stacks2026foundations} verifies the exact current statements.}
\GSfield{Present status}{As of November 2026, the theorem and proof mechanism are standard; inclusion of sheaves into presheaves is generally only left exact.}
\GSfield{Assessment}{literal statement: correct with missing hypotheses; publication match: exact; verification: read the result; development: none; confidence: high}
\end{GSitem}
\begin{GSitem}{1971-265.5}{E}{Cohomology on a site as a derived section functor}
\GSfield{Kind; source}{definition; \GSpage{P000433}}
\GSfield{The notes}{The definition is \(H^q(T,X,F)=R^qi(F)(X)\), where \(i\) includes sheaves in presheaves.}
\GSfield{Statement}{For an object \(X\) of a small site and an abelian sheaf \(F\), define \(H^q(X,F)=R^q\Gamma(X,-)(F)\). Abelian sheaves have enough injectives. Evaluation of presheaves at \(X\) is exact, so computing an injective resolution also gives \(H^q(X,F)=(R^qi(F))(X)\).}
\GSfield{Reconstruction}{Enough injectives and the exactness of evaluation explain why the notebook's seemingly unusual inclusion-functor definition agrees with ordinary site cohomology.}
\GSfield{Fate}{The homological construction predates the talk \citep{grothendieck1957homological}; the site version and enough-injectives theorem are verified in \citet[Tags 01DL and 03NT]{stacks2026foundations}.}
\GSfield{Present status}{As of November 2026, the definition is correct for abelian sheaves on a site, with standard size conventions.}
\GSfield{Assessment}{literal statement: correct; publication match: exact; verification: read the result; development: none; confidence: high}
\end{GSitem}
\begin{GSitem}{1971-265.6}{E}{Continuity of a functor of sites}
\GSfield{Kind; source}{claim; \GSpage{P000433-P000434}}
\GSfield{The notes}{A morphism of topologies is defined as a functor carrying coverings to coverings; precomposition \(F\mapsto F\circ f\) is asserted to carry sheaves to sheaves.}
\GSfield{Statement}{If \(f:T_1\to T_2\) carries covering families to covering families and preserves the fibre products that occur in their descent diagrams, then precomposition carries sheaves on \(T_2\) to sheaves on \(T_1\). More generally continuity is defined by that sheaf-preservation condition. Cover preservation alone does not justify substituting \(f(U_i\times_UU_j)\) for \(f(U_i)\times_{f(U)}f(U_j)\).}
\GSfield{Reconstruction}{The missing fibre-product compatibility is supplied. The direction and notation in the notes are retained as precomposition, without identifying this automatically with the exact inverse-image functor of a geometric morphism of topoi.}
\GSfield{Fate}{The distinction is part of \citet[Expos\'e III]{sga41972topos}. Its 1963--64 circulation predates the seminar; the volume appeared in 1972.}
\GSfield{Present status}{As of November 2026, the construction is standard under continuity hypotheses, while the displayed definition is insufficient as written.}
\GSfield{Assessment}{literal statement: correct with missing hypotheses; publication match: partial; verification: publisher metadata checked; development: none; confidence: high}
\end{GSitem}
\begin{GSitem}{1971-265.7}{E}{Grothendieck and Leray spectral sequences}
\GSfield{Kind; source}{theorem; \GSpage{P000433-P000435}}
\GSfield{The notes}{For \(A\xrightarrow F B\xrightarrow G C\), the notes state \(R^pG\,R^qF\Rightarrow R^{p+q}(GF)\) ``under certain conditions'', and present the site Leray sequence as a special case.}
\GSfield{Statement}{Let \(A,B,C\) be abelian categories, with enough injectives in \(A,B\). Suppose \(F,G\) are additive left exact functors and \(F\) carries injectives to \(G\)-acyclic objects. For \(M\in A\), there is a first-quadrant spectral sequence \(E_2^{p,q}=R^pG(R^qF(M))\Rightarrow R^{p+q}(GF)(M)\). For a morphism of topoi \(f\), this yields \(H^p(Y,R^qf_*F)\Rightarrow H^{p+q}(X,F)\).}
\GSfield{Reconstruction}{The acyclicity assumption is the decisive missing condition. The modern Leray formulation uses the direct-image functor \(f_*\); the manuscript's opposite site-functor conventions must be translated before comparing symbols.}
\GSfield{Fate}{The composite-functor spectral sequence is due to the earlier homological formalism \citep{grothendieck1957homological}; its hypotheses and convergence are verified in \citet[Tag 015L, Lemma 13.22.2]{stacks2026foundations}. The earlier site formulation is also explicit in \citet[Chapter II, Theorem 4.11]{artin1962topologies}, read in the 1962 primary notes.}
\GSfield{Present status}{As of November 2026, the spectral sequence is established. Exactness of precomposition on presheaves alone does not replace the injective-acyclicity check on sheaves.}
\GSfield{Assessment}{literal statement: correct with missing hypotheses; publication match: exact; verification: read the result; development: none; confidence: high}
\end{GSitem}
\begin{GSitem}{1971-265.8}{E}{Group objects and the classifying site}
\GSfield{Kind; source}{construction; \GSpage{P000434-P000435}}
\GSfield{The notes}{A group object is described by a functorial group structure on \(h_G(Y)\), or by multiplication, inversion and a unit. The new site \(TG\) consists of \(G\)-objects, equivariant maps and underlying coverings; its final object is called a classifying space.}
\GSfield{Statement}{In a category with finite products, a group object is an object \(G\) equipped with the group-law diagrams. Yoneda identifies this with a group structure on every \(\operatorname{Hom}(Y,G)\), natural in \(Y\). For a group object in a topos \(T\), the topos of \(G\)-objects is its classifying topos \(B_TG\). Its final object is the base's final object with trivial action; the topology and equivariance, not a nontrivial underlying terminal set, encode classifying cohomology.}
\GSfield{Reconstruction}{The classifying topos is distinguished from a topological classifying space. For a discrete group over the topos of sets, abelian objects are \(G\)-modules and derived global sections recover group cohomology.}
\GSfield{Fate}{The construction belongs to \citet[Expos\'e IV]{sga41972topos}. It is the conceptual language needed for the subsequent comparison, not a new construction of ordinary \(BG\) from a terminal object.}
\GSfield{Present status}{As of November 2026, the classifying-topos interpretation is standard; the notebook's terminal-object phrase is too compressed if read as a statement about ordinary spaces.}
\GSfield{Assessment}{literal statement: correct with missing hypotheses; publication match: partial; verification: publisher metadata checked; development: none; confidence: high}
\end{GSitem}
\begin{GSitem}{1971-265.9}{E}{Constant sheaves and the local coefficients in a fibration sequence}
\GSfield{Kind; source}{definition; \GSpage{P000435}}
\GSfield{The notes}{The constant sheaf associated with \(A\) is correctly obtained from the constant presheaf. A later display simplifies the \(E_2\) page to \(H^p(BG)\otimes H^q(G)\).}
\GSfield{Statement}{The constant sheaf \(\underline A\) is the sheafification of the constant presheaf with value \(A\). Let \(F\to E\to B\) be a Serre fibration with connected CW base and fibre of CW type, and let \(k\) be a field. Its cohomology spectral sequence has \(E_2^{p,q}=H^p(B;\mathcal H^q(F;k))\), where monodromy defines the coefficient system. If monodromy on each \(H^q(F;k)\) is trivial and each such vector space is finite dimensional, then \(E_2^{p,q}\cong H^p(B;k)\otimes_k H^q(F;k)\).}
\GSfield{Reconstruction}{The fibration version is a precise model of the analogy drawn in the notes. Trivial local coefficients and tensor-product hypotheses are supplied; no general tensor-product formula for arbitrary sites or coefficient rings is asserted. Finite-dimensional fibre cohomology in every degree is one explicit sufficient condition for the tensor identification.}
\GSfield{Fate}{Constant sheaves and site Leray sequences are already present in SGA 4 \citep{sga41972topos}; the classical fibration sequence is the topological analogy used by the speaker.}
\GSfield{Present status}{As of November 2026, the constant-sheaf construction is correct. The tensor notation needs the additional conditions expressly alluded to but not specified in the notes.}
\GSfield{Assessment}{literal statement: correct with missing hypotheses; publication match: partial; verification: publisher metadata checked; development: none; confidence: high}
\end{GSitem}
\begin{GSitem}{1971-265.10}{E}{Quillen’s mod-l cohomology of a finite general linear group}
\GSfield{Kind; source}{theorem; \GSpage{P000436}}
\GSfield{The notes}{The notes write \(H^*(GL_n(\mathbb F_q);\mathbb F_\ell)=S[c_1,\ldots,c_n]\otimes\Lambda[c'_1,\ldots,c'_n]\), with degrees \(2i\) and \(2i-1\), but do not state the relation between \(q\) and \(\ell\). One coefficient field is also written \(\mathbb F_q\).}
\GSfield{Statement}{For an odd prime \(\ell\neq\operatorname{char}\mathbb F_q\), let \(r\) be the multiplicative order of \(q\) modulo \(\ell\), and \(m=\lfloor n/r\rfloor\). Quillen's description is \[H^*(GL_n(\mathbb F_q);\mathbb F_\ell)\cong\mathbb F_\ell[c_r,c_{2r},\ldots,c_{mr}]\otimes\Lambda(e_r,e_{2r},\ldots,e_{mr}),\] with \(|c_{jr}|=2jr\), \(|e_{jr}|=2jr-1\). The displayed notebook formula is the case \(r=1\), namely \(\ell\mid q-1\). The prime \(2\) requires separate cases.}
\GSfield{Reconstruction}{The odd-prime restrictions and the order \(r\) supply the absent hypotheses. The coefficient of every cohomology group in this formula is \(\mathbb F_\ell\).}
\GSfield{Fate}{The formula, including the odd-prime algebra qualification, is explicit in \citet[Section 2, p. 50]{quillen1971cohomology}, announced at ICM 1970 and printed in 1971. The fuller journal account appeared in November 1972 \citep{quillen1972finitefields}.}
\GSfield{Present status}{As of November 2026, the formula with these qualifications is established. The unrestricted printed formula is false, for example when \(\ell\nmid|GL_n(\mathbb F_q)|\), when all positive-degree cohomology vanishes.}
\GSfield{Assessment}{literal statement: false as written; publication match: exact; verification: read the result; development: none; confidence: medium}
\end{GSitem}
\begin{GSitem}{1971-265.11}{E}{Finite generation through a faithful unitary representation}
\GSfield{Kind; source}{method; \GSpage{P000436}}
\GSfield{The notes}{A faithful irreducible representation \(G\to U\) and the induced \(BG\to BU\) are invoked to explain finite generation of group cohomology.}
\GSfield{Statement}{A finite group has a faithful finite-dimensional unitary representation, for example its regular representation. It need not have a faithful irreducible representation. The classical finite-generation theorem says that \(H^*(G;\mathbb F_\ell)\) is a finitely generated algebra for every finite group \(G\); the topological proof uses \(G\subset U(N)\) and the fibration \(U(N)/G\to BG\to BU(N)\).}
\GSfield{Reconstruction}{Irreducibility is removed; it is false for a noncyclic finite abelian group. The coefficient ring and finiteness of \(G\) are stated explicitly.}
\GSfield{Fate}{Finite generation is already attributed to Venkov and Evens in \citet[Section 1, p. 47]{quillen1971cohomology}. A modern independent proof appears in \citet[Theorem 2.4]{rouquier2014finite}. The faithful embedding need not be irreducible; the regular representation supplies an unconditional embedding of any finite group.}
\GSfield{Present status}{As of November 2026, finite generation is established. It gives no automatic explicit set of Quillen generators or a Brauer-lift comparison.}
\GSfield{Assessment}{literal statement: false as written; publication match: exact; verification: read the result; development: none; confidence: medium}
\end{GSitem}
\begin{GSitem}{1971-265.12}{E}{Comparing finite-field and complex algebraic groups}
\GSfield{Kind; source}{claim; \GSpage{P000436}}
\GSfield{The notes}{Passing from \(GL_n(\mathbb F_q)\) to \(GL_n(\overline{\mathbb F}_q)\), the notes say the resulting classifying object, properly interpreted, equals \(BU\). Bernstein adds that representations over the algebraic closure and over \(\mathbb C\) are described by highest weights.}
\GSfield{Statement}{The comparison concerns étale or prime-to-characteristic completed homotopy/cohomology of algebraic groups and their classifying objects, not an equality between the discrete classifying space of \(GL_n(\overline{\mathbb F}_q)\) and the ordinary topological \(BU(n)\). For a split connected reductive group over an algebraically closed field, simple rational representations are indexed by dominant highest weights. In positive characteristic their dimensions, characters and semisimplicity properties need not agree with characteristic zero.}
\GSfield{Reconstruction}{The topology, finite rank and coefficient restrictions are restored. The highest-weight observation is separated from the much stronger and false assertion that the two representation categories coincide.}
\GSfield{Fate}{Quillen's comparison strategy and its prime-to-characteristic coefficients were already announced in \citet[Section 2]{quillen1971cohomology}. For Bernstein's aside, \citet[Theorem 22.2 and Sections 22.3--22.5]{milne2017groups} states the characteristic-independent highest-weight parametrization for split reductive groups. It does not identify the dimensions or characters of all corresponding simple modules across characteristics.}
\GSfield{Present status}{As of November 2026, these qualified comparisons are established; neither literal equality of spaces nor characteristic-independent ordinary representation theory follows from the notebook.}
\GSfield{Assessment}{literal statement: correct with missing hypotheses; publication match: partial; verification: read the result; development: none; confidence: medium}
\end{GSitem}

\GStalk{266}{1972-04-10}{A. Yu. Geronimus}{P000437-P000438}{The homotopy category and localization at quasi-isomorphisms}
Geronimus builds the derived category from complexes and homotopies, emphasizes its universal property, and explicitly rejects identifying maps merely because they induce the same map on cohomology. A roof construction follows, although the notebook mistakenly calls the resulting limit projective. The examples begin in nonnegative degrees. A crossed-out reprise on P000439 belongs to this discussion but adds no distinct live assertion.

\begin{GSitem}{1971-266.1}{E}{Complexes, cohomology, degree and homotopy}
\GSfield{Kind; source}{definition; \GSpage{P000437}}
\GSfield{The notes}{A complex \(0\to C_0\xrightarrow{d_0}C_1\to\cdots\) has \(H^q=\ker d_q/\operatorname{im}d_{q-1}\). Maps of degree \(\tau\) are introduced and a homotopy formula \(f-g=d\alpha+\alpha d'\) is written without consistent indices or signs.}
\GSfield{Statement}{For a cochain complex \(C^\bullet\) in an abelian category, \(d^{n+1}d^n=0\) and \(H^n(C)=\ker d^n/\operatorname{im}d^{n-1}\). Degree-zero chain maps \(f,g:C\to D\) are homotopic if \(f-g=d_Dh+hd_C\), with \(h^n:C^n\to D^{n-1}\). In the graded Hom complex a degree-\(\tau\) map has differential \(d_Df-(-1)^\tau fd_C\); the corresponding signed boundary formula governs higher-degree homotopies.}
\GSfield{Reconstruction}{The degree-zero statement makes the notebook's intended homotopy unambiguous. The graded sign convention is explicitly supplied rather than treating the unsigned formula as correct in every degree.}
\GSfield{Fate}{These are the classical chain-complex definitions, incorporated in the derived-category formalism before 1972; see \citet[Chapter I]{hartshorne1966residues} and \citet[Chapter 12]{may1999concise}.}
\GSfield{Present status}{As of November 2026, the definitions are standard; homotopic maps induce equal maps on cohomology.}
\GSfield{Assessment}{literal statement: correct with missing hypotheses; publication match: partial; verification: indirect; development: none; confidence: high}
\end{GSitem}
\begin{GSitem}{1971-266.2}{E}{The homotopy category is generally not abelian}
\GSfield{Kind; source}{definition; \GSpage{P000437}}
\GSfield{The notes}{Morphisms are maps modulo null-homotopic maps. The notes collect morphisms of all degrees into a direct sum, then stress that the resulting category is not abelian because of homotopy equivalence.}
\GSfield{Statement}{The ordinary homotopy category \(K(A)\) has complexes as objects and homotopy classes of degree-zero chain maps as morphisms. It is additive and triangulated, but generally not abelian. The groups \(\operatorname{Hom}_{K(A)}(C,D[\tau])\) form a graded enhancement; collecting every degree changes the morphism convention and should not be confused with ordinary \(K(A)\).}
\GSfield{Reconstruction}{The text's all-degrees convention is separated from the standard category used in the subsequent localization. The abelian category of actual complexes and chain maps must also be distinguished from its homotopy category.}
\GSfield{Fate}{The distinction is already in \citet[Chapter I]{hartshorne1966residues}; the precise conventions are verified in \citet[Tag 03T3]{stacks2026foundations}.}
\GSfield{Present status}{As of November 2026, the non-abelian homotopy-category statement is established, with the specified convention for morphisms.}
\GSfield{Assessment}{literal statement: correct with missing hypotheses; publication match: exact; verification: read the result; development: none; confidence: high}
\end{GSitem}
\begin{GSitem}{1971-266.3}{E}{The derived category and its universal property}
\GSfield{Kind; source}{construction; \GSpage{P000437-P000438}}
\GSfield{The notes}{The functor \(K(A)\to D(A)\) sends cohomology isomorphisms to isomorphisms, is universal with this property, and in fact a chain map becomes invertible exactly when it induces cohomology isomorphisms.}
\GSfield{Statement}{Let \(S\) be the quasi-isomorphisms in \(K(A)\). The derived category is \(D(A)=K(A)[S^{-1}]\), with the localization universal property: every functor inverting \(S\) factors uniquely up to natural isomorphism through \(D(A)\). The cohomology functors descend to \(D(A)\); a map in \(K(A)\) is a quasi-isomorphism if and only if its localized image is invertible.}
\GSfield{Reconstruction}{The relevant set-theoretic universe and existence of the localization are tacit; the bounded-below setting of the talk suffices in the standard module and sheaf categories.}
\GSfield{Fate}{The construction was available in the 1963--64 Harvard seminar published as \citet[Chapter I]{hartshorne1966residues}; \citet[Tag 03T3]{stacks2026foundations} verifies the current convention.}
\GSfield{Present status}{As of November 2026, this is the standard derived category, and the universal property recorded in the notes is accurate.}
\GSfield{Assessment}{literal statement: correct; publication match: exact; verification: read the result; development: none; confidence: high}
\end{GSitem}
\begin{GSitem}{1971-266.4}{E}{Derived morphisms retain information invisible to cohomology}
\GSfield{Kind; source}{remark; \GSpage{P000438}}
\GSfield{The notes}{The notes emphatically say that \(D(A)\) is not obtained by declaring \(f\sim g\) when \(f^*=g^*\); it has more morphisms, obtained by formally adjoining inverses.}
\GSfield{Statement}{In general the functor \(H^*:D(A)\to\operatorname{Gr}(A)\) is not faithful. For modules, \(\operatorname{Hom}_{D(R)}(M,N[1])\cong\operatorname{Ext}^1_R(M,N)\). For example \(R=\mathbb Z\), \(M=N=\mathbb Z/p\), this group is \(\mathbb Z/p\), whereas a morphism \(M\to N[1]\) induces zero on each cohomology group because the nonzero groups lie in different degrees.}
\GSfield{Reconstruction}{\ours{} The Ext example is an editorial illustration of the correct warning in the notebook; it is not presented as an example supplied by the speaker.}
\GSfield{Fate}{This is an elementary consequence of the earlier derived-functor formalism \citep[Chapter I]{hartshorne1966residues}.}
\GSfield{Present status}{As of November 2026, the warning is correct and essential: cohomology detects quasi-isomorphisms without determining every morphism.}
\GSfield{Assessment}{literal statement: correct; publication match: partial; verification: indirect; development: annotated; confidence: high}
\end{GSitem}
\begin{GSitem}{1971-266.5}{E}{Roofs and the direction of the limit}
\GSfield{Kind; source}{construction; \GSpage{P000438}}
\GSfield{The notes}{Morphisms from \(C\) to \(D\) are represented by diagrams \(C\leftarrow C'\to D\), with the left arrow a cohomology isomorphism. The notes call the set of such morphisms a projective limit.}
\GSfield{Statement}{A morphism in \(D(A)\) is represented by a roof \(C\xleftarrow{s}C'\xrightarrow{f}D\), with \(s\) a quasi-isomorphism, modulo the calculus-of-fractions equivalence by common refinements in \(K(A)\). The Hom set is a filtered colimit of the corresponding Hom sets over quasi-isomorphic replacements, with the appropriate opposite indexing category. It is not the inverse limit of all roof diagrams.}
\GSfield{Reconstruction}{The homotopy-category equivalence relation and the direction of the colimit are supplied. The word \ru{проект. предел} is visible in the scan and is treated as a mathematical emendation, not an English translation mistake.}
\GSfield{Fate}{Roofs are the standard localization construction of \citet[Chapter I]{hartshorne1966residues}; they implement the universal property described correctly earlier on the same page.}
\GSfield{Present status}{As of November 2026, the roof construction is established, while the unqualified projective-limit statement needs correction.}
\GSfield{Assessment}{literal statement: false as written; publication match: partial; verification: indirect; development: annotated; confidence: high}
\end{GSitem}
\begin{GSitem}{1971-266.6}{E}{When an additive functor has a total derived functor}
\GSfield{Kind; source}{construction; \GSpage{P000438}}
\GSfield{The notes}{The concluding line sends \(f:A\to A'\) to \(D(f):D(A)\to D(A')\), without stating exactness or resolution hypotheses.}
\GSfield{Statement}{An exact additive functor preserves quasi-isomorphisms and induces a functor on derived categories. For a left exact additive functor \(F:A\to A'\), with enough injectives in \(A\), the bounded-below right derived functor is \(RF:D^+(A)\to D^+(A')\), computed by applying \(F\) to an injective resolution. A general additive functor has no automatic induced functor on the unlocalized cohomology classes.}
\GSfield{Reconstruction}{The two cases, exact functor and right derived functor, are distinguished. Kazhdan's next intervention supplies the missing injective-resolution mechanism.}
\GSfield{Fate}{The construction is classical; \citet[Tag 03T3, Definition 64.6.4]{stacks2026foundations} verifies exactly the bounded-below formulation.}
\GSfield{Present status}{As of November 2026, derived functors are standard under these hypotheses; the unrestricted displayed notation suppresses necessary assumptions.}
\GSfield{Assessment}{literal statement: correct with missing hypotheses; publication match: exact; verification: read the result; development: none; confidence: high}
\end{GSitem}

\GStalk{267}{1972-04-10}{D. A. Kazhdan}{P000439}{Computing the derived category with injective complexes}
Kazhdan supplies the practical model omitted from the preceding abstract localization: enough injectives allow one to work with injective complexes. The adjacent complex-bordism lecture starts only after this short intervention. The omitted boundedness qualification is recoverable from Geronimus's nonnegative complexes on P000437.

\begin{GSitem}{1971-267.1}{E}{Bounded-below injective complexes model the derived category}
\GSfield{Kind; source}{theorem; \GSpage{P000439}}
\GSfield{The notes}{If \(A\) has enough injectives, the category of complexes in \(\operatorname{Inj}(A)\) is said to be equivalent to the derived category of \(A\).}
\GSfield{Statement}{For an abelian category \(A\) with enough injectives, the natural functor \(K^+(\operatorname{Inj}A)\to D^+(A)\) is an equivalence. Morphisms on the left are chain maps modulo homotopy, not arbitrary chain maps. For unbounded complexes, termwise injectivity is insufficient in general; one uses \(K\)-injective resolutions in categories where they exist.}
\GSfield{Reconstruction}{Bounded-below complexes and passage to homotopy classes are restored. This is the precise result compatible with the preceding nonnegative complexes.}
\GSfield{Fate}{The result belongs to the pre-1972 derived-category theory \citep{hartshorne1966residues}; \citet[Tag 03T3, Lemma 64.6.3]{stacks2026foundations} states the exact equivalence.}
\GSfield{Present status}{As of November 2026, the bounded-below theorem is established. Its unbounded analogue needs the stronger \(K\)-injective condition.}
\GSfield{Assessment}{literal statement: correct with missing hypotheses; publication match: exact; verification: read the result; development: none; confidence: high}
\end{GSitem}

\GStalk{268}{1972-04-10}{Jack Morava}{P000439-P000445}{Complex bordism, formal groups and the emerging height picture}
Morava starts from geometric bordism and interprets the complex-bordism module of a space as a sheaf on the scheme of formal group laws. Coordinate changes, height strata, periodic theories and division-algebra stabilizers give an early form of chromatic geometry. Lazard and Quillen supply the older foundations; Morava's geometric programme is contemporary research, with a November 1972 unpublished announcement and later published developments. The strong claim that every fibre in one dense open set is integral \(K\)-theory requires substantial repair.

\begin{GSitem}{1971-268.1}{E}{Unoriented, oriented and stably complex bordism}
\GSfield{Kind; source}{definition; \GSpage{P000439-P000440}}
\GSfield{The notes}{Closed compact smooth manifolds are equivalent when they form the boundary of a compact manifold. Unoriented bordism has \([M]+[M]=0\). A weak complex manifold includes a complex bundle \(\xi\) and a fixed stable real-bundle isomorphism with the tangent bundle.}
\GSfield{Statement}{Unoriented bordism uses \(\partial W\cong M\sqcup N\) and disjoint-union addition. Oriented bordism requires \(\partial W\cong M\sqcup(-N)\). A stable complex structure on a smooth manifold is a complex structure on its stabilized tangent bundle, up to the standard stable homotopy equivalence; its choice induces an orientation and must extend across a bordism. It is data, not merely the existence of some isomorphism of underlying real bundles.}
\GSfield{Reconstruction}{Boundary orientations and stable homotopies of the chosen structure are supplied. The cylinder proves exponent two only in unoriented bordism, not in complex bordism.}
\GSfield{Fate}{These geometric definitions precede the talk and are the input of \citet{connerfloyd1966relation} and the complex-bordism theory used in \citet{quillen1969formal}.}
\GSfield{Present status}{As of November 2026, all three bordism theories are standard; keeping the extra structures prevents the false transfer of the unoriented exponent-two assertion.}
\GSfield{Assessment}{literal statement: correct with missing hypotheses; publication match: exact; verification: publisher metadata checked; development: none; confidence: high}
\end{GSitem}
\begin{GSitem}{1971-268.2}{E}{Representatives by complex manifolds: a necessary qualification}
\GSfield{Kind; source}{claim; \GSpage{P000440}}
\GSfield{The notes}{The notes assert that every weak complex manifold is bordant to a complex manifold.}
\GSfield{Statement}{A safe precise result is that the complex bordism ring is generated, as a ring over \(\mathbb Z\), by classes of nonsingular projective varieties; hence each class is an integral linear combination of such classes. This does not mean that every class, with its specified stable complex orientation, is represented by a nonnegative disjoint union of complex manifolds. In dimension zero, \(MU_0\cong\mathbb Z\), and a negatively oriented point represents \(-1\), which no ordinary compact complex zero-manifold represents.}
\GSfield{Reconstruction}{\ours{} The generated-ring statement is the verified replacement. The zero-dimensional counterexample diagnoses the universal literal wording; no assertion about the sharp positive-dimensional realization theorem is inferred from this brief line.}
\GSfield{Fate}{The use of nonsingular projective generators is explicit in the proof of \citet[Theorem 2]{quillen1969formal}, relying on the earlier complex-bordism calculation.}
\GSfield{Present status}{As of November 2026, the ring-generation result is established. The unqualified literal claim has a zero-dimensional counterexample.}
\GSfield{Assessment}{literal statement: false as written; publication match: partial; verification: read the result; development: annotated; confidence: high}
\end{GSitem}
\begin{GSitem}{1971-268.3}{E}{The complex bordism coefficient ring}
\GSfield{Kind; source}{theorem; \GSpage{P000441}}
\GSfield{The notes}{The coefficient ring is written \(\Omega_*^U=\mathbb Z[t_0,t_1,\ldots]\), with \(\deg t_i=2i\). Products are Cartesian products, and inverses require changing the stable complex structure appropriately.}
\GSfield{Statement}{The complex bordism coefficient ring is \(MU_*\cong\mathbb Z[x_1,x_2,\ldots]\), where \(|x_i|=2i\) for \(i\geq1\); \(MU_0=\mathbb Z\) and \(MU_{2i+1}=0\). There is no independent degree-zero polynomial generator. Multiplication is induced by Cartesian product, with the product stable complex structure.}
\GSfield{Reconstruction}{The subscript zero is removed as an independent generator; if \(t_0\) was intended to denote the point, it must equal the unit rather than a polynomial variable.}
\GSfield{Fate}{The Milnor--Novikov coefficient-ring calculation predates the talk; \citet[Theorem 2 and its proof]{quillen1969formal} identifies it with Lazard's polynomial ring and uses the projective generators.}
\GSfield{Present status}{As of November 2026, the polynomial presentation is established; the notebook's extra degree-zero variable is erroneous if read literally.}
\GSfield{Assessment}{literal statement: false as written; publication match: exact; verification: read the result; development: none; confidence: high}
\end{GSitem}
\begin{GSitem}{1971-268.4}{E}{Bordism of a space and the module-classification programme}
\GSfield{Kind; source}{construction; \GSpage{P000441-P000442}}
\GSfield{The notes}{A cycle over \(X\) is \([M,\xi,\alpha,f]\) with \(f:M\to X\), and a bordism extends all the data. The direct sum \(\Omega_*^U(X)\) is an \(\Omega_*^U\)-module. A problem asks to describe its module structure.}
\GSfield{Statement}{For a CW space \(X\), geometric complex bordism defines the homology theory \(MU_*(X)\). External products with coefficient manifolds make it a graded \(MU_*\)-module. The broad request to describe this module has no specified class of spaces or proposed classification invariant in the notes; it becomes a programme rather than a single quantified problem.}
\GSfield{Reconstruction}{The covariant theory here is homology, despite later notation resembling \(K^*(X)\). No claim is made that all graded modules arise as \(MU_*(X)\), or that height data alone classify them.}
\GSfield{Fate}{Module structures were already a focus of complex-bordism theory before 1972. \citet[Introduction and Section 3]{johnsonwilson1975operations} makes the later height-local structure results precise and explicitly builds on Morava's 1971 unpublished work.}
\GSfield{Present status}{As of November 2026, the module construction is standard, while the broad classification request is not a sufficiently bounded open problem to receive a new problem identifier.}
\GSfield{Assessment}{literal statement: correct with missing hypotheses; publication match: partial; verification: read the result; development: none; confidence: high}
\end{GSitem}
\begin{GSitem}{1971-268.5}{E}{Modules as quasicoherent sheaves and their fibres}
\GSfield{Kind; source}{construction; \GSpage{P000442}}
\GSfield{The notes}{Ignoring most grading, the polynomial ring is viewed as the coordinate ring of an infinite affine space, and \(MU_*(X)\) as a sheaf; even and odd degrees are retained separately.}
\GSfield{Statement}{Put \(L=MU_*\) and \(M=MU_*(X)\). The associated quasicoherent sheaf \(\widetilde M\) on \(\operatorname{Spec}L\) has fibre \(M\otimes_L\kappa(\mathfrak p)\) at a prime \(\mathfrak p\). Forgetting to parity preserves a \(\mathbb Z/2\)-grading but discards the full weight information. Since \(L\) has infinitely many polynomial generators, this affine scheme is not noetherian.}
\GSfield{Reconstruction}{The residue-field tensor product and the distinction between a stalk and its fibre are supplied. Fibres are vector spaces over varying residue fields, so they cannot literally all be the same integral \(K\)-group.}
\GSfield{Fate}{The module-to-sheaf construction is older affine algebraic geometry. Morava's use of it for formal groups is made explicit in \citet{morava1972moduli} and developed in \citet[Section 0.1]{morava1985noetherian}.}
\GSfield{Present status}{As of November 2026, the construction is standard and gives the correct language for the following geometric claims, provided the grading and base field are tracked.}
\GSfield{Assessment}{literal statement: correct with missing hypotheses; publication match: exact; verification: read the result; development: none; confidence: high}
\end{GSitem}
\begin{GSitem}{1971-268.6}{E}{The dense-open K-theory claim and the Conner–Floyd comparison}
\GSfield{Kind; source}{claim; \GSpage{P000442}}
\GSfield{The notes}{A theorem attributed to Morava and Quillen says that every fibre over a dense open set is \(K^0(X)\oplus K^1(X)\); the open is identified with inverting \(t_1=[\mathbb {CP}^1]\).}
\GSfield{Statement}{For a finite CW complex, the Conner--Floyd comparison identifies \(MU_*(X)\otimes_{MU_*}KU_*\) with \(KU_*(X)\), where the coefficient map is the multiplicative formal group law and \(KU_*=\mathbb Z[\beta^{\pm1}]\). Inverting a degree-two polynomial generator in the universal ring does not impose that coefficient map. A geometric fibre at \(\mathfrak p\) is instead a vector space over \(\kappa(\mathfrak p)\); even at \(X=\mathrm{pt}\), it is not the integral group \(K^0(\mathrm{pt})=\mathbb Z\).}
\GSfield{Reconstruction}{\ours{} The comparison is stated in the covariant homology convention of the lecture. Passing to cohomology requires the corresponding finite-complex comparison, not an unmentioned identification of homology and cohomology. The height-one locus is defined prime by prime through the Hasse invariant, not by one universal \(t_1\) in all characteristics.}
\GSfield{Fate}{The correctly delimited comparison predates the talk \citep{connerfloyd1966relation}; the formal-group interpretation is explained by \citet{quillen1969formal}. Later height-local theory implements Morava's broader idea, but does not vindicate the literal fibre claim.}
\GSfield{Present status}{As of November 2026, the Conner--Floyd theorem is established. The notebook's global dense-open fibre statement is false as written and should not be counted as an unsolved conjecture.}
\GSfield{Assessment}{literal statement: false as written; publication match: partial; verification: indirect; development: annotated; confidence: high}
\end{GSitem}
\begin{GSitem}{1971-268.7}{B}{Formal coordinate changes and equivariant bordism data}
\GSfield{Kind; source}{construction; \GSpage{P000442-P000444}}
\GSfield{The notes}{The group \(\Gamma(R)\) consists of \(f(t)=a_1t+a_2t^2+\cdots\), with \(a_1\) invertible. It acts by \(F^f(x,y)=f^{-1}F(f(x),f(y))\), and the bordism sheaf is called equivariant, ``constant on orbits''.}
\GSfield{Statement}{Formal substitutions with invertible linear term form a group under composition and act on one-dimensional commutative formal group laws by the displayed formula. Strict coordinate changes have \(a_1=1\). The flat Hopf algebroid \((MU_*,MU_*MU)\) encodes formal group laws and strict isomorphisms; the grading records the additional scaling. The comodule \(MU_*(X)\) therefore has descent data. Equivariance identifies fibres along chosen isomorphisms while allowing nontrivial stabilizer actions; it does not make the sheaf canonically constant on each orbit.}
\GSfield{Reconstruction}{The strict/full-coordinate distinction and stabilizer action are supplied. A moduli stack with its grading is the accurate modern home for the informal quotient picture.}
\GSfield{Fate}{The coordinate action is recorded in Morava's unpublished announcement communicated on 22 November 1972 \citep{morava1972moduli}. Its comodule geometry appears in the published \citet[Sections 0.1 and 1.0]{morava1985noetherian}; Johnson--Wilson's 1975 paper already develops substantial parts of the programme.}
\GSfield{Present status}{As of November 2026, the comodule and descent interpretation is established. The word ``constant'' is valid only in the weak sense of isomorphic fibres after choosing transport.}
\GSfield{Assessment}{literal statement: correct with missing hypotheses; publication match: partial; verification: read the result; development: none; confidence: high}
\end{GSitem}
\begin{GSitem}{1971-268.8}{B}{Height classifies geometric orbits in fixed positive characteristic}
\GSfield{Kind; source}{theorem; \GSpage{P000443}}
\GSfield{The notes}{Orbits \(\mathcal D_{(p,h)}\) for primes \(p\) and heights \(h\), together with an infinite-height stratum, are proposed. The theory is said to depend on describing these orbits.}
\GSfield{Statement}{Over a fixed algebraically closed field \(k\) of characteristic \(p>0\), one-dimensional commutative formal groups are classified up to isomorphism by height \(h\in\{1,2,\ldots,\infty\}\), defined by the first nonzero term of \([p]_F(x)\), of degree \(p^h\); infinite height means \([p]_F=0\). The closure of the height-\(h\) stratum consists of heights at least \(h\). In characteristic zero all such groups become isomorphic to the additive group over a field. These are fibrewise orbit statements, not a claim that all points over \(\operatorname{Spec}\mathbb Z\) lie in just three orbits.}
\GSfield{Reconstruction}{The algebraic closure and fixed characteristic are indispensable. Arithmetic forms over smaller fields require additional descent data.}
\GSfield{Fate}{The elementary height classification rests on older formal-group theory \citep{lazard1955formal}. The geometric stratification is exactly \citet[Theorem 1]{morava1972moduli}, an unpublished two-page announcement rather than a 1972 journal article; later published support includes \citet[Corollary 1.10]{johnsonwilson1975operations}.}
\GSfield{Present status}{As of November 2026, the height stratification and invariant-prime picture are established. The orbit geometry must be read with its stated base field.}
\GSfield{Assessment}{literal statement: correct with missing hypotheses; publication match: exact; verification: read the result; development: none; confidence: high}
\end{GSitem}
\begin{GSitem}{1971-268.9}{B}{Morava K-theories and the advertised periodicity}
\GSfield{Kind; source}{construction; \GSpage{P000443}}
\GSfield{The notes}{The notes announce functors \(K(h)\), with periodicity \(2(p^h-1)\), and describe an infinite-height fibre as ordinary integral cohomology.}
\GSfield{Statement}{For each prime \(p\) and integer \(h\geq1\), Morava \(K\)-theory is a homology theory with coefficient ring \(K(h)_*=\mathbb F_p[v_h^{\pm1}]\), \(|v_h|=2(p^h-1)\). Multiplication by \(v_h\) is the stated periodicity. Its construction cannot in general be replaced by the naive tensor product \(MU_*(X)\otimes_{MU_*}K(h)_*\). The additive formal group is represented by ordinary cohomology with the chosen coefficients; an infinite-height point in characteristic \(p\) gives a mod-\(p\) formal group, not automatically integral cohomology.}
\GSfield{Reconstruction}{The coefficient field, grading and distinction between a constructed homology theory and an underived fibre are supplied. The notebook's \(\mathbb Z_p\) notation is not assumed to distinguish \(p\)-adic integers from integers modulo \(p\).}
\GSfield{Fate}{\citet[Introduction, Sections 2--3]{johnsonwilson1975operations} constructs these theories and gives the coefficient ring and period explicitly. Its references identify Morava's Columbia notes as unpublished in 1971; this is earlier circulation distinct from the 1975 publication.}
\GSfield{Present status}{As of November 2026, Morava \(K\)-theories and the periodicity are established. Their heightwise relation to complex bordism has precise hypotheses, not the universal fibre identity suggested by the sketch.}
\GSfield{Assessment}{literal statement: correct with missing hypotheses; publication match: partial; verification: read the result; development: none; confidence: high}
\end{GSitem}
\begin{GSitem}{1971-268.10}{B}{Height stabilizers and p-adic division algebras}
\GSfield{Kind; source}{theorem; \GSpage{P000443}}
\GSfield{The notes}{Adams-type operations are indexed by units in a \(p\)-adic division algebra of dimension \(h^2\), described as the stabilizer of a point of a height orbit.}
\GSfield{Statement}{For a height-\(h\) one-dimensional formal group over an algebraic closure of \(\mathbb F_p\), its automorphism group is the unit group \(\mathcal O_D^\times\) of the maximal order in the central division algebra \(D/\mathbb Q_p\) of invariant \(1/h\), with \(\dim_{\mathbb Q_p}D=h^2\). The full group \(D^\times\) allows quasi-isogenies and is larger. Stabilizer actions on height-theoretic objects require the chosen coefficient extension and realization; the notebook does not construct every asserted operation.}
\GSfield{Reconstruction}{Units of the maximal order are distinguished from all nonzero division-algebra elements. The stabilizer identification is verified separately from the stronger claim of an action on a particular integral or mod-\(p\) cohomology theory.}
\GSfield{Fate}{The precise stabilizer is \citet[Theorem 2]{morava1972moduli}; the November 1972 announcement remained unpublished and was later made available from archival material. Morava's published comodule theory gives the broader realization framework \citep{morava1985noetherian}.}
\GSfield{Present status}{As of November 2026, the stabilizer identification is established. A literal \(D^\times\)-action by automorphisms of the formal group would be too large.}
\GSfield{Assessment}{literal statement: correct with missing hypotheses; publication match: partial; verification: read the result; development: none; confidence: high}
\end{GSitem}
\begin{GSitem}{1971-268.11}{E}{Formal group laws and the Lazard–Quillen universal ring}
\GSfield{Kind; source}{definition; \GSpage{P000443-P000444}}
\GSfield{The notes}{A one-dimensional commutative formal group law satisfies symmetry, a zero identity and associativity; the inverse exists automatically. The functor of such laws is represented by \(\Omega_U\).}
\GSfield{Statement}{For a commutative ring \(R\), a formal group law is \(F(x,y)\in R[[x,y]]\) with \(F(x,0)=x\), \(F(0,y)=y\), associativity and commutativity. The linear terms are \(x+y\), so the inverse series is determined recursively. The functor of laws in a chosen coordinate is represented by the Lazard ring \(L\cong\mathbb Z[x_1,x_2,\ldots]\); Quillen identifies its universal law with the complex orientation law on \(MU_*\). The representing scheme classifies laws with coordinates, not isomorphism classes of formal groups.}
\GSfield{Reconstruction}{The chosen coordinate is made explicit. Modding out coordinate changes leads to a stack rather than replacing the universal law ring by a ring of isomorphism classes.}
\GSfield{Fate}{Representability is Lazard's 1955 theorem \citep{lazard1955formal}; the topological identification is \citet[Theorem 2]{quillen1969formal}, communicated in May and published in November 1969.}
\GSfield{Present status}{As of November 2026, the universal-property statement is established, with the distinction between laws and isomorphism classes.}
\GSfield{Assessment}{literal statement: correct; publication match: exact; verification: read the result; development: none; confidence: high}
\end{GSitem}
\begin{GSitem}{1971-268.12}{E}{Invariant differentials determine the logarithm over a Q-algebra}
\GSfield{Kind; source}{claim; \GSpage{P000445}; D. A. Kazhdan}
\GSfield{The notes}{Kazhdan's intervention says that in characteristic zero the invariant differential gives coordinates for the formal group; the coefficient of each power can be prescribed.}
\GSfield{Statement}{Over a \(\mathbb Q\)-algebra \(R\), a normalized differential \(\omega=a(x)\,dx\), \(a(0)=1\), has a unique integral \(\ell(x)\in R[[x]]\) with \(\ell(0)=0\), \(\ell'(0)=1\). The law \(F(x,y)=\ell^{-1}(\ell(x)+\ell(y))\) has invariant differential \(\omega\). Thus normalized invariant differentials and commutative one-dimensional formal group laws are in bijection over \(\mathbb Q\)-algebras. Characteristic zero alone is insufficient for the integration argument: denominators may be absent, as in \(\mathbb Z\).}
\GSfield{Reconstruction}{The \(\mathbb Q\)-algebra hypothesis replaces the weaker phrase characteristic zero. The claim is about a chosen coordinate and a normalized differential.}
\GSfield{Fate}{The formula and the necessary rationalization are explicit in \citet[Section 1, formulas (1)--(2)]{quillen1969formal}. Morava returns to it on 14 April.}
\GSfield{Present status}{As of November 2026, the rational classification is established; integral formal groups retain information not removed by a logarithm over the original ring.}
\GSfield{Assessment}{literal statement: correct with missing hypotheses; publication match: exact; verification: read the result; development: none; confidence: high}
\end{GSitem}

\GStalk{269}{1972-04-14}{I. M. Gelfand}{P000445-P000446}{Dyson’s missed opportunities: eta identities and contractions}
Gelfand reports two examples from Dyson's Gibbs Lecture of 17 January 1972: the relation between powers of the Dedekind eta function and root systems, and contractions of kinematical groups. The manuscript preserves the explicit five-variable formula for Ramanujan's tau function and singles out the exponent 26 as unexplained. The missing leading factor in the eta product is a mathematical emendation, while the exceptional exponent has a documented later explanation. The final contraction remark occupies only the top of P000446, before Morava's heading.

\begin{GSitem}{1971-269.1}{E}{Ramanujan’s tau function and the missing leading factor}
\GSfield{Kind; source}{definition; \GSpage{P000445}}
\GSfield{The notes}{The notebook writes \(\prod_{m=1}^{\infty}(1-x^m)^{24}=\eta^{24}(x)=\sum_{n\geq1}\tau(n)x^n\). The first lower index is faint in transcription, but the scan and Dyson's original use \(m=1\).}
\GSfield{Statement}{For \(\tau\) in the complex upper half-plane set \(q=e^{2\pi i\tau}\) and \(\eta(\tau)=q^{1/24}\prod_{m\geq1}(1-q^m)\). Ramanujan's integers are defined by \[\Delta(q)=\eta(\tau)^{24}=q\prod_{m\geq1}(1-q^m)^{24}=\sum_{n\geq1}\tau(n)q^n.\] The product without \(q\) has constant term 1 and coefficients \(\tau(n+1)\) at \(q^n\).}
\GSfield{Reconstruction}{The variable \(q\), the upper-half-plane parameter and the leading factor are restored. The notebook conflates the eta function with its product part; this is not repaired by changing the convention for \(\tau(1)=1\).}
\GSfield{Fate}{The normalization appears explicitly in \citet[equation (1), p. 636]{dyson1972missed}. Dyson delivered the lecture in January 1972; the journal text appeared in September. The missing factor is therefore a notebook error, not a later change of definition.}
\GSfield{Present status}{As of November 2026, the normalization is standard. Its correction follows already by comparing constant terms and is also required by the independent coefficient verification in this chapter.}
\GSfield{Assessment}{literal statement: false as written; publication match: exact; verification: read the result; development: annotated; confidence: high}
\end{GSitem}
\begin{GSitem}{1971-269.2}{E}{Dyson’s five-variable lattice formula for tau}
\GSfield{Kind; source}{theorem; \GSpage{P000445}}
\GSfield{The notes}{The notes sum a ten-factor Vandermonde product over integers \(a,b,c,d,e\) satisfying residues \(1,2,3,4,5\pmod5\), zero sum and square sum \(10n\), with denominator \(1!2!3!4!\).}
\GSfield{Statement}{For every positive integer \(n\), put \(v=(a,b,c,d,e)\) and \(V(v)=\prod_{1\leq i<j\leq5}(v_i-v_j)\). Then \[\tau(n)=\frac1{288}\!\sum_{\substack{a\equiv1,\ b\equiv2,\ c\equiv3,\ d\equiv4,\ e\equiv0\ (5)\\ a+b+c+d+e=0\\ a^2+b^2+c^2+d^2+e^2=10n}}\! V(a,b,c,d,e).\] Each coordinate is labelled by its residue class; the sum is not divided by an additional permutation factor.}
\GSfield{Reconstruction}{Residue 5 is written as residue 0, and the denominator is evaluated. The formula has been checked editorially by three exact-integer algorithms for \(1\leq n\leq100\); the computational appendix marks this as \oursComputed{} and does not infer an all-\(n\) proof from the finite check.}
\GSfield{Fate}{The exact identity is \citet[equation (2), p. 636]{dyson1972missed}, already in the January 1972 lecture. Its Lie-theoretic setting is the affine denominator identities discussed by \citet{macdonald1972affine}.}
\GSfield{Present status}{As of November 2026, the identity is established. The notebook's residue order, square-sum scale and denominator agree with the primary formula; the error lies in the preceding generating function, not in this lattice sum.}
\GSfield{Assessment}{literal statement: correct; publication match: exact; verification: read the result; development: annotated; confidence: high}
\end{GSitem}
\begin{GSitem}{1971-269.3}{E}{Eta powers and products over positive roots}
\GSfield{Kind; source}{claim; \GSpage{P000445}}
\GSfield{The notes}{The exponents \(3,8,10,14,15,21,24,26,28,35,36\) are compared with dimensions of simple Lie algebras. The numerator of the tau formula is described as a product over ``simple roots (positive)'' of \(SL(5)\).}
\GSfield{Statement}{The root system \(A_4\) has ten positive roots \(e_i-e_j\), \(i<j\), and only four simple roots. Evaluating the product over all positive roots gives the ten-factor Vandermonde in the formula above. The Lie-algebra dimensions \(\dim A_1=3\), \(\dim A_2=8\), \(\dim C_2=10\), \(\dim A_4=24\) illustrate the ordinary affine-denominator specializations. The exponent 26 is exceptional to this dimension list and must be considered separately.}
\GSfield{Reconstruction}{``Simple'' is deleted before ``positive roots''. A denominator-identity family is identified, without asserting that every eta-power identity is uniquely determined by the dimension of a simple Lie algebra.}
\GSfield{Fate}{\citet[pp. 636--639]{dyson1972missed} describes the root-system discovery and Macdonald's general proof. The journal account \citep{macdonald1972affine} appeared in 1972, but the connection was already part of the January lecture summarized here.}
\GSfield{Present status}{As of November 2026, the root-system interpretation is established. Counting ten positive roots versus four simple roots provides an immediate internal check on the manuscript's wording.}
\GSfield{Assessment}{literal statement: false as written; publication match: partial; verification: read the result; development: none; confidence: high}
\end{GSitem}
\begin{GSitem}{1971-269.4}{B}{Why does the exponent 26 occur?}
\GSfield{Kind; source}{question; \GSpage{P000445}}
\GSfield{The notes}{The exponent 26 is doubly underlined; Gelfand records that it ``drops out'' of the simple-Lie-algebra dimension pattern and remains unclear.}
\GSfield{Statement}{Explain the eta-power identity with exponent 26 within the root-system framework, although 26 is not the dimension of a simple complex Lie algebra. The question concerns the structural source of that identity, not the existence of a simple Lie algebra of dimension 26.}
\GSfield{Reconstruction}{The structural formulation follows the explicit contrast with the surrounding dimension list. No unspecified coefficient identity beyond the one discussed by Dyson is invented.}
\GSfield{Fate}{\citet[Section I, pp. 200--201]{monastyrskii1980appendix} explains the exponent 26 through the dual-root construction associated with \(F_4\). The unequal root lengths lead to a specialization different from the ordinary dimension-52 pattern. This supplies a documented explanation of the notebook's exceptional exponent.}
\GSfield{Present status}{As of November 2026, the exceptional exponent is explained within affine root-system identities. It is not an unresolved problem and does not require a nonexistent 26-dimensional simple Lie algebra.}
\GSfield{Assessment}{literal statement: correct; publication match: exact; verification: read the result; development: none; confidence: high; problem \hyperref[prob:GP-1971-41]{GP-1971-41}}
\end{GSitem}
\begin{GSitem}{1971-269.5}{E}{Limits of the kinematical groups}
\GSfield{Kind; source}{method; \GSpage{P000446}}
\GSfield{The notes}{The second example consists of \(c\to0\), \(c\to\infty\) and \(R\to\infty\), called limiting transitions in the kinematical group.}
\GSfield{Statement}{A Lie-algebra contraction chooses invertible linear changes of basis depending on a parameter and takes a finite limit of the transformed structure constants. In the Poincar\'e algebra, normalize boosts by \([K_i,H]=P_i\), \([K_i,P_j]=\delta_{ij}H/c^2\) and \([K_i,K_j]=-\epsilon_{ijk}J_k/c^2\). The limit \(c\to\infty\) gives the Galilei brackets. For \(c\to0\), first set \(C_i=c^2K_i\); then \([C_i,P_j]=\delta_{ij}H\) survives while \([C_i,H]\) and \([C_i,C_j]\) vanish, giving the Carroll limit. Vanishing curvature \(R^{-2}\) contracts a de Sitter algebra to the Poincar\'e algebra.}
\GSfield{Reconstruction}{The changes of basis and bracket normalizations are supplied to make the shorthand meaningful; limits of groups without specifying coordinates or rescaling are ambiguous. We do not reconstruct a complete classification from this one line.}
\GSfield{Fate}{These are precisely the three processes discussed by \citet[Section 4, pp. 642--643]{dyson1972missed}, which cites Bacry and L\'evy-Leblond's 1968 analysis. Both the contraction method and these examples predate the seminar.}
\GSfield{Present status}{As of November 2026, the contractions are established examples. The excerpt contains no separate conjecture about physical applicability or an unproved classification.}
\GSfield{Assessment}{literal statement: correct with missing hypotheses; publication match: partial; verification: read the result; development: none; confidence: high}
\end{GSitem}

\GStalk{270}{1972-04-14}{Jack Morava}{P000446-P000449}{Formal groups, invariant differentials and homogeneous orbits}
Kirillov first recapitulates the geometric bordism construction and the picture of a bordism module over the coefficient scheme. The discussion then clarifies geometric fibres, coordinate changes, characteristic-zero logarithms and Mishchenko's projective-space coefficients. It ends with a compressed homogeneous-space assertion. The last conclusion is at the top of P000449; P000450--P000451 belong to Gabrielov's 25 September 1972 lecture, outside this record.

\begin{GSitem}{1971-270.1}{E}{Kirillov’s recap of bordism cycles over a space}
\GSfield{Kind; source}{definition; \GSpage{P000446}; A. A. Kirillov}
\GSfield{The notes}{Under the heading ``Kirillov (recapitulation)'', a cycle is a quadruple \((M,\xi,\alpha,f)\): \(M\) is a compact weakly complex manifold, \(\alpha\) is a specified stable tangent-bundle isomorphism, and \(f:M\to X\). A diagram gives the extending bundle and map on a bordism.}
\GSfield{Statement}{For a CW space \(X\), a geometric complex-bordism cycle consists of a closed smooth manifold with a chosen stable complex structure and a continuous map to \(X\). Two cycles are equivalent when their disjoint union with the opposite bordism structure bounds a compact manifold whose stable complex structure and map extend the boundary data. Disjoint union and Cartesian products give \(MU_*(X)\) its graded \(MU_*\)-module structure.}
\GSfield{Reconstruction}{``Weak complex set'' in the translation is read as a stably complex manifold. The stable isomorphism is part of the structure and must extend; existence of a complex bundle with the same stable real class alone does not specify a cycle.}
\GSfield{Fate}{This repeats the established geometric construction in 1971-268.1 and 1971-268.4; its pre-seminar context is \citet{connerfloyd1966relation}. The recurrence is an exposition, not a second proposed theory.}
\GSfield{Present status}{As of November 2026, the geometric model and module action are standard. This entry records Kirillov's intervention separately from Morava's later assertions.}
\GSfield{Assessment}{literal statement: correct with missing hypotheses; publication match: exact; verification: publisher metadata checked; development: none; confidence: high}
\end{GSitem}
\begin{GSitem}{1971-270.2}{E}{The coefficient ring and the affine sheaf attached to bordism}
\GSfield{Kind; source}{theorem; \GSpage{P000446}; A. A. Kirillov}
\GSfield{The notes}{The coefficient theorem is now written \(A=\Omega_*^U(*)=\mathbb Z[t_1,t_2,\ldots]\), \(\deg t_i=2i\), with \(t_1\leftrightarrow\mathbb {CP}^1\). The module \(M=\Omega_*^U(X)\) gives a sheaf over \(\mathbb A=\operatorname{Spec}A\).}
\GSfield{Statement}{Complex bordism has coefficient ring \(MU_*\cong\mathbb Z[x_1,x_2,\ldots]\), with \(|x_i|=2i\), and a degree-two generator may be chosen as \([\mathbb {CP}^1]\). Any module \(M\) defines a quasicoherent sheaf \(\widetilde M\) on \(\operatorname{Spec}MU_*\), after forgetting or separately recording its grading. The infinitely generated polynomial coefficient ring is not noetherian.}
\GSfield{Reconstruction}{This page correctly starts the polynomial generators at index 1, unlike P000441. The nonnoetherian qualification matters when comparing the later finite-type homogeneous-space assertion with the infinite formal-group scheme.}
\GSfield{Fate}{The coefficient theorem and identification with the universal formal-group ring are explicit in \citet[Theorem 2]{quillen1969formal}; the associated-sheaf construction is standard affine algebraic geometry. Morava's equivariant version is developed in \citet[Section 1.0]{morava1985noetherian}.}
\GSfield{Present status}{As of November 2026, the coefficient and sheaf constructions are established. They do not by themselves identify the fibres with familiar integral cohomology theories.}
\GSfield{Assessment}{literal statement: correct; publication match: exact; verification: read the result; development: none; confidence: high}
\end{GSitem}
\begin{GSitem}{1971-270.3}{B}{The three orbit types and the proposed fibre theories}
\GSfield{Kind; source}{claim; \GSpage{P000446-P000447}}
\GSfield{The notes}{There is said to be one dense open orbit, one closed orbit and countably many intermediate orbits; their fibres are labelled \(K(X)\), \(H^*(X,\mathbb Z)\), and \(K(n)(X,\mathbb Z_p)\). The intermediate labels also involve the prime \(p\).}
\GSfield{Statement}{Over an algebraically closed field of a fixed positive characteristic \(p\), coordinate-change orbits of one-dimensional commutative formal groups are indexed by heights \(1,2,\ldots,\infty\): height 1 is open, infinite height is closed, and finite higher heights lie between them. The positive finite-height cohomology theories are Morava \(K(h)\), with coefficients \(\mathbb F_p[v_h^{\pm1}]\), \(|v_h|=2(p^h-1)\). These are not a literal classification into three orbits over \(\operatorname{Spec}\mathbb Z\), nor are all their geometric fibres integral \(K\)-groups or integral ordinary cohomology.}
\GSfield{Reconstruction}{The orbit theorem is restricted to a fixed algebraically closed base field. The labels of homology theories are separated from raw tensor-product fibres, repeating the corrections in 1971-268.6--1971-268.9. No claim is made that the fibre alone defines a generalized homology theory for every space.}
\GSfield{Fate}{The geometric orbit theorem is \citet[Theorem 1]{morava1972moduli}, an unpublished November 1972 announcement. The coefficient rings and realization of the periodic theories are treated in \citet[Introduction and Sections 2--3]{johnsonwilson1975operations}, citing Morava's unpublished 1971 work.}
\GSfield{Present status}{As of November 2026, the height picture and the periodic theories are established. The global integral-fibre formulation remains false as written, and is not an open problem.}
\GSfield{Assessment}{literal statement: false as written; publication match: partial; verification: read the result; development: none; confidence: high}
\end{GSitem}
\begin{GSitem}{1971-270.4}{E}{The fibre M/MJ and the residue field}
\GSfield{Kind; source}{construction; \GSpage{P000447}}
\GSfield{The notes}{A point \((a_1,a_2,\ldots)\) gives a purported maximal ideal \(J=(t_1-a_1,t_2-a_2,\ldots)\). The fibre is \(M/MJ\), called a vector space and then more accurately an \(A/J\)-module.}
\GSfield{Statement}{For an ideal \(J\subset A\), \(M/JM\cong M\otimes_A A/J\) is an \(A/J\)-module. It is a vector space when \(J\) is maximal. The geometric fibre at a prime \(\mathfrak p\) is \(M\otimes_A\kappa(\mathfrak p)\), not in general \(M/\mathfrak p M\). For \(A=\mathbb Z[t_1,t_2,\ldots]\) and integers \(a_i\), the displayed evaluation ideal has quotient \(\mathbb Z\), so it is prime but not maximal; a residue characteristic must also be imposed to obtain the corresponding finite-characteristic point.}
\GSfield{Reconstruction}{The correction already begun in the notebook is completed. Points lie in \(\operatorname{Spec}A\), not in the module \(M\); an ideal is a subset of \(A\), not a subset of the scheme \(\mathbb A\).}
\GSfield{Fate}{This is the ordinary affine-scheme meaning of fibre. The same coefficient-scheme viewpoint underlies \citet[Section 1.0]{morava1985noetherian}; that later account clarifies the geometry rather than introducing the elementary tensor-product identity.}
\GSfield{Present status}{As of November 2026, the construction is standard. Keeping the base ring prevents a false vector-space assertion and clarifies the varying coefficient fields of the height strata.}
\GSfield{Assessment}{literal statement: false as written; publication match: partial; verification: read the result; development: none; confidence: high}
\end{GSitem}
\begin{GSitem}{1971-270.5}{E}{Full and strict formal coordinate changes}
\GSfield{Kind; source}{definition; \GSpage{P000447}}
\GSfield{The notes}{The group consists of \(g(x)=a_1x+a_2x^2+\cdots\). Both \(a_1=1\) and the condition that \(a_1\) be a unit occur. Laws are equivalent under \(F_1=\varphi^{-1}F_2(\varphi(x),\varphi(y))\).}
\GSfield{Statement}{For any commutative ring \(R\), series \(g\in xR[[x]]\) with invertible linear coefficient form a group under substitution. The inverse is found recursively. The subgroup with \(g'(0)=1\) consists of strict coordinate changes; over \(\mathbb Z\), the full group permits \(g'(0)=\pm1\). Isomorphism of one-dimensional formal group laws means an invertible series intertwining the laws, equivalently the displayed conjugation formula.}
\GSfield{Reconstruction}{The two linear-coefficient conditions specify different groups, not interchangeable conventions. The ambiguous margin phrase translated as ``same-form groups'' means isomorphic formal groups in the displayed formula.}
\GSfield{Fate}{Formal changes of coordinate are part of Lazard's pre-seminar theory \citep{lazard1955formal}. Their full/strict distinction and grading are explicit in \citet[Sections 1.0.2--1.0.4]{morava1985noetherian}.}
\GSfield{Present status}{As of November 2026, the construction is established over arbitrary commutative rings, with invertibility required only for the linear coefficient.}
\GSfield{Assessment}{literal statement: correct with missing hypotheses; publication match: exact; verification: read the result; development: none; confidence: high}
\end{GSitem}
\begin{GSitem}{1971-270.6}{F}{Formal vector fields and the ambiguous dual of differentials}
\GSfield{Kind; source}{remark; \GSpage{P000447}}
\GSfield{The notes}{The margin names \(W_1\), the algebra of formal vector fields, and \(W'_1\), a ``dual algebra of differentials'', without specifying the duality or its operations.}
\GSfield{Statement}{Over a characteristic-zero field \(k\), the formal vector fields are \(W_1=k[[x]]\partial_x\), with \([f\partial_x,g\partial_x]=(fg'-gf')\partial_x\). The module of formal differentials \(k[[x]]\,dx\) is dual to this rank-one module over \(k[[x]]\). This module dual is distinct from the continuous \(k\)-linear dual of the infinite-dimensional Lie algebra. The Lie algebra of origin-fixing formal substitutions is the subalgebra \(xk[[x]]\partial_x\).}
\GSfield{Reconstruction}{These are precise nearby constructions, not an identification of the unexplained prime on \(W'_1\). No Lie bracket on the differential module, topology for a continuous dual, or claim about a coadjoint representation is supplied by the notes.}
\GSfield{Fate}{The distinctions are explicit in the modern formal-geometry treatment \citet[Section 2.1 and Corollary 2.5.18]{gorbounov2020chiral}. Searches for formal vector fields, dual differentials and Gelfand--Fuks terminology recover that framework, but no exact additional theorem can be assigned to this isolated line.}
\GSfield{Present status}{As of November 2026, the elementary objects are standard, but the intended meaning of ``dual algebra'' remains unidentified in this occurrence. This is a limitation of the record, not evidence of an undeveloped mathematical idea.}
\GSfield{Assessment}{literal statement: not determinable; publication match: partial; verification: read the result; development: none; confidence: medium}
\end{GSitem}
\begin{GSitem}{1971-270.7}{E}{Associativity and the characteristic-zero commutativity theorem}
\GSfield{Kind; source}{claim; \GSpage{P000447-P000448}}
\GSfield{The notes}{A law \(F=x+y+\alpha xy+\beta x^2y+\delta xy^2+\cdots\) is said to have associativity as its sole requirement. On the next page, over a field of characteristic zero, it is declared isomorphic to addition and hence commutative.}
\GSfield{Statement}{Let \(R\) be a \(\mathbb Q\)-algebra and let \(F\in R[[x,y]]\) satisfy \(F(x,0)=x\), \(F(0,y)=y\) and associativity. Writing \(f(x)=\partial_2F(x,0)\), define \(\ell(x)=\int_0^x dt/f(t)\). Then \(\ell(F(x,y))=\ell(x)+\ell(y)\); thus \(F\) is isomorphic to the additive law and is commutative. Over arbitrary rings, associativity and the displayed identity terms define a possibly noncommutative one-dimensional law; commutativity cannot be silently added.}
\GSfield{Reconstruction}{\ours{} Differentiating associativity in its third variable at zero gives \(f(F(x,y))=\partial_2F(x,y)f(y)\). Consequently the derivative in \(y\) of \(\ell(F(x,y))\) is \(\ell'(y)\); the value at \(y=0\) gives the stated identity. This checks the notebook's logarithm argument without assuming commutativity first. The integral uses division by all positive integers.}
\GSfield{Fate}{Characteristic-zero logarithms are classical formal-group theory \citep{lazard1955formal}; the normalized logarithm and rational coefficient requirement are explicit in \citet[Section 1]{quillen1969formal}. The derivation above is our check of the slightly stronger associative formulation used by the notes.}
\GSfield{Present status}{As of November 2026, the characteristic-zero field assertion is established. Its unrestricted-ring extrapolation is invalid; integral and positive-characteristic formal groups require separate hypotheses.}
\GSfield{Assessment}{literal statement: correct with missing hypotheses; publication match: partial; verification: read the result; development: annotated; confidence: high}
\end{GSitem}
\begin{GSitem}{1971-270.8}{E}{Invariant differentials, pullback and the obstruction to integration}
\GSfield{Kind; source}{construction; \GSpage{P000448}}
\GSfield{The notes}{The invariant differential is \(a(x)dx=dx/f(x)\); its integral gives the inverse coordinate. The substitution rule is \(a(x)\mapsto a(\varphi(x))\varphi'(x)\). The text says any series \(a(x)=1+a_1x+\cdots\) determines a formal group, and notes that \(\int x^{p-1}dx\) does not exist in characteristic \(p\).}
\GSfield{Statement}{Over a \(\mathbb Q\)-algebra, every \(a(x)\in1+xR[[x]]\) has the unique normalized primitive \(\ell(x)=x+\sum_{n\geq1}a_nx^{n+1}/(n+1)\). The law \(\ell^{-1}(\ell(x)+\ell(y))\) has normalized invariant differential \(a(x)dx\). For a strict change \(\varphi'(0)=1\), pullback sends its coefficient to \(a(\varphi(x))\varphi'(x)\). For a full coordinate change, the normalized differential acquires the additional factor \(1/\varphi'(0)\). In characteristic \(p\), the derivative of every power series has zero coefficient at \(x^{p-1}\), so that monomial has no primitive.}
\GSfield{Reconstruction}{The scope remains the characteristic-zero field in the preceding theorem, or more generally a \(\mathbb Q\)-algebra. Arbitrary pullback one-forms need not remain normalized. The isolated dot above the second inverse-coordinate symbol is not interpreted as an extra derivative; that reading would contradict the displayed integral.}
\GSfield{Fate}{The invariant differential and logarithm formulas are in \citet[Section 1, equations (1)--(2)]{quillen1969formal}, before the seminar. The substitution factor follows from the usual pullback of one-forms and the choice of normalization.}
\GSfield{Present status}{As of November 2026, the rational correspondence is established. The universal Lazard ring does not remove the positive-characteristic integration obstruction or make every normalized differential over an arbitrary ring arise by this logarithm formula.}
\GSfield{Assessment}{literal statement: correct with missing hypotheses; publication match: partial; verification: read the result; development: none; confidence: high}
\end{GSitem}
\begin{GSitem}{1971-270.9}{E}{Mishchenko’s projective-space coefficients}
\GSfield{Kind; source}{theorem; \GSpage{P000448}}
\GSfield{The notes}{Morava identifies the coefficients of \(a(x)\) with \([\mathbb {CP}^1],[\mathbb {CP}^2],\ldots\), citing Mishchenko; the page also records a suggestion by Kirillov.}
\GSfield{Statement}{For the formal group law of the canonical complex orientation of \(MU\), its normalized invariant differential and rational logarithm are \[\omega_F(x)=\left(\sum_{n\geq0}[\mathbb {CP}^{n}]x^n\right)dx,\qquad \log_F(x)=\sum_{n\geq0}\frac{[\mathbb {CP}^{n}]}{n+1}x^{n+1},\] with \([\mathbb {CP}^{0}]=1\). The differential has integral bordism coefficients; the logarithm is taken in \(MU_*\otimes\mathbb Q\), with grading interpreted in the usual complex-oriented convention.}
\GSfield{Reconstruction}{The denominator belongs in the logarithm, not in the differential. A complex orientation and its coordinate are fixed; arbitrary polynomial generators of \(MU_*\) are not being equated with every projective-space class.}
\GSfield{Fate}{\citet[Corollary to Theorem 1]{quillen1969formal} explicitly credits Mishchenko and gives the logarithm formula. Its differential is exactly the coefficient assertion recorded here, already published in 1969.}
\GSfield{Present status}{As of November 2026, the formula is established. It connects geometric projective-space classes with the formal-group logarithm and supplies a precise bridge between the two parts of the lecture.}
\GSfield{Assessment}{literal statement: correct; publication match: exact; verification: read the result; development: none; confidence: high}
\end{GSitem}
\begin{GSitem}{1971-270.10}{E}{When a transitive group-scheme action is a homogeneous quotient}
\GSfield{Kind; source}{claim; \GSpage{P000448-P000449}}
\GSfield{The notes}{The final theorem assumes a smooth \(G\)-scheme \(X\) over \(\mathbb Z\), a flat group scheme \(G\), noetherian hypotheses, and transitivity of \(G(K)\) on \(X(K)\) for an algebraically closed field. The continuation asserts \(X\cong G/S\), with \(S\) a stabilizer.}
\GSfield{Statement}{A precise field version is the following. Let \(k\) be algebraically closed, \(G\) a smooth affine algebraic group over \(k\), and \(X\) a nonempty separated reduced finite-type \(k\)-scheme on which \(G(k)\) acts transitively. For \(x\in X(k)\), the scheme-theoretic stabilizer \(S=G_x\) is closed and the orbit map identifies \(X\) with the faithfully flat quotient \(G/S\). The stabilizer need not be reduced. Over a base such as \(\operatorname{Spec}\mathbb Z\), a corresponding quotient description needs a global section and a faithfully flat orbit map. Geometric-point transitivity alone does not supply those data.}
\GSfield{Reconstruction}{The verified statement is explicitly restricted to a field. The manuscript's quantifier before \(K\) is crossed out, and geometric-fibre transitivity does not itself provide a chosen global section or a globally defined stabilizer over \(\mathbb Z\). Nor do finite-type noetherian hypotheses automatically apply to the infinite-dimensional formal-coordinate group used earlier.}
\GSfield{Fate}{The orbit and quotient theorem is established algebraic-group theory; see \citet[Propositions 7.11--7.17, especially Remark 7.16(c)]{milne2017groups}. This gives a verified precise core, not verification of the stronger relative formulation exactly as abbreviated in the notebook.}
\GSfield{Present status}{As of November 2026, the field theorem and the faithfully flat orbit criterion are established. The intended full relative theorem remains underdetermined by the cancelled quantifier and missing section data.}
\GSfield{Assessment}{literal statement: correct with missing hypotheses; publication match: partial; verification: read the result; development: none; confidence: medium}
\end{GSitem}

\GStalk{271}{}{-}{P000449-P000450}{Catalogue heading: 1972}
The catalogue marks the start of a new year. No mathematical assertion belongs to this heading. P000449 ends Morava's 14 April 1972 discussion; P000450 begins Gabrielov's 25 September 1972 lecture, continued on P000451. These adjacent pages explain the supplied overlapping range but do not turn Gabrielov's talk into record 271.


\GStalk{286}{}{-}{P000477-P000478}{Catalogue heading: YEAR 1973}
This catalogue heading separates calendar-year material and carries no mathematical talk. P000477 is the end of Buchstaber's 20 November 1972 discussion of formal groups; P000478--P000479 belong to Buchstaber on 12 February 1973. The heading's supplied page range includes both sides of that boundary. Neither adjoining lecture is assigned mathematical occurrences under record 286.


\GStalk{288}{1973-02-19}{R. L. Dobrushin}{P000480}{Nelson’s Euclidean Markov-field construction}
The manuscript explicitly dates this short record 19 February 1973, although the catalogue supplied with the task says 19 February 1972. The YEAR1973 heading and adjacent Buchstaber dates support a catalogue typo. It is retained among the 45 requested records but uses the corrected academic-year prefix 1972. Gelfand distinguishes quantizing given equations from finding physically adequate equations; Dobrushin then identifies Nelson's probabilistic Euclidean-field method and says that Glimm--Jaffe results can ``almost'' be recovered. Only P000480 belongs to this talk: P000479 and P000481 are parts of different Buchstaber lectures.

\begin{GSitem}{1972-288.1}{F}{Gelfand’s programme of quantizing existing equations}
\GSfield{Kind; source}{problem; \GSpage{P000480}; I. M. Gelfand}
\GSfield{The notes}{Gelfand's first task is \ru{проквантовать существующ. ур-я}: quantize the existing equations.}
\GSfield{Statement}{No specific classical equation, phase space, quantization rule, dimension or required quantum theory is supplied. The line therefore states a research programme; it cannot be converted into a definite existence or uniqueness question without adding its mathematical object and success criterion.}
\GSfield{Reconstruction}{The neighbouring Nelson discussion suggests constructive quantum field theory as context, but does not determine whether Gelfand intended a particular polynomial interaction or a canonical quantization prescription.}
\GSfield{Fate}{Contemporary context is supplied by \citet[Introduction and Section 1]{cartier1974quantum}, whose November 1972 lecture discusses explicit scalar-field equations and the gap between formal equations and construction. That is thematic evidence, not an exact match for the unspecified equation in this occurrence.}
\GSfield{Present status}{As of November 2026, the problem cannot be identified at the level required for an open-problem entry. No claim about its solution, insolubility or lack of later work follows from this shorthand.}
\GSfield{Assessment}{literal statement: not determinable; publication match: thematic only; verification: indirect; development: none; confidence: low}
\end{GSitem}
\begin{GSitem}{1972-288.2}{F}{Gelfand’s programme of finding physically adequate equations}
\GSfield{Kind; source}{problem; \GSpage{P000480}; I. M. Gelfand}
\GSfield{The notes}{The second task is to write equations that really correspond to physics; no model or observable is named.}
\GSfield{Statement}{This is a request to choose a physically adequate dynamical model. The notes contain no mathematical class of equations, data to be reproduced, accuracy criterion, or existence assertion. A precise modern formulation cannot be supplied from this record alone.}
\GSfield{Reconstruction}{The distinction from the first task is preserved: constructing a quantum theory for a specified classical model and selecting that model are different questions. Neither is identified with a later named conjecture.}
\GSfield{Fate}{The model-construction discussion in \citet[Section 1.5]{cartier1974quantum} provides contemporary background. Searches centered on Dobrushin, Nelson and the listed date did not identify a particular equation or claim corresponding to Gelfand's second programme.}
\GSfield{Present status}{As of November 2026, the intended mathematical problem remains unidentified. This is not evidence that a definite 1973 conjecture remains open.}
\GSfield{Assessment}{literal statement: not determinable; publication match: thematic only; verification: unresolved lead; development: none; confidence: low}
\end{GSitem}
\begin{GSitem}{1972-288.3}{E}{Euclidean-invariant random fields and the spatial Markov property}
\GSfield{Kind; source}{definition; \GSpage{P000480}}
\GSfield{The notes}{Dobrushin writes random quantities \(\xi_t\), indexed by \(t\in\mathbb R^{\nu+1}\), satisfying Euclidean invariance and the Markov property.}
\GSfield{Statement}{A random generalized scalar field is a continuous linear map \(T:\mathcal S(\mathbb R^d)\to L^0(\Omega,\mathbb P)\), where \(L^0\) carries convergence in probability. Euclidean invariance means equality of joint laws after every Euclidean isometry. Let \(\mathcal F(U)\) be generated by field tests supported in an open region \(U\), and choose a specified boundary sigma-field \(\mathcal F(\partial U)\). The spatial Markov property says that, for integrable \(\mathcal F(U)\)-measurable \(Z\), \(\mathbb E[Z\mid\mathcal F(U^c)]=\mathbb E[Z\mid\mathcal F(\partial U)]\). This concerns conditional independence of spatial regions, not merely a one-parameter Markov process.}
\GSfield{Reconstruction}{Generalized fields replace pointwise random variables because the latter need not exist. The boundary sigma-field is part of the definition: in the Sobolev formulation below it is generated by the extended field evaluated on supported \(H^{-1}\) distributions. It must not automatically be replaced by an intersection of neighbourhood sigma-fields.}
\GSfield{Fate}{This is Nelson's framework, already described in \citet[Section 3.1 and note (3), p. 27]{cartier1974quantum}, a November 1972 lecture published in 1974. Nelson's own bibliography dates the initial Berkeley exposition to August 1971, with proceedings publication in 1973 \citep{nelson1973quantum}.}
\GSfield{Present status}{As of November 2026, the framework is established. The notebook's two labels identify the method but do not specify the analytic regularity needed for a reconstruction theorem.}
\GSfield{Assessment}{literal statement: correct with missing hypotheses; publication match: partial; verification: read the result; development: none; confidence: high}
\end{GSitem}
\begin{GSitem}{1972-288.4}{E}{Nelson reconstruction with explicit analytic assumptions}
\GSfield{Kind; source}{method; \GSpage{P000480}}
\GSfield{The notes}{The main author of a probabilistic method for constructing objects satisfying the Wightman axioms is identified as Nelson. Only Euclidean invariance and a Markov property are written.}
\GSfield{Statement}{One precise version, in four Euclidean dimensions, starts with a real random tempered distribution \(T\) whose law is invariant under all Euclidean isometries, including reflection. Assume that \(T\) extends continuously to \(H^{-1}(\mathbb R^4)\). For closed \(F\), let \(\mathcal F(F)\) be generated by \(T(u)\) with \(u\in H^{-1}\) supported in \(F\), and impose the spatial Markov property with these sigma-fields. Conditional expectation and time translation then give a strongly continuous self-adjoint contraction semigroup \(e^{-tH}\), \(H\geq0\), on \(L^2(\Omega,\mathcal F(\{t=0\}),\mathbb P)\). Put \(\phi_0(u)=T(u\otimes\delta_0)\). Assume in addition that some integer \(m\geq0\) satisfies \[\int |\phi_0(u)|\,|\Phi|^2\,d\mathbb P<\infty\quad\text{for all }u\in\mathcal S(\mathbb R^3),\ \Phi\in D(H^m).\] Nelson's construction then yields smeared fields on the smooth domain of \(H\); restricting to the cyclic vacuum subspace gives temperedness, Poincar\'e covariance, the spectrum condition and local commutativity. Uniqueness of the invariant vacuum is not asserted without an additional ergodicity or clustering hypothesis.}
\GSfield{Reconstruction}{The dimension-four restriction, Sobolev extension, supported-distribution boundary fields and energy-domain condition are taken from the contemporary precise account, not from the two notebook bullets. The displayed integrability condition contains the first power of \(|\phi_0(u)|\); this was checked visually on p. 20. The statement is not the assertion that Euclidean invariance plus an arbitrary Markov property automatically implies all Wightman axioms.}
\GSfield{Fate}{The construction was presented in Nelson's 1971 Berkeley lecture and published in \citet{nelson1973quantum,nelson1973construction}. The detailed formulation above is read in \citet[Sections 3.2--3.4, pp. 19--23, and note (3)]{cartier1974quantum}, whose lecture predates this corrected February 1973 occurrence. The later Osterwalder--Schrader approach has its own hypotheses and a documented 1975 correction \citep{osterwalderschrader1975euclidean}.}
\GSfield{Present status}{As of November 2026, Nelson reconstruction is established with explicit analytic assumptions. This entry identifies and delimits the method; it does not claim that the notebook specifies a new interacting model or verifies those assumptions for one.}
\GSfield{Assessment}{literal statement: correct with missing hypotheses; publication match: partial; verification: read the result; development: none; confidence: high}
\end{GSitem}
\begin{GSitem}{1972-288.5}{F}{The unspecified recovery of Glimm–Jaffe results}
\GSfield{Kind; source}{claim; \GSpage{P000480}}
\GSfield{The notes}{The notebook says that one can ``almost'' obtain the results of Glimm and Jaffe by the probabilistic method, without naming those results.}
\GSfield{Statement}{No interaction polynomial, spacetime dimension, cutoff regime, limit or axiom verified is supplied. Consequently there is no definite theorem to compare with the Glimm--Jaffe constructions. A plausible historical context is the use of Euclidean methods in two-dimensional interacting scalar-field theory, but that context does not identify the word ``almost'' mathematically.}
\GSfield{Reconstruction}{No full construction in four dimensions or assertion of all Wightman axioms for a specified nonlinear equation is inferred. The modern statement remains a delimited identification problem.}
\GSfield{Fate}{\citet[Section 3.5, pp. 23--24]{cartier1974quantum} discusses Nelson's relation to truncated two-dimensional Glimm--Jaffe Hamiltonians and Nelson symmetry. Its note added in November 1973 records further measure constructions, showing why talk date and publication date must be separated. These are plausible nearby results, not an exact match for the unqualified notebook line.}
\GSfield{Present status}{As of November 2026, the exact referenced result and omitted qualification remain unidentified in this audit. Established constructive-field literature exists; the absence of a precise notebook statement gives no basis for class C or D.}
\GSfield{Assessment}{literal statement: not determinable; publication match: thematic only; verification: indirect; development: none; confidence: low}
\end{GSitem}

\section{Developments}\label{dev:1971}
The work below is an editorial verification of three precise statements in the
notebook.  It makes no claim of a new general theorem or of a solution to an
unresolved problem in the problem index.  Results proved here are explicitly
marked; historical identities retain their attribution.  The decision about
research development of the class C and D items is recorded in the overview
and report.

\subsection{The grading check for Atiyah's de Rham example}
\label{verification:grading}

On \GSpage{P000373}, the notes use \(d+\delta\) on differential forms as an
example in the discussion of spectral flow and an obstruction to an invariant
metric.  The full de Rham operator needs to be distinguished from the odd
signature operator.  The following elementary check concerns the full operator
and does not change the general spectral-flow/index theorem recorded in the
preceding pages.
The general index identification used below is
\citet[\S7, Theorem 7.4 and pp.94--95]{atiyahpatodisinger1976spectral3}.

\begin{proposition}[Vanishing for a graded loop]\ours{}
Let \(D_t\), \(0\leq t\leq1\), be a smooth family of self-adjoint first-order
elliptic operators on a fixed Hermitian vector bundle over a closed smooth
manifold.  Regard them as operators on a fixed Hilbert space with the same
Sobolev domain.  Suppose \(D_1=UD_0U^{-1}\), where \(U\) is the unitary gluing
operator, and suppose a fixed unitary involution \(\Gamma\) preserves the
domains, commutes with \(U\), and satisfies
\[
\Gamma D_t\Gamma=-D_t\qquad(0\leq t\leq1).
\]
Then the spectral flow of the resulting closed family is zero.  In particular,
the full de Rham family \(d+\delta_{g_t}\), glued by the action of a
diffeomorphism and identified with a fixed Hilbert space in a
degree-preserving way, has zero spectral flow.
\end{proposition}

\begin{proof}
Spectral flow is unchanged by a unitary conjugacy of the closed family, and
replacing every operator by its negative reverses the oriented crossing
count.  These properties follow first for a family with isolated regular
crossings, and then for arbitrary families by perturbation and homotopy
invariance.  For a closed family the dimensions of the endpoint kernels
agree, so no endpoint correction survives.  Consequently,
\[
\operatorname{sf}(D_t)=\operatorname{sf}(\Gamma D_t\Gamma)
=\operatorname{sf}(-D_t)=-\operatorname{sf}(D_t).
\]
The spectral flow is an integer, so it is zero.

For the last assertion, let \(\Gamma\) act by \((-1)^k\) on forms of degree
\(k\).  Both \(d\) and \(\delta_{g_t}\) reverse this parity, and a
diffeomorphism preserves degree.  The positive changes of density and of the
fibre metric used to identify the varying \(L^2\) spaces are block diagonal
by degree and hence commute with \(\Gamma\).  All the required identities
therefore hold after those identifications.
\end{proof}

\paragraph{Independent check and scope.}
For geometric bundle gluing, the same conclusion can be checked by the
suspended operator.  On the mapping torus use a unitary connection along the circle which preserves
\(\Gamma\), and put \(P=\nabla_t+D_t\).  Integration by parts gives
\(P^*=-\nabla_t+D_t\), while
\(\Gamma P\Gamma=\nabla_t-D_t=-P^*\).  Unitary conjugation preserves index
and the adjoint changes its sign, whence \(\operatorname{ind}P=0\).  This is
an independent algebraic check of the vanishing once the historical
spectral-flow/index identification is used.  It explains why an unspecified
reference to \(d+\delta\) cannot, by itself, supply a nonzero obstruction.
It does not assert that every natural self-adjoint elliptic family has zero
spectral flow, nor does it identify a signature grading that the notes do not
specify.

\subsection{An exact finite audit of Dyson's coefficient formula}
\label{verification:dyson}

The product printed on \GSpage{P000445} needs the normalization
\[
\Delta(q)=q\prod_{m=1}^{\infty}(1-q^m)^{24}
=\sum_{n=1}^{\infty}\tau(n)q^n.
\]
The leading \(q\) is already necessary on comparison of constant terms:
the unshifted product has constant coefficient one and the displayed series
has constant coefficient zero.  This is a mathematical emendation of the
notebook formula, not an inference about an invisible pen stroke.
The historical all-coefficient formula is
\citet[pp.636--637, equations (1)--(3)]{dyson1972missed}.

\begin{proposition}[Finite coefficient agreement]\oursComputed{}
For every integer \(1\leq n\leq100\), the coefficient of \(q^n\) in the
corrected product agrees exactly with
\[
\sum_{\substack{
 (a,b,c,d,e)\equiv(1,2,3,4,0)\pmod5\\
 a+b+c+d+e=0\\a^2+b^2+c^2+d^2+e^2=10n}}
\frac{\displaystyle\prod_{1\leq i<j\leq5}(v_i-v_j)}{1!2!3!4!},
\qquad v=(a,b,c,d,e).
\]
This assertion reports a finite exact computation and is not a proof of the
all-\(n\) identity.
\end{proposition}

\paragraph{Algorithms and exhaustiveness.}
The file \texttt{code/check\_dyson.py} first multiplies the finite binomial
expansions of \((1-q^m)^{24}\), retaining terms only through the requested
degree.  Factors with larger \(m\) cannot change those coefficients.  A second
algorithm computes the coefficients \(a_n\) of the unshifted product from
its formal logarithmic derivative:
\[
a_0=1,\qquad
n a_n=-24\sum_{j=1}^{n}\sigma_1(j)a_{n-j},\qquad
\tau(n)=a_{n-1}.
\]
The third algorithm enumerates the first four lattice coordinates with the
specified residues and sets the fifth equal to the negative of their sum.
For a run through \(N\), every coordinate is bounded in absolute value by
\(\lfloor\sqrt{10N}\rfloor\); the squared-length condition then selects
the coefficient.  Thus the enumeration omits no qualifying tuple.  All
arithmetic is integral, including explicit divisibility checks for the
denominator \(1!2!3!4!=288\).

\begin{center}
\begin{tabular}{r r r r}
\toprule
\(n\) & \(\tau(n)\) & \(n\) & \(\tau(n)\)\\
\midrule
1&1&7&-16744\\
2&-24&8&84480\\
3&252&9&-113643\\
4&-1472&10&-115920\\
5&4830&11&534612\\
6&-6048&12&-370944\\
\bottomrule
\end{tabular}
\end{center}

\paragraph{Second check and limitation.}
The complete comparison was run separately with \(N=32\) and \(N=100\);
all three arrays agreed in both runs.  The two output files in
\texttt{audit/} record every checked coefficient and the number of contributing
lattice tuples.  Agreement also checks the residue ordering and the sign of
the ten-factor Vandermonde as copied.  The verification supports the
normalization correction; it does not convert finite evidence into an
identity for every integer, and it has no bearing on the separate question
about the exponent \(26\).
An independent review also checked the labelled residue classes, the absence
of a factorial permutation divisor, the enumeration bound, and the recurrence;
it found no error.

\subsection{Recovering the eighty-million lattice bound}
\label{verification:mass}

The numerical remark on \GSpage{P000398} becomes identifiable after restoring
positive definiteness, integrality and evenness.  For positive-definite even
unimodular lattices of rank \(n=8k>0\), the classical mass formula is
\[
m_n=\sum_{[L]}\frac{1}{|\operatorname{Aut}(L)|}
=\frac{|B_{n/2}|}{n}\prod_{j=1}^{n/2-1}\frac{|B_{2j}|}{4j},
\]
where \(t/(e^t-1)=\sum_{j\geq0}B_jt^j/j!\).
The formula and comparison bound are recorded by
\citet[\S\S2, 6]{king2003mass}.

\begin{proposition}[Exact mass evaluation]\oursComputed{}
Evaluation in rational arithmetic gives
\[
m_{32}=40309232.93560208\ldots,\qquad
\lceil 2m_{32}\rceil=80618466.
\]
Thus the mass formula alone proves that there are at least \(80618466\)
isometry classes in this genus.  It does not give their exact number.
\end{proposition}

\begin{proof}
The program \texttt{code/check\_mass.py} evaluates the displayed historical
formula using exact fractions; its full numerator and denominator are retained
in \texttt{audit/mass\_check.json}.  Every positive-rank lattice admits the two
distinct isometries \(I\) and \(-I\), so each summand of the mass is at most
\(1/2\).  The number of summands is therefore at least \(2m_{32}\), and
integrality gives the ceiling.
\end{proof}

\paragraph{Independent arithmetic check.}
The Bernoulli numbers were computed both from
\(\sum_{j=0}^{r}\binom{r+1}{j}B_j=0\) for \(r\geq1\), with \(B_0=1\),
and from
\(B_r=\sum_{j=0}^{r}(-1)^j j!S(r,j)/(j+1)\), where \(S(r,j)\) are Stirling
numbers of the second kind.  All values through \(B_{30}\) agreed.
As normalization checks, the same computation gives
\(m_8=1/696729600\) and
\[
m_{16}=\frac{1}{2(696729600)^2}+\frac{1}{2^{15}16!},
\]
the reciprocal-automorphism sum for \(E_8\oplus E_8\) and \(D_{16}^{+}\).
All checks passed.  This recovers the notebook's rough magnitude; the later
stronger lower bound belongs to the cited literature and is not recomputed
here.

\section{Problem index}\label{problems:1971}
Earlier-year occurrence identifiers are retained as text; current-year identifiers link to their entries.

\subsection{GP-1970-26: Homotoping a distribution to a foliation}\label{prob:GP-1970-26}
For a smooth codimension-\(n\) distribution \(D\subset TM\), determine whether it can be homotoped through distributions of the same rank to an integrable distribution. Bott’s necessary condition is the vanishing over \(\mathbb R\) of every Pontryagin monomial of \(TM/D\) whose cohomological degree exceeds \(2n\).
\par\noindent\textit{Occurrences:} 1970-206.1, 1970-206.2, \hyperref[item:1971-243.17]{1971-243.17}.
\par\noindent\textit{Status (2026-11):} As of November 2026, the Bott obstruction is established. The present occurrence supplies its formal-vector-field interpretation; the checked sources do not turn this necessary condition into a general sufficiency criterion \citep{bott1970integrability,bott1976characteristic}.

\subsection{GP-1971-01: Critical exponents of finite-dimensional lp norms}\label{prob:GP-1971-01}
For fixed \(n\) and \(1\le p<\infty\), over the real and complex scalar fields separately, does every norm-one contraction on \(\ell_p^n\) that preserves norm through a bounded number of powers have spectral radius one? Equivalently, is \(q(\ell_p^n)<\infty\)?
\par\noindent\textit{Occurrences:} \hyperref[item:1971-233.5]{1971-233.5}.
\par\noindent\textit{Status (2026-11):} As of November 2026, the verified result is finiteness for every rational \(p\), already proved before the talk by \citet[Theorem 1, Corollaries 2--3]{kirzhnertabachnikov1971critical}. The exact later fate for irrational \(p\) and optimal bounds was not established in this search; absence of a verified successor is not asserted to mean current openness.

\subsection{GP-1971-02: Growth of the minimum critical exponent over all norms}\label{prob:GP-1971-02}
For a fixed scalar field let \(M(n)=\min_{\dim E=n}q(E)\). Must \(\liminf_{n\to\infty}M(n)=\infty\)?
\par\noindent\textit{Occurrences:} \hyperref[item:1971-233.6]{1971-233.6}.
\par\noindent\textit{Status (2026-11):} As of November 2026, the Hilbert bound \(M(n)\le n\) is verified, but neither the notebook’s claimed square-root subsequence construction nor the precise later resolution of the lower-growth question was identified. The similarly named 1995 exponential-number invariant does not answer this question.

\subsection{GP-1971-03: Maximizing a power norm under a spectral-radius constraint}\label{prob:GP-1971-03}
For a specified finite-dimensional normed space \(E\), determine \(C(E,r,\rho)=\max\{\|T^r\|:\|T\|\le1,\ r(T)\le\rho\}\), including sharp bounds and extremizers, for \(r\ge1\), \(0\le\rho<1\).
\par\noindent\textit{Occurrences:} \hyperref[item:1971-233.8]{1971-233.8}.
\par\noindent\textit{Status (2026-11):} As of November 2026, compactness gives existence of the maximum and the critical exponent characterizes when it is strictly below one for every fixed \(\rho<1\). This does not evaluate the maximum. The exact state of the optimization problem for arbitrary norms was not verified; specific Pták operator results must be matched separately.

\subsection{GP-1971-04: Geometric computation of spectral flow}\label{prob:GP-1971-04}
For a closed family of self-adjoint elliptic operators, compute its signed eigenvalue flow from the family symbol and geometric gluing data; in first order, identify it with the index of the suspended mapping-torus operator.
\par\noindent\textit{Occurrences:} \hyperref[item:1971-237.4]{1971-237.4}.
\par\noindent\textit{Status (2026-11):} Solved in the elliptic setting by the topological-index formula and suspension equality of \citet[Theorem 7.4, pp.~94--95]{atiyahpatodisinger1976spectral3}. The notebook supplies September 1971 circulation evidence; publication followed in 1976. As of November 2026 the required hypotheses and distinction from arbitrary nonlocal unitary gluing are standard.

\subsection{GP-1971-20: Normalized universal relations among elementary matrices}\label{prob:GP-1971-20}
Let \(R=\mathbb Z\langle X\rangle\) be the free associative unital ring, augmented to \(\mathbb Z\) by \(X\mapsto0\). Does \(\ker(K_2(R)\to K_2(\mathbb Z))\) vanish? Equivalently, is every universally valid stable elementary matrix relation whose Steinberg lift specializes to the identity under augmentation a consequence of the Steinberg relations? The normalization is necessary; arbitrary constant relations include the nonzero \(K_2(\mathbb Z)\).
\par\noindent\textit{Occurrences:} \hyperref[item:1971-247.5]{1971-247.5}, \hyperref[item:1971-253.4]{1971-253.4}, \hyperref[item:1971-256.5]{1971-256.5}.
\par\noindent\textit{Status (2026-11):} As of November 2026: resolved in the stated augmented formulation. Vaserstein announces the answer in the March 1972 notebook and the free-ring theorem in \citet[second theorem]{vaserstein1973foundations}; \citet[Theorem 28.1 and Corollary 28.1]{vaserstein1976foundations} gives its precise reduced formulation. The February notebook omits the augmentation condition, so the match is qualified rather than an assertion that every unnormalized relation is trivial. No exact earlier recurrence was found in the supplied prior index.

\subsection{GP-1971-21: Agreement of the early higher algebraic K-theory constructions}\label{prob:GP-1971-21}
For every associative unital ring \(A\) and every \(n\ge1\), do the stable simplicial constructions of Volodin and Vaserstein and the free-resolution constructions of Swan and Gersten agree naturally with \(\pi_nBGL(A)^+\)? One concrete consequence asked in the notes is the Quillen-theory suspension isomorphism \(K_{n+1}(\Sigma A)\cong K_n(A)\). Polynomial-homotopy Karoubi--Villamayor groups are excluded without additional regularity.
\par\noindent\textit{Occurrences:} \hyperref[item:1971-253.14]{1971-253.14}, \hyperref[item:1971-256.3]{1971-256.3}, \hyperref[item:1971-257.2]{1971-257.2}.
\par\noindent\textit{Status (2026-11):} As of November 2026: resolved for the specified stable constructions and unital rings. The spring 1972 seminar announcement is documented by \citet[introduction, p.97]{vaserstein1976foundations}; \citet[first theorem]{vaserstein1973foundations} publishes the comparison, with the detailed Quillen and free-resolution arguments in the 1976 paper, \S\S27--29. Record 257's catalogue attribution is conditional and does not create an additional independent priority claim. No exact earlier problem identifier was located in the supplied index.

\subsection{GP-1971-40: Unitary equivalence modulo trace class from spectral data}\label{prob:GP-1971-40}
For bounded self-adjoint operators \(A,B\) on separable Hilbert spaces, characterize from their spectral measures and multiplicities whether \(A-UBU^*\in\mathcal S_1\) for some unitary \(U\).
\par\noindent\textit{Occurrences:} \hyperref[item:1971-261.2]{1971-261.2}.
\par\noindent\textit{Status (2026-11):} As of November 2026, the subject has a dedicated later classification literature and the purely singular diagonalization case is verified indirectly. This audit has not independently checked the complete necessary-and-sufficient criterion for a prescribed pair of spectral types, and does not assert that equality of absolutely continuous parts alone suffices.

\subsection{GP-1971-41: The exceptional exponent 26 in the eta identities}\label{prob:GP-1971-41}
Explain the \(\eta^{26}\) identity in root-system terms despite the absence of a simple complex Lie algebra of dimension 26.
\par\noindent\textit{Occurrences:} \hyperref[item:1971-269.4]{1971-269.4}.
\par\noindent\textit{Status (2026-11):} As of November 2026, resolved in the affine-root framework: the dual-root specialization for \(F_4\) is described by \citet[Section I, pp. 200--201]{monastyrskii1980appendix}.

\section{Corrections to the transcription and translation}\label{corr:1971}
The evidence distinguishes a transcription correction from a mathematical emendation of wording faithfully preserved in the notebook.
\subsection*{P000358-P000364 (both; provenance): catalogue span for record 229}
\noindent\textit{Transcription reads.} P000358-P000360\par
\noindent\textit{Critical reading.} Actual record 229: P000360-P000364; P000358-P000359 are carryover from record 228.\par
\noindent\textit{Evidence.} Heading on P000360 and continuation through P000364; previous\_index identifies 1970-228.7--.13.\par
\subsection*{P000360 (both; faithful error): musical map}
\noindent\textit{Transcription reads.} alpha\_X(Y)=omega(X,Y), with partial\_q mapped to dp\par
\noindent\textit{Critical reading.} For omega=dp wedge dq under this contraction, partial\_q maps to -dp and partial\_p to dq.\par
\noindent\textit{Evidence.} Scan contains both the contraction definition and opposite coordinate signs; faithful mathematical inconsistency.\par
\subsection*{P000361 (both; faithful error): Poisson bracket}
\noindent\textit{Transcription reads.} dF=2 alpha\_xi\_F and {F,G}=xi\_F G=2 omega(xi\_F,xi\_G)\par
\noindent\textit{Critical reading.} These equations require a sign change in one convention.\par
\noindent\textit{Evidence.} Evaluate dG(xi\_F)=2 omega(xi\_G,xi\_F); normalized wedge does not remove the sign discrepancy.\par
\subsection*{P000363 (both; ambiguity): Kähler polarization and sphere remark}
\noindent\textit{Transcription reads.} F intersect conjugate F=0 implies Kähler; sphere admits no polarization\par
\noindent\textit{Critical reading.} Require positivity for Kähler; sphere obstruction concerns nonsingular real polarizations.\par
\noindent\textit{Evidence.} Scan context separates real and complex polarizations; S\textasciicircum{}2 has a complex Kähler polarization.\par
\subsection*{P000365 (both; faithful error): flag variety}
\noindent\textit{Transcription reads.} \(B\backslash G\) called projective space\par
\noindent\textit{Critical reading.} \(B\backslash G\) is a projective flag variety; generally not projective space.\par
\noindent\textit{Evidence.} Scan confirms the wording; for n>2 the dimensions and cohomology distinguish it.\par
\subsection*{P000365 (both; provenance): record 230 heading}
\noindent\textit{Transcription reads.} Separate year heading inferred from broad page range\par
\noindent\textit{Critical reading.} Retain catalogue heading; no separate handwritten heading established on P000365.\par
\noindent\textit{Evidence.} P000365 begins directly with 13/IX-71 and Gelfand on Springer.\par
\subsection*{P000366 (both; ambiguity): Frölicher–Nijenhuis sign}
\noindent\textit{Transcription reads.} d\_L=d i\_L+(-1)\textasciicircum{}p i\_L d\par
\noindent\textit{Critical reading.} Retain formula as a sign convention: d\_L=(-1)\textasciicircum{}p[i\_L,d].\par
\noindent\textit{Evidence.} Direct graded-commutator calculation; no transcription emendation justified.\par
\subsection*{P000367 (both; faithful error): critical-exponent sharpness at q=1}
\noindent\textit{Transcription reads.} exists B with ||B||=||B\textasciicircum{}(q-1)||=1 and spectral radius<1\par
\noindent\textit{Critical reading.} The predecessor witness clause applies for q>1; define q=1 by the universal implication.\par
\noindent\textit{Evidence.} For dimension1 every norm-one scalar has spectral radius1, so the written witness cannot exist.\par
\subsection*{P000367 (both; ambiguity): dimension sign}
\noindent\textit{Transcription reads.} dim E <= infinity (provisional)\par
\noindent\textit{Critical reading.} Use finite dimension for the critical-exponent discussion; exact handwritten sign remains uncertain.\par
\noindent\textit{Evidence.} The scan preserves overlapping strokes; published formulation and all examples are finite-dimensional.\par
\subsection*{P000367 (both; ambiguity): direct-sum dimension}
\noindent\textit{Transcription reads.} ell\_4\textasciicircum{}2 direct-sum ell\_8\textasciicircum{}2 has dimension12\par
\noindent\textit{Critical reading.} The literal standard notation gives dimension4; intended index convention is unresolved.\par
\noindent\textit{Evidence.} Scan has superscript2 and subscripts4,8; no unique justified replacement.\par
\subsection*{P000367 (both; faithful error): lp state of knowledge}
\noindent\textit{Transcription reads.} unknown except p=2\textasciicircum{}m\par
\noindent\textit{Critical reading.} Already known for every rational p by Kirzhner–Tabachnikov, May–June1971.\par
\noindent\textit{Evidence.} Full primary paper, Theorem1 and Corollaries2–3; catalogue talk date13September1971.\par
\subsection*{P000368 (both; ambiguity): Radon planes}
\noindent\textit{Transcription reads.} all planes\par
\noindent\textit{Critical reading.} all affine hyperplanes for the displayed formula\par
\noindent\textit{Evidence.} The scan coefficients count projective directions of codimension-one affine subspaces.\par
\subsection*{P000368-P000369 (both; provenance): source order of records 234 and 235}
\noindent\textit{Transcription reads.} 234 followed by 235 in catalogue number order\par
\noindent\textit{Critical reading.} 235 on P000368 precedes 234 on P000369.\par
\noindent\textit{Evidence.} Both headings visually inspected.\par
\subsection*{P000369 (both; faithful error): oscillatory converse}
\noindent\textit{Transcription reads.} every distribution with wavefront contained in a Lagrangian has the displayed symbol integral\par
\noindent\textit{Critical reading.} Require membership in a specified Lagrangian distribution class, not only wavefront containment.\par
\noindent\textit{Evidence.} Scan explicitly says every; Hörmander1971 includes amplitude estimates.\par
\subsection*{P000369 (both; faithful error): wavefront test}
\noindent\textit{Transcription reads.} arbitrary compactly supported test function with dpsi prescribed at one point\par
\noindent\textit{Critical reading.} Use sufficiently local cutoffs and an open conic neighbourhood of the covector.\par
\noindent\textit{Evidence.} Scan is faithful; without localization singularities elsewhere invalidate the criterion.\par
\subsection*{P000370 (both; faithful error): Godbillon–Vey deformation claim}
\noindent\textit{Transcription reads.} stable under deformations\par
\noindent\textit{Critical reading.} Invariant under foliated concordance; not under unrestricted smooth families of foliations.\par
\noindent\textit{Evidence.} Scan wording is legible; Thurston1972p513 gives continuously variable values.\par
\subsection*{P000370 (both; faithful error): cohomology of W1}
\noindent\textit{Transcription reads.} all cohomology zero except H3=R\par
\noindent\textit{Critical reading.} In positive degrees only H3=R; H0=R as well.\par
\noindent\textit{Evidence.} Trivial-coefficient degree-zero cohomology and formal-vector-field computation.\par
\subsection*{P000370 (both; faithful error): foliation transition maps}
\noindent\textit{Transcription reads.} phi\_j phi\_i\par
\noindent\textit{Critical reading.} phi\_j=psi\_ji composed with phi\_i on an overlap\par
\noindent\textit{Evidence.} The raw displayed composition has incompatible domains; scan confirms shorthand.\par
\subsection*{P000371 (both; faithful error): circle conjugation}
\noindent\textit{Transcription reads.} D\_t=-i d/dx+t and U=e\textasciicircum{}{ix}, with D\_(t+1)=U D\_t U\textasciicircum{}{-1}\par
\noindent\textit{Critical reading.} Use U=e\textasciicircum{}{-ix}, or D\_(t+1)=U\textasciicircum{}{-1}D\_tU.\par
\noindent\textit{Evidence.} Direct differentiation gives e\textasciicircum{}{ix}(-i d/dx+t)e\textasciicircum{}{-ix}=D\_t-1.\par
\subsection*{P000372 (both; faithful error): eigenvalue derivative bound}
\noindent\textit{Transcription reads.} |lambda prime|<C|lambda|\par
\noindent\textit{Critical reading.} For a suitable elliptic family use |lambda prime|<=C(1+|lambda|).\par
\noindent\textit{Evidence.} Faithful to scan; lambda(t)=t in preceding circle example contradicts the original bound at t=0.\par
\subsection*{P000372 (both; ambiguity): finite eigenvalue orbit}
\noindent\textit{Transcription reads.} integral\_0\textasciicircum{}N lambda(t)dt=0\par
\noindent\textit{Critical reading.} This equality is the condition for a periodic scalar solution, not automatic for every finite orbit.\par
\noindent\textit{Evidence.} The diagonal ODE has monodromy exp(-integral lambda).\par
\subsection*{P000373 (both; faithful error): Schubert indices}
\noindent\textit{Transcription reads.} a\_1<=...<=a\_n<=N\par
\noindent\textit{Critical reading.} 0<=a\_1<=...<=a\_n<=N-n when jumps are a\_i+i.\par
\noindent\textit{Evidence.} All jump positions must lie between1andN; scan indexing formula fixes this convention.\par
\subsection*{P000373 (both; faithful error): de Rham example for nonzero spectral flow}
\noindent\textit{Transcription reads.} d+delta on forms\par
\noindent\textit{Critical reading.} Ungraded d+delta has parity symmetry and zero loop spectral flow; a nontrivial signature example needs a specified modified operator.\par
\noindent\textit{Evidence.} Scan contains no restriction; Gamma D Gamma=-D.\par
\subsection*{P000374 (both; faithful error): basic-form condition}
\noindent\textit{Transcription reads.} vanishes when all arguments are tangent to the fibre\par
\noindent\textit{Critical reading.} Must vanish when any one argument is vertical.\par
\noindent\textit{Evidence.} The written weaker condition does not imply descent; invariant curvature construction requires horizontality.\par
\subsection*{P000374 (both; ambiguity): uncertain curvature factor}
\noindent\textit{Transcription reads.} theta\_js wedge conjugate Theta\_ks (case uncertain)\par
\noindent\textit{Critical reading.} Mathematically requires theta\_js wedge conjugate theta\_ks.\par
\noindent\textit{Evidence.} Scan does not settle case; degrees1+1=2 settle the intended formula, not the paleography.\par
\subsection*{P000375 (both; faithful error): Schur form integral}
\noindent\textit{Transcription reads.} integral over Omega(a)=1 with a\_i=N-n-lambda\_i\par
\noindent\textit{Critical reading.} Use a closed complementary Schubert cycle dual to the cohomology basis, with consistent partition conventions.\par
\noindent\textit{Evidence.} Form degree2|lambda| and stated cell dimension2(n(N-n)-|lambda|) generally disagree.\par
\subsection*{P000375 (both; faithful error): ring hypothesis}
\noindent\textit{Transcription reads.} K any ring, even noncommutative\par
\noindent\textit{Critical reading.} The displayed functor calculus requires commutative unital K, or a separately specified noncommutative replacement.\par
\noindent\textit{Evidence.} Vinogradov1972section1 explicitly imposes commutativity; the two scalar actions need it.\par
\subsection*{P000376 (both; ambiguity): jet algebra by scalar-action generators}
\noindent\textit{Transcription reads.} J\textasciicircum{}n(K) is the algebra generated by k and k+ on Diff\_n K\par
\noindent\textit{Critical reading.} Universal principal parts are (K tensor\_R K)/I\textasciicircum{}(n+1); equality with a particular operator image requires faithfulness.\par
\noindent\textit{Evidence.} Universal-property reconstruction; unclear lower index in scan remains unresolved.\par
\subsection*{P000377 (both; ambiguity): Hamilton–Jacobi repeated commutator}
\noindent\textit{Transcription reads.} delta\_f repeated m times is the Hamilton–Jacobi left side\par
\noindent\textit{Critical reading.} In the symbol convention without i, this is multiplication by m! sigma\_m(A)(df); division requires m! invertible.\par
\noindent\textit{Evidence.} Direct repeated-product-rule calculation, with scalar operator and order specified.\par
\subsection*{P000379 (both; faithful error): prime-field example}
\noindent\textit{Transcription reads.} K=Z\_2 gives D\_2 nonzero\par
\noindent\textit{Critical reading.} For K=F\_2 and the stated additive derivation functors, every D\_i for i>0 vanishes.\par
\noindent\textit{Evidence.} Visual scan and derivation of1 imply the literal claim is false; neighbouring group ring is a distinct possible intended example.\par
\subsection*{P000379 (both; faithful error): two-coordinate expansion}
\noindent\textit{Transcription reads.} b(k\_2) partial\_x k\_1 and b(k\_1) partial\_x k\_2\par
\noindent\textit{Critical reading.} The b coefficients of the vector field multiply partial\_y, not partial\_x.\par
\noindent\textit{Evidence.} Visual scan confirms the repeated x; comparison with a partial\_x+b partial\_y fixes the mathematical emendation.\par
\subsection*{P000379 (both; faithful error): variance of D\_k(M)}
\noindent\textit{Transcription reads.} D\_k(M) are k-forms with values in M\par
\noindent\textit{Critical reading.} D\_k(M) are multiderivations; the representing Lambda\_k are covariant form modules.\par
\noindent\textit{Evidence.} Definitions, duality and Vinogradov2015section5.2.\par
\subsection*{P000380 (both; ambiguity): group-ring comparison}
\noindent\textit{Transcription reads.} de Rham cohomology of group rings equals group cohomology\par
\noindent\textit{Critical reading.} An exact complex and coefficient convention are required; the ordinary exterior Kähler complex fails this assertion already for F\_2[C\_2] in degree2.\par
\noindent\textit{Evidence.} Direct ordinary-complex calculation; no exact published alternative identified.\par
\subsection*{P000381 (both; faithful error): ordinary forms over smooth functions}
\noindent\textit{Transcription reads.} all universal Lambda\_i are ordinary differential forms\par
\noindent\textit{Critical reading.} The equality with ordinary smooth forms uses the geometric module category.\par
\noindent\textit{Evidence.} Vinogradov2015Proposition7.2 and Example7.1 distinguish unrestricted algebraic modules.\par
\subsection*{P000383 (both; faithful error): resolution terminology}
\noindent\textit{Transcription reads.} Spencer resolution for an arbitrary ring\par
\noindent\textit{Critical reading.} Spencer complex; exactness requires further hypotheses.\par
\noindent\textit{Evidence.} The scan later says complex; Vinogradov1972section3 explicitly denies general exactness.\par
\subsection*{P000384 (both; faithful error): Lemma1 composition type}
\noindent\textit{Transcription reads.} D composed with f: M\_1 to M\_2\par
\noindent\textit{Critical reading.} D composed with f: M\_1 to N.\par
\noindent\textit{Evidence.} Visual scan confirms M\_2; domain and codomain of the displayed maps force N.\par
\subsection*{P000384 (both; ambiguity): Lemma2 corollary}
\noindent\textit{Transcription reads.} A-A(1) is a derivation\par
\noindent\textit{Critical reading.} Subtract the map k maps to kA(1).\par
\noindent\textit{Evidence.} A(1) is a module element, so the unqualified subtraction is shorthand; constant-map subtraction is not additive.\par
\subsection*{P000384-P000385 (both; ambiguity): commutator sign}
\noindent\textit{Transcription reads.} delta\_k A=kA-Ak\par
\noindent\textit{Critical reading.} This is the negative of P376 delta\_k; either convention gives the same order filtration.\par
\noindent\textit{Evidence.} Visual scan and product-rule calculation; convention change is harmless here.\par
\subsection*{P000385 (both; ambiguity): Lemma4 order}
\noindent\textit{Transcription reads.} evaluation is of order n\par
\noindent\textit{Critical reading.} of order at most n.\par
\noindent\textit{Evidence.} Zero modules and order collapse show exact equality is not automatic.\par
\subsection*{P000385 (both; faithful error): Lemma5 proof1}
\noindent\textit{Transcription reads.} restriction to a submodule K\par
\noindent\textit{Critical reading.} Precomposition with the homomorphism K to M, k maps to km; it need not be injective.\par
\noindent\textit{Evidence.} The displayed map has kernel the annihilator of m.\par
\subsection*{P000387 (both; ambiguity): finite-dimensional integration}
\noindent\textit{Transcription reads.} one can pass to a group when the Lie algebra is finite-dimensional\par
\noindent\textit{Critical reading.} Local integration; global integration requires additional completeness assumptions.\par
\noindent\textit{Evidence.} The incomplete vector field on an interval disproves the unrestricted global reading; Bernstein--Rosenfeld introduction explicitly says perhaps only locally.\par
\subsection*{P000388 (both; faithful error): equivalence of formal coordinates}
\noindent\textit{Transcription reads.} all derivatives of the transition map are zero\par
\noindent\textit{Critical reading.} The transition has identity infinite jet; equivalently the coordinate difference is flat.\par
\noindent\textit{Evidence.} An invertible coordinate transition has invertible first derivative, so it cannot have every derivative zero; visually confirmed on the scan.\par
\subsection*{P000389 (both; faithful error): Chevalley--Eilenberg exponent}
\noindent\textit{Transcription reads.} \(s+t=1\)\par
\noindent\textit{Critical reading.} A convention-consistent exponent, for example \(s+t\) with the left Maurer--Cartan convention.\par
\noindent\textit{Evidence.} The scan has two horizontal strokes. The equality is faithfully transcribed but does not define a sign; an intended opposite-action exponent with plus one is possible.\par
\subsection*{P000389-P000390 (both; ambiguity): inverse-limit manifold definition}
\noindent\textit{Transcription reads.} arbitrary onto finite-dimensional inverse system; global union of function rings\par
\noindent\textit{Critical reading.} Use the jet-tower category with compatible affine-bundle stages and locally finite-coordinate smooth functions.\par
\noindent\textit{Evidence.} Bernstein--Rosenfeld 1973, Definition 1.2 and Example 1.1, give the missing hypotheses and explicitly exclude arbitrary inverse systems.\par
\subsection*{P000390 (both; faithful error): projective-limit Lie algebra}
\noindent\textit{Transcription reads.} projective limit of finite-dimensional Lie algebras\par
\noindent\textit{Critical reading.} Inverse limit as a topological vector space, with continuous Lie bracket.\par
\noindent\textit{Evidence.} Constant formal vector fields lower vanishing order, so the finite coefficient truncations are not Lie algebra quotients.\par
\subsection*{P000391 (both; faithful error): relative-complex cohomology}
\noindent\textit{Transcription reads.} cohomology of the quotient in the classical case\par
\noindent\textit{Critical reading.} A characteristic map to de Rham cohomology, not a general isomorphism.\par
\noindent\textit{Evidence.} The additive real line has one-dimensional Lie algebra cohomology in degree one and zero de Rham cohomology there.\par
\subsection*{P000391-P000392 (both; ambiguity): independence of point in fibre}
\noindent\textit{Transcription reads.} independent of the chosen point\par
\noindent\textit{Critical reading.} Connected fibres suffice; with disconnected structure group require invariance under all components.\par
\noindent\textit{Evidence.} Bernstein--Rosenfeld 1973, Proposition 3.2 and Remark 3.2.\par
\subsection*{P000392 (both; faithful error): degrees of characteristic generators}
\noindent\textit{Transcription reads.} all \(\psi_i\) have degree \(2n\)\par
\noindent\textit{Critical reading.} \(\deg\psi_i=2i\); only products of degree greater than \(2n\) are killed.\par
\noindent\textit{Evidence.} Scan inspected; Bernstein--Rosenfeld Theorem 10.2 and Bott 1976 section 2 give the graded ring.\par
\subsection*{P000392 (both; ambiguity): marginal map}
\noindent\textit{Transcription reads.} \(M\to S(M)\)\par
\noindent\textit{Critical reading.} If a section is meant, use \(M\to S(M)/GL_n\), unless a global framing is assumed.\par
\noindent\textit{Evidence.} The preceding marginal word is illegible; a full formal-frame section induces a global tangent frame. Do not silently identify the marginal word.\par
\subsection*{P000394 (both; ambiguity): Bott vanishing boundary}
\noindent\textit{Transcription reads.} prose appears to say degree \(\geq2n\)\par
\noindent\textit{Critical reading.} strictly greater than \(2n\)\par
\noindent\textit{Evidence.} The preceding formula clearly has a strict inequality; the prose comparison stroke is ambiguous on the scan. The verified theorem uses the strict inequality.\par
\subsection*{P000394 (both; provenance): Gelfand--Fuks reference date}
\noindent\textit{Transcription reads.} Izvestiya 1971\par
\noindent\textit{Critical reading.} Russian original 1970; English printing bears 1971 copyright.\par
\noindent\textit{Evidence.} Primary original metadata and full English PDF; received 4 November 1969.\par
\subsection*{P000394 (both; provenance): end of record 243}
\noindent\textit{Transcription reads.} whole page attached to Gelfand formal-geometry record\par
\noindent\textit{Critical reading.} End record 243 before the separate Arnold heading.\par
\noindent\textit{Evidence.} The heading is explicit in the Russian transcription and visually checked on the scan; the Arnold discussion belongs to the next record.\par
\subsection*{P000396 (both; faithful error): last sentence, maximal form parameter}
\noindent\textit{Transcription reads.} \(\Lambda_{\max}\): the quadratic form is recovered from \(|u|_q\)\par
\noindent\textit{Critical reading.} Recovered from the hermitianized pairing \(b_q\), not from the diagonal refinement alone.\par
\noindent\textit{Evidence.} Scan inspected: the diagonal length symbol is present. For trivial involution, epsilon=-1 and Lambda=A, all lengths are zero but alternating pairings need not be zero.\par
\subsection*{P000397 (both; faithful error): hyperbolic scalar-product matrix}
\noindent\textit{Transcription reads.} \(\left(\begin{smallmatrix}0&1\\-1&0\end{smallmatrix}\right)\)\par
\noindent\textit{Critical reading.} For general symmetry use \(\left(\begin{smallmatrix}0&1\\\varepsilon&0\end{smallmatrix}\right)\), up to block order.\par
\noindent\textit{Evidence.} The preceding page fixes arbitrary epsilon. The displayed alternating matrix is only the epsilon=-1 specialization.\par
\subsection*{P000398 (both; faithful error): field example}
\noindent\textit{Transcription reads.} \(K_0O(A)=\mathbb Z\) after the field cases epsilon=+/-1\par
\noindent\textit{Critical reading.} Valid for the alternating case; symmetric forms require the field and its Witt theory.\par
\noindent\textit{Evidence.} Scan confirms the blanket Z. Over the real field positive and negative indices give Grothendieck-Witt group Z\textasciicircum{}2.\par
\subsection*{P000398 (both; faithful error): integral classification by dimension and signature}
\noindent\textit{Transcription reads.} Dimension and signature classify forms over Z.\par
\noindent\textit{Critical reading.} Separate stable classes from isometry classes and specify parity and definiteness.\par
\noindent\textit{Evidence.} Scan inspected. Positive definite even unimodular forms already supply distinct isometry classes with the same rank and signature.\par
\subsection*{P000400 (both; faithful error): four-factor elementary-matrix display}
\noindent\textit{Transcription reads.} Four factors with adjacent roots are equated to \(1\); first two signs overwritten.\par
\noindent\textit{Critical reading.} Retain sign ambiguity; with inverse signs the adjacent-root commutator is \(1+ab e_{ik}\), not \(1\).\par
\noindent\textit{Evidence.} Scan inspected. The final identity is visible; multiplication for distinct i,j,k gives the nonzero ik term.\par
\subsection*{P000400 (both; ambiguity): last field example}
\noindent\textit{Transcription reads.} \(K_1(\mathbb C)=\mathbb C\) with unresolved upper mark\par
\noindent\textit{Critical reading.} \(K_1(\mathbb C)=\mathbb C^\times\).\par
\noindent\textit{Evidence.} Scan inspected: upper mark remains ambiguous; determinant determines the required multiplicative group.\par
\subsection*{P000402 (both; faithful error): inverse construction before the finite-matrix argument}
\noindent\textit{Transcription reads.} Cokernel class alone represents the Laurent boundary.\par
\noindent\textit{Critical reading.} Use \([\operatorname{coker}(t^ng)]-nr[A]\).\par
\noindent\textit{Evidence.} Dependence on the clearing power is immediate on the identity matrix; the normalization is explicitly supplied in the next Kirillov occurrence.\par
\subsection*{P000403 (both; faithful error): last line of the truncated elementary relation}
\noindent\textit{Transcription reads.} L is over the finite-matrix ideal.\par
\noindent\textit{Critical reading.} \(L-1\) is in the finite-matrix ideal; \(L=1+F\) represents a stable automorphism.\par
\noindent\textit{Evidence.} Scan inspected: the conclusion follows from Psi\_2(L)=1, which means L-1 lies in the kernel.\par
\subsection*{P000404 (both; faithful error): third bilinear relation and later Weil-operator condition}
\noindent\textit{Transcription reads.} With \(\Psi=(-1)^{n(n+1)/2}\int\), the notes write \(i^{p-q}\Psi(x,\bar x)>0\).\par
\noindent\textit{Critical reading.} Keep Psi and replace by \(i^{q-p}\Psi(x,\bar x)>0\), or change the initial exponent to \(n(n-1)/2\).\par
\noindent\textit{Evidence.} Scan inspected: both source exponents are clear. The later Psi(x,C bar x)>0 agrees with the first correction; the two conventions differ by (-1)\textasciicircum{}n.\par
\subsection*{P000405 (both; provenance): Kirillov date heading}
\noindent\textit{Transcription reads.} Catalogue metadata: 28 February 1972.\par
\noindent\textit{Critical reading.} Notebook heading: 27/II. Catalogue date retained with this qualification.\par
\noindent\textit{Evidence.} Scan inspected: the second digit has the angular form of 7. No silent harmonization to the expected Monday date.\par
\subsection*{P000406 (both; faithful error): Gelfand scalar-extension diagram}
\noindent\textit{Transcription reads.} \(B\otimes_A(N\oplus M)=A^k\)\par
\noindent\textit{Critical reading.} \(B\otimes_A(N\oplus M)=B^k\).\par
\noindent\textit{Evidence.} The same paragraph describes a B-linear automorphism of B\textasciicircum{}k; extending N+M=A\textasciicircum{}k changes the free scalar ring to B.\par
\subsection*{P000409 (both; faithful error): initial Hodge filtration}
\noindent\textit{Transcription reads.} \(F^p=\bigoplus_{p\prime>p}H^{p\prime,q}\)\par
\noindent\textit{Critical reading.} In the decreasing convention use \(p\prime\ge p\).\par
\noindent\textit{Evidence.} Scan inspected: the strict sign is present. A shifted convention would require corresponding changes to all subsequent indices.\par
\subsection*{P000409 (both; ambiguity): smoothness gloss}
\noindent\textit{Transcription reads.} Smooth means local triviality in the analytic sense.\par
\noindent\textit{Critical reading.} Local projection near each point is valid; local triviality of the entire family is differentiable, not generally holomorphic.\par
\noindent\textit{Evidence.} Scan inspected. The wording does not distinguish the two scopes of local triviality.\par
\subsection*{P000410 (both; faithful error): last step and exact sequence of the filtration}
\noindent\textit{Transcription reads.} \(F^n=R^nf_*\Omega^0\); \(0\to F^{p-1}\to F^p\to E_\infty^{p,n-p}\to0\)\par
\noindent\textit{Critical reading.} \(F^n=f_*\Omega^n\); \(0\to F^{p+1}\to F^p\to R^{n-p}f_*\Omega^p\to0\).\par
\noindent\textit{Evidence.} Scan inspected. This is the consistent decreasing filtration associated with the displayed E\_1\textasciicircum{}{p,q}=R\textasciicircum{}q f\_*Omega\textasciicircum{}p.\par
\subsection*{P000411 (both; faithful error): description of K\_0 as exponentially decaying solutions}
\noindent\textit{Transcription reads.} K\_0 consists of solutions decaying exponentially in time near the boundary.\par
\noindent\textit{Critical reading.} Distinguish local or compressed evolution from conserved full energy; exponential decay needs additional hypotheses.\par
\noindent\textit{Evidence.} Scan inspected. The same page explicitly conserves the global energy, precluding nonzero decay in that norm.\par
\subsection*{P000412 (both; faithful error): Birman functional-calculus trace-ideal line}
\noindent\textit{Transcription reads.} \(\varphi(H)-\varphi(H_0)\) nuclear without restrictions on \(\varphi\).\par
\noindent\textit{Critical reading.} Specify an admissible invariance-principle function or use the resolvent trace-class criterion.\par
\noindent\textit{Evidence.} A constant function satisfies the displayed condition for arbitrary pairs, so omitted hypotheses cannot be dropped.\par
\subsection*{P000412 (both; faithful error): two-space wave-operator condition 1}
\noindent\textit{Transcription reads.} \(D(H)=D(H_0)\)\par
\noindent\textit{Critical reading.} For different spaces require \(J D(H_0)\subseteq D(H)\), with the localized commutator condition.\par
\noindent\textit{Evidence.} Scan inspected: the equality is retained after introducing different spaces. Belopolskii--Birman Theorem 5.1 formulates domain compatibility through the identification.\par
\subsection*{P000413 (both; ambiguity): derived-functor author name}
\noindent\textit{Transcription reads.} Swan; Gershtein\par
\noindent\textit{Critical reading.} Swan; Gersten (S. M. Gersten).\par
\noindent\textit{Evidence.} The named construction matches Gersten’s 1971 higher K-functor papers; the phonetic notebook form is normalized in commentary.\par
\subsection*{P000413 (both; faithful error): free-presentation recursion}
\noindent\textit{Transcription reads.} \(K_{n+1}(A)=K_n(A\prime\prime)\), where \(A\prime\prime\) was called the presentation kernel.\par
\noindent\textit{Critical reading.} The published dimension shift uses \(D(A)=F(A)\times_A F(A)\), or relative groups, not an unqualified absolute theory of the kernel.\par
\noindent\textit{Evidence.} Vaserstein 1976, §29, Theorem 29.1 gives the precise doubled fibre-product construction. General excision for arbitrary ideals must not be assumed.\par
\subsection*{P000414 (both; ambiguity): relative pullback diagram}
\noindent\textit{Transcription reads.} The source of the inclusion into the cone is labelled mu B (possibly mB or MB).\par
\noindent\textit{Critical reading.} Leave the letter uncertain; it cannot be the previously defined quotient mu B with its evident inclusion.\par
\noindent\textit{Evidence.} Scan inspected. The published relative-ring model is A times\_B F(B); the notebook diagram and its shifted indexing are not identified silently.\par
\subsection*{P000415 (both; faithful error): Laurent splitting and relative-group caption}
\noindent\textit{Transcription reads.} K\_2 relative is trivial.\par
\noindent\textit{Critical reading.} The evaluation map is split, but its relative kernel contains the K\_1 and Nil contributions.\par
\noindent\textit{Evidence.} Scan inspected. The same displayed Laurent decomposition already contains a K\_1 summand.\par
\subsection*{P000415 (both; ambiguity): Whitehead matrix factorization, two blank entries}
\noindent\textit{Transcription reads.} Third upper-right and fourth lower-left entries are dots.\par
\noindent\textit{Critical reading.} They are a\textasciicircum{}{-1}-1 and -a, respectively.\par
\noindent\textit{Evidence.} Direct multiplication restores diag(a,a\textasciicircum{}{-1}) for any invertible a; this is a mathematical completion, not a claim that the notebook wrote these entries.\par
\subsection*{P000415 (both; faithful error): general plus-construction description}
\noindent\textit{Transcription reads.} For any group with perfect commutator subgroup, cell attachments turn BG into an H-space.\par
\noindent\textit{Critical reading.} The plus construction exists; an H-space multiplication needs further structure, available for stable GL through block sum.\par
\noindent\textit{Evidence.} Scan inspected. The line explicitly says any group and then H-space. Quillen ICM §3 supplies the stable-linear categorical structure.\par
\subsection*{P000416 (both; faithful error): quotient for compression}
\noindent\textit{Transcription reads.} \(\ell A/\mu A\)\par
\noindent\textit{Critical reading.} With P000414 notation, \(\ell A/MA=\mu A\).\par
\noindent\textit{Evidence.} The denominator must be the finite-matrix ideal; mu A already denotes the quotient. The displayed letter is preserved in the diplomatic source.\par
\subsection*{P000416 (both; ambiguity): residue commutator and rank caption}
\noindent\textit{Transcription reads.} Trace of [p\_? f,p\_? g], labelled finite rank.\par
\noindent\textit{Critical reading.} Specify the projections and the space on which the trace is taken; on A use [P M\_f|\_A,P M\_g|\_A]. On K the general construction uses the finite-potent trace.\par
\noindent\textit{Evidence.} Scan inspected. Subscripts are not sufficiently certain to emend diplomatically. Tate 1968, §§1–2 distinguishes the relevant trace notions.\par
\subsection*{P000417 (both; ambiguity): monotonicity condition in the simplex category}
\noindent\textit{Transcription reads.} \(\varphi(1)\le\cdots\le\varphi(k)\) for \([k]=\{0,\ldots,k\}\)\par
\noindent\textit{Critical reading.} Include \(\varphi(0)\le\varphi(1)\).\par
\noindent\textit{Evidence.} The displayed objects include zero; Segal 1968 §1 uses all order-preserving maps. The initial digit is overwritten and remains qualified in the diplomatic source.\par
\subsection*{P000419 (both; provenance): Boardman–Vogt publication label}
\noindent\textit{Transcription reads.} Boardman–Vogt 1968 unpublished.\par
\noindent\textit{Critical reading.} A recognition announcement was published in November 1968; full details may be the intended unpublished reference.\par
\noindent\textit{Evidence.} Opened Boardman–Vogt 1968 primary article, pp.1117–1122, communicated 24 May 1968; its first page promises full details elsewhere.\par
\subsection*{P000419 (both; ambiguity): cohomology grading}
\noindent\textit{Transcription reads.} \(K_{\mathcal C}^{-k}(Y)=\lim_i[Y,\Omega^iB^{k+i}\mathcal C]\)\par
\noindent\textit{Critical reading.} Fix a spectrum-level convention before using the sign: standard \(E^j(Y)=\varinjlim_i[\Sigma^iY,E_{i+j}]\).\par
\noindent\textit{Evidence.} Scan inspected. B C is initially the nerve, so both the level shift and the opposite cohomological sign need declaration.\par
\subsection*{P000419 (both; ambiguity): finite-set category}
\noindent\textit{Transcription reads.} Category of finite sets, with BC equal to the disjoint union of B Sigma\_n.\par
\noindent\textit{Critical reading.} Specify bijections as morphisms.\par
\noindent\textit{Evidence.} The displayed classifying spaces force the groupoid interpretation; all set maps would give a category with a terminal object.\par
\subsection*{P000421 (both; ambiguity): K\_i(GL) and cone contractibility}
\noindent\textit{Transcription reads.} \(K_iA=\pi_{i-1}GL(A)\); \(GL(\ell A)\) is contractible.\par
\noindent\textit{Critical reading.} Homotopy is taken in the specified simplicial realization, not in the discrete group.\par
\noindent\textit{Evidence.} The page explicitly introduces a simplicial scheme, and Vaserstein 1976 §§7–15 makes that structure essential.\par
\subsection*{P000421 (both; faithful error): free-kernel degree shift}
\noindent\textit{Transcription reads.} \(K_i(A)=K_{i-1}(L\prime)\) for a free presentation.\par
\noindent\textit{Critical reading.} Use the relative homotopy fibre; the published recursion uses the doubled fibre product F(A) ×\_A F(A).\par
\noindent\textit{Evidence.} Scan inspected; Vaserstein 1976 §29, Theorem 29.1. Free acyclicity does not supply unrestricted excision.\par
\subsection*{P000422 (both; faithful error): Karoubi comparison hypothesis}
\noindent\textit{Transcription reads.} No zero divisors — Karoubi also coincide.\par
\noindent\textit{Critical reading.} Regular noetherian is a sufficient hypothesis; integral domain alone is false.\par
\noindent\textit{Evidence.} The scan plainly says no zero divisors. Hesselholt–Nikolaus Theorem A gives nonzero K\_2 for a graded cuspidal domain, while polynomial-homotopy K\_2 contracts to that of its finite base field.\par
\subsection*{P000422-P000423 (both; provenance): speaker and date at the Quillen comparison block}
\noindent\textit{Transcription reads.} Catalogue attributes a separate talk to Kazhdan on 20 March.\par
\noindent\textit{Critical reading.} Preserve the catalogue metadata conditionally: this block has no separate dated Kazhdan heading; it may continue Vaserstein.\par
\noindent\textit{Evidence.} P422 and P423 scans inspected. Gelfand 20/III is explicit below the P423 divider; Kazhdan Proof begins only on P424.\par
\subsection*{P000423 (both; ambiguity): finite-field odd degree and field exponent}
\noindent\textit{Transcription reads.} Overwritten \(K_{2i\pm1}(F_q)\), with right side provisionally \(F_{q^i}^*\).\par
\noindent\textit{Critical reading.} The verified theorem is \(K_{2i-1}(F_q)\cong\mathbb Z/(q^i-1)\); equivalently \(K_{2i+1}(F_q)\cong\mathbb Z/(q^{i+1}-1)\).\par
\noindent\textit{Evidence.} Scan inspected; the overwritten sign remains ambiguous. Quillen ICM §3 Theorems 1–2 determine the mathematical indexing.\par
\subsection*{P000423 (both; ambiguity): homology and K\_0 components}
\noindent\textit{Transcription reads.} \(H_iGL(\mu A)=H_iGL(MA)\).\par
\noindent\textit{Critical reading.} On the left use the identity component of the simplicial realization; on the right use stable group homology.\par
\noindent\textit{Evidence.} The scan has an explicit identity-component annotation in the right margin. Vaserstein 1973 and Weibel IV Definition 1.1.1 separate K\_0 components.\par
\subsection*{P000423-P000424 (both; provenance): chronology of the finite-field result}
\noindent\textit{Transcription reads.} The prominent full publication is Quillen 1972.\par
\noindent\textit{Critical reading.} The theorem was announced at the 1970 ICM and printed in its 1971 proceedings, before Gelfand’s March 1972 report.\par
\noindent\textit{Evidence.} Quillen Cohomology of groups, §3 Theorems 1–2, pp.50–51, primary text opened; Annals metadata independently checked.\par
\subsection*{P000424 (both; provenance): record 260 supplied page interval}
\noindent\textit{Transcription reads.} P000423-P000427\par
\noindent\textit{Critical reading.} P000424\par
\noindent\textit{Evidence.} The Fuchs heading and both assertions are confined to P000424, between Kazhdan passages.\par
\subsection*{P000424-P000425 (both; faithful error): Quillenization and H-space}
\noindent\textit{Transcription reads.} K\_i(G\_n)=pi\_i(B\_n); B\_n is an H-space\par
\noindent\textit{Critical reading.} Use the stable group GL(F\_q) and its plus construction; fixed rank requires a stability range.\par
\noindent\textit{Evidence.} The notes use finite n throughout and then say large n. The stable construction is necessary for all degrees.\par
\subsection*{P000424-P000426 (both; provenance): record 259 supplied page interval}
\noindent\textit{Transcription reads.} P000423-P000425\par
\noindent\textit{Critical reading.} P000424-P000426\par
\noindent\textit{Evidence.} Kazhdan starts after the Gelfand record 258 continuation on P000424 and resumes after Fuchs on P000425; the Lang-map discussion ends on P000426.\par
\subsection*{P000425 (both; faithful error): fixed subgroup represented by fibre}
\noindent\textit{Transcription reads.} [K(X)]\textasciicircum{}k=[X,E psi\textasciicircum{}k]\par
\noindent\textit{Critical reading.} Maps into the homotopy fibre also contain a null-homotopy; use the long exact sequence.\par
\noindent\textit{Evidence.} The scanned formula is literal; X=S\textasciicircum{}1 and k>2 gives a nonzero Z/(k-1) on the right and zero on the left.\par
\subsection*{P000426 (both; faithful error): unipotent value}
\noindent\textit{Transcription reads.} phi\_mu=1 on unipotents\par
\noindent\textit{Critical reading.} The displayed sum is n on rank-n unipotents; the reduced value is zero.\par
\noindent\textit{Evidence.} All eigenvalues of a unipotent matrix are 1, and mu(1)=1.\par
\subsection*{P000427 (both; provenance): record 261 supplied page interval}
\noindent\textit{Transcription reads.} P000425-P000428\par
\noindent\textit{Critical reading.} P000427\par
\noindent\textit{Evidence.} Birman begins at the top of P000427; the 27 March Kazhdan heading starts a different record below.\par
\subsection*{P000427-P000428 (both; faithful error): Artin–Mazur complex and completion}
\noindent\textit{Transcription reads.} An ordinary complex C\_X; the compatible functor is representable\par
\noindent\textit{Critical reading.} Retain the étale homotopy pro-object and profinite comparison; specify the pro-category.\par
\noindent\textit{Evidence.} The notes omit the ambient pro-homotopy category required by Artin–Mazur.\par
\subsection*{P000427-P000428 (both; provenance): record 262 supplied page interval}
\noindent\textit{Transcription reads.} P000426-P000429\par
\noindent\textit{Critical reading.} P000427-P000428\par
\noindent\textit{Evidence.} Kazhdan begins at the bottom of P000427; completion continues only to the first separator on P000428.\par
\subsection*{P000428 (both; provenance): record 264 supplied page interval}
\noindent\textit{Transcription reads.} P000427-P000431\par
\noindent\textit{Critical reading.} P000428\par
\noindent\textit{Evidence.} Gelfand’s reading plan lies between Kazhdan’s completion and Birman’s continuation on P000428; it precedes record 263 in source order.\par
\subsection*{P000428-P000430 (both; provenance): record 263 supplied page interval}
\noindent\textit{Transcription reads.} P000427-P000429\par
\noindent\textit{Critical reading.} P000428-P000430\par
\noindent\textit{Evidence.} Birman begins below Gelfand’s plan on P000428 and ends before the 3 April Geronimus heading on P000430.\par
\subsection*{P000430-P000436 (both; provenance): record 265 supplied page interval}
\noindent\textit{Transcription reads.} P000429-P000438\par
\noindent\textit{Critical reading.} P000430-P000436\par
\noindent\textit{Evidence.} The dated 3 April heading begins on P000430; P000437 explicitly starts 10 April Geronimus.\par
\subsection*{P000431 (both; faithful error): sheaf equalizer}
\noindent\textit{Transcription reads.} Direct sums over arbitrary covering families\par
\noindent\textit{Critical reading.} Products over the covering family and the pairwise fibre products\par
\noindent\textit{Evidence.} The scan shows direct-sum signs. For infinite covers the sheaf axiom uses products.\par
\subsection*{P000433 (both; faithful error): continuity of site functor}
\noindent\textit{Transcription reads.} Cover-preserving implies precomposition preserves sheaves\par
\noindent\textit{Critical reading.} Add preservation of the relevant fibre products or the general continuity condition\par
\noindent\textit{Evidence.} Applying the sheaf equalizer requires identifying the overlaps after applying the functor.\par
\subsection*{P000436 (both; faithful error): Quillen cohomology formula}
\noindent\textit{Transcription reads.} All primes and all q give n polynomial and n exterior generators\par
\noindent\textit{Critical reading.} For odd ell use indices divisible by ord\_ell(q); the displayed indices require ell dividing q-1\par
\noindent\textit{Evidence.} The scan supplies no condition on ell. Prime 2 also has separate cases.\par
\subsection*{P000436 (both; faithful error): degree of c\_i}
\noindent\textit{Transcription reads.} c\_i in H\textasciicircum{}{2i}(G,F\_q)\par
\noindent\textit{Critical reading.} c\_i in H\textasciicircum{}{2i}(G,F\_ell)\par
\noindent\textit{Evidence.} The surrounding formula and companion exterior generators use F\_ell; mathematical coefficient correction.\par
\subsection*{P000436 (both; faithful error): unitary embedding}
\noindent\textit{Transcription reads.} faithful irreducible representation\par
\noindent\textit{Critical reading.} faithful finite-dimensional unitary representation\par
\noindent\textit{Evidence.} A finite noncyclic abelian group has no faithful irreducible representation.\par
\subsection*{P000437 (both; faithful error): maps of arbitrary degree}
\noindent\textit{Transcription reads.} f-g=d alpha+alpha d-prime without graded signs\par
\noindent\textit{Critical reading.} For degree zero use d\_D h+h d\_C; for degree tau use the graded Hom differential\par
\noindent\textit{Evidence.} The displayed arbitrary degree requires a consistent signed and indexed convention.\par
\subsection*{P000437-P000438 (both; provenance): record 266 supplied page interval}
\noindent\textit{Transcription reads.} P000436-P000440\par
\noindent\textit{Critical reading.} P000437-P000438\par
\noindent\textit{Evidence.} The 10 April Geronimus lecture occupies P000437-P000438. Its short reprise on P000439 is crossed out and is not a new live assertion.\par
\subsection*{P000438 (both; faithful error): roof construction limit}
\noindent\textit{Transcription reads.} projective limit\par
\noindent\textit{Critical reading.} Filtered colimit of Hom sets over quasi-isomorphic replacements, modulo roof equivalence\par
\noindent\textit{Evidence.} The scan visibly says projective limit; derived localization formally adds inverses and uses the calculus of fractions.\par
\subsection*{P000439 (both; provenance): record 267 supplied page interval}
\noindent\textit{Transcription reads.} P000438-P000440\par
\noindent\textit{Critical reading.} P000439\par
\noindent\textit{Evidence.} Kazhdan’s injective-complex intervention is at the top of P000439, before the crossed-out Geronimus reprise and Morava.\par
\subsection*{P000439-P000445 (both; provenance): record 268 supplied page interval}
\noindent\textit{Transcription reads.} P000438-P000446\par
\noindent\textit{Critical reading.} P000439-P000445\par
\noindent\textit{Evidence.} Morava starts below the crossed-out passage on P000439 and ends with Kazhdan’s differential remark at the top of P000445; the 14 April Gelfand heading follows.\par
\subsection*{P000440 (both; faithful error): complex representative}
\noindent\textit{Transcription reads.} Every weakly complex manifold is cobordant to a complex manifold\par
\noindent\textit{Critical reading.} Use the verified statement that complex bordism is generated integrally by nonsingular projective varieties; unrestricted actual-representative statement fails in dimension zero\par
\noindent\textit{Evidence.} The negatively stably complex point is -1 in MU\_0, unlike any complex zero-manifold.\par
\subsection*{P000441-P000444 (both; faithful error): coefficient ring}
\noindent\textit{Transcription reads.} Z[t\_0,t\_1,...], degree t\_i=2i\par
\noindent\textit{Critical reading.} Z[t\_1,t\_2,...]; t\_0, if used for the point, is the unit\par
\noindent\textit{Evidence.} MU\_0=Z, so a free degree-zero polynomial generator cannot occur.\par
\subsection*{P000442 (both; faithful error): dense open fibres}
\noindent\textit{Transcription reads.} Every fibre over D(t\_1) equals integral K\textasciicircum{}0(X) plus K\textasciicircum{}1(X)\par
\noindent\textit{Critical reading.} State the Conner–Floyd base change through the multiplicative formal group; geometric fibres use their residue fields\par
\noindent\textit{Evidence.} Already at a point the geometric fibre is a field, not Z. Inverting t\_1 alone does not give the K-theory coefficient map.\par
\subsection*{P000443 (both; faithful error): stabilizer}
\noindent\textit{Transcription reads.} Units of a p-adic division algebra\par
\noindent\textit{Critical reading.} Units of its maximal order for automorphisms; the whole division-algebra unit group is for quasi-isogenies\par
\noindent\textit{Evidence.} Morava 1972, Theorem 2 identifies the stabilizer with E\_n units.\par
\subsection*{P000445 (both; faithful error): differential classification}
\noindent\textit{Transcription reads.} Characteristic zero ring\par
\noindent\textit{Critical reading.} Q-algebra for unrestricted termwise integration and arbitrary differential coefficients\par
\noindent\textit{Evidence.} The logarithm construction divides by positive integers.\par
\subsection*{P000445 (both; faithful error): eta-product generating series}
\noindent\textit{Transcription reads.} Product (1-x\textasciicircum{}m)\textasciicircum{}24 equals sum from n=1 of tau(n)x\textasciicircum{}n\par
\noindent\textit{Critical reading.} Insert x before the product; use eta(tau)\textasciicircum{}24 with x=exp(2 pi i tau)\par
\noindent\textit{Evidence.} Scan confirms the omitted factor. Dyson 1972, equation (1) includes x; the unshifted product has constant term 1.\par
\subsection*{P000445 (both; faithful error): root product}
\noindent\textit{Transcription reads.} all simple roots (positive) of SL(5)\par
\noindent\textit{Critical reading.} all positive roots of A4\par
\noindent\textit{Evidence.} The written numerator has 10 factors; A4 has 10 positive roots but 4 simple roots.\par
\subsection*{P000445-P000446 (both; provenance): record 269 supplied page interval}
\noindent\textit{Transcription reads.} P000444-P000447\par
\noindent\textit{Critical reading.} P000445-P000446\par
\noindent\textit{Evidence.} Gelfand’s 14 April heading starts on P000445; the contraction line ends before Morava’s heading on P000446.\par
\subsection*{P000446 (both; faithful error): fibre point}
\noindent\textit{Transcription reads.} Description of M\_x, x in M\par
\noindent\textit{Critical reading.} The point x lies in Spec A, not in the module M\par
\noindent\textit{Evidence.} The next page describes the ideal corresponding to the point.\par
\subsection*{P000447 (both; faithful error): formal substitutions}
\noindent\textit{Transcription reads.} a\_1=1 and a\_1 an arbitrary unit treated together\par
\noindent\textit{Critical reading.} Distinguish strict coordinate changes from all coordinate changes\par
\noindent\textit{Evidence.} Scan shows both conditions; over Z the units are +1 and -1.\par
\subsection*{P000447 (en; translation): isomorphism margin}
\noindent\textit{Transcription reads.} These are same-form groups\par
\noindent\textit{Critical reading.} These are isomorphic formal groups\par
\noindent\textit{Evidence.} The displayed conjugation formula defines isomorphism; the Russian marginal abbreviation is difficult but agrees with this mathematical reading.\par
\subsection*{P000447 (both; faithful error): maximal evaluation ideal over Z}
\noindent\textit{Transcription reads.} J=(t\_i-a\_i) is maximal and M/MJ is a vector space\par
\noindent\textit{Critical reading.} For integers a\_i the quotient is Z; impose a residue characteristic for a maximal evaluation ideal, or base change to a field\par
\noindent\textit{Evidence.} The manuscript itself qualifies vector space as A/J-module.\par
\subsection*{P000448 (both; faithful error): differential substitution and normalization}
\noindent\textit{Transcription reads.} a maps to a(phi) phi prime for all invertible phi while a(0)=1\par
\noindent\textit{Critical reading.} This preserves normalization for strict phi; for general phi divide by phi prime(0)\par
\noindent\textit{Evidence.} Chain rule gives the pullback; its constant coefficient is phi prime(0).\par
\subsection*{P000448 (both; ambiguity): second inverse-coordinate formula}
\noindent\textit{Transcription reads.} Dot above phi inverse before integral a dx\par
\noindent\textit{Critical reading.} Use the inverse coordinate itself as the primitive; retain the dot as ambiguous manuscript evidence\par
\noindent\textit{Evidence.} Scan has an isolated mark; interpreting it as differentiation is incompatible with the preceding formula.\par
\subsection*{P000449-P000451 (both; provenance): record 270/271 boundary}
\noindent\textit{Transcription reads.} Supplied record 270 range includes P000450-P000451; record 271 is labelled 1972\par
\noindent\textit{Critical reading.} Record 270 ends at P000449; record 271 is a catalogue heading at the P000449/P000450 boundary\par
\noindent\textit{Evidence.} P000450 is headed 25/IX-72 Gabrielov. P000449 contains the final G/S assertion and closing dots only.\par
\subsection*{P000477-P000479 (both; provenance): record 286 boundary}
\noindent\textit{Transcription reads.} Supplied YEAR 1973 heading spans P000477-P000479\par
\noindent\textit{Critical reading.} Treat 286 as a heading at the P000477/P000478 boundary; assign no mathematical items to it\par
\noindent\textit{Evidence.} P000477 continues the 20 November 1972 Buchstaber lecture; P000478 is headed 12 February 1973, Buchstaber and continues on P000479.\par
\subsection*{P000479-P000481 (both; provenance): record 288 supplied page span}
\noindent\textit{Transcription reads.} Assign the entire supplied range P000479-P000481 to Dobrushin\par
\noindent\textit{Critical reading.} Only P000480 belongs to Dobrushin. P000479 continues Buchstaber’s 12 February 1973 lecture; P000481 begins his 26 February 1973 lecture\par
\noindent\textit{Evidence.} Independent dates and speaker headings on the scans identify the three boundaries.\par
\subsection*{P000480 (both; provenance): record 288 date}
\noindent\textit{Transcription reads.} Catalogue: 19 Feb 1972 / 1972-02-19\par
\noindent\textit{Critical reading.} 19 February 1973 / 1973-02-19; academic-year item prefix 1972\par
\noindent\textit{Evidence.} Manuscript heading explicitly 19/II-73, following YEAR 1973 and 12 February 1973, and preceding 26 February 1973. Catalogue date is preserved as a provenance exception.\par
